% Framework v3.069: five-core axial/decay same-action matching; conditional only.
% Framework source amendment v3.068: scoped geometry/operator additions, physical occurrence unchanged.
% Paper source Paper7A-QCDKernel-v1_179.tex: v3.040 S3 literature mathematical/conditional comparison; physical EM packet unresolved.
% Framework S3 finite-response interface v3.035: physical Gauss/source-response definition, no additional operator or action assumption.
% Framework v3.027: exact 600-cell/Hurwitz return channels, Cartan rank, native clutch, Clifford constraints and conditional physical gates.
% Framework v3.026: common/relative tight-frame stability and conditional same-owner five-core transport.
% Framework v3.025: scoped QCD nonlinear-action sensitivity and temporal-transfer positivity; no physical owner promotion.
% Paper7A-QCDKernel-v1_174: polyhedral full-2O working carrier and exact 2I meson/baryon geometry; physical occurrence remains conditional.
% Paper-series v2.381: complete algebraic exceptional torsion atlas and all-cocycle twist-stability; original claims retained.
% Paper-series v2.380: unified bibliography, exact exceptional double/length/gauge mathematics in I-B; previous theorem labels retained.
% Paper-series v2.379: full equivariant Dirac/Hodge/eta/torsion atlas in Paper I-B; paper-root references synchronized.
% Paper-series v2.378: 2026 exceptional literature consolidation and reference synchronization. The mathematical additions are located in Paper I-B.
% Paper-series v2.377: reference-only rebind to current Paper I-B v1.1; original mathematics retained except new Paper I-B comparison chapter.
% Framework v2.370: derived exceptional-root/Hurwitz comparison in I-A/I-C; other papers reference-only synchronized.
% Framework v2.368: preserve prior II-B convention/unital fixes and verified II-D corrections; current references synchronized.
% Framework v2.366: exact chamber--transporter/Hurwitz integration; physical action and hadron decay gates remain conditional.
% Framework v2.366: scoped Hurwitz/600-cell hadron-source integration; physical APBI/decay frontier remains open.
% Framework v2.358: current 30-root source reference synchronization; selectively updated III-E, III-F, VI-B, VI-C.
% Framework v2.357: normal-jet/Clifford/hadron-shell observable proof additions at scoped mathematical, conditional, or QCD-source levels.
% Framework v2.356: new source-owned conditional results and source-registry synchronization.
% Framework v2.350: repaired Paper VII mathematics; exact finite identities separated from analytic/physical certificates.
% Framework v2.349: integrated exact conditional mathematical closures; G0 and GP1--GP3 retained as axioms.
% Framework v2.345: F40 native three-readout and source-visibility integration.
% Framework v2.344: Twisted Hurwitz C3/Petersen bridge, native baryon source and same-action curvature firewall.
% Framework v2.338: exact Petersen 120-carrier, oriented two-source and source-ownership results. No hadron mass prediction.
% Framework v2.327: exact two-core graded transport, conditional meson spin source, native mass-spectrum firewall.
% Framework v2.323: Petersen finite strong-source coherence, exact multi-core mirror spectra, QCD action-incidence and continuum firewalls. Existing source theorems preserved.
% Framework v2.275: exact positive dense-twin word cubature, finite repeatability and eighth-moment classification, nested sampling frames, joint action and POVM-jet exports; preserves v2.274.
% Framework v2.274: two-scalar positive convolution certificates, first-leaked exceptional coset detection, branch-symmetric twin mixing, and finite endpoint moments; all v2.273 branches preserved.
% Framework v2.273: restored missing Problem 3 theorems from v2.268; exact exceptional frames, twin geometry, convolution repeatability and projector/occurrence firewalls; merged on v2.272 without overwriting independent branches.
% Framework v2.272: merged v2.270 covariant native-readout mathematics with v2.271 Problem 4 finite-action Hopf hypercharge/holonomy corrections; no universal bulk-gauge uniqueness claim.
% Independent branches: v2.270 readout/guard exports and v2.270 K16 action-consumed charge closure; all prior scoped assumptions retained.
% Framework v2.267: exact Paper II-A converse/observable certificates, III-F action tests, V-B right orbit and V-C native gauge scope; prior v2.266 preserved.
% Framework v2.266: exact Paper II-C outcome-count/readout converse results and scoped downstream imports; 30-root references synchronized.
% Framework v2.265: II-C coordinate-stage constitution and exact first-symbol host inheritance; prior mathematical domains retained.
% Framework v2.264: II-B reverse and action-bridge integration; preserves prior mathematical domains; active 30-root references synchronized.
% Framework v2.263: Paper I-D adds bridge-controlled root-carry parity and source-relative old-cap cancellation; other root mathematical bodies retained; active 30-root references synchronized.
% Framework v2.262: merges Paper II-F three-gap closure from v2.260 into v2.261, preserving newer II-B/II-D repairs; 30-root references synchronized.
% Framework v2.261: merged Paper II-B GP2-level/geometric repairs; preserves v2.260 II-D repairs; active references synchronized.
% Framework v2.260: Paper II-D signed RS, finite-symbol typing, and conditional curvature/Jordan repairs; references synchronized.
% Framework v2.259: conditional native SU(2) atlas and explicit later discharge; v2.258 mathematical bodies retained elsewhere.
% Framework v2.258: preserves parallel II-A/dynamics update and adds complete readout/frequency closures.
% Framework v2.257: exact minimal complete position readout and arbitrary-preparation frequency law.
% Framework v2.256: exact native inertia and sharp bounded ordered-source closure; preserves concurrent v2.255 work.
% Framework v2.255: preserves parallel detector integration and adds exact finite readout closures.
% Framework v2.254: shared bibliography/reference sync; substantive Pasterski detector-Ward additions only IV-B/IV-C/VII-C.
% Framework v2.253: typed raw coordinate defect in I-A and synchronized VI-A import; active references synchronized.
% Framework v2.252: optional shell-frame source-return calculation in III-G; active references synchronized.
% Framework v2.251: forward common-action execution, joint source jets and scoped coupled control.
% Framework v2.251: exact full-heat variations, fermion source jets and explicit physical completion frontier.
% Framework v2.251: certified interacting source jets, pole stability and scoped completion proofs.
% Framework v2.251: interacting finite-source evaluation and coupled local vacuum classification.
% Framework v2.251: exact finite source execution and native causal common-parent bridge.
% Framework v2.251: constructive parent and scoped source-incidence audit; see release report.
% Framework v2.251: common-action two-jet, corrected source/response closure and synchronized references.
% Framework v2.244: scoped rooted source-order Schur diagnostics and synchronized paper references.
% Framework v2.242: signed rooted phase lift, complete transverse return and exact family cubic.
% Framework v2.241: proved rooted conormal geometry, conditional cap realization and scoped action tests.
% Framework v2.239: coherent 30-paper reference rebinding; selective new proofs on established owners.
% Framework v2.236: source registry synchronized; existing mathematical body retained except declared insertions.
\documentclass[12pt,a4paper]{article}
\PassOptionsToPackage{hidelinks}{hyperref}

% Paper-specific support packages not supplied by the shared series macro file.
\usepackage{microtype}
\usepackage{etoolbox}

% Macros-v2_203: v3.077 source-current bindings for unconditionality source integration. No mathematical macro changes.
%% Macros-v2_201: downstream full-carrier geometry update, five current root bindings, no mathematical macro change.
%% Macros-v2_200: Paper III-F v1.280 current-source binding for framework v3.074; mathematical macro definitions unchanged.
%% Macros-v2_198: current project v3.072 III-D and bibliography binding; mathematical macros unchanged.
%% Current project v3.070 source-reference bindings; previous macro bodies unchanged.
% Current II-D reference root for PapersMinimal v3.063; mathematical macros otherwise unchanged.
% Macros-v2_188: re-based current paper identities and A--G family; no change to geometry.
% Macros-v2_184: active Paper III-E v1.200/III-F v1.267 binding for framework v3.031; mathematical macros unchanged.
% Macros-v2_183: updated shared cross-reference bindings for Paper1B/III-E/III-F v3.030; mathematical macros unchanged.
% Macros-v2_182: active Paper II-F/V-A cross-references for framework v3.029; definitions unchanged.
% Paper-series v2.381: new I-B torsion/modular results, bibliography and current source roots.
% Paper-series v2.380: new exceptional bibliography and synchronized paper references.
% Paper-series v2.379 current 31-paper macro/aux registry.
% Paper-series v2.378: current 31-paper AUX/bibliography registry updated for Paper I-B v1.2 literature additions.
% Framework v2.371: current-source reference registry refreshed; old theorem labels retained.
% Framework v2.370: current 30-paper source/AUX root rebinding; native mathematical macros unchanged.
% Framework v2.368: synchronized merged 30-root registry; mathematical notation unchanged.
% Framework v2.358: active 30-paper source registry synchronized; existing mathematical macros unchanged.
% Framework v2.357: selective new source-owned certificates; active reference roots updated for five revised papers.
% Framework v2.356: active 30-root registry synchronized to the 30 papers in this release.
% Framework v2.275: exact positive dense-twin word cubature, finite repeatability and eighth-moment classification, nested sampling frames, joint action and POVM-jet exports; preserves v2.274.
% Framework v2.274: filename and root-source synchronization.
% Framework v2.273: restored missing Problem 3 theorems from v2.268; exact exceptional frames, twin geometry, convolution repeatability and projector/occurrence firewalls; merged on v2.272 without overwriting independent branches.
% Framework v2.272: merged v2.270 covariant native-readout mathematics with v2.271 Problem 4 finite-action Hopf hypercharge/holonomy corrections; no universal bulk-gauge uniqueness claim.
% Independent branches: v2.270 readout/guard exports and v2.270 K16 action-consumed charge closure; all prior scoped assumptions retained.
% Framework v2.269: synchronized 30 active AUX roots; mathematics and bibliography preserved from Macros-v2_161.
% Framework v2.266: Paper II-C converse-derived outcome resources; II-E, III-C, IV-A/B, V-A and VI-B downstream mathematical closures; 30-root registry synchronized.
% Framework v2.265: II-C coordinate-stage constitution and exact first-symbol host inheritance; prior mathematical domains retained.
% Framework v2.264: paper II-B reverse/audit and IV-A observable-minimal exchange; synchronized external references to 30 active roots.
% Framework v2.263: Paper I-C adds bridge-controlled root-carry parity and source-relative old-cap cancellation; other root mathematical bodies retained; active 30-root references synchronized.
% Framework v2.262: integrates Paper II-F radius, Casimir and readout results; active 30-root references synchronized.
% Framework v2.261: merged Paper II-B GP2-level/geometric repairs; preserves v2.260 II-D repairs; active references synchronized.
% Framework v2.260: Paper II-D signed RS, finite-symbol typing, and conditional curvature/Jordan repairs; references synchronized.
% Framework v2.259: conditional native SU(2) atlas and explicit later discharge; v2.258 mathematical bodies retained elsewhere.
% Framework v2.258: preserves parallel II-A/dynamics update and adds complete readout/frequency closures.
% Framework v2.257: exact minimal complete position readout and arbitrary-preparation frequency law.
% Framework v2.256: exact native inertia and sharp bounded ordered-source closure; preserves concurrent v2.255 work.
% Framework v2.255: preserves parallel detector integration and adds exact finite readout closures.
% Framework v2.254: Pasterski bibliography and synchronized 30-paper external-reference registry.
% Framework v2.253: typed raw coordinate defect in I-A and synchronized VI-A import; active references synchronized.
% Framework v2.252: synchronized source registry; optional III-G shell-frame calculation.
% Framework v2.251: forward common-action execution, joint source jets and scoped coupled control.
% Framework v2.251: exact full-heat variations, fermion source jets and explicit physical completion frontier.
% Framework v2.251: certified interacting source jets, pole stability and scoped completion proofs.
% Framework v2.251: synchronized active source registry; mathematical notation unchanged.
% Framework v2.251: synchronized active source registry; mathematical notation unchanged.
% Framework v2.251: synchronized active source registry; mathematical notation unchanged.
% Framework v2.251: synchronized source registry; no mathematical notation changes.
% Framework v2.244: synchronized 30-paper AUX registry and rooted phase-source action-order tests.
% Framework v2.242: signed rooted phase lift, complete transverse return and exact family cubic.
% Framework v2.241: proved rooted conormal geometry, conditional cap realization and scoped action tests.
% Framework v2.239: all 30 active source roots synchronized; no mathematical macro changes.
% Framework v2.238: Paper IV A/B/C split and 30-root reference registry.
% Framework v2.237: Paper VII-D technical split and 28-root reference registry.
% Framework v2.236: Paper VI-C technical-volume reference added; all current root bindings advanced.
% Framework v2.235: 26-root AUX and bibliography bindings synchronized; mathematical macros unchanged.
% BEGIN INLINE SOURCE: Paper00-TeXBuildCompatibility-v1_0.tex
% Rendering-only repairs for inherited long appendices and bitmap text fonts.
% Integer counters through 26 retain the original labels; later values use AA,
% AB, ... . This file changes no action coefficient, state or theorem premise.
\RequirePackage{alphalph}
\makeatletter
\let\BHPrimitiveAlph\@Alph
\newalphalph{\BHExtendedAlph}{\BHPrimitiveAlph}{26}
\let\@Alph\BHExtendedAlph
\AtBeginDocument{%
  \@ifpackageloaded{microtype}{\microtypesetup{expansion=false}}{}%
}
\makeatother

% END INLINE SOURCE: Paper00-TeXBuildCompatibility-v1_0.tex

% Framework v2.164: current reference/bibliography profile; mathematical definitions retained.
% Framework v2.162: finite metric, functional Hilbert and scoped Ward integration.
% Framework v2.161: inherits Macros-v2_57.tex exactly except the current external-reference binding.
% Macros v2.57 / framework v2.156: current reference registry only; mathematical definitions unchanged.
% Framework v2.143 / Macros v2.55: current external-reference bindings; mathematical notation unchanged.
% Framework v2.110 / Macros v2.54: current external-reference registry and bibliography; mathematical notation inherited.
% Framework v2.67 / Macros-v2_53: adds Born-action-cell and joint-EH readout symbols; absolute unit realization is metrological; MATCH_{J/scal}^{EH} remains the sole G24 normalization bridge.
% Framework v2.35 / Macros-v2_52: grouped paper numbering, active AUX registry, and Born-v6_20 bibliography binding synchronized; notation unchanged.
% Framework v2.26: active-source AUX and bibliography bindings synchronized; notation unchanged.
% Framework v2.09: reviewed early-reciprocity integration; scoped physical imports.
% Framework v2.08: Born--Euler mathematics and conditional complete-action potential criterion.
% Framework v2.07: active AUX registry includes Paper I-B v1.2; shared bibliography Born-v6_18.
% Macros-v2_46 / framework v2.04: mathematical atlas and downstream realization correction.
% Macros-v2_46 / framework v2.03: scoped structure integration and current dependency bindings.
% Macros-v2_46 / framework v2.01: shared source bindings; prior revision comments are historical.
% v2.42 / framework v1.99: Synchronized RG/AS owner and bibliography bindings.
% v2.41 / framework v1.99: active manuscript AUX bindings and geometric-closure bibliography synchronized.
% Release-v1.80 centralized source synchronization for scalar-circle/comparison-band integration.
% Release v1.79: Hurwitz hierarchy and kinetic-normalized action response integration; physical mass matching remains conditional.
% Release v1.78: approved M8 closure-batch integration; physical and analytic domains retained.
% Current v1.78 registry: all active TeX roots are synchronized below.
% Earlier release comments are historical scope records, not current status.
% Release-v1.77 synchronization: native same-owner Noether closure, generator--flow/history typing, analytic-certification firewalls, scalaron/K16, LF/QCD, quantum-Ward, and RAR boundary-owner transport are registered; shared mathematical notation is unchanged.
% Release-v1.76 synchronization: registry points to the v1.76 active paper roots.
% Release-v1.76 synchronization: registers the local A+Phi+g Noether/constraint closure and dynamics-frontier propagation; shared mathematical notation is unchanged.
% Release-v1.75 synchronization: registry points to the v1.75 active paper roots.
% Release-v1.75 synchronization: integrates the joint-Dirac Einstein bulk real form and the local odd-scalar heat-coefficient action on explicit analytic first-jet/interior domains; v1.74 Yang--Mills closure is retained. Fermionic reality, pure-metric curvature-squared/exact-nonlocal terms, noncanonical boundaries, physical vacuum selection, and actual solution-history tokens remain open.
% Release-v1.74 synchronization: registry points to the v1.74 active paper roots and retains the v1.73 theorem labels.
% Release-v1.74 synchronization: canonical Macros-v2_34.tex; generated local nonlinear Yang--Mills action, constrained connection-field phase space, source-free Gauss momentum map, and Legendre closure are integrated on their proved domains. The remaining non-gauge ACT_cl sectors and actual solution-history selection remain open.
% Release-v1.71 synchronization: closes the former JOINT_x/JOINT_5
% connection-occurrence gap by canonical structure-group extension of the
% nonflat Born connection, faithful associated representation, and the local
% cotangent--Clifford operator constructor.  Actual fluctuation members,
% actions, vacua, solutions, calibrations, and analytic QFT domains remain
% separately typed.  The active-AUX registry is advanced once.
% Release-v1.71 canonical macro and active-AUX registry.
% ======================================================================
% Macros-v2_39.tex -- canonical publication macro file
% Release v1.71 retains the v1.70 null-ray/QCD boundaries and additionally
% removes the duplicate structural connection actuality premise.  It does not
% claim global bundle topology, select a gauge/Higgs member, or promote the
% finite construction to an interacting continuum theory.
% Release v1.68 constructs the finite light-front total-fraction operators and
% parton number/momentum operator measures, proves their simplex and Hausdorff-limit
% theorems, and separates positive cut measures from scheme-renormalized
% light-ray distributions.  It also separates finite real Schur poles, stable
% norm-resolvent/Riesz pole-residue limits, and continued resonance
% pole/residue limits.  The Paper-VII-C scattering package is claim-indexed:
% local QFT and spectral claims import no S-matrix or infrared premise, while
% asymptotic amplitudes, LSZ, S-matrix poles, completeness, and cross sections
% import only their named massive or massless/long-range clauses.
% Release v1.67 factorizes BND48_H into a same-action Hessian/Schur
% descendant and still-open typed boundary-trace, critical-lift, and
% refinement-limit conditions.  It proves the finite-collar Dirichlet-to-
% Neumann law sqrt(y)tanh(L sqrt(y)), its nonuniform critical limit, and the
% finite-return ultra-infrared turnover.  The Paper-I-A Hardy constant mode is
% an exact geometric near-hit, not yet the physical metric boundary mode.
% The current local f(R) Hessian remains massive and therefore still refutes
% a RAR prediction on the minimal action branch.
% Release v1.66 derives the transverse depth-six normalizer bridge from the
% already generated C6 Reynolds image and Paper-I-A binary-dihedral normalizer.
% The bridge is the unique least 2D3-stable scalar-containing realization of
% the faithful quotient 2D3/C6=C2, has trace defect six and hence K=48, but is
% transverse to the nested active carrier.  Exact K=44 and K=52 contraction,
% semigroup, and energy-midpoint identities are registered as structural
% readouts rather than projector events.  The active E48 carrier gap remains
% unchanged.  A same-action boundary-symbol/Calderon intertwiner is isolated
% as the necessary and still unproved owner bridge to the static DtN sort;
% representation equivalence alone neither derives RAR nor creates a K=48
% action junction, pole, occupation, or cosmological transition.
% Release v1.65 derives the native E_48-centered carrier mass gap from the
% GP1--GP3-generated active spectrum, distinguishes its carrier-free interval
% from a particle located at E_48, and withdraws the unsupported promotion to
% a positive continuum Yang--Mills mass gap.  The latter remains a valid
% conditional transfer theorem on an enlarged continuum branch, but it is not
% an active framework prediction or an unfinished native-gap proof.
% Release v1.64 synchronizes Derivation Batches 14--22.  It registers the
% GP1--GP3 heat-profile orbit, same-Hessian electroweak response and pole
% ownership, K16/K48 action-response corrections, perturbation state and
% non-Gaussianity typing, finite-vacuum/gap/QCD limit repairs, factorized
% position-area ownership, the actual static metric response, and the ordered
% Wilson/string-tension limits.  Strict continuum convergence, positive
% physical gap, surviving higher cumulants, and positive tension remain open.
% Release v1.61 closes TIMEWIDTH16 as an independent direct-production datum.
% JUNC_{2|7} continues to own the finite intrinsic microscopic profile/width,
% while the actual same-action background/production history determines its
% physical-time pullback, support, and duration.  The current Paper-I-A--VII-C AUX
% registry is advanced once; shared mathematical notation is unchanged.
% Release v1.60 derives the final-octet primitive-remnant depletion, canonical
% Hopf-current representative, exact Hopf descent to dOmega/(4pi), and
% closed-FLRW a^{-3} neutral-current shape.  It retypes K=48 as the onset
% of the final one-unit geometric transfer and records cDM/RAR as distinct
% physical projections of that common neutral parent while preserving
% occurrence, Hilbert-owner, stress-matching, and static-kernel premises.
% Release v1.59 transports the v1.58 exceptional-shell/K16 closure through
% Papers III-F and VI-B--VII-C.  It replaces Paper VI-B's obsolete mixed-Hessian
% output typing by the third-variation pair vertex and Bogoliubov occupation
% support, derives the post-seed Higgs-pair feedback vertex and its exact
% alpha--delta kernel factors, propagates the unique K16/K48 action-junction
% split, identifies the QCD delta count with the upstream 2O shell atom only
% as a common-scalar numerical compatibility, and narrows THERM16 to dynamical
% approach/beta selection because the finite Gibbs/KMS target is already
% unique on a supplied simple observable algebra once its Hamiltonian and
% inverse temperature are supplied; direct-sum central weights remain typed.
% Release v1.58 registers the exceptional nearest-shell closure, the derived
% transverse K=16 action event, the unique interior action junction, the
% type-correct ordered pair vertex, the thin-layer squeeze law, and the
% radiation-like/matter-like kinematic split.  REG_0, transfer magnitude,
% physical layer width, and thermalization remain explicitly conditional.
% Release v1.57 registers the active-depth-six obstruction, the native
% K=48 carrier desert, the typed 7|9 owner reduction, and the dynamical
% reheating handoff across Papers III-B--VI-B.
% Release v1.56 registers the least-support n_2 inflation dressing of the
% persistent n_7 Laurent response, the analytic global-convexity and positive
% endpoint-measure proofs, the derived inflation core and K=16 coefficient
% freeze, and the no-particle dark-sector split into static RAR plus the
% conditional homogeneous effective-curvature response at reflected K=48.
% There is no n_6 action, scalaron-DM mass, or second Einstein--Hilbert term.
% Release v1.55 registers the protected-octet resolved/residual reflection
% q_sp+Z_W=8, its j<->8-j and K<->64-K scale covariance, and the derived
% rank-one-exit quadratic grade C_1^2/U=C_2 in Papers III-C, III-E, and VIII.
% The reflection is a typed trace/scale theorem, not a time reversal or a
% physical event-occurrence theorem; reheating transfer remains conditional.
% Release v1.54 registers the exact positive-part Laurent indices, the
% active-domain/terminal-endpoint action split, the no-second-EH firewall, and
% the q_DM=6/K_DM=48 scalaron readout typing across Papers III-D/E/F, VI,
% VIII, and IX.
% Release v1.53 registers the common nonnegative Laurent obstacle projection
% for n_1^+ and n_7^+, discharges the three former exit/freeze assumptions,
% and advances the Paper-III-D/E/F, VIII, and IX consumers.  Shared notation
% and bibliography ownership are unchanged.
% Release v1.51 registers the PQTI derivation batch and its updated
% Paper-III-A/C/D/E/F, VIII, and IX consumers.  Shared notation and
% bibliography ownership are unchanged.
% Release v1.50 advances the Paper-III-B AUX target to v1.40 after the
% rigorous audit repair; shared notation and bibliography ownership are unchanged.
% Release v1.49 advances the Paper-II-E/F AUX target to v3.14 after the
% minimal audit repair; shared notation and bibliography ownership are unchanged.
% Release v1.46 synchronizes Paper II-A's minimal frozen constructor-origin
% theorem and its five dependent AUX targets.  Shared notation is unchanged.
% Release v1.45 updates the Paper-II-B AUX target to v2.23 after the
% GP3-rigidity and appendix-placement repair; shared notation is unchanged.
% Release v1.44 advances the centralized Paper-I-A--VII-C AUX registry after the
% Paper-I-A v5.16 protected-completion scope repair.  Every current
% paper imports this one file; no paper-level opt-in or \externaldocument list
% is permitted.  The bibliography remains centralized here through Born-v6_13.
% Shared bibliography, theorem infrastructure, colors, notation alignment,
% and current Paper-I-A--VII-C external-reference registration.
% This file centralizes common imports, bibliography sorting, colors, theorem
% infrastructure, and the notation shared by the complete series.
% v2.9 (release v1.43) made the Paper-I-A--VII-C AUX registry automatic and
% centralized the bibliography and all external-reference ownership.
% v2.8 release-registry refresh (v1.33) registers the rooted generated-parent
% cascade and the POL/REG/free-law/11+4/Born-group closure revisions.  No
% shared mathematical notation or Paper-II-A fundamental principle is changed.
% v2.8 release-registry refresh (v1.32) registers the universal generated-
% continuation revisions in Papers II-A, III-A, III-C, IV, and VIII.  No
% shared mathematical notation or Paper-II-A fundamental principle is changed.
% v2.8 release-registry refresh (v1.30) registers the CRP4/projective-state
% source revisions and their direct registry-only consumers.  No shared
% mathematical notation or Paper-II-A axiom is changed.
% v2.7 release-registry refresh (v1.27) updates only the current external
% Paper-I-A--VII-C basenames after the placement/assumption synchronization; no
% shared mathematical notation or Paper-II-A axiom is changed.
% v2.7 synchronizes the geometric defect-repair revision: Paper-II-A v3.42,
% Paper-III-C v1.26, Paper-III-D v1.27, Paper-III-E v1.32, and Paper-III-F
% v1.34.  It registers the common trace-defect/spectral-filtration theorem,
% its owner-typed response pullback, and the reduced GRED and TERM interfaces.
% No shared mathematical notation is changed.
% v2.6 synchronizes Paper-III-E v1.31's cone-type repair, minimal-response
% classification, and explicit classical action-role boundary, together with
% Paper-III-F v1.32's downstream response-atlas synchronization.  No shared
% mathematical notation is changed.
% v2.5 synchronizes the Paper-III-D v1.26 faithful-trace-duality proof,
% exact ordered-owner joint locus, and typed alternative-presentation
% bookkeeping repair, together with its Paper-III-E and closure-frontier
% references.  No shared mathematical notation is changed.
% v2.4 synchronizes the Paper-III-C v1.25 orbit/process/owner repair, its
% downstream Papers III-B--F, V, VI, VIII, and IX revisions, and Born-v6_13.
% v2.3 synchronizes the Paper-II-A v3.41 witness/support/principal inheritance
% layer, the Paper-III-A v1.24 joint-realizability/readout firewall, all dependent
% host-rate revisions, and the current Paper-I-A--VII-B source registry.
% v2.2 synchronizes the generated active-realization, determinant-depth,
% dimension-weighted hierarchy, and typed cosmological-event interfaces.
% v2.1 synchronizes the exceptional-support interface revisions: Paper I-A
% v5.9, Paper II-A v3.30, Paper II-E/F v3.7, and Paper IV v4.16.
% v2.0 introduces the stage-safe G0 occurrence notation, synchronizes the
% active Paper-I-A--IV source registry (current rigorous Paper-III closure:
% III-A v1.17, III-B v1.22, III-C--E v1.17, III-F v1.14, IV v4.14), and makes
% external-reference self-skip depend on logical paper role rather than a
% versioned job name.
% ======================================================================

% ======================================================================
% macros.tex — Shared macros for the Born Geometry series
% ======================================================================

% ======================================================================
% 1. PACKAGES, BIBLIOGRAPHY, COLORS, AND UNIT DEFINITIONS
% ======================================================================
% Core packages shared across the complete Paper-I-A--VII-C series.  Loading them here makes
% the macro file a genuine series preamble while remaining compatible with existing
% paper preambles that load the same packages before \input.
\usepackage{amsmath}
\usepackage{amssymb}
\usepackage{amsthm}
\usepackage{mathrsfs}
\usepackage{array}
\usepackage{booktabs}
\usepackage{xcolor}
\makeatletter
\@ifpackageloaded{geometry}{}{\usepackage{geometry}\geometry{margin=1in}}
\makeatother
\usepackage{mdframed}

% Shared support packages required by the macros below and by most papers.
\usepackage{ragged2e}
\usepackage{longtable}
\usepackage{tabularx}
\usepackage{enumitem}
\usepackage[most]{tcolorbox}

% Series-wide numeric bibliography convention using the portable BibTeX path.
% All Papers I-A--VII-C use the ordinary \cite interface.
\usepackage[numbers,sort&compress]{natbib}
\providecommand{\printbibliography}{%
  \bibliographystyle{plainnat}%
  \bibliography{Born-v6_60}%
}

% Shared colors.  \providecolor avoids clobbering a paper-specific override.
\providecolor{postulatecolor}{RGB}{230,242,255}
\providecolor{compare}{RGB}{230,245,230}
\providecolor{definitioncolor}{RGB}{230,242,255}

\providecolor{openriskcolor}{RGB}{255,235,235}
\providecolor{bridgecolor}{RGB}{240,240,255}
\providecolor{gold}{RGB}{255,215,0}

% ======================================================================
% 2. CUSTOM ENVIRONMENTS
% ======================================================================
% Common table column types used in all papers.  Guarded so that a paper may
% define its own versions before loading this file.
\makeatletter
\@ifundefined{NC@find@L}{%
  \newcolumntype{L}[1]{>{\RaggedRight\arraybackslash}p{#1}}%
}{}
\@ifundefined{NC@find@C}{%
  \newcolumntype{C}[1]{>{\Centering\arraybackslash}p{#1}}%
}{}
\@ifundefined{NC@find@Y}{%
  \newcolumntype{Y}{>{\RaggedRight\arraybackslash}X}%
}{}
\makeatother

% ======================================================================
% 3. THEOREM STYLES AND ENVIRONMENTS
% ======================================================================
% Custom underlined theorem style (useful for remarks or notes)
\newtheoremstyle{underlined}
  {3pt}   % Space above
  {3pt}   % Space below
  {}      % Body font (empty means normal/roman text, typical for remarks)
  {}      % Indent amount
  {}      % Theorem head font (can be \itshape, \bfseries, etc.)
  {.}     % Punctuation after theorem head
  { }     % Space after theorem head
  {\underline{\thmname{#1}\thmnumber{ #2}}\thmnote{ (#3)}} % Theorem head formatting

% Standard plain theorem style environments.  Each environment is guarded so
% the shared macro file can be loaded after a paper-specific theorem setup.
\theoremstyle{plain}
\makeatletter
\@ifundefined{theorem}{\newtheorem{theorem}{Theorem}[section]}{}
\@ifundefined{lemma}{\newtheorem{lemma}{Lemma}[section]}{}
\@ifundefined{proposition}{\newtheorem{proposition}{Proposition}[section]}{}
\@ifundefined{corollary}{\newtheorem{corollary}{Corollary}[section]}{}
\makeatother

% Definition style environments (typically non-italicized body)
\theoremstyle{definition}
\makeatletter
\@ifundefined{definition}{\newtheorem{definition}{Definition}[section]}{}
\@ifundefined{postulate}{\newtheorem{postulate}{Postulate}[section]}{}
\@ifundefined{closure}{\newtheorem{closure}{Closure}[section]}{}
\makeatother

% Applying the custom underlined style
\theoremstyle{underlined}
\makeatletter
\@ifundefined{remark}{\newtheorem{remark}{Remark}[section]}{}
\@ifundefined{convention}{\newtheorem{convention}{Convention}[section]}{}
\@ifundefined{principle}{\newtheorem{principle}{Principle}[section]}{}
\makeatother

% Framed proof environment.
% Every \begin{proof} ... \end{proof} is automatically framed/shaded.
% This avoids mdframed/trivlist nesting, which can drop the top border after
% theorem-like environments.
\makeatletter
\renewenvironment{proof}[1][\proofname]{%
  \par\addvspace{6pt}%
  \noindent
  \begin{tcolorbox}[%
    enhanced,
    breakable,
    colback=black!4,
    colframe=black!35,
    boxrule=0.35pt,
    arc=1pt,
    left=6pt,
    right=6pt,
    top=4pt,
    bottom=4pt,
    boxsep=0pt,
    before skip=0pt,
    after skip=6pt,
    parbox=false
  ]%
  \pushQED{\qed}%
  \normalfont\itshape #1\@addpunct{.}\enspace
  \normalfont\ignorespaces
}{%
  \unskip\nobreak\hfill\popQED
  \end{tcolorbox}%
  \par\addvspace{2pt}%
}
\makeatother

\mdfdefinestyle{postulatebox}{
  backgroundcolor=yellow!10,
  linecolor=black,
  linewidth=0.4pt,
  skipabove=\baselineskip,
  skipbelow=\baselineskip,
  nobreak=true
}

% ======================================================================
% PROOF-READABILITY HELPERS
% ======================================================================

\newcommand{\ProofPlan}[1]{%
  \par\smallskip
  \noindent\textbf{Proof plan. }#1%
  \par\smallskip
}

\newlist{proofsteps}{enumerate}{2}
\setlist[proofsteps,1]{
  label=\textbf{(\alph*)},
  ref=\alph*,
  leftmargin=2.2em,
  itemsep=0.45\baselineskip,
  topsep=0.35\baselineskip
}
\setlist[proofsteps,2]{
  label=\textbf{(\roman*)},
  ref=\roman*,
  leftmargin=2.0em,
  itemsep=0.25\baselineskip,
  topsep=0.2\baselineskip
}


% ======================================================================
% 4. NUMBER SYSTEMS AND ALPHABETS
% ======================================================================
\providecommand{\RR}{\mathbb{R}} % Real numbers
\providecommand{\PP}{\mathbb{P}} % Projective space / Probability
\providecommand{\CC}{\mathbb{C}} % Complex numbers
\providecommand{\ZZ}{\mathbb{Z}} % Integers
\providecommand{\NN}{\mathbb{N}} % Natural numbers
\providecommand{\HH}{\mathbb{H}} % Quaternions / Hyperbolic space
\providecommand{\LFRayGeom}{\mathsf{LF\!\mbox{-}RAY\!\mbox{-}GEOM}}
\providecommand{\LFRep}{\mathsf{LF\!\mbox{-}REP}}
\providecommand{\LFState}{\mathsf{LF\!\mbox{-}STATE}}
\providecommand{\LFLim}{\mathsf{LF\!\mbox{-}LIM}}
\providecommand{\LFScat}{\mathsf{LF\!\mbox{-}SCAT}}
% Number systems / projective spaces
\providecommand{\CP}{\mathbb{CP}}
\providecommand{\RP}{\mathbb{RP}}

% ======================================================================
% 5. GROUPS AND LIE THEORY
% ======================================================================
\providecommand{\SO}{\ensuremath{\mathrm{SO}}}     % Special Orthogonal group
\providecommand{\SU}{\ensuremath{\mathrm{SU}}}     % Special Unitary group
\providecommand{\Spin}{\ensuremath{\mathrm{Spin}}} % Spin group
\providecommand{\CS}{\ensuremath{\mathrm{CS}}}     % Chern-Simons (capitalized)
\providecommand{\cs}{\ensuremath{\mathrm{cs}}}     % Chern-Simons (lowercase)

% ======================================================================
% 6. COMMON MANIFOLDS AND TOPOLOGY
% ======================================================================
\providecommand{\Sone}{S^1}                 % Standard circle

% ======================================================================
% 7. MATH OPERATORS (Declared via \DeclareMathOperator)
% ======================================================================
\DeclareMathOperator{\End}{End}   % Endomorphisms
\DeclareMathOperator{\Aut}{Aut}   % Automorphisms
\DeclareMathOperator{\spec}{spec} % Lowercase spectrum
\DeclareMathOperator{\tr}{tr}     % Lowercase trace (if you already have \Tr)
\DeclareMathOperator{\supp}{supp} % Support
\DeclareMathOperator{\TrA}{Tr_{A}} % Partial trace over subsystem A
\DeclareMathOperator{\TrB}{Tr_{B}} % Partial trace over subsystem B
% Lie/algebra operators
\DeclareMathOperator{\Ad}{Ad}
\DeclareMathOperator{\Sym}{Sym}
\DeclareMathOperator{\su}{su}
% ======================================================================
% 8. STANDARD OPERATOR-LIKE SYMBOLS (Safe definitions via \providecommand)
% ======================================================================
\providecommand{\vol}{\operatorname{vol}}   % Volume form
\providecommand{\Vol}{\operatorname{Vol}}   % Total volume
\providecommand{\Ric}{\operatorname{Ric}}   % Ricci tensor
\providecommand{\Rm}{\operatorname{Rm}}     % Riemann tensor
\providecommand{\Tr}{\operatorname{Tr}}     % Uppercase trace
\providecommand{\Spec}{\operatorname{Spec}} % Uppercase spectrum
\providecommand{\diag}{\operatorname{diag}} % Diagonal matrix
\providecommand{\rank}{\operatorname{rank}} % Rank
\providecommand{\scal}{\mathcal R} % Scalar curvature
\providecommand{\Cl}{\operatorname{Cl}}     % Clifford algebra
\providecommand{\diam}{\operatorname{diam}} % Diameter

% ======================================================================
% 9. FUNCTIONS, TRIGONOMETRY, AND CALCULUS
% ======================================================================

% Codifferential naming conventions
% (Avoid \ddag: commonly defined by physics.sty and symbol packages)
\providecommand{\codiff}{\mathrm{d}^{\dagger}} % Codifferential operator

% Quantum mechanics bra-ket notation (needed for basis tables)
\providecommand{\ket}[1]{\lvert #1 \rangle}
\providecommand{\bra}[1]{\langle #1 \rvert}
\providecommand{\braket}[2]{\langle #1 \mid #2 \rangle}

% ======================================================================
% 10. PHYSICAL SCALES AND CONSTANTS (Paper 1c / Born Geometry)
% ======================================================================
\providecommand{\Rp}{\ensuremath{R_{\mathrm P}}}     % matched common EH curvature scale (not defined from observed G)
\providecommand{\hBorn}{\ensuremath{h_{\mathrm B}}} % physical representative 2\pi\hbar of one primitive Born scalar action cell
\providecommand{\UEHjoint}{\ensuremath{U_{J}}} % joint-Dirac EH curvature-coordinate scale 8 kappa_J
\providecommand{\Gread}{\ensuremath{G_{\rm scal}}} % scalar-action coordinate c_op^3/(h_B U); physical matching explicit
\providecommand{\GreadJoint}{\ensuremath{G_{J,\mathrm{raw}}}} % raw ordinary-trace coordinate = G_scal/96
\providecommand{\Rds}{R_{\mathrm{dS}}}    % de Sitter Curvature
\providecommand{\Mpl}{M_{\mathrm{Pl}}}    % Planck mass
\providecommand{\lp}{\ell_{\mathrm{Pl}}}  % Planck length
\providecommand{\Pmax}{P_{\mathrm{max}}}  % Maximum momentum
\providecommand{\Mplbar}{\bar{M}_{\mathrm{Pl}}} % Reduced Planck mass
\providecommand{\eff}{\mathrm{eff}}       % Subscript for effective quantities
\providecommand{\Cosmo}{\ensuremath{\Lambda}}                % Cosmological constant (use \Cosmo in source)
\providecommand{\LCDM}{\ensuremath{\Cosmo\mathrm{CDM}}}  % \LambdaCDM (typesets with \Cosmo)
\providecommand{\Lcut}{\ensuremath{\Lambda_{\mathrm{cut}}}}  % Cutoff scale (brace when used as script)
\providecommand{\Rcs}{\ensuremath{R_{\mathrm{cs}}}}      % Configuration-space radius
\providecommand{\CosmoEff}{\ensuremath{\Cosmo_{\eff}}}   % Effective cosmological constant

% Disambiguation: \Rmom is the momentum-space (S^3) radius. 
% Reserve plain R for the 4D curvature scalar in f(R).
\providecommand{\Rmom}{\ensuremath{R_{\mathrm{mom}}}} 
\providecommand{\Punit}{P_{0}} % Momentum unit set by \Rmom (= \hbar/\Rmom)
\providecommand{\Rvec}{\ensuremath{R_{\mathrm{vec}}}} % invariant vector-space radius in the reduced round sector
% Specific numeric values and expressions
\providecommand{\Vlogcap}{L_{\mathrm{Hopf}}^{(1)}}
\providecommand{\Hopf}{\operatorname{Hopf}}
% Scalar phase (S^1) sector + universal cover notation (Paper I-A)
\providecommand{\tper}{T_\star}      % Time period/scale
\providecommand{\Eper}{E_\star}      % Energy period/scale
\providecommand{\ttilde}{\tilde t}   % Universal cover time
\providecommand{\Etilde}{\tilde E}   % Universal cover energy

% ======================================================================
% 10A. SERIES-WIDE LOGICAL / NOTATIONAL DISAMBIGUATIONS
% ======================================================================
% Use these symbols to avoid collisions among: (i) the Paper--III
% no-insertion/geometric-closure class, (ii) Paper--I-A trimer diagonals,
% (iii) Paper--V-A/V-C family-space diagonals, (iv) scalar-index running,
% and (v) protected shell-clock labels.
\providecommand{\Rprim}{\ensuremath{\mathsf R_{\mathrm{prim}}^{\Sigma}}} % primitive finite structural readouts
\providecommand{\Rloop}{\ensuremath{\mathsf R_{\circ}^{\Sigma}}} % closed-loop readout role
\providecommand{\Rbdry}{\ensuremath{\mathsf R_{\partial}^{\Sigma}}} % boundary-conditioned readout role
\providecommand{\Rgen}{\ensuremath{\mathsf R_{\mathrm{fin}}^{\Sigma}}} % finite generated closure of primitive roles
\providecommand{\PhysOcc}{\operatorname{PhysOcc}}
\providecommand{\Actual}{\operatorname{Actual}}
\providecommand{\GeomLic}{\operatorname{GeomLic}}
\providecommand{\InvDist}{\operatorname{InvDist}}
\providecommand{\Dcl}{\ensuremath{\mathsf{D}_1}}      % unmarked positive Dirac-even closed spectral-readout form class
\providecommand{\Dheat}{\ensuremath{\mathsf D_{\mathrm{heat}}}} % one-parameter heat-semigroup subclass
\providecommand{\DtriI}{\ensuremath{D^{\mathrm{tri}}_1}} % Paper--I-A cyclic-trimer diagonal
\providecommand{\DtriII}{\ensuremath{D^{\mathrm{tri}}_2}}
\providecommand{\DfamI}{\ensuremath{D^{\mathrm{fam}}_1}} % Family-space diagonal selector
\providecommand{\DfamII}{\ensuremath{D^{\mathrm{fam}}_2}}
\providecommand{\alpharun}{\ensuremath{\alpha_{\mathrm{run}}}} % scalar spectral-index running
\providecommand{\Ksh}{\ensuremath{K_{\mathrm{sh}}}} % protected shell-clock coordinate
\providecommand{\ddef}{\ensuremath{\delta_{\mathrm{def}}}} % universal defect scalar after Paper--III realization

% ======================================================================
% 10B. SERIES-WIDE LOGICAL STATUS TAGS
% ======================================================================
% Upright small-caps wrapper for audit/status labels.  It prevents LaTeX from
% trying to use unavailable small-caps italic font shapes inside theorem bodies
% such as TU/lmr/m/scit.
\providecommand{\AuditLabel}[1]{}
\renewcommand{\AuditLabel}[1]{{\normalfont\scshape #1}}

% These tags are bookkeeping labels for consistency audits. They do not
% create new mathematical hypotheses; they record how a datum is being used
% in the proof chain.
\providecommand{\StatusPrimitive}{}
\renewcommand{\StatusPrimitive}{\AuditLabel{Primitive}}
\providecommand{\StatusDerived}{}
\renewcommand{\StatusDerived}{\AuditLabel{Derived}}
\providecommand{\StatusImported}{}
\renewcommand{\StatusImported}{\AuditLabel{Imported}}
\providecommand{\StatusConditional}{}
\renewcommand{\StatusConditional}{\AuditLabel{Conditional}}
\providecommand{\StatusBridge}{}
\renewcommand{\StatusBridge}{\AuditLabel{Bridge-derived}}
\providecommand{\StatusComparison}{}
\renewcommand{\StatusComparison}{\AuditLabel{Comparison-only}}

% ======================================================================
% 10C. SERIES-WIDE CLAIM-WORDING AUDIT LABELS
% ======================================================================
% These are textual audit labels used to keep modal language consistent
% across the series.  They introduce no new mathematical assumptions.
\providecommand{\ClaimForced}{}
\renewcommand{\ClaimForced}{\AuditLabel{forced}}
\providecommand{\ClaimUnique}{}
\renewcommand{\ClaimUnique}{\AuditLabel{unique}}
\providecommand{\ClaimCanonical}{}
\renewcommand{\ClaimCanonical}{\AuditLabel{canonical}}
\providecommand{\ClaimMinimal}{}
\renewcommand{\ClaimMinimal}{\AuditLabel{minimal}}
\providecommand{\ClaimClosed}{}
\renewcommand{\ClaimClosed}{\AuditLabel{closed}}
\providecommand{\ClaimPrediction}{}
\renewcommand{\ClaimPrediction}{\AuditLabel{prediction}}

% ======================================================================
% 10D. SERIES-WIDE THEOREM-CHAIN AUDIT LABELS
% ======================================================================
% These labels mark the direction of load-bearing theorem chains. They are
% audit annotations only; they introduce no new mathematical assumptions.
\providecommand{\ChainInput}{}
\renewcommand{\ChainInput}{\AuditLabel{Chain input}}
\providecommand{\ChainOutput}{}
\renewcommand{\ChainOutput}{\AuditLabel{Chain output}}
\providecommand{\ChainGuard}{}
\renewcommand{\ChainGuard}{\AuditLabel{Chain guard}}
\providecommand{\ChainNoBackflow}{}
\renewcommand{\ChainNoBackflow}{\AuditLabel{No backflow}}

% ======================================================================
% 10E. SERIES-WIDE LOCAL THEOREM-AUDIT LABELS
% ======================================================================
% These labels mark a statement-level dependency audit around individual
% theorem/definition nodes. They are bookkeeping labels only.
\providecommand{\LocalAudit}{}
\renewcommand{\LocalAudit}{\AuditLabel{Local audit}}
\providecommand{\UsesOnly}{}
\renewcommand{\UsesOnly}{\AuditLabel{Uses only}}
\providecommand{\Exports}{}
\renewcommand{\Exports}{\AuditLabel{Exports}}
\providecommand{\DoesNotUse}{}
\renewcommand{\DoesNotUse}{\AuditLabel{Does not use}}

% ======================================================================
% 10F. SERIES-WIDE NUMERICAL-ROW AUDIT LABELS
% ======================================================================
% These labels classify individual rows in numerical, benchmark, status,
% and comparison tables.  They are audit annotations only and do not change
% any theorem, value, formula, or comparison number.
\providecommand{\RowAudit}{}
\renewcommand{\RowAudit}{\AuditLabel{Numerical-row audit}}
\providecommand{\RowTheorem}{}
\renewcommand{\RowTheorem}{\AuditLabel{Theorem-fixed row}}
\providecommand{\RowBridge}{}
\renewcommand{\RowBridge}{\AuditLabel{Bridge-fixed row}}
\providecommand{\RowRepresentative}{}
\renewcommand{\RowRepresentative}{\AuditLabel{Representative row}}
\providecommand{\RowComparison}{}
\renewcommand{\RowComparison}{\AuditLabel{Comparison row}}
\providecommand{\RowCalibrated}{}
\renewcommand{\RowCalibrated}{\AuditLabel{Calibrated row}}
\providecommand{\RowEstimate}{}
\renewcommand{\RowEstimate}{\AuditLabel{Estimate row}}
\providecommand{\RowOpen}{}
\renewcommand{\RowOpen}{\AuditLabel{Open row}}
\providecommand{\RowNoSelector}{}
\renewcommand{\RowNoSelector}{\AuditLabel{No selector}}

% ======================================================================
% 10G. TABLE HELPERS
% ======================================================================
% Table helpers used by Paper V-C.  The unused readout-label preset has been
% removed; only active table helpers remain.
\providecommand{\gpsection}[1]{%
  \addlinespace[0.55em]
  \multicolumn{3}{@{}l}{\textsc{#1}}\\[-0.25em]
  \midrule
}
\providecommand{\mathcell}[1]{%
  \(\begin{gathered}#1\end{gathered}\)%
}

\newcommand{\DerBy}[1]{%
  \begin{tabular}[t]{@{}l@{}}#1\end{tabular}%
}

\newcommand{\TableNotes}[1]{%
  \par\medskip
  \begin{minipage}{0.94\linewidth}
  \footnotesize\emph{Notes.} #1
  \end{minipage}\par
}
% ======================================================================
% 10H. PAPER VI-A / QCD NOTATION
% ======================================================================
% Shared symbols used by Paper VI-A and downstream QCD/Yang--Mills handoffs.
\providecommand{\Pqcd}{\mathfrak P_{\rm VIII}^{\rm QCD}}
\providecommand{\Iqcd}{\mathfrak I_{\rm QCD}^{(7)}}
\providecommand{\muQCD}{\mu_{\rm QCD}^{\rm int}}
\providecommand{\LambdaQCD}{\Lambda_{\rm QCD}}
\providecommand{\LambdaBH}{\Lambda_{\rm BH}^{\rm QCD}}
\providecommand{\HQCD}{\mathcal H_{\rm QCD}^{\rm BH}}
\providecommand{\HQCDfin}{H_{\rm QCD}^{\rm fin}}
\providecommand{\Hraw}{\mathcal H_{\rm raw}}
\providecommand{\HphysQCD}{\mathcal H_{\rm phys}^{\rm QCD}}
\providecommand{\Vcol}{V_C}
\providecommand{\Ps}{P_{\rm sing}}
\providecommand{\PiOct}{\Pi_8}
\providecommand{\MSbar}{\overline{\rm MS}}
\providecommand{\YM}{\mathrm{YM}}
\providecommand{\Ktri}{K_{\triangle}}
\providecommand{\KYM}{K_{\YM}^{\rm BH}}
\providecommand{\aBH}{a_{\rm BH}}
\providecommand{\epsBH}{\varepsilon_{\rm BH}}
\newcommand{\RowConstructive}{\textnormal{\textsc{Constructive}}}
\newcommand{\RowObligation}{\textnormal{\textsc{Bridge obligation}}}
\newcommand{\RowTarget}{\textnormal{\textsc{Target row}}}
\newcommand{\RowValidation}{\textnormal{\textsc{Validation row}}}
\providecommand{\slashed}[1]{\not{#1}}

% ======================================================================
% 10I. PAPER VI-B / CONTINUOUS-QFT NOTATION
% ======================================================================
% Shared symbols used by Paper VI-B and any downstream continuous-QFT or
% Yang--Mills handoff paper.  All definitions are guarded so individual
% papers may override them locally before loading this file if needed.
\providecommand{\A}{\mathfrak A}
\providecommand{\Hphys}{\mathcal H_{\rm phys}}
\providecommand{\QCD}{\mathrm{QCD}}
\providecommand{\BH}{\mathrm{BH}}
\providecommand{\loc}{\mathrm{loc}}

\providecommand{\RBH}{\mathcal R_{\rm BH}}
\providecommand{\DG}{\mathfrak D_G}
\providecommand{\PW}{\mathcal{P\!W}} % Peter--Weyl spectral-symbol subspace
\providecommand{\Toe}{\operatorname{Toe}} % Toeplitz compression/representative
\providecommand{\Hol}{\operatorname{Hol}} % Holonomy

% ======================================================================
% 11. LATE-LOADED PACKAGES (Must be at the end of the preamble)
% ======================================================================
% Load late packages only if a paper has not already loaded them, avoiding
% option clashes in paper-specific preambles.
\makeatletter
\@ifpackageloaded{hyperref}{}{\usepackage{hyperref}}
\makeatother
% --- Series-wide disambiguations (Pass 3) ---
\providecommand{\diracdelta}{\delta_{\rm D}} % Dirac delta distribution (delta reserved for hierarchy)
\providecommand{\Hg}{H_{\mathrm g}} % Hodge harmonic projector (keep H for Hubble)
\providecommand{\Gr}{G_{\mathrm r}} % Green operator (keep G for Newton)
\providecommand{\FD}{\mathrm{D}} % functional/variational derivative symbol (delta reserved)
\providecommand{\Id}{\mathbf{1}} % identity/Kronecker tensor (avoid \delta_{ij})


% BEGIN INLINE SOURCE: Paper00-ActiveExternalReferences-v2_130.tex
% Framework v2.191: current reference binding; existing definitions retained.
% Framework v2.190: current22-root registry with native clock-evolution proofs.
% Framework v2.189: exact current22-root external registry, including retained v2.188 hardening owners.
% Framework v2.187: current registry synchronization only; source release v2.186 preserved.
% Framework v2.186: authoritative current 22-paper registry.
% Framework v2.182: all current 22 logical paper AUX bindings.
% Framework v2.170: current I-B v1.43 external AUX binding.
% Framework v2.167: exact current 22-root registry after CP--BL2 integration.
% Framework v2.166: exact current 22-root compact-release AUX registry.
% Framework v2.165: authoritative current root bindings.
% Framework v2.164: exact current 22 manuscript AUX bindings.
% Framework v2.162: finite metric, functional Hilbert and scoped Ward integration.
% Framework v2.161: current 22 manuscript AUX bindings; source only.
% Framework v2.156 / registry v2.113: synchronized current 22 manuscript roots; source-only registry.
% Framework v2.143 / registry v2.111: active manuscript AUX basenames; no generated AUX or PDF included.
% ======================================================================
% Central external-reference imports, numbers only, no PDF jumping
% ======================================================================

% xr-hyper is loaded centrally so that papers may retain their existing
% \usepackage{xr-hyper} lines without redefining \externaldocument after this
% registry is installed.  If plain xr was loaded earlier, retain it.
\makeatletter
\@ifpackageloaded{xr-hyper}{}{%
  \@ifpackageloaded{xr}{}{\usepackage{xr-hyper}}%
}
\makeatother

% Make external-document registration idempotent.  The registry identifies a
% logical paper by its Paper<number><optional letter> prefix.  Prefix-qualified
% imports remain independent.  Current paper roots contain no local
% \externaldocument calls; this guard also makes accidental duplicates harmless.
\ExplSyntaxOn
\prop_new:N \g_bh_external_refs_prop
\cs_new_eq:NN \BHOriginalExternalDocument \externaldocument

\cs_new_protected:Npn \bh_external_document:nnnn #1#2#3#4
 {
  \tl_set:Nn \l_tmpa_tl {#3}
  \regex_replace_once:nnN
   { \A (?: xref/ )? (Paper[0-9]+[A-G]?) .* \Z }
   { \1 }
   \l_tmpa_tl
  \tl_set:Nx \l_tmpb_tl { #1 | \tl_use:N \l_tmpa_tl }
  \prop_if_in:NVTF \g_bh_external_refs_prop \l_tmpb_tl
   {
    \PackageInfo{BornHopf}
     {Skipping~duplicate~external~refs~for~#3}
   }
   {
    \prop_gput:NVn \g_bh_external_refs_prop \l_tmpb_tl {#3}
    \IfNoValueTF{#4}
     { \BHOriginalExternalDocument[#1][#2]{#3} }
     { \BHOriginalExternalDocument[#1][#2]{#3}[#4] }
   }
 }

\RenewDocumentCommand{\externaldocument}{O{} O{} m o}
 { \bh_external_document:nnnn {#1}{#2}{#3}{#4} }

\NewDocumentCommand{\BHExternalRefs}{m}
 {
  % Do not import the aux file of the logical paper currently being compiled,
  % even when the job and registry use different version suffixes.
  \tl_set:Nx \l_tmpa_tl { \c_sys_jobname_str }
  \regex_replace_once:nnN
   { \A (Paper[0-9]+[A-G]?) .* \Z }
   { \1 }
   \l_tmpa_tl
  \tl_set:Nn \l_tmpb_tl { #1 }
  \regex_replace_once:nnN
   { \A (Paper[0-9]+[A-G]?) .* \Z }
   { \1 }
   \l_tmpb_tl
  \str_if_eq:eeTF { \tl_use:N \l_tmpa_tl } { \tl_use:N \l_tmpb_tl }
   {
    \PackageInfo{BornHopf}
     {Skipping~external~refs~for~#1~because~this~is~the~current~document}
   }
   {
    \IfFileExists{xref/#1.aux}
     { \externaldocument[][nocite]{xref/#1}[#1.pdf] }
     {
      \IfFileExists{#1.aux}
       { \externaldocument[][nocite]{#1} }
       {
        \PackageInfo{BornHopf}
         {External~aux~file~#1.aux~not~found;~no~labels~imported}
       }
     }
   }
 }
\ExplSyntaxOff

% External references are automatic.  This file is the sole owner of both the
% bibliography and the current Paper-I-A--VII-D AUX registry (30 roots including I-C).  Missing auxiliaries
% are informational on the first pass and are populated by the source-only
% two-pass series build.  Basenames deliberately omit temporary download/copy
% suffixes such as ``(1)'' or ``-1''.
  % Paper I-A: round S3 compendium.
  \BHExternalRefs{Paper1A-Geometry-v5_181}

  % Paper I-B: complete exceptional geometry and finite spectral structures.
  \BHExternalRefs{Paper1B-ExceptionalGeometry-v1_27}
  % Paper I-D: technical companion of I-A through I-C.
  \BHExternalRefs{Paper1D-TechnicalMaterial-v1_69}
  % Paper I-C: reciprocal M8 compendium.
  \BHExternalRefs{Paper1C-M8Geometry-v1_146}

\ifdefined\BHMathOnly
  % Papers I-A and I-B import only the mathematical atlas labels.
\else
  % Paper II: five physical-mathematical main papers plus one technical proof volume.
  \BHExternalRefs{Paper2A-GeometricConstitution-v3_227}
  \BHExternalRefs{Paper2B-MomentumCarrier-v2_147}
  \BHExternalRefs{Paper2C-ReciprocalLocalization-v2_177}
  \BHExternalRefs{Paper2D-SymmetrySpectra-v2_145}
  % The former II-F export labels remain in II-E; current II-F is the technical volume.
  \BHExternalRefs{Paper2E-KinematicExports-v3_140}
  % Paper II-F: complete technical proof and calculation owner for II-A--II-E.
  \BHExternalRefs{Paper2F-TechnicalMaterial-v1_71}

  % Paper III: Dynamics and Intrinsic Geometric Constitution (six components).
  \BHExternalRefs{Paper3A-GeneratedDynamics-v1_181}
  \BHExternalRefs{Paper3B-SpectralAction-v1_192}
  \BHExternalRefs{Paper3C-ProtectedHierarchy-v1_196}
  \BHExternalRefs{Paper3D-RunningResponse-v1_168}
  \BHExternalRefs{Paper3E-GravityAction-v1_204}
  \BHExternalRefs{Paper3F-CompleteDynamics-v1_281}
  % Paper III-G: complete preserved technical proof owner.
  \BHExternalRefs{Paper3G-TechnicalMaterial-v1_58}

  % Paper IV-A: generated finite quantum mechanics.
  \BHExternalRefs{Paper4A-FiniteQuantumMechanics-v4_217}
  % Paper IV-B: complete finite quantum processes and physical response.
  \BHExternalRefs{Paper4B-QuantumResponse-v1_55}
  % Paper IV-C: extended proofs and finite technical calculations.
  \BHExternalRefs{Paper4C-TechnicalMaterial-v1_50}

  % Paper V-A: The Standard Model.
  \BHExternalRefs{Paper5A-FiniteMatterGeometry-v3_219}

  % Paper V-B: Finite Spectrum and Flavor Closure.
  \BHExternalRefs{Paper5B-TerminalSpectrum-v5_202}

  % Paper V-C: Finite Interaction Boundaries of the Standard Model.
  \BHExternalRefs{Paper5C-InteractionResponse-v2_185}
  % Paper V-D: preserved technical proof and calculation owner.
  \BHExternalRefs{Paper5D-TechnicalMaterial-v1_58}

  % Paper VI-A: Gravity and Closed Cosmology.
  \BHExternalRefs{Paper6A-GravityCosmology-v2_268}

  % Paper VI-B: Geometric Dark Sector and Cosmological Readouts.
  \BHExternalRefs{Paper6B-CosmologyReadouts-v1_215}

  % Paper VI-C: complete technical proof and comparison owner.
  \BHExternalRefs{Paper6C-TechnicalMaterial-v1_53}

  % Paper VII-A: Finite QCD Kernel and Nonperturbative Closure.
\BHExternalRefs{Paper7A-QCDKernel-v1_185}

  % Paper VII-B: Hadronic, Thermal, and Nuclear QCD Readouts.
  \BHExternalRefs{Paper7B-QCDReadouts-v1_169}

  % Paper VII-C: Continuous Quantum Field Theory.
  \BHExternalRefs{Paper7C-ContinuousQFT-v2_208}

  % Paper VII-D: complete technical material for VII-A through VII-C.
  \BHExternalRefs{Paper7D-TechnicalMaterial-v1_60}

\fi

% END INLINE SOURCE: Paper00-ActiveExternalReferences-v2_130.tex



\providecommand{\LambdacomparisonBH}{\Lambda_{\rm BH}^{\rm comp}}
\providecommand{\deltageom}{\delta_{\rm geom}}
\providecommand{\alphasseed}{\alpha_s^{\rm seed}}
\providecommand{\Hseedzero}{\mathsf H_{\rm seed}^{0}}


% Paper-specific layout refinements.
\setlength{\tabcolsep}{2pt}
\AtBeginEnvironment{longtable}{\footnotesize}
\ifcsname c@biburlnumpenalty\endcsname\setcounter{biburlnumpenalty}{9000}\fi
\ifcsname c@biburllcpenalty\endcsname\setcounter{biburllcpenalty}{9000}\fi
\ifcsname c@biburlucpenalty\endcsname\setcounter{biburlucpenalty}{9000}\fi


\title{\textbf{Born--Hopf Geometry in a Closed Universe VII-A: Finite QCD Kernel}\\[0.4em]
\large Hurwitz geometry, color singlets, strong action, baryon spin--flux and finite source theorems}
\author{Peter Kugelmann\thanks{\texttt{pkugelmann@proton.me}}}
\date{Version 1.185 --- 20 November 2026 \\ Framework v3.077}
\setcounter{tocdepth}{2}

% Current release citation-only xref imports (no material DAG dependency).
\BHExternalRefs{Paper1A-Geometry-v5_175}
\BHExternalRefs{Paper1B-ExceptionalGeometry-v1_27} % v3.026 five-core theorem
\BHExternalRefs{Paper7D-TechnicalMaterial-v1_56}
\begin{document}
% Publication typesetting: controlled prose line-breaking, without scaling mathematics.
\emergencystretch=2em
\maketitle
\begin{abstract}
Paper VII-A owns the finite Born--Hopf QCD kernel: Hurwitz/24-cell
 geometry, generated $SU(3)_c$ color singlets, exact finite Hamiltonian, Cartan
 and Hopf clock, same-action heat/gauge source, chiral/topological response,
 and the complete three-face baryon spin--flux geometry. Finite operator
 results, physical scales, realized states, uniform Yang--Mills gap, and positive
 asymptotic string tension are distinct conditions. The derived Hurwitz split
 grammar is preserved as geometry. Full auxiliary matrices, bounds, numerical
 tests and conditional comparators are retained in Paper VII-D. Paper VII-B
 owns observable hadronic readouts and Paper VII-C bounded-momentum continuum
 reconstruction.
\end{abstract}
\tableofcontents
\listoftables
\clearpage

\section{Geometric input, strong-sector constitution, and source ownership}
\label{p10:sec:v237-a1}
This paper owns the finite QCD kernel and imports physical charged-source data without selecting a second quark-mass member. The source and action conditions must be checked before an operator can be interpreted as a hadronic observable.

% Source Paper VII-A, original line 81; inherited proof bodies unchanged.
% Former logical part: Framework Prefix and Finite Strong-Kernel Inputs; integrated into the new thematic argument.
\label{p10:part:finite-strong-sector-closure}

\noindent\textbf{Part orientation.}
This opening part fixes the inherited framework prefix, the Paper--VII-A theorem scope, and the finite QCD input ledger.  It records what is native to the Born--Hopf strong layer, what is imported from Papers~I-A--VI-B, and what is merely a later bridge or validation map.


% Source Paper VII-A, original line 87; inherited proof bodies unchanged.
\subsection{Framework status of Paper~VII-A}
\label{p10:sec:scope}

Paper~VII-A is the native QCD layer over the inherited framework prefix
\begin{equation}
\label{p10:eq:framework-prefix-fle9}
\mathfrak F_{\le 6B}
:=
\mathfrak F_{1A}^{\rm ref}
+
\mathfrak F_{2}^{\rm kin}
+
\mathfrak F_{3}^{\rm dyn}
+
\mathfrak F_{4}^{\rm q}
+
\mathfrak F_{5A}^{\rm SM}
+
\mathfrak F_{5B}^{\rm readout}
+
\mathfrak F_{5C}^{\rm int}+\mathfrak F_{6A}^{\rm cos,grav}+\mathfrak F_{6B}^{\rm cos,read}.
\end{equation}
Here Paper~I-A \cite{Kugelmann2026a} supplies reference geometry, Paper~II compact kinematics, Paper~III dynamics, Paper~IV finite quantum closure, Paper~V-A the Standard--Model carrier, Paper~V-B the finite spectrum/readout layer, Paper~V-C the Standard--Model interaction layer, Paper~VI-A cosmological gravity, and Paper~VI-B cosmological readouts.  Paper~VII-A begins with this inherited prefix and constructs the finite strong kernel together with rigorously typed chiral, anomaly, and topological computation interfaces.

\paragraph{Scope and ownership.}
Every theorem in Paper~VII-A is read over the inherited prefix \(\mathfrak F_{\le 6B}\) together with the hypotheses explicitly printed in that theorem.  Conditional source members, regulator domains, and continuum/order-of-limits hypotheses are not silently absorbed into the prefix.  Paper~V-A exclusively owns selection of the color carrier,
\(SU(3)_c\), the triplet/antitriplet/singlet matter assignments, and the adjoint
octet.  Paper~V-C owns the GAUGE/CP/QRG interface families and the singlet handoff; Paper~VII-A may compile them parametrically or specialize to an explicitly declared source member.  Paper~VII-A owns the finite strong compilation datum, exact singlet projection inside the compiled QCD space, finite Hamiltonian/transfer families, commutant and Hodge computations, and the definitions of gap, confinement, chiral, anomaly, and topological response targets.  It does not own a numerical strong wall, mass gap, string tension, decay constant, condensate, or susceptibility until the corresponding spectral/correlator theorem is discharged.  Paper~VII-B owns hadron spectroscopy, distributions,
fragmentation, thermal response, nuclear transfer, and numerical validation;
Paper~VII-C owns the continuum/local-field reconstruction.

\begin{longtable}{>{\raggedright\arraybackslash}p{0.23\textwidth} >{\raggedright\arraybackslash}p{0.31\textwidth} >{\raggedright\arraybackslash}p{0.34\textwidth}}
\caption[Paper-VII-A ownership boundary]{Paper~VII-A ownership boundary.}\label{p10:tab:paper-x-ownership-boundary}\\
\toprule
Layer & Fixed object & Logical role \\
\midrule
\endfirsthead
\toprule
Layer & Fixed object & Logical role \\
\midrule
\endhead
Papers~I-A--VI-B & compact carrier, finite quantum calculus, Standard--Model carrier, QCD interface, and cosmological realization & inherited theorem prefix \\
Paper~VII-A & finite strong compiler over the imported \(SU(3)_c\) carrier, \(\Ps\), finite Hamiltonian/transfer families, commutant/Hodge data, and chiral/anomaly/topology response functionals & finite-QCD kernel plus explicit computation targets \\
Paper~VII-B & hadron poles and residues, PDFs/DAs/GPDs/FFs, thermal traces, nuclear transfer, and comparison tables & execution of the fixed source and transfer interface \\
Paper~VII-C & admissible continuum limit, OS/Wightman reconstruction, Euclidean/local-field presentation & continuum handoff without coefficient backflow \\
\bottomrule
\end{longtable}
\TableNotes{``Native'' means fixed by the inherited finite datum before comparison.  Standard representation, index, chiral, anomaly, or spectral theorems may be applied after native fixation.  Lattice, PDG, \(\overline{\rm MS}\), collider, heavy-ion, nuclear, and cosmological data are validation maps only.}

\begin{longtable}{>{\raggedright\arraybackslash}p{0.23\textwidth} >{\raggedright\arraybackslash}p{0.65\textwidth}}
\caption[Finite QCD status classes]{Finite QCD status classes.}\label{p10:tab:finite-qcd-status-classes}\\
\toprule
Status & Meaning \\
\midrule
\endfirsthead
\toprule
Status & Meaning \\
\midrule
\endhead
Finite datum & exact finite object constructed from the inherited carrier and operator calculus \\
Analytic bound & rigorous finite geometric/spectral lower or boundary statement, with physical promotion only when its typing hypotheses are proved \\
Kernel readout & defined spectral, transfer, current, chiral, anomaly, or topology functional; its numerical value may remain a computation target \\
Execution problem & deterministic finite matrix, Schur, trace, transition-measure, or cumulant evaluation assigned to Paper~VII-B \\
Validation comparison & external data used only after native fixation \\
\bottomrule
\end{longtable}
\TableNotes{A standard bridge theorem may be applied only after its printed hypotheses are satisfied.  Defining a finite functional is not a claim that its value, continuum limit, or physical normalization has already been derived.}


% Source Paper VII-A, original line 155; inherited proof bodies unchanged.
\subsubsection{Momentum carrier versus reciprocal-position unfolding}
\label{p10:sec:mp-mx-separation}

Paper~VII-A uses two logically distinct size notions.  The first is the primitive compact momentum multiplicity inherited from Papers~I-A--VI-B.  The second is the large reciprocal-position/comparison unfolding used whenever ordinary spatial locality, large loops, scattering separation, or continuum charts are discussed.

\begin{definition}[Fixed momentum carrier and large position unfolding]
\label{p10:def:fixed-mp-large-mx}
The compact momentum carrier multiplicity is fixed once and for all by the upstream Born--Hopf closure:
\begin{equation}
\label{p10:eq:fixed-mp-one}
m_p=1.
\end{equation}
It is the one primitive compact spectral cell from which the color, shell, and strong-pivot data are read.  By contrast, \(m_x\) denotes the reciprocal-position/comparison unfolding.  Large laboratory locality, large Wilson contours, four-volumes \(V_4\), radial nuclear variables, and Paper-VII-C continuum charts are read in regimes
\begin{equation}
\label{p10:eq:large-mx-position-unfolding}
m_x\gg1,
\end{equation}
with \(m_p\) still fixed by \eqref{p10:eq:fixed-mp-one}.  The finite level symbol \(L\) used below labels an admissible finite readout, basis, cell-complex, or comparison resolution.  It may approximate or organize the \(m_x\)-side unfolding, but it never replaces \(m_p\), never generates additional primitive compact momentum cells, and never changes the already derived color multiplicities \(3,8,9\).
\end{definition}

Whenever a reciprocal carrier comparison is used, the unequal level/window is
a retained mark.  Paper~II-A,
Proposition~\ref{p2a:prop:typed-born-branch-transport}, therefore supplies
the physical permission and typing for transport to a separately certified
reciprocal comparison presentation.  Paper~I-C,
Theorem~\ref{p1b:thm:unequal-exchange}, then supplies the concrete projector
identity transporting the carrier and operator together.  An unequal window
is not automatically an invariant subspace of exchange.  This typing is
applied before the finite readout; it neither enlarges the compact momentum
seed nor assumes that the complete QCD interaction, source, mass packet, or
transfer kernel has an additional unmarked Born symmetry.


\begin{theorem}[No \(m_p\)-flow in Paper~VII-A]
\label{p10:thm:no-mp-flow-paper-x}
Every native QCD coefficient derived in Paper~VII-A is computed with the compact momentum carrier fixed at \(m_p=1\).  Refinement labels \(L\), Wilson areas \(A_L(C)\), topological four-volumes \(V_4\), scattering radii, and Paper-VII-C comparison charts are \(m_x\)-side localization or validation devices.  They may be taken large only after the finite strong datum has been fixed, and they cannot alter the fixed geometric number \(\deltageom=(9\cdot8)^{-1/2}\), \(SU(3)_c\), or the exact singlet projector.  They also cannot by themselves alter the already derived common-heat gauge normalization, select a low-energy scale identification, transfer time, shell response, or infrared energy scale.
\end{theorem}
\begin{proof}
The inherited prefix fixes the compact carrier, color triplet, adjoint octet, and exact singlet representation data, while the common-heat gauge normalization is inherited as a theorem at its native matching datum, while any low-scale coupling coordinate, the strong-CP source member, and the quark mass operator retain the scale/matching or conditional statuses stated in Section~\ref{p10:sec:imported-interface}.  Enlarging \(m_p\) would change the microscopic carrier on which both the structural and conditional interfaces are defined and would therefore constitute backflow into Papers~I-A--VI-B.  By construction, the later labels \(L\), \(A_L(C)\), \(V_4\), and the charts \(\chi_L\) only enumerate finite matrices, compare them with large-position readouts, or export them to the Paper-VII-C continuum theorem.  GP2 therefore permits them only as \(m_x\)-side comparison structure, not as new compact-momentum input.
\end{proof}

\textbf{Reading convention.}\par\smallskip
Whenever Paper~VII-A says ``large loop'', ``area'', ``four-volume'',
``radial block'', or ``finite-volume spectrum'', the spatial or macroscopic
size is on the reciprocal-position side.  The same convention applies to any
``continuum'' or ``refinement'' comparison.  The compact momentum/spectral
generator remains the already-closed \(m_p=1\) Born--Hopf cell.

\begin{corollary}[Reciprocal unfolding does not generate an extensive quark regulator]
\label{p10:cor:v2180-mx-unfolding-no-quark-constructor}
The generated finite Hurwitz/GW quark regulator and the reciprocal-position
$m_x$ unfolding are distinct constructions.  Existing Born/reciprocal
transport may transport an already licensed constructor word only on a target
presentation that independently carries that transported word with its origin
certificate.  Consequently
\[
 \boxed{\text{finite Hurwitz quark regulator}+m_x\text{ unfolding}
 \ \not\Rightarrow\ \text{cofinal extensive quark regulator}}
\]
without an additional extensive fermion/CAR-and-covariant-shift constructor.
This statement does not alter the fixed $m_p=1$ regulator, its finite topology,
or any already derived finite QCD readout.
\end{corollary}

\begin{proof}
Proposition~\ref{p2a:prop:typed-born-branch-transport} requires the reciprocal
branch to contain the transported licensed constructor with its own frozen
origin certificate.  The unequal-exchange theorem transports a compressed
operator that is already defined; it does not create a quark/CAR occurrence or
covariant shift on a carrier for which no such physical word has been
constituted.  Definition~\ref{p10:def:fixed-mp-large-mx} independently types
$m_x$ as reciprocal-position/comparison unfolding.
\end{proof}

% Source Paper VII-A, original line 228; inherited proof bodies unchanged.
\subsubsection{Imported QCD interface from Paper~V-C}
\label{p10:sec:imported-interface}

Paper~VII-A separates theorem-level color geometry from the conditional physical QCD source member.  The structural interface is
\begin{equation}
\label{p10:eq:imported-qcd-interface}
\Iqcd:=\left(\mathcal C_{\rm QCD}^{\rm alg},\Vcol,\PiOct,\Ps,\deltageom,\alphasseed,\mathcal T_q,\bar\theta_{\rm QCD},\muQCD[\alphasseed],\mathcal H_{\rm QCD}^{\rm prim}\right).
\end{equation}
The color carrier and its adjoint decomposition are theorem-level:
\begin{equation}
\label{p10:eq:color-pair-import}
C\cong\CC^3,\qquad G_c=SU(C)\cong SU(3)_c,\qquad
\Vcol=C\otimes C^\ast\cong\mathbf1\oplus\mathbf8.
\end{equation}
The associated geometric count defect is
\begin{equation}
\label{p10:eq:strong-pivot-import}
\deltageom=(9\cdot8)^{-1/2}=\frac1{6\sqrt2}.
\end{equation}
Current Paper~V-C now separates two gauge-normalization branches.  On the
source-native common-heat branch imported from Papers~III-F and V-A,
\begin{equation}
\label{p10:eq:qcd-heat-source-branch}
 k_i g_i^2=4\pi^2,
 \qquad
 \alpha_s^{\rm heat}=\frac\pi2
\end{equation}
at the native common-heat matching datum.  The scale-free count family
\begin{equation}
 \alpha_s^{(t)}=(d_A d_P)^{-t},\qquad t>0,
\end{equation}
remains only a broader kinetic-class calibration chart; its historical
\begin{equation}
\label{p10:eq:qcd-source-branch}
\Hseedzero:\qquad \alphasseed=\deltageom
\end{equation}
is not a native count theorem.  Paper~VII-A therefore never derives
\(\alphasseed=\deltageom\).  A numerical coupling at a named external scale is
obtained only after the common-heat matching scale and RG transport have been
typed; until then \(\alphasseed\) denotes an explicit comparison-scale
coordinate rather than a second gauge-normalization owner.

Likewise, current Papers~V-A--V-B distinguish the diagnostic finite split packet from the actual fermion-source member.  We retain
\begin{equation}
\mathcal T_q=(6=3+3,\ 16=7+9,\ \kappa_u,\kappa_d,\ \text{conditional coordinate representatives})
\end{equation}
only as support/comparison data.  If an actual nonzero G0 quark-source member occurs, the already-generated injective source-owner map fixes its \((H_u,H_d)\) member, and G16 RG/threshold transport supplies the running microscopic mass/source operator at the execution scale:
\begin{equation}
\label{p10:eq:qrg-quark-mass-operator-needed}
\texttt{FERMION\_PRETERMINAL\_PARENT\_CONTRACTION}
\longmapsto(H_u,H_d)
\xrightarrow{\;\mathrm{G16}_{\rm RG}\;}M_q^{\rm phys}(\mu).
\end{equation}
This is an inherited member-plus-transport handoff, not a new QCD mass constructor.  A G16 pole continuation/isolation/residue certificate is consumed only by a pole or mass-shell claim.  No Paper--V-B split coordinate is promoted to a physical threshold by notation.

The sharp one-loop comparison branch remains useful.  On \(\Hseedzero\), with the same conditional coordinate representatives used in Papers~V-B--V-C, it gives
\begin{equation}
\label{p10:eq:qcd-interface-scale-import}
\muQCD\big|_{\Hseedzero,\,\mathcal T_q^{\rm coord}}
=0.2843536970\,\mathrm{GeV}.
\end{equation}
This is a comparison coordinate, not a theorem-level \(\Lambda_{\rm QCD}\), hadron mass, confinement scale, or transfer-time normalization.

\begin{definition}[Paper-V-B/Paper-V-C numerical import convention for Paper VII-A]
\label{p10:def:paper-vi-numerical-import-convention}
Numerical Paper--V-B split coordinates may be inserted only into explicitly labelled diagnostic/comparison formulas.  They do not become eigenvalues of a physical quark-mass operator by notation.  The historical member \(\alphasseed=\deltageom\) retains only its broad-branch comparison status; the source-native common-heat normalization is instead \eqref{p10:eq:qcd-heat-source-branch}.  The physical \(\bar\theta_{\rm QCD}=0\) statement retains its Paper--V-C/G9 conditions.  No hadronic or lattice datum may be used to choose either source member inside Paper~VII-A.
\end{definition}

\begin{remark}[Geometric defect versus physical strong seed]
\label{p10:rem:geometric-defect-versus-strong-seed}
The number \(\deltageom\) is a theorem-level normalized count and, by
Corollary~\ref{p10:cor:exceptional-shell-ancestry-of-qcd-delta}, the same exact
positive scalar as \(\omega_{2O}\).  This common ancestry does not make it the
determinant/gauge owner of the source-native common-heat branch.  The latter
fixes the boundary gauge normalization by \eqref{p10:eq:qcd-heat-source-branch};
a separately written \(\alphasseed\) is only a coupling coordinate at a named
comparison scale after scale/RG matching.  The identity
\(g_s=\sqrt{4\pi\alphasseed}\) follows from the definition
\(\alpha_s=g_s^2/(4\pi)\), but no transition probability or finite transfer
amplitude is thereby fixed.  In particular, \(\sqrt{\deltageom}\) is not
promoted to a universal QCD transition amplitude.
\end{remark}

\begin{remark}[Quark split packet versus physical mass operator]
\label{p10:rem:paper-vi-quark-split-versus-p8-spurions}
The local \(6=3+3\) carrier split and its coefficient ledger remain valid finite support data.  Physical threshold running, chiral spurions, and quark-mass ratios consume the inherited source-member/RG handoff \eqref{p10:eq:qrg-quark-mass-operator-needed}.  Paper~VII-A does not manufacture a mass operator from the coordinate packet and does not select a second physical mass member; it consumes the derived charged preterminal source member after the applicable G16 RG/threshold transport.
\end{remark}

\begin{longtable}{>{\raggedright\arraybackslash}p{0.21\textwidth} >{\raggedright\arraybackslash}p{0.21\textwidth} >{\raggedright\arraybackslash}p{0.24\textwidth} >{\raggedright\arraybackslash}p{0.22\textwidth}}
\caption[Imported QCD interface inventory]{Imported QCD interface inventory after source-type repair.}\label{p10:tab:qcd-interface-inventory}\\
\toprule
Object & Source & Status & Paper-VII-A role \\
\midrule
\endfirsthead
\toprule Object & Source & Status & Paper-VII-A role \\ \midrule
\endhead
\(C\cong\CC^3\), \(SU(3)_c\) & Paper~V-A & \RowTheorem{} & color carrier and gauge group \\
\(\mathbf1\oplus\mathbf8\) & Paper~V-A & \RowTheorem{} & singlet/adjoint decomposition \\
\(\deltageom\) & exceptional-shell scalar plus finite QCD count & \RowTheorem{} & common geometric scalar; normalized count only \\
\(\alpha_s^{\rm heat}\) & Papers~III-F--V-C & \RowTheorem{} & source-native common-heat gauge boundary at its matching datum \\
\(\alphasseed\) at a named external scale & RG/scale matching & \RowBridge{} & comparison-scale coupling coordinate after transport \\
\(\Ps\) & Paper~V-C & \RowTheorem{} & exact singlet projector \\
\(\mathcal T_q\) & Paper~V-B & \RowBridge{} & carrier/coefficient packet, not a mass operator \\
\(M_q^{\rm phys}(\mu)\) & derived charged preterminal source + G16 RG/threshold transport & \RowBridge{} & inherited running mass/source operator; no independent mass selector \\
\(\muQCD[\alphasseed]\) & perturbative comparison & \RowBridge{} & conditional comparison crossing \\
\bottomrule
\end{longtable}

\begin{definition}[Paper-VII-A logical-status classes]
\label{p10:def:logical-status-classes}
\RowTheorem{} denotes a proved finite/mathematical result; \RowBridge{} a typed conditional interface; \RowConstructive{} a deterministic finite computation; \RowObligation{} a theorem or computation still needed before physical promotion; \RowTarget{} a native quantity defined by the kernel but not yet evaluated; and \RowValidation{} an external comparison datum.  A computation target is not silently promoted to a numerical prediction.
\end{definition}

\begin{theorem}[Inherited structural QCD interface and explicit conditional coordinates]
\label{p10:thm:inherited-qcd-interface}
The completed prefix fixes \(C\cong\CC^3\), \(SU(3)_c\), \(\mathbf1\oplus\mathbf8\), \(\deltageom\), the exact singlet projector \(\Ps\), and the source-native common-heat gauge boundary at its matching datum.  The charged preterminal noncentral source member is already generated upstream on the rooted charged branch.  Its named execution-scale RG/threshold transport remains a dependent calculation.  A coupling at a named external scale and any absolute nonperturbative scale retain their stated matching interfaces.  Paper~V-C fixes the native finite trivial sector character and rooted strong-CP invariant; the auxiliary topological Fourier source supplies cumulants rather than a new physical CP member. Regulated physical transport remains claim-scoped.  Consequently no theorem in Paper~VII-A may infer a physical amplitude, threshold, decay constant, gap, string tension, or topological susceptibility from a finite count alone.
\end{theorem}
\begin{proof}
The finite color/interface objects are direct outputs of the upstream carrier and representation theorems.  Paper~V-A now generates the charged preterminal noncentral finite source member on the rooted branch; Paper~V-B/G16 transport that already owned member to the execution scale.  Thus Paper~VII-A inherits a running mass/source operator once those two inputs are supplied, without a second mass constructor.  The strong-coupling scale and other dynamical readouts retain their separately stated conditions.  This is ordinary dependency propagation and introduces no new assumption.
\end{proof}


% Source Paper VII-A, original line 349; inherited proof bodies unchanged.
\subsubsection{Geometric spine and imported \texorpdfstring{\(S^3\)}{S³} theorem bank}
\label{p10:sec:geometric-spine-s3-bank}

Standard strong-interaction uses of \(S^3\) include it appears in Skyrmion baryons, instanton boundary topology, finite-volume Yang--Mills on \(S^3\times S^1\), and Hopf/Faddeev--Niemi knot sectors~\cite{Skyrme1961,AtiyahSinger1963,tHooft1976,Unsal2007,FaddeevNiemi1997,VakulenkoKapitanskii1979}.  The role in the present framework is different.  Paper~VII-A does not introduce \(S^3\) as a spatial regulator or as an after-the-fact soliton target.  The upstream carrier is the compact momentum/spectral carrier.  Therefore the external theorem bank is imported only at the level of topology, harmonic analysis, heat-kernel positivity, and compact transfer arguments; no Skyrmion mass, finite-volume transition temperature, phenomenological knot energy, or lattice scale is imported.

\begin{theorem}[Geometric QCD spine]
\label{p10:thm:geometric-qcd-spine}
The theorem-level finite QCD spine is
\begin{equation}
\label{p10:eq:geometric-qcd-spine}
S^3_{\rm mom}\simeq SU(2)
\longrightarrow 2T\;(24\text{-cell})
\longrightarrow C\cong\CC^3
\longrightarrow C\otimes C^\ast=\mathbf1\oplus\mathbf8
\longrightarrow SU(3)_c,
\end{equation}
with normalized geometric defect
\begin{equation}
\label{p10:eq:delta-from-color-pair-and-octet}
\deltageom=(9\cdot8)^{-1/2}=\frac1{6\sqrt2}.
\end{equation}
The map from this spine to a physical strong seed, quark-mass spectrum, or absolute infrared energy is not part of the finite count theorem and remains separately typed.
\end{theorem}
\begin{proof}
The 24-cell/Hurwitz carrier and the local trimer realize the Paper--V-A color triplet.  The standard identity \(\End(C)=\mathbf1\oplus\End_0(C)\) gives dimensions \(1+8\), hence the pair/adjoint count \(9,8\) and the normalized number \(\deltageom\).  None of these representation identities contains a dynamical energy or source-normalization map, so no such map is inferred.
\end{proof}

\begin{corollary}[Exceptional-shell ancestry of the QCD normalized count]
\label{p10:cor:exceptional-shell-ancestry-of-qcd-delta}
The QCD normalized count is exactly the smooth exceptional-shell scalar
already owned by the Paper--III determinant calculation:
\begin{equation}
\label{p10:eq:qcd-delta-exceptional-shell-ancestry}
 \boxed{
 \deltageom=\omega_{2O}=\frac1{6\sqrt2},
 \qquad
 \deltageom^{-2}
 =\operatorname{rank}(E_6)h(E_6)
 =6\cdot12
 =8\cdot9
 =72.}
\end{equation}
Thus the exceptional-shell and QCD constructions have a common exact scalar:
the former supplies the analytic determinant owner and its binary-octahedral
certificate, while the latter supplies the independent color-pair/adjoint
\(9\)-by-\(8\) certificate.  This equality closes the numerical ancestry of
\(\deltageom\); it does not select \(\alphasseed\), a gauge source, a
transition amplitude, a transfer time, or a nonperturbative energy scale.
\end{corollary}

\begin{proof}
Paper~III-C, Corollary~\ref{p3:cor:delta-e6-root-denominator-strengthened},
gives \(\omega_{2O}=1/(6\sqrt2)\) and
\(\omega_{2O}^{-2}=\operatorname{rank}(E_6)h(E_6)=72\).  Theorem~\ref{p10:thm:geometric-qcd-spine}
gives \(\deltageom^{-2}=9\cdot8=72\).  Positivity fixes the common positive
square root.  The owner firewall follows because equality of scalar values
does not identify the determinant germ, gauge source, or transfer operator.
\end{proof}

\begin{proposition}[\texorpdfstring{$D_4/A_2$}{D4/A2} reading of the finite QCD color support]
\label{p10:prop:d4-a2-reading-of-qcd-color-support}
The geometric QCD spine of Theorem~\ref{p10:thm:geometric-qcd-spine} has the equivalent root-system reading
\[
\begin{gathered}
\Phi(D_4)\longleftrightarrow 2T
\supset C_6^{\rm link}=T_+\sqcup T_-,
\qquad C=\CC^{T_+},\\
2D_3/\{\pm1\}\cong W(A_2),
\qquad
\mathfrak{su}(3)_c=\mathfrak{su}(C).
\end{gathered}
\]
Consequently the finite interacting color octet used by Paper~VII-A is the \(A_2\)-adjoint decomposition on the trimer support
\[
\mathfrak{su}(3)_c\otimes_{\mathbb R}\CC
\cong\mathfrak{sl}_3(\CC)
=\mathfrak h_{A_2}^{\CC}
\oplus
\bigoplus_{\alpha\in\Phi(A_2)}\mathfrak g_\alpha,
\qquad
\dim_{\CC}\mathfrak h_{A_2}^{\CC}+|\Phi(A_2)|=2+6=8,
\]
where the root spaces belong to the complexified Lie algebra and the real compact form is the eight-real-dimensional $\mathfrak{su}(3)_c$, coming from the trimer selected by the vertex-figure geometry of the \(D_4\)/Hurwitz \(24\)-cell.  The center-triality used in the confinement and hadron-selection theorems below is the center grading of this \(SU(3)_c\) closure,
\[
Z(SU(3)_c)\cong\ZZ_3,
\]
not the \(S_3\) diagram triality of \(D_4\).
\end{proposition}

\begin{proof}
Paper~I-A identifies the Hurwitz \(24\)-cell with the \(D_4\) root shell, records the local cube-link trimer and \(A_2\) defect-plane geometry, and proves \(2D_3/\{\pm1\}\cong W(A_2)\).  Paper~V-A applies this finite Weyl skeleton to the visible color support: the selected axis-removed six-cycle is the local vertex-figure hexagon, not a literal great \(D_4\) rank-two root subsystem; one alternating trimer gives \(C=\CC^{T_+}\), and the no-insertion unitary closure gives \(\mathfrak{su}(C)\cong\mathfrak{su}(3)_c\).  The adjoint decomposition of \(\mathfrak{su}(3)\) has a two-dimensional Cartan plane and six root spaces, hence the \(2+6=8\) octet used in \(C\otimes C^\ast=\mathbf1\oplus\mathbf8\).

The center statement is a role separation.  The \(D_4\) triality symmetry permutes the three triality octet frames of the finite support geometry.  After the color support has been selected and closed to \(SU(C)\), the center acting on fundamental quarks is instead the scalar subgroup \(\{1,\omega,\omega^2\}\subset SU(C)\).  This is the \(\ZZ_3\) center grading used in the finite QCD singlet projection and confinement arguments; it is not an additional \(D_4\) diagram-symmetry assumption.
\end{proof}


\begin{theorem}[Momentum-\(S^3\) harmonic basis]
\label{p10:thm:momentum-s3-harmonic-basis}
Because the native carrier is a compact momentum carrier with \(S^3\simeq SU(2)\), every scalar carrier field admits a Peter--Weyl/hyperspherical expansion~\cite{Peter1927,Weyl1946}
\begin{equation}
\label{p10:eq:s3-harmonic-expansion}
\psi(p)=\sum_{\ell,m,n}\psi_{\ell mn}Y_{\ell mn}(p),
\qquad
\Delta_EY_{\ell mn}=\ell(\ell+2)Y_{\ell mn},
\end{equation}
with degeneracy
\begin{equation}
\label{p10:eq:s3-degeneracy}
d_\ell=(\ell+1)^2 .
\end{equation}
Under a finite admissible readout cutoff \(\ell\le L\), the scalar momentum basis used for computation is finite, with
\begin{equation}
\label{p10:eq:s3-cutoff-mode-count}
N_L=\sum_{\ell=0}^L(\ell+1)^2=\frac{(L+1)(L+2)(2L+3)}6 .
\end{equation}
Consequently the one-body mode carrier is finite. The full QCD matrix
problem additionally consumes the actual finite occupation/capacity carrier
of Theorem~\ref{p10:thm:v2211-finite-occupation-assembly}; unrestricted
bosonic Fock space is not made finite by a one-body mode cutoff.  The cutoff label \(L\) in this theorem is an auxiliary finite readout/comparison resolution; it is not a change of the primitive compact momentum multiplicity \(m_p=1\) from Definition~\ref{p10:def:fixed-mp-large-mx}.
\end{theorem}
\begin{proof}
The identification \(S^3\simeq SU(2)\) makes the harmonic decomposition an immediate instance of Peter--Weyl theory.  Paper~I-A's representation--harmonic theorem gives \(\Delta_EY_{\ell mn}=\ell(\ell+2)Y_{\ell mn}\) with degeneracy \((\ell+1)^2\); equivalently \(\Delta_{S^3(\Rmom)}=\Rmom^{-2}\Delta_E\).  A finite admissible readout cutoff leaves finitely many modes in the computation basis.  This truncation is part of the \(L\)-level matrix readout and not a physical flow of \(m_p\).  Color and singlet projection act on this finite mode algebra as internal
finite tensor factors. On the separately retained finite occupation carrier
the regulated QCD transfer is therefore a finite matrix.\qedhere
\end{proof}

\begin{theorem}[Connected null flow cannot permute a finite point shell]
\label{p10:thm:connected-flow-finite-shell-no-go}
Let \(F\) be a finite discrete set and let \(G\) be a connected topological
group.  Every continuous homomorphism
\[
 \rho:G\longrightarrow\operatorname{Sym}(F)
\]
is trivial.  In particular, neither the compact clock--Hopf \(U(1)\) action
nor its elapsed \(\RR\)-cover can act nontrivially by exact permutations of
the 24 Hurwitz points.
\end{theorem}

\begin{proof}
The continuous image of connected \(G\) is connected, whereas the finite
permutation group is discrete.  Its image is therefore a singleton, and a
homomorphism sends the identity to the identity permutation.
\end{proof}

\begin{corollary}[Exact harmonic action and shell-intertwining boundary]
\label{p10:cor:harmonic-null-action-shell-boundary}
Let \(\mathcal H_{x,L}^{\rm PW}\subset L^2(S_x^3)\) be a finite sum of
complete round-Laplacian eigenspaces and let \(\Pi_{x,L}^{\rm PW}\) be its
orthogonal projection.  The position-side clock--Hopf right action preserves
this carrier, so
\[
 [\Pi_{x,L}^{\rm PW},U_H(s)]=0,
 \qquad
 U_{H,L}(s):=U_H(s)|_{\mathcal H_{x,L}^{\rm PW}}
\]
is an exact finite-dimensional unitary representation.  On the compact
\(U(1)\) quotient its dimensionless Hopf generator has integer weights.  The
elapsed \(\RR\)-cover uses the lifted one-parameter representation; the full
physical null generator, which also contains the Hamiltonian, is not thereby
forced to have integer spectrum.

This supplies an admissible \emph{kinematic} carrier for a null-action
certificate.  It does not identify \(\mathcal H_{x,L}^{\rm PW}\) with the native
Hurwitz-shell QCD carrier.  Any use of shell sampling must declare a map
\[
 S_L:\mathcal H_{x,L}^{\rm PW}\longrightarrow
 \CC^{24}\otimes\mathcal K_L
\]
and prove the exact or controlled approximate intertwining relation required
by the claimed QCD readout.  This map includes the already generated
position--momentum comparison \(\iota_1^{-1}:S_x^3\to S^3_{\rm mom}\) whenever
the native Hurwitz points are read on the momentum parent.  The 24-cell may
remain the internal support,
stencil, quadrature, or regulator without carrying a nontrivial continuous
point-permutation action.
\end{corollary}

\begin{proof}
Round isometries commute with the round Laplacian and hence preserve every
complete spectral subspace.  The restricted action is therefore unitary and
finite dimensional.  Compact \(U(1)\) representations decompose into integer
characters.  The remaining statements follow from
Theorem~\ref{p10:thm:connected-flow-finite-shell-no-go} and from the fact that
sampling is a map out of the harmonic carrier, not an action on the point
set.
\end{proof}


\begin{corollary}[Exact finite geometric sampling interfaces]
\label{p10:cor:v203-exact-geometric-sampling}
Paper~II-C, Corollary~\ref{p2c:cor:v204-sampled-realization}, transports
the already realized intrinsic class to its explicit rank-fifteen Hurwitz
fibre presentation. Paper~I-A supplies its exact coordinate matrices and
scalar compiler in Theorem~\ref{p1:thm:v203-scalar-observable-compiler}
and Corollary~\ref{p1:cor:v203-quadratic-top-block-correction}; the latter
uses the separate rank-fourteen scalar band.
\end{corollary}
\begin{proof}
Apply the II-C realization result and its explicit Paper-I-A isometries.
Their domains and exact symbol matrices are those of the stated carriers.
\end{proof}
For operators already acting on either such carrier, Paper~III-A,
Theorem~\ref{p3a:thm:v204-action-law-transport}, transports the actual
action and evolution, and Paper~IV,
Corollary~\ref{p4:cor:v204-quantum-sampled-transport}, transports any
existing one-particle, species and Gauss structure. No additional
realization, Stone theorem, state postulate or symmetry assumption is
introduced by this change of presentation.

In Definition~\ref{p10:def:finite-null-representation-certificate}, each
already established clause on a common carrier is therefore preserved
exactly.  The v2.52 null-flow theorem removes the former need to identify an
independently supplied R1 generator with the positive-polar generator: once a
native equal-time Hopf-equivariant good-field seed exists on the actual QCD
carrier, the same \(H_L,P_{H,L}\) pair constructs the elapsed generator
uniquely.  The remaining G35 tests are therefore construction of that native
seed, proof of the actual-carrier positive-polar inequality, exact R3
physical/Fock compatibility, and the color/Gauss/carrier intertwiners.  This
is not a fresh justification of $\mathcal H_1$ or $\mathsf C_D$.  The
dimensionless Levi--Civita comparison of
Corollary~\ref{p1:cor:v203-geometric-positive-polar} does not by itself
construct the native QCD seed or prove the R3 equality.  A larger spectral
carrier still needs its own representation map.

\textbf{Import discipline.}
Theorems on \(S^3\) topology, Peter--Weyl decomposition, heat kernels on compact connected spaces, Hopf charge, and Faddeev--Skyrme stability may be imported as mathematical theorem types.  Numerical Skyrmion masses, finite-volume transition temperatures on spatial \(S^3\times S^1\), phenomenological knot energies, and lattice scales are comparison data and cannot select any Paper~VII-A coefficient.
\begin{theorem}[External theorem-bank import map]
\label{p10:thm:external-theorem-bank-import-map}
The external QCD literature may be used in Paper~VII-A only through
hypothesis-matched theorem types.  The admissible imports are
\begin{enumerate}[label=(\roman*), leftmargin=2.0em]
\item reflection positivity and positive transfer-matrix existence for
Wilson-type Euclidean gauge regulators
\cite{OsterwalderSeiler1978,MenottiPelissetto1987};

\item strong-coupling character-expansion area-law mechanisms for Wilson loops
\cite{Wilson1974,KogutSusskind1975,OsterwalderSeiler1978};

\item admissible finite topological sectors and regulated lattice index
technology \cite{Luscher1982Topology,GinspargWilson1982,Neuberger1998};

\item near-zero Dirac spectral universality applies after the symmetry class and
topology are fixed.  This import is used only in that matched setting
\cite{LeutwylerSmilga1992,Verbaarschot1994};

\item center-vortex and center-disorder mechanisms as comparison routes for the
exact \(Z_3\) center grading and conditional area-law route
\cite{tHooft1978,Greensite2003};

\item finite-volume comparison control once a uniform mass gap is established
\cite{Luscher1986Volume};

\item Hopf/Faddeev energy bounds once the Born--Hopf color-direction reduction
yields a positive Faddeev--Skyrme functional
\cite{FaddeevNiemi1997,VakulenkoKapitanskii1979}.
\end{enumerate}
None of these imports may supply an absolute nonperturbative scale, \(f_\pi\), \(\Sigma_q\),
\(M_N\), \(T_c\), \(\sigma_{\YM}\), PDF shapes, or hadronic matrix elements as
numerical inputs.
\end{theorem}
\begin{proof}
Each cited theorem has hypotheses independent of the physical values of QCD observables: positivity of a Euclidean regulator, compact/admissible gauge-field topology, representation-theoretic center charge, or universal microscopic spectral correlations after symmetry class and index are fixed.  When the Born--Hopf finite carrier satisfies the corresponding hypotheses, the theorem may be imported.  The theorem's numerical consequences in ordinary lattice simulations or effective models are comparison data, not selectors for the Born--Hopf coefficients.\qedhere
\end{proof}


\begin{theorem}[Readout-chronology criterion for QCD completion]
\label{p10:thm:qcd-readout-chronology}
A Paper~VII-A QCD kernel quantity is native to the Born--Hopf framework only if every coefficient, kernel, truncation order, boundary scale, and representative state used in its construction is fixed before consulting hadron masses, lattice values, scattering data, PDF fits, fragmentation data, heavy-ion data, or cosmological likelihoods.  Downstream QCD data may test, qualify, or rule out a readout, but may not select it.
\end{theorem}
\begin{proof}
The upstream series fixes the finite carrier and the QCD interface before any QCD comparison is made.  If a hadronic datum is used to choose an absolute strong scale, \(\Sigma_q\), \(f_\pi\), a spectrum kernel, a PDF shape, a fragmentation map, or a thermal equation of state, then the resulting object is not a consequence of the finite carrier; it is a phenomenological calibration.  This would move the object outside the Paper~VII-A finite-kernel domain.\qedhere
\end{proof}

\begin{theorem}[Finite regulated pure-glue operator layer is inherited; full-QCD mass layer is conditional]
\label{p10:thm:finite-qcd-operator-layer-inherited}
The finite state/operator calculus together with \(SU(3)_c\) and \(\Ps\) is sufficient to form the regulated pure-glue operator algebra and its singlet projection.  A full quark Hamiltonian is additionally well-defined once a typed physical mass operator \(M_q^{\rm phys}\) and the canonically normalized coupling at the execution scale are supplied; on the source-native branch that coupling is obtained by RG/scale transport from the common-heat boundary.  No hadron, lattice, PDF, fragmentation, thermal, or nuclear datum is required for either definition.
\end{theorem}
\begin{proof}
The pure-glue statement uses only finite gauge variables, the adjoint representation, and the finite quantum calculus.  The quark mass term requires the actual G0 quark-source member transported to the execution scale by the G16 RG/threshold map; the diagnostic split packet is not that member.  Thus the pure-glue layer is unconditional within the finite prefix, while numerical massive full-QCD execution is conditional on the occurring source member and its transport.
\end{proof}


% Source Paper VII-A, original line 6116; inherited proof bodies unchanged.
\subsection{physical quark-mass handoff}
\label{px:sec:post-v125-qmass}

The physical quark-mass input required by QCD is now typed more narrowly than
a raw Paper-V-B split packet.  On the generic low-energy branch define
\[
 H_u=Y_uY_u^*,\qquad H_d=Y_dY_d^*.
\]
The physical quark flavor object is their simultaneous-conjugacy class under
$U(3)_{Q_L}$; the two spectra provide the six quark singular values and the
relative spectral frame provides CKM.  Right-handed weak-basis frames are
presentation and are not QCD threshold inputs.

Accordingly the historical `X-QMASS` structural/member interface is no longer
an independent QCD frontier.  The weighted source-owner/family geometry is
generated upstream, and its actual numerical member consumes the generated rooted charged-preterminal
\texttt{FERMION\_PRETERMINAL\_PARENT\_CONTRACTION} source member (on its certified hierarchy arrow), rather than a retired same-action SVD assembly prerequisite.  The only mass-handoff
certificate consumed by ordinary threshold/chiral-spurion execution is the
finite-to-continuum scheme/RG matching-and-running map on the relevant scale
interval.  The historical terminal-row lookup and comparison-exponent atlas
remain nonphysical diagnostics.  A pole certificate is added only for a pole
claim.



% BEGIN TRANSITIVE-QCD-MASS-HANDOFF v1.64

% Source Paper VII-A, original line 6142; inherited proof bodies unchanged.
\subsection{transitive quark-mass handoff into full QCD}
\label{p10:sec:v164-transitive-qmass-handoff}

\begin{theorem}[No independent full-QCD quark-mass constructor]
\label{p10:thm:v164-no-independent-qmass-constructor}
On a branch with an actual nonzero G0 quark-source member and a certified G16
RG/threshold transport to scale \(\mu\), the finite full-QCD constructor
consumes \(M_q^{\rm phys}(\mu)\) as the transported member of the already
generated microscopic source-owner family.  The QCD compiler does not select a
second flavor member and does not require a physical pole certificate unless
the target observable is itself a pole/mass-shell observable.
\end{theorem}
\begin{proof}
Paper~V-A's source-owner map is injective on the generated support slots;
Paper~V-B retypes the actual member to G0 and proves the RG/pole premise split.
The finite QCD Hamiltonian already accepts the whole inherited mass/source
family.  Composing these maps fixes the member inserted at scale \(\mu\)
without any additional section selector.  Pole continuation data have a
different target type and therefore enter only the corresponding pole branch.
\end{proof}

\begin{corollary}[G34 transitive status]
\label{p10:cor:v164-g34-transitive-status}
The structural/member part of G34 is
\[
 \boxed{\texttt{DERIVED\_DEPENDENT\_WHOLE\_FAMILY\_TRANSPORT\_CLOSED}}.
\]
For member-sensitive numerical QCD predictions the live mass-side premises are
only the actual \texttt{FERMION\_PRETERMINAL\_PARENT\_CONTRACTION} and the consumed
\(\mathrm{G16}_{\rm RG}\) matching/running certificate.  G16 pole data remain
claim-scoped.
\end{corollary}
% END TRANSITIVE-QCD-MASS-HANDOFF v1.64


% Source Paper VII-A, original line 6581; inherited proof bodies unchanged.
\subsection{Effective rooted mass-source handoff}
\label{p10:sec:v2109-rooted-source-handoff}
The charged mass-source constructor is already closed on the occurring
rooted charged/action branch.  Import the complete mass/incidence sibling
packet and preterminal source of
Theorem~\ref{p5:thm:v397-charged-preterminal-srcval-realization} and
Corollary~\ref{p5:cor:v397-charged-preterminal-assembly-closure}:
\[
 R_s^{\rm ch}=\sum_r m_{s,r}^{\rm fin}\iota_{s,r},\qquad
 D_{F,s}^{\rm ch}
 =\Sigma_{\sigma_s}\otimes P_{Q_s}\otimes R_s^{\rm ch}+\mathrm{h.c.},
 \quad s\in\{u,d\}.
\]
The dependency is this inherited finite source followed by the applicable
\(\mathrm{G16}_{\rm scheme/RG}\) transport to \(M_q^{\rm phys}(\mu)\).
The direct finite coefficients and their family incidences remain siblings;
they are not redefined as singular values of $R_s^{\rm ch}(R_s^{\rm ch})^\dagger$.
The former noncentral SVD/terminal-binding gap is superseded on this branch,
not a live QCD constructor or a second source-member occurrence gate.
Actual scheme/threshold/running evaluation remains dependent execution.
Pole certification is consumed only by a pole claim; state/history, ordered
chiral limits and extensive bounds remain only where the readout uses them.


% BEGIN FRAMEWORK v2.228: CHARGED ORIGIN AND CHARACTER TESTS

% Source Paper VII-A, original line 6606; inherited proof bodies unchanged.
\subsection{Current charged-input origin in the strong kernel}
\label{p7a:sec:v228-charged-origin}
The noncentral quark-source input in the strong kernel consumes the
hierarchy-arrow certificate of
Definition~\ref{p5:def:v228-hierarchy-origin}.
Proposition~\ref{p5:prop:v228-charged-origin-scope} supplies assembly and
source-value handoff on that generated-packet domain. The kernel and its
complete reductions remain defined on their supplied actual source/state
domain; a retained scalar mass coefficient or a geometric family host does
not evaluate the source arrows. Hadronic pole claims retain their full
kernel, stationary support and physical visibility requirements.

% END FRAMEWORK v2.228: CHARGED ORIGIN AND CHARACTER TESTS


% Source Paper VII-A, original line 7819; inherited proof bodies unchanged.
\subsection{charged finite-source import}
\label{p10:sec:v168-charged-source-import}
The upstream charged family/source occurrence boundary is closed on the rooted
charged branch.  Paper~VII-A therefore treats the charged finite quark
mass/source operator as an inherited theorem output and retains only the G16
running/threshold transport needed to evaluate it at a QCD execution scale.
No QCD mass constructor or second charged source-member selector is introduced.

% END CHARGED-SOURCE-IMPORT v1.68


% BEGIN CHARGED-COMPARISON-IMPORT v1.69

% Source Paper VII-A, original line 7831; inherited proof bodies unchanged.
\subsection{charged finite-comparison import}
\label{p10:sec:v169-charged-comparison-import}
The microscopic quark source imported from Paper~V-B now includes the derived
rooted charged comparison character.  In particular, its relative
$24/25$-grading is fixed before QCD running.  The QCD kernel may transport this
owned source through its existing renormalization and threshold machinery but
cannot reselect the charged polar class or use running to manufacture the
finite family-relative comparison factor.
% END CHARGED-COMPARISON-IMPORT v1.69


% BEGIN EXCEPTIONAL-ROOT-INTEGRATION v2.173

% BEGIN INLINE SOURCE: Paper7A-OverlapSourceGram-Appendix-v1_0.tex


\section{Hurwitz carrier, color singlets, and executable finite Hamiltonian}
\label{p10:sec:v237-a2}
The Hurwitz/24-cell incidence, $SU(3)_c$ color carrier, singlet projection and finite strong algebra are geometric and operator constructions. Their exact split grammar is retained; finite arithmetic alone does not choose QCD matrix elements.

% Source Paper VII-A, original line 630; inherited proof bodies unchanged.
\subsection{Finite QCD compiler and transfer}
\label{p10:sec:strong-sector-ledger}

The imported interface and the finite quantum calculus determine a finite
strong-sector matrix problem.  Paper~VII-A closes that problem by specifying its
basis, admissible connection domain, singlet projection, Hamiltonian, transfer,
and source algebra.  Physical hadron, distribution, thermal, and nuclear
execution is deferred to Paper~VII-B.

\begin{definition}[Paper-VII-A QCD compiler endpoint]
\label{p10:def:qcd-compiler-endpoint}
Paper~VII-A reaches its kernel-definition endpoint when the structural finite compilation datum, exact singlet projection, compiler maps, and normalized source algebra have been specified, with every unresolved physical interface carried as an explicit typed argument.  Paper~VII-B evaluates hadronic, distributional, thermal, and nuclear consequences only on branches for which those arguments have been supplied and may not modify the compiler.
\end{definition}

\begin{definition}[Finite QCD compilation datum]
\label{p10:def:finite-qcd-compilation-datum}
At finite level \(L\), define
\begin{equation}
\label{p10:eq:finite-qcd-compilation-datum}
\mathfrak C_L^{\rm QCD}:=
\bigl(G_L,\mathcal A_{L,{\rm adm}},\mathcal B_L^{\rm raw},C,F_L,\Ps^{(L)},\alphasseed,M_{q,L}^{\rm phys},\mathcal O_L\bigr),
\end{equation}
where \(\alphasseed\) and \(M_{q,L}^{\rm phys}\) are present only on branches on which those physical interfaces have been supplied.  The pure-glue datum is obtained by deleting the quark entries.
\end{definition}

\begin{definition}[Compiled finite strong matrices and transfer family]
\label{p10:def:compiled-finite-strong-matrices}
The compiled Hamiltonian is
\begin{equation}
\label{p10:eq:compiled-hamiltonian-matrix}
H_L:=\left[\Ps^{(L)}\bigl(H_{\YM,L}+H_{q,L}+H_{{\rm int},L}\bigr)\Ps^{(L)}\right]\big|_{\operatorname{Ran}\Ps^{(L)}},
\end{equation}
with the quark terms omitted on the pure-glue branch. This is the exact global-singlet finite matrix. Its identification with the locally gauge-physical Hamiltonian additionally requires a well-defined gauge-invariant finite truncation and the reducing local Gauss projector of Theorem~\ref{p10:thm:v350-local-gauss-projector}. Diagonalize directly on the stated global-singlet carrier,
\begin{equation}
\label{p10:eq:compiled-energy-eigenvalues}
H_L\Psi_{a,L}=E_{a,L}\Psi_{a,L}.
\end{equation}
For every \(t>0\) the normalized Euclidean transfer family is
\begin{equation}
\label{p10:eq:compiled-transfer-matrix}
\widehat K_L(t):=\exp[-t(H_L-E_{0,L})],
\end{equation}
so \(\widehat K_L(t)\) is positive self-adjoint and has vacuum eigenvalue one.  No preferred \(t\) is needed to define the Hamiltonian spectrum.
\end{definition}

\begin{theorem}[Finite QCD matrix compiler]
\label{p10:thm:finite-qcd-matrix-compiler}
For every fully typed compilation datum \(\mathfrak C_L^{\rm QCD}\), basis enumeration, gauge/quark matrix assembly, singlet projection, diagonalization, source differentiation, commutant computation, and formation of \(\widehat K_L(t)\) are deterministic finite operations.  The compiler introduces no additional phenomenological coefficient.  If \(\alphasseed\) has not been transported to the execution scale, or if the actual G0 quark-source member and its G16 RG transport have not supplied \(M_q^{\rm phys}(\mu)\), the corresponding full-QCD numerical execution is undefined rather than guessed.
\end{theorem}
\begin{proof}
All listed operations are finite linear algebra on the supplied matrices.  Determinism does not imply that an omitted physical input has been derived: the compiler therefore treats missing RG/threshold transport as a typed obligation.  It does not create a replacement source/mass selector.  This separates mathematical executability from physical closure.
\end{proof}

\begin{longtable}{>{\raggedright\arraybackslash}p{0.22\textwidth} >{\raggedright\arraybackslash}p{0.36\textwidth} >{\raggedright\arraybackslash}p{0.30\textwidth}}
\caption[Finite QCD compiler proof-discharge ledger]{Finite QCD compiler proof-discharge ledger.}\label{p10:tab:finite-qcd-compiler-proof-discharge}\\
\toprule
Compiler stage & Fixed input consumed & Output produced by finite operation \\
\midrule
\endfirsthead
\toprule
Compiler stage & Fixed input consumed & Output produced by finite operation \\
\midrule
\endhead
Basis and domain enumeration & Hurwitz shell, admissible refinement level, finite color/quark support, and exact upstream projectors & Finite raw basis \(\mathcal B_L^{\rm raw}\) and admissible domain \(\mathcal A_{L,{\rm adm}}\). \\
Matrix assembly & \(SU(3)_c\), the coupling at the execution scale (transported from the source-native common-heat boundary or supplied only on an explicitly labelled comparison branch), a typed \(M_q^{\rm phys}\) when quarks are present, covariant shifts, an admissible gauge-action representative, and the source/operator list & Finite gauge, quark, interaction, source, current, Wilson, topological, and thermal matrices. \\
Global singlet projection & Exact diagonal/global singlet projector and conserved charge/projector labels & Globally invariant finite matrix problem on \(\operatorname{Ran}\Ps^{(L)}\); physical local-Gauss reduction requires its own certificate. \\
Transfer and spectrum & Positive reflection-compatible finite Hamiltonian block & \(H_L\), \(\widehat K_L(t)=e^{-t(H_L-E_0)}\), eigenvectors, eigenvalues, and finite energy differences. \\
Observable extraction & Gauge-invariant operator list and finite eigenbasis & Matrix elements, source derivatives, Schur poles, Wilson-loop expectations, susceptibilities, and finite traces. \\
Comparison sequence & Native values already fixed at finite level & Only the prescribed \(m_x\)-unfolding/refinement presentation remains; no lattice scale or hadron value selects a coefficient. \\
\bottomrule
\end{longtable}
\TableNotes{The table expands Theorem~\ref{p10:thm:finite-qcd-matrix-compiler}.  Each row is a deterministic finite algebraic or spectral operation on data fixed before comparison.}


% Source Paper VII-A, original line 703; inherited proof bodies unchanged.
\subsubsection{Hurwitz-shell momentum regulator and native QCD transfer}
\label{p10:subsec:hurwitz-shell-momentum-regulator-native-qcd-transfer}

The finite strong compiler is most naturally read in transfer-first form.  The primitive object is not a Euclidean position lattice and not a Fourier transform of a chosen lattice action.  It is the selected Hurwitz/\(24\)-cell momentum-shell complex on the compact Born--Hopf carrier.  Reciprocal-position behaviour can be discussed only after a fully typed finite Hamiltonian/transfer representative has been fixed, and then only through the \(m_x\)-side pairing/unfolding appropriate to Wilson areas, volumes, and continuum comparison.

\begin{definition}[Hurwitz-shell transfer complex]
\label{p10:def:hurwitz-shell-transfer-complex}
Let
\begin{equation}
\label{p10:eq:hurwitz-shell-vertex-set}
V_H:=2T\subset S^3_{\rm mom}
\end{equation}
be the \(24\) Hurwitz units.  The primitive shell graph has edges
\begin{equation}
\label{p10:eq:hurwitz-shell-edge-set}
E_H:=\{(q,q')\in V_H^2:q\ne q',\ \langle q,q'\rangle=1/2\}.
\end{equation}
Together with the triangular faces and octahedral three-cells of the \(24\)-cell boundary this defines the finite cellular complex
\begin{equation}
\label{p10:eq:hurwitz-shell-cell-complex}
\mathcal K_H=(V_H,E_H,F_H,C_H,\partial).
\end{equation}
The one-step Euclidean transfer complex is \(\mathcal K_H^\tau\), obtained by adjoining a transfer layer to \(\mathcal K_H\).  It has spatial triangular Hurwitz faces and temporal ribbon faces over the Hurwitz edges.  The complex \(\mathcal K_H\) is a compact momentum-shell regulator; it is not a primitive spacetime lattice.
\end{definition}

\begin{theorem}[Hurwitz-shell adjacency spectrum]
\label{p10:thm:hurwitz-shell-adjacency-spectrum}
Let \(A_H\) be the adjacency matrix of \((V_H,E_H)\).  Then
\begin{equation}
\label{p10:eq:hurwitz-adjacency-spectrum}
\operatorname{Spec}(A_H)=
8^{(1)}\oplus4^{(4)}\oplus0^{(9)}\oplus(-2)^{(8)}\oplus(-4)^{(2)}.
\end{equation}
Equivalently, the graph Laplacian \(L_H=8I-A_H\) has spectrum
\begin{equation}
\label{p10:eq:hurwitz-laplacian-spectrum}
\operatorname{Spec}(L_H)=
0^{(1)}\oplus4^{(4)}\oplus8^{(9)}\oplus10^{(8)}\oplus12^{(2)}.
\end{equation}
Thus the primitive shell algebra has the fixed decomposition
\begin{equation}
\label{p10:eq:hurwitz-shell-band-decomposition}
\CC^{24}=\mathcal V_1\oplus\mathcal V_4\oplus\mathcal V_9\oplus\mathcal V_8\oplus\mathcal V_2.
\end{equation}
In particular the adjacent native \(9\)- and \(8\)-band dimensions define the geometric normalized count
\begin{equation}
\label{p10:eq:delta-from-hurwitz-eight-nine-bands}
\deltageom=(\dim\mathcal V_9\dim\mathcal V_8)^{-1/2}=\frac1{\sqrt{9\cdot8}}=\frac1{6\sqrt2}.
\end{equation}
\end{theorem}
\begin{proof}
This is the Bose--Mesner spectrum of the regular \(24\)-cell/Hurwitz nearest-neighbour graph imported from Paper~I-A.  The graph is \(8\)-regular.  Subtracting the adjacency spectrum from \(8\) gives the Laplacian spectrum.  The idempotent ranks are fixed by the corresponding association scheme.  The resulting \(8\)-by-\(9\) number is a geometric invariant of this decomposition; identifying it with a physical coupling requires the separately typed source map described in Section~\ref{p10:sec:imported-interface}.
\end{proof}

\begin{lemma}[Hurwitz cube-link second moment]
\label{p10:lem:hurwitz-cube-link-second-moment}
At each vertex \(q\in V_H\), the link \(N(q)\) is a cube.  In a local tangent frame its eight directions may be written
\begin{equation}
\label{p10:eq:hurwitz-cube-link-directions}
\hat n=\frac1{\sqrt3}(\pm1,\pm1,\pm1).
\end{equation}
They satisfy
\begin{equation}
\label{p10:eq:hurwitz-cube-link-moment-zero}
\sum_{\hat n\in N(q)}\hat n_a=0,
\end{equation}
\begin{equation}
\label{p10:eq:hurwitz-cube-link-second-moment}
\sum_{\hat n\in N(q)}\hat n_a\hat n_b=\frac83\delta_{ab}.
\end{equation}
\end{lemma}
\begin{proof}
The signs in the cube link cancel all odd and off-diagonal contributions.  Each coordinate square contributes \(1/3\), and there are eight cube vertices.  Hence the diagonal second moment is \(8/3\), giving \eqref{p10:eq:hurwitz-cube-link-second-moment}.
\end{proof}

\begin{theorem}[Local Hurwitz realization of the imported Paper--V-A color carrier]
\label{p10:thm:hurwitz-trimer-color-carrier}
Choosing an oriented body diagonal in the cube link gives the canonical finite split
\begin{equation}
\label{p10:eq:hurwitz-link-trimer-split}
N(q)=\{u_+(q),u_-(q)\}\sqcup T_+(q)\sqcup T_-(q),
\qquad |T_+(q)|=|T_-(q)|=3.
\end{equation}
The oriented-axis stabilizer is \(2C_3\), while forgetting the axial orientation gives
\begin{equation}
\label{p10:eq:local-binary-dihedral-color-skeleton}
N_{2O}(2C_3)=2D_3,
\qquad
2D_3/2C_3\cong C_2,
\qquad
2D_3/\{\pm1\}\cong W(A_2).
\end{equation}
Define the local trimer span
\begin{equation}
\label{p10:eq:local-trimer-color-carrier}
C_q:=\CC^{T_+(q)}\cong\CC^3.
\end{equation}
There is a canonical color-frame identification, unique up to the finite
Hurwitz/cube symmetry,
\[
\iota_q:C_q\xrightarrow{\sim}C,
\]
where \(C\) is the color carrier already selected in Paper~V-A.  Transporting the
Paper--V-A color theorem through \(\iota_q\) gives
\begin{equation}
\label{p10:eq:local-trimer-adjoint}
\begin{gathered}
\operatorname{End}(C_q)\cong\mathbf1\oplus\mathbf8,\\
\operatorname{End}_0(C_q)\cong\mathfrak{sl}_3(\CC),\\
\mathfrak{su}(C_q)\subset\operatorname{End}_0(C_q)
\ \text{is its compact real form}.
\end{gathered}
\end{equation}
Thus the Hurwitz cube link realizes the imported color carrier locally; it does
not select color, the gauge group, or the adjoint octet anew.
\end{theorem}
\begin{proof}
The body diagonal selects two axial endpoints and two opposite triangular
triples, so the local transverse span has dimension three.  Paper~I-A supplies the
binary-dihedral completion and projective \(A_2\) Weyl action.  Paper~V-A has
already assigned one visible trimer to the color carrier and closed its
sector-preserving traceless unitary action to \(SU(3)_c\).  Mapping the local
trimer basis to that imported carrier transports the Paper--V-A decomposition
\(\operatorname{End}(C)=\mathbf1\oplus\mathbf8\) and
\(\operatorname{End}_0(C)=\mathfrak{su}(3)_c\) to \(C_q\).  A different oriented
diagonal changes only the finite color frame by Hurwitz/cube symmetry and
therefore supplies an equivalent local realization, not a new physical
projector or a second color-selection theorem.
\end{proof}

\begin{corollary}[QCD realization of the Paper--V-B split carrier]
\label{p10:cor:paper6-split-carrier-realization}
The Hurwitz-shell transfer complex realizes, on the QCD side, the local carrier
shape used by the Paper~V-B quark split packet.  Removing the axial pair in
\eqref{p10:eq:hurwitz-link-trimer-split} gives the axis-removed support
\begin{equation}
\label{p10:eq:paper6-axis-removed-six-support}
C_6(q):=T_+(q)\sqcup T_-(q),
\qquad
|C_6(q)|=6=3+3.
\end{equation}
After the sector orientation is fixed, \eqref{p10:eq:paper6-axis-removed-six-support}
is the local finite-carrier meaning of the Paper~V-B notation
\begin{equation}
\label{p10:eq:paper6-six-as-two-trimers}
6=3_u+3_d.
\end{equation}
This corollary fixes only the local split support.  The global complement
\begin{equation}
\label{p10:eq:paper6-global-complement-not-local-link}
\mathcal G_{16}=\mathcal G_7\sqcup\mathcal G_9
\end{equation}
and the compensated endpoint coefficients
\begin{equation}
\label{p10:eq:paper6-endpoint-coefficients-not-local-link}
\kappa_u=\frac1{\sqrt2},
\qquad
\kappa_d=\frac1{8\sqrt2}
\end{equation}
remain Paper~V-A/Paper~V-B shell-complement data.  They are not consequences of the
local cube link alone.
\end{corollary}
\begin{proof}
Theorem~\ref{p10:thm:hurwitz-trimer-color-carrier} gives the axial pair and the
two transverse trimers in the Hurwitz cube link.  Deleting the axial pair leaves
the unique six-point two-trimer support, so the local carrier shape is exactly
the one used by the Paper~V-B split packet.  The complement
\(\mathcal G_{16}=\mathcal G_7\sqcup\mathcal G_9\) and the pair
\((\kappa_u,\kappa_d)\) require the first global shell complement and the
compensated shell operator, hence belong to the upstream shell-flow derivation
rather than to the local cube-link identity.
\end{proof}

\begin{theorem}[Hurwitz-shell principal symbol and fixed-scale spinorial remainder]
\label{p10:thm:hurwitz-shell-dirac-compression}
Let $S_H$ be point sampling of a $C^3$ smooth spinor in local spin frames at
the $24$ Hurwitz vertices.  Denote spin parallel transport from $q'$ back
to $q$ along its geodesic by $P_{q\leftarrow q'}$.  For the eight tangent
cube directions $u_{qq'}=\widehat{\operatorname{Im}(q^{-1}q')}$ and
$\ell_H=\pi R_{\rm mom}/3$, define
\begin{equation}
\label{p10:eq:hurwitz-dirac-compression}
(D_{H,0}S_H\psi)(q):=\frac{3}{8\ell_H}\sum_{q'\sim q}
 c(u_{qq'}) P_{q\leftarrow q'}\psi(q').
\end{equation}
Then the covariant Taylor expansion has the correct round $S^3$
Dirac principal symbol, and its fixed-shell remainder satisfies
\begin{equation}
\label{p10:eq:v350-spinorial-fixed-remainder}
\left\|(D_{H,0}S_H\psi)(q)-(S_HD_{S^3}\psi)(q)\right\|
 \le \frac{\ell_H^2}{2}
 \sup_{q'\sim q,\,0\le s\le\ell_H}
 \left\|\nabla_{\dot\gamma_{qq'}}^3\psi(\gamma_{qq'}(s))\right\|.
\end{equation}
The comparison uses covariant parallel transport and the identical
Clifford sign convention on both sides.  Since $\ell_H/R_{\rm mom}=\pi/3$
is fixed, this is not an asymptotic convergence theorem and does not
assert equality on every smooth spinor.
\end{theorem}
\begin{proof}
In the parallel-transported normal frame about $q$ the eight directions
satisfy $\sum u_a=0$, $\sum u_a^iu_a^j=(8/3)\delta^{ij}$ and
$\sum u_a^iu_a^ju_a^k=0$.  Covariant Taylor expansion through second
order gives the zeroth term zero, the first term exactly
$\gamma^i\nabla_i\psi$ after multiplication by $3/(8\ell_H)$, and the
quadratic term zero by antipodal cube symmetry.  In each direction,
the integral order-three remainder has norm at most
$\ell_H^3\sup\|\nabla_{\dot\gamma}^3\psi\|/6$.
Summing eight terms and multiplying by $3/(8\ell_H)$ gives the stated
$\ell_H^2/2$ bound.  The round-space geodesic angle is fixed at
$\pi/3$ on the normalized shell, so no vanishing-step limit is inferred.
\end{proof}

\begin{theorem}[Hurwitz-shell gauge-covariant finite compiler]
\label{p10:thm:hurwitz-shell-qcd-transfer-compiler}
On the Hurwitz cell complex, assign \(SU(3)\) link variables \(U_{qq'}\) and triangular holonomies
\begin{equation}
\label{p10:eq:hurwitz-triangular-face-holonomy}
U_f=U_{q_0q_1}U_{q_1q_2}U_{q_2q_0}.
\end{equation}
The quark stencil
\begin{equation}
\label{p10:eq:hurwitz-gauge-covariant-dirac}
(D_H[U]\psi)(q)=\frac{3}{8\ell_H}\sum_{q'\sim q}c(u_{qq'})P_{q\leftarrow q'}U_{qq'}\psi(q')+M_q^{\rm phys}\psi(q)
\end{equation}
is gauge covariant whenever the physical mass operator is supplied.  On a supplied finite occupation/flux realization with a self-adjoint
same-owner Hamiltonian, an admissible gauge action and this stencil compile
to a finite self-adjoint \(H_H^{\rm QCD}\), and hence to the positive family
\begin{equation}
\label{p10:eq:hurwitz-shell-qcd-transfer}
\widehat K_H(t)=\exp[-t(H_H^{\rm QCD}-E_{0,H})],\qquad t>0.
\end{equation}
The 24-cell geometry fixes the incidence and covariance structure, but does not by itself fix a unique heat time, transition amplitude, or physical energy unit.
\end{theorem}
\begin{proof}
Gauge covariance follows from \(\psi_q\mapsto g_q\psi_q\) and \(U_{qq'}\mapsto g_qU_{qq'}g_{q'}^{-1}\).  The licensed self-adjoint matrix assembly and exponentiation give the
transfer family. Pointwise positivity of a face density alone does not prove
that its temporal integral kernel equals this positive transfer; that
same-owner identification retains its reflection/positive-type certificate.
Distinct admissible gauge actions can share the carrier and symmetries, so
uniqueness of their coefficients is not inferred from geometry.
\end{proof}

\begin{corollary}[No primitive position-lattice input]
\label{p10:cor:no-primitive-position-lattice-input}
The data \(a\mathbb Z^4\), lattice spacing \(a\), Wilson coupling \(\beta_{\rm lat}\), hopping parameter \(\kappa\), Wilson parameter, box volume, and scale-setting prescription are not primitive QCD data in Paper~VII-A.  They are downstream reciprocal-position/comparison-chart conventions after the compact Hurwitz-shell transfer has been fixed.
\end{corollary}
\begin{proof}
The primitive regulator is \(\mathcal K_H\subset S^3_{\rm mom}\) together with the gauge-covariant finite compiler and Dirac stencil above.  Coordinate locality, large Wilson areas, and continuum-looking fields are obtained only by the \(m_x\)-side Born/Fourier unfolding.  Introducing a position grid or scale-setting rule before \eqref{p10:eq:hurwitz-shell-qcd-transfer} is fixed would violate the readout-chronology criterion.
\end{proof}


% Source Paper VII-A, original line 930; inherited proof bodies unchanged.
\subsubsection{Computed dynamical sectors and the finite strong commutant}
\label{p10:subsec:no-branch-strong-algebra}

\begin{definition}[Finite strong generated algebra]
\label{p10:def:finite-strong-generated-algebra}
Let \(\mathfrak A_{L,{\rm str}}\subset\End(\mathcal H_L)\) be the unital \(*\)-algebra generated by the actually compiled gauge-covariant shifts, curvature matrices, quark/Fock operators when present, source operators, and exact upstream projectors.  Its commutant is \(\mathfrak A'_{L,{\rm str}}\).
\end{definition}

\begin{lemma}[Dynamical component projectors are computable outputs]
\label{p10:lem:primitive-exact-block-shell-connectivity}
The support graph of the compiled finite shifts decomposes uniquely into connected components.  The orthogonal projector onto each component is obtained from that graph and commutes with every component-preserving generated operator.  Hence disconnected dynamical sectors are outputs of the compiled QCD algebra; they are not forbidden merely because no upstream label named them.
\end{lemma}
\begin{proof}
Connected components of a finite graph are unique.  Every generator that preserves edges and local multiplication preserves the component span, so its component projector commutes with the generated algebra.  No external selector is introduced.
\end{proof}

\begin{lemma}[Shell matrix units on a connected generated component]
\label{p10:lem:shell-graph-matrix-units}
On any connected support component on which the compiled covariant shifts contain reversible transports along every graph edge, the path products and vertex projectors generate all shell matrix units \(|x\rangle\langle y|\) on that component.
\end{lemma}
\begin{proof}
Choose a path from \(y\) to \(x\); composing the reversible transports and sandwiching with endpoint projectors yields the corresponding matrix unit after undoing the transported fiber action.
\end{proof}

\begin{lemma}[Local color--Fock algebra test]
\label{p10:lem:local-color-fock-matrix-units}
For any fixed finite fiber, the associative algebra generated by the compiled color/Fock matrices is a finite matrix algebra whose irreducible decomposition is obtained constructively by computing its commutant.  It equals the full endomorphism algebra precisely when that commutant is scalar.
\end{lemma}
\begin{proof}
This is the finite-dimensional double-commutant/Burnside criterion over \(\CC\).  It is an algorithmic test on the actual compiled matrices rather than an assumed irreducibility property.
\end{proof}

\begin{lemma}[Gauge-singlet rank-one closure criterion]
\label{p10:lem:gauge-singlet-rank-one-closure}
On a singlet block \(\mathcal H_{\beta,L}\), the gauge-invariant generated algebra contains every rank-one operator if and only if its restricted commutant is scalar.  Thus rank-one closure is equivalent to a checkable irreducibility certificate.
\end{lemma}
\begin{proof}
If all rank-one operators are present, the algebra is \(\End(\mathcal H_{\beta,L})\) and its commutant is scalar.  Conversely, a unital complex matrix algebra with scalar commutant acts irreducibly, and Burnside's theorem gives the full endomorphism algebra.
\end{proof}

\begin{theorem}[Wedderburn decomposition of the finite strong algebra]
\label{p10:thm:finite-strong-algebra-centralizer}
For every finite compiled block there is a unitary decomposition
\begin{equation}
\label{p10:eq:finite-strong-algebra-centralizer}
\mathcal H_L\cong\bigoplus_j\bigl(\CC^{m_j}\otimes\CC^{n_j}\bigr),\qquad
\mathfrak A_{L,{\rm str}}\cong\bigoplus_j\bigl(I_{m_j}\otimes M_{n_j}(\CC)\bigr),
\end{equation}
with
\begin{equation}
\mathfrak A'_{L,{\rm str}}\cong\bigoplus_j\bigl(M_{m_j}(\CC)\otimes I_{n_j}\bigr).
\end{equation}
The minimal central projectors and multiplicity factors are therefore computable dynamical sector data.  A scalar commutant occurs exactly when there is one irreducible block with multiplicity one.
\end{theorem}
\begin{proof}
Apply the structure theorem for finite-dimensional \(C^*\)-algebras to the norm-closed algebra generated by the finite matrices.  The displayed form and commutant are the standard Wedderburn decomposition.
\end{proof}

\begin{corollary}[Irreducibility is a certificate, not a no-insertion axiom]
\label{p10:cor:no-hidden-dynamical-bypass-projector}
A nontrivial projector commuting with the compiled strong algebra is admissible exactly when it occurs in \(\mathfrak A'_{L,{\rm str}}\).  Therefore absence of a dynamical bypass, topological quench sector, or residue-decoupling block must be established by the computed commutant or a stronger positivity/irreducibility theorem; it cannot follow from the slogan ``no new upstream label.''
\end{corollary}
\begin{proof}
By Theorem~\ref{p10:thm:finite-strong-algebra-centralizer}, every commuting projector lies in the explicitly computed commutant.  Conversely, every orthogonal projector in that commutant defines a reducing subspace for the generated algebra.  Thus irreducibility is equivalent to the absence of nontrivial projections in the commutant and is decidable from the compiled matrices.
\end{proof}

\begin{theorem}[Gauge-coupling normalization does not fix a transfer amplitude]
\label{p10:thm:quadratic-strong-seed-normalization}
On any supplied seed branch,
\begin{equation}
\alpha_s=\frac{g_s^2}{4\pi}=\alphasseed
\qquad\Longrightarrow\qquad
g_s=\sqrt{4\pi\alphasseed}.
\end{equation}
No universal finite transition amplitude or probability follows from this identity alone.  Such an amplitude is a matrix element of the compiled transfer/source operator and must be computed from that operator.
\end{theorem}
\begin{proof}
The equality is the definition of \(\alpha_s\).  Multiplying an independent finite transition operator by a nonzero scalar changes its matrix elements while leaving the gauge-covariant derivative coupling fixed, so the latter cannot determine the former without a typed operator identity.
\end{proof}

\begin{theorem}[Scalar-shell response as a computed spectral functional]
\label{p10:thm:primitive-scalar-shell-coefficient-fixed}
Let \(P_{{\rm sh},L}\ne0\) be a specified scalar adjoint-curvature orthogonal projector obtained by spectral calculus from the finite curvature/Hodge operator, and let \(H_L=H_L^*\) on the finite carrier. Define the finite shell Rayleigh energy diagnostic
\begin{equation}
\label{p10:eq:primitive-scalar-shell-coefficient-fixed}
\gamma_{{\rm sh},L}:=
\inf_{\substack{\psi\in P_{{\rm sh},L}\mathcal H_L\\\|\psi\|=1}}
\langle\psi,(H_L-E_{0,L})\psi\rangle\ge0.
\end{equation}
For transfer time \(t>0\), define the Rayleigh-gap diagnostic and, separately, the actual compressed-transfer loss on the nonzero shell:
\begin{equation}
\begin{aligned}
\epsilon_{{\rm sh},L}^{\rm Ray}(t)&:=1-e^{-t\gamma_{{\rm sh},L}},\\
\epsilon_{{\rm sh},L}^{\rm tr}(t)&:=1-\bigl\|P_{{\rm sh},L}e^{-t(H_L-E_{0,L})}P_{{\rm sh},L}\big|_{\operatorname{Ran}P_{{\rm sh},L}}\bigr\|.
\end{aligned}
\end{equation}
They coincide when \([P_{{\rm sh},L},H_L]=0\).  Without that reducing-subspace certificate, the Rayleigh functional is not the actual projected-transfer contraction.
These quantities are outputs of the compiled Hamiltonian.  In particular, neither \(24/15\) nor \(\sqrt{\deltageom}\) is used as a shell amplitude.  If a diagnostic ratio to a supplied positive physical seed is desired, one may report \(c_{{\rm sh},L}^{\rm Ray}(t):=\sqrt{\epsilon_{{\rm sh},L}^{\rm Ray}(t)/\alphasseed}\); it is a Rayleigh readout, not a physical transition amplitude or input.
\end{theorem}
\begin{proof}
The nonzero finite-dimensional shell has a compact unit sphere, so the Rayleigh infimum is finite, nonnegative, and basis independent. The compressed exponential is a positive contraction because \(H_L-E_{0,L}\ge0\). If \(P_{{\rm sh},L}\) reduces \(H_L\), its restricted spectrum gives \(\|P e^{-t(H_L-E_{0,L})}P|_{\operatorname{Ran}P}\|=e^{-t\gamma_{{\rm sh},L}}\). For noncommuting \(P\) and \(H_L\) this last identity does not follow from Rayleigh--Ritz. The current Paper~I-A source separates 24/15 geometric counts from amplitudes.
\end{proof}


% Source Paper VII-A, original line 1029; inherited proof bodies unchanged.
% Former logical part: Finite Color Datum, Singlet Projection, and Executable Transcript; integrated into the new thematic argument.
\label{p10:part:finite-color-executable-transcript}


\noindent\textbf{Part orientation.}
This part turns the finite QCD ledger into the explicit computable datum: the singlet basis, sparse Hamiltonian factorization, source dictionary, transfer/readout maps, and admissible color-connection normalization.  It is the finite operator/compiler layer before the strong clock and nonperturbative readouts are evaluated.


% Source Paper VII-A, original line 1036; inherited proof bodies unchanged.
\subsection{Executable finite-QCD transcript}
\label{p10:sec:executable-finite-qcd-transcript}

This section gives the constructive compiler interface: the canonical finite
singlet basis, sparse-entry factorization, source-operator dictionary, and
spectral extraction rules.  Together they make
\[
\mathfrak C_L^{\rm QCD}\longmapsto\mathfrak M_L^{\rm QCD}
\]
an explicit finite computation.


% Source Paper VII-A, original line 1047; inherited proof bodies unchanged.
\subsubsection{Finite singlet basis and color projection}
\label{p10:subsec:finite-singlet-basis-transcript}

At refinement/readout level \(L\), let \(R_{q,L}\), \(R_{\bar q,L}\), and \(R_{g,L}\) denote the finite non-colour mode spaces: shell, spin/polarization, flavour, and exact upstream labels, but not colour.  Put
\[
\mathcal Q_L:=R_{q,L}\otimes \CC^3,
\qquad
\bar{\mathcal Q}_L:=R_{\bar q,L}\otimes (\CC^3)^*,
\qquad
\mathcal G_L:=R_{g,L}\otimes {\rm ad}(\CC^3).
\]
For fixed occupation numbers \((n_q,n_{\bar q},n_g)\), define the raw finite Fock block
\begin{equation}
\label{p10:eq:raw-qcd-occupation-block}
\mathcal H_{L}^{\rm raw}(n_q,n_{\bar q},n_g)
=
\wedge^{n_q}\mathcal Q_L
\otimes
\wedge^{n_{\bar q}}\bar{\mathcal Q}_L
\otimes
\operatorname{Sym}^{n_g}\mathcal G_L .
\end{equation}
The ordered monomial basis of \eqref{p10:eq:raw-qcd-occupation-block} is part of the transcript row \(\mathcal B_L^{\rm raw}\).  Its raw dimension is
\begin{equation}
\label{p10:eq:raw-qcd-occupation-dimension}
N_L^{\rm raw}(n_q,n_{\bar q},n_g)
=
\binom{3r_{q,L}}{n_q}
\binom{3r_{\bar q,L}}{n_{\bar q}}
\binom{8r_{g,L}+n_g-1}{n_g},
\end{equation}
where \(r_{q,L}=\dim R_{q,L}\), \(r_{\bar q,L}=\dim R_{\bar q,L}\), and \(r_{g,L}=\dim R_{g,L}\).

\begin{theorem}[Finite occupation assembly and its retained boundary defect]
\label{p10:thm:v2211-finite-occupation-assembly}
The fixed-$n_g$ spaces above are kinematic summands, not automatically
reducing Hamiltonian sectors. An actual finite QCD occupation carrier consumes
the inherited finite occupation/capacity projector $P_{{\rm occ},L}$ of its
Paper-IV representative, or an explicitly licensed finite occupation sum.
Finite one-body mode count alone does not make unrestricted bosonic Fock
space finite. On that supplied finite-rank carrier,
\[
 H_{{\rm occ},L}=P_{{\rm occ},L}H_L^{\rm parent}P_{{\rm occ},L}
 \big|_{\operatorname{ran}P_{{\rm occ},L}}
\]
is self-adjoint whenever the parent matrix is self-adjoint on the licensed
comparison domain. If the projector commutes with the color action and each
retained conserved charge, these symmetries survive compression. A fixed
occupation label is a conserved block only after its commutator with the
actual interaction has been checked.

For a single bosonic mode truncated to $0,\ldots,N$, its actual annihilator
$a_N$ obeys exactly
\[
 [a_N,a_N^\dagger]=I-(N+1)|N\rangle\langle N|.
\]
If the capacity/source chart changes, the derivative of the compressed
parent includes all three terms
$P'HP+PH'P+PHP'$, expressed in the same transported finite frame. No
boundary CCR or source contact is discarded by calling the mode count finite.
\end{theorem}
\begin{proof}
Each licensed occupation summand has the finite binomial dimension above;
a finite sum is finite. Orthogonal compression preserves adjoints, and
commuting projectors preserve the stated symmetry restrictions. On number
states below $N$, the commutator is one; at $N$ it is $-N$, yielding the
displayed boundary term. The source identity is the ordinary product rule
in the complete canonical-frame pullback. It supplies a constructor and
test once the upstream capacity is supplied; it selects no new physical cap.
\end{proof}

\begin{definition}[Finite colour-singlet projector]
\label{p10:def:finite-colour-singlet-projector}
On \(\mathcal H_{L}^{\rm raw}(n_q,n_{\bar q},n_g)\), let
\[
U_L(g)
=
\wedge^{n_q}(1\otimes g)
\otimes
\wedge^{n_{\bar q}}(1\otimes g^*)
\otimes
\operatorname{Sym}^{n_g}(1\otimes{\rm Ad}_g),
\qquad g\in SU(3)_c .
\]
The exact \emph{diagonal/global} colour-singlet projector is
\begin{equation}
\label{p10:eq:finite-colour-singlet-projector}
\Ps^{(L)}=P_{\rm sing}^{(L)}:=\int_{SU(3)}U_L(g)\,dg .
\end{equation}
Its image \(\mathcal H_{L,\rm glob\,sing}^{\rm QCD}:=P_{\rm sing}^{(L)}\mathcal H_L^{\rm raw}\) is a global singlet trial block, not automatically the local gauge-physical space. When the actual vertexwise link/matter representation and invariant truncation exist, Theorem~\ref{p10:thm:v350-local-gauss-projector} defines \(P_{\rm Gauss}\), and only \(\operatorname{Ran}P_{\rm Gauss}\subseteq\mathcal H_{L,\rm glob\,sing}^{\rm QCD}\) is certified as the locally gauge-invariant block. Any use of ``physical'' for this global projector, earlier or later in this paper, requires that additional local-Gauss reduction.
\end{definition}

\begin{theorem}[Finite colour-singlet basis transcript]
\label{p10:thm:finite-colour-singlet-basis-transcript}
The projector \eqref{p10:eq:finite-colour-singlet-projector} is Hermitian, idempotent, and fixed before comparison.  The singlet dimension in an occupation sector is
\begin{equation}
\label{p10:eq:singlet-dimension-character-formula}
\begin{aligned}
d_{\rm sing}(L;n_q,n_{\bar q},n_g)
&=[z_q^{n_q}z_{\bar q}^{n_{\bar q}}z_g^{n_g}]
\int_{SU(3)}
\det(1+z_q g)^{r_{q,L}}\\
&\quad\times
\det(1+z_{\bar q}g^*)^{r_{\bar q,L}}
\det(1-z_g{\rm Ad}_g)^{-r_{g,L}}\,dg .
\end{aligned}
\end{equation}
In particular, the centre of \(SU(3)_c\) imposes the exact triality law
\begin{equation}
\label{p10:eq:qcd-triality-law}
n_q-n_{\bar q}\equiv0\pmod 3 .
\end{equation}
Thus baryon number \(B=(n_q-n_{\bar q})/3\) is integer on every singlet block.
\end{theorem}
\begin{proof}
The Haar average of a unitary representation is the orthogonal projection onto the invariant subspace, giving Hermiticity and idempotence.  Taking the trace of the averaged representation gives the character integral.  The generating function follows from exterior powers of \(\mathbf 3\) and \(\bar{\mathbf 3}\) and symmetric powers of \(\mathbf 8\).  If \(z\in Z(SU(3))\), a raw sector transforms by \(z^{n_q-n_{\bar q}}\).  Invariance under all centre elements is equivalent to \eqref{p10:eq:qcd-triality-law}.
\end{proof}

The primitive invariant colour tensors used later are therefore fixed:
\[
q\bar q:\frac{1}{\sqrt3}\delta_i{}^j,
\qquad
qqq:\frac{1}{\sqrt6}\epsilon_{ijk},
\qquad
gg:\frac{1}{\sqrt8}\delta_{ab},
\qquad
ggg:\frac{1}{\sqrt{24}}f_{abc},\ \sqrt{\frac3{40}}d_{abc}.
\]
No alternative colour projector or phenomenological singlet weight is admissible.


% Source Paper VII-A, original line 1177; inherited proof bodies unchanged.
\subsubsection{Sparse Hamiltonian entry factorization}
\label{p10:subsec:sparse-entry-factorization}

Let \(e_A,e_B\in\mathcal B_L^{\rm raw}\).  Every normal-ordered elementary QCD monomial in the finite basis has a matrix element of the form
\begin{equation}
\label{p10:eq:elementary-entry-factorization}
\langle e_A|O_v|e_B\rangle
=
\eta_v(A,B)\,N_v(A,B)\,S_v^{(L)}(A,B)\,C_v(A,B)\,F_v(A,B)\,\kappa_v .
\end{equation}
Here \(\eta_v\) is the fermionic reordering sign, \(N_v\) the bosonic occupation square-root factor, \(S_v^{(L)}\) the native shell/mode overlap, \(C_v\) the colour tensor contraction, \(F_v\) the spin/flavour contraction, and \(\kappa_v\) the fixed clock, mass, or coupling coefficient.

\begin{theorem}[Polynomial canonical QCD sparse-entry factorization]
\label{p10:thm:finite-qcd-sparse-entry-factorization}
On the polynomial canonical Yang--Mills/quark branch, every nonzero entry
of the singlet Hamiltonian
\[
H_L=P_{\rm sing}^{(L)}(H_{\YM,L}+H_{q,L}+H_{{\rm int},L})P_{\rm sing}^{(L)}
\]
can be written as a finite sum
\begin{equation}
\label{p10:eq:singlet-entry-factorization}
(H_L)_{rs}
=
\sum_v
\eta_v\,N_v\,S_v^{(L)}\,C_v\,F_v\,\kappa_v .
\end{equation}
The only colour tensors appearing in \eqref{p10:eq:singlet-entry-factorization} are generated by
\begin{equation}
\label{p10:eq:allowed-colour-tensors}
\delta_i{}^j,
\qquad
(T^a)_i{}^j,
\qquad
f_{abc},
\qquad
d_{abc},
\qquad
f_{abe}f_{cde},
\qquad
T^a\otimes T^a,
\end{equation}
and their Haar-singlet contractions.  The coefficient list is fixed by
\begin{equation}
\label{p10:eq:allowed-qcd-entry-coefficients}
\kappa_v\in
\{1,M_{q,L}^{\rm phys},g_s,g_s^2\},
\qquad
\frac{g_s^2}{4\pi}=\alphasseed .
\end{equation}
Therefore no hadron mass, string tension, fitted PDF, fragmentation function, thermal table, lattice scale, or collider datum can appear as a matrix-entry coefficient.
\end{theorem}
\begin{proof}
The finite Yang--Mills quadratic term carries colour \(\delta_{ab}\), the cubic term carries \(g_sf_{abc}\), and the quartic term carries \(g_s^2f_{abe}f_{cde}\), up to the finite shell/polarization contractions.  The free quark term carries \(\delta_i{}^j\) and the typed physical mass matrix \(M_{q,L}^{\rm phys}\).  The quark--gluon vertex carries \(g_s(T^a)_i{}^j\), and current--current or Coulomb blocks carry \(g_s^2T^a\otimes T^a\).  Finite Fock algebra supplies only the signs and occupation factors.  Projection by \(P_{\rm sing}^{(L)}\) is a fixed finite linear combination of raw entries and therefore preserves factorization as a finite sum.
\end{proof}

The executable sparse row format is consequently
\begin{equation}
\label{p10:eq:executable-sparse-row-format}
(i,j;{\rm block},v,\eta_v,N_v,S_v^{(L)},C_v,F_v,\kappa_v,{\rm value}).
\end{equation}
This is the constructive content of the polynomial finite-matrix audit
row. For the complete heat or generated multiface action, the exact
Haar/character compiler of Theorem~\ref{p10:thm:v199-haar-interaction-compiler}
is the corresponding dictionary: its fixed representation weights,
nonlinear vertices and complete source jets are retained as entries of the
actual action member. They are not replaced by the polynomial coefficient
list, and are not independently fitted to downstream data.


% Source Paper VII-A, original line 1246; inherited proof bodies unchanged.
\subsubsection{Finite source-operator dictionary}
\label{p10:subsec:finite-source-operator-dictionary}

The Hamiltonian alone gives spectra, but not residues, decay constants, distribution moments, or thermal response.  The transcript must therefore include a finite operator dictionary.

\begin{definition}[Finite QCD source-operator dictionary]
\label{p10:def:finite-source-operator-dictionary}
The minimal native source list is
\begin{equation}
\label{p10:eq:minimal-source-operator-dictionary}
\mathcal O_L^{\rm min}
=
\left\{\begin{gathered}
O_G^{0^{++}},O_G^{0^{-+}},O_G^{2^{++}},O_S^T,O_P^a,
O_V^{\mu,a},O_A^{\mu,a},J_{\rm em}^{\mu},\\
O_B,W_R(C),\nu_L,O_{q,n},O_{g,n},O_{{\rm DA},n},B,Q,S,P_A
\end{gathered}\right\}.
\end{equation}
Each element is a finite normal-ordered, gauge-singlet operator built from the same \(a,a^\dagger,b,b^\dagger,d,d^\dagger\) algebra and then projected by \(P_{\rm sing}^{(L)}\).
\end{definition}

The gluonic sources are represented by
\[
O_G^{0^{++}}=P_{\rm sing}^{(L)}\sum_x{\rm tr}\bigl(F_{\mu\nu}F^{\mu\nu}\bigr)P_{\rm sing}^{(L)},
\qquad
O_G^{0^{-+}}=P_{\rm sing}^{(L)}\sum_x{\rm tr}\bigl(F_{\mu\nu}\widetilde F^{\mu\nu}\bigr)P_{\rm sing}^{(L)},
\]
and the tensor source is the traceless spatial stress component.  The quark bilinear sources are
\[
O_{\Gamma,T}^{(L)}=P_{\rm sing}^{(L)}\sum_x\bar q(x)\Gamma Tq(x)P_{\rm sing}^{(L)},
\]
with \(\Gamma=1,i\gamma_5,\gamma^\mu,\gamma^\mu\gamma_5\) for scalar, pseudoscalar, vector, and axial readouts.  Baryon sources have the determinant colour tensor
\[
O_B^\alpha=P_{\rm sing}^{(L)}\sum_x\epsilon_{ijk}\bigl(q_i^T C\Gamma_1q_j\bigr)\Gamma_2q_k^\alpha P_{\rm sing}^{(L)}.
\]
Distribution operators are Wilson-line-completed twist-two moments,
\[
O_{q_f,n}^{(L)}=P_{\rm sing}^{(L)}\bar q_f\slashed n(i n\!\cdot\!D)^{n-1}q_f P_{\rm sing}^{(L)}-{\rm traces},
\]
\[
O_{g,n}^{(L)}=P_{\rm sing}^{(L)}F^{+\alpha}(i n\!\cdot\!D)^{n-2}F^+{}_{\alpha}P_{\rm sing}^{(L)}-{\rm traces}.
\]

\begin{theorem}[Current and spin-two descendants of the native twist-two source family]
\label{p10:thm:v264-twist-two-current-stress-owner}
Let \(n\) be the generated null direction used in the finite source dictionary, and let
\(T_{\mu\nu,L}^{\rm QCD}\) denote the stress tensor obtained by metric variation of the same complete finite QCD action/source owner.  With the source normalizations of Definition~\ref{p10:def:finite-source-operator-dictionary}, the first two members of the twist-two tower satisfy
\begin{equation}
\label{p10:eq:v264-n1-current-owner}
 O_{q_f,1}^{(L)}
 =P_{\rm sing}^{(L)}\,n_\mu J_f^\mu\,P_{\rm sing}^{(L)},
\end{equation}
and, on every forward physical singlet matrix element for which the finite source insertions are defined,
\begin{equation}
\label{p10:eq:v264-n2-forward-stress-owner}
 \boxed{
 \sum_f\langle h,L|O_{q_f,2}^{(L)}|h,L\rangle
 +\langle h,L|O_{g,2}^{(L)}|h,L\rangle
 =\langle h,L|T_{L}^{nn}|h,L\rangle,}
 \qquad
 T_L^{nn}:=n^\mu n^\nu T_{\mu\nu,L}^{\rm QCD}.
\end{equation}
Total-derivative and improvement representatives that differ off forward give no extra forward matrix element.  Hence the charge and total-momentum normalizations of the first two moments are descendants of the native current/metric source owner and do not require an exact finite one-body/Fock realization of \(P^+\).
\end{theorem}
\begin{proof}
For \(n=1\), Definition~\ref{p10:def:finite-source-operator-dictionary} removes all covariant derivatives, so the quark twist-two insertion is exactly the null contraction of the normalized vector current, giving \eqref{p10:eq:v264-n1-current-owner}.  For \(n=2\), varying the complete quark-plus-gluon action with respect to the metric gives the sum of the quark and gluon symmetric spin-two pieces.  Contracting twice with a null vector removes pure trace terms because \(n^2=0\).  Forward matrix elements annihilate total derivatives, while source normalization is inherited from the same metric/action owner by Theorem~\ref{p10:thm:no-independent-source-normalization}.  This yields \eqref{p10:eq:v264-n2-forward-stress-owner}.  No species projector, one-body plus generator, or second-quantized fraction operator enters the argument.
\end{proof}

\begin{remark}[Exact finite LFRep is a stronger readout realization]
\label{p10:rem:v264-lfrep-stronger-than-twist-two-owner}
The source-owned light-ray/twist-two hierarchy and Theorem~\ref{p10:thm:v264-twist-two-current-stress-owner} are defined before any exact finite species/Fock equality.  Clause~\textnormal{(R3)} of Definition~\ref{p10:def:finite-null-representation-certificate} is still required when a claim asks for the stronger finite commuting fraction-simplex or positive number/momentum POVM realization.  It is not a premise for the current, stress-tensor, or higher twist-two source operators themselves.
\end{remark}

The thermal charges are
\[
B=\frac{1}{3}\bigl(N_q-N_{\bar q}\bigr),
\qquad
Q=\sum_f Q_f(N_f-N_{\bar f}),
\qquad
S=-(N_s-N_{\bar s}).
\]

\begin{theorem}[Source normalization and dynamical residues]
\label{p10:thm:no-independent-source-normalization}
Once a source coupling convention in \(H_L[J]\) is fixed, its finite operator is the derivative \(O_A=\partial H_L[J]/\partial J_A|_{J=0}\) and requires no further arbitrary normalization.  However the matrix elements
\(\langle\Psi_m|O_A|\Psi_n\rangle\), decay constants, form factors, PDF moments, and susceptibilities are dynamical spectral readouts and are not fixed by the operator's Ward or charge normalization alone.
\end{theorem}
\begin{proof}
Differentiation fixes the operator associated with a normalized source.  Matrix elements additionally depend on the eigenvectors of \(H_L\); varying the dynamics while preserving the same current algebra changes those residues.  Thus source normalization and dynamical residue normalization are logically distinct.
\end{proof}


\begin{definition}[Native QCD generating functional with sources]
\label{p10:def:native-qcd-generating-functional}
For a finite source packet
\[
J=(v_\mu,a_\mu,s,p,J_{\rm em},J_B,J_{\rm heavy},\ldots),
\]
let \(D_H[U;J]\) denote the Hurwitz-shell Dirac operator and Wilson/source insertions
obtained by the corresponding covariant deformation of the transfer compiler.  The
native finite generating functional is
\begin{equation}
\label{p10:eq:native-qcd-generating-functional}
Z_H[J]
=
\operatorname{Tr}_{\Ps\mathcal H_{\rm raw}}
K_H^{\rm QCD}[J],
\qquad
W_H[J]=\log Z_H[J].
\end{equation}
On the strictly positive finite exponential-transfer branch
\(K_H^{\rm QCD}[J]=e^{-tH_H[J]}\), with fixed \(t>0\) and fixed invariant carrier, the normalized Hamiltonian insertion and the connected log-source responses are
\begin{equation}
\label{p10:eq:native-source-operator-derivative}
\mathcal O_A^{\rm BH}:=\left.\partial_{J_A}H_H[J]\right|_{J=0}
=-\left.\frac1t\frac{\delta\log K_H^{\rm QCD}[J]}{\delta J_A}\right|_{J=0},
\end{equation}
\begin{equation}
\label{p10:eq:native-connected-source-derivatives}
\mathcal C_{A_1\cdots A_n}^{\rm src}
:=\left.\frac{\delta^n W_H[J]}
{\delta J_{A_1}\cdots\delta J_{A_n}}\right|_{J=0}.
\end{equation}
In particular \(\mathcal C_A^{\rm src}=-t\langle\mathcal O_A^{\rm BH}\rangle\). For \(n\ge2\), these derivatives are \emph{Duhamel-ordered} connected source responses: nonlinear source-contact jets and noncommuting insertions are retained. They are not, without extra source-linear/time-order normalization, ordinary unweighted products of the \(\mathcal O_A^{\rm BH}\). The complete Hessian is Theorem~\ref{p10:thm:v350-dynamical-source-hessian}.
\end{definition}

\begin{definition}[Complete finite QCD source-transfer object]
\label{p10:def:complete-finite-qcd-source-transfer-datum}
The complete finite Paper-VII-A source-transfer kernel datum is
\begin{equation}
\label{p10:eq:complete-finite-qcd-source-transfer-object}
\mathfrak S_{\rm QCD}^{\rm BH}
:=
\left(
\mathfrak C_L^{\rm QCD},
\mathfrak M_L^{\rm QCD},
\mathcal O_L,
W_H[J]
\right).
\end{equation}
It contains the finite state space, singlet Hamiltonian and transfer, kernel spectrum,
normalized source operators, topological operator, and source-generating
functional required by the kernel, chiral, anomaly, and topology sectors.  It is
not the completed physical readout datum: Paper~VII-B obtains that datum by executing
the fixed spectral, source, Schur, transition, trace, and cumulant maps on
\(\mathfrak S_{\rm QCD}^{\rm BH}\).
\end{definition}

\begin{theorem}[Finite source-transfer closure and dependency criterion]
\label{p10:thm:master-finite-nonperturbative-qcd-closure}
\label{p10:thm:no-observational-input-covered-qcd-sectors}
Every Paper-VII-A kernel quantity is a fixed algebraic, spectral, transfer, source,
or topology functional of \(\mathfrak S_{\rm QCD}^{\rm BH}\).  Every Paper-VII-B
execution quantity is obtained from the same object by finite diagonalization,
source differentiation, Schur/Feshbach extraction, transition spectral measure,
finite trace, or transfer cumulant.  Hadron, lattice, distribution, thermal,
nuclear, collider, or cosmological data may validate these outputs but may not
select a coefficient, projector, matrix entry, scale, state, or source
normalization.
\end{theorem}
\begin{proof}
The compiler fixes the finite matrices and the source dictionary before
comparison.  The listed extraction operations are deterministic functions of
those objects.  Any observational selector would change the dependency graph
and violate the readout-chronology theorem.
\end{proof}

\begin{theorem}[Source-interface universality for downstream QCD]
\label{p10:thm:source-interface-universality}
Every downstream QCD observable supported by the Paper-VII-A interface is obtained from one of three native
operations:
\begin{enumerate}[label=\textnormal{(\roman*)},leftmargin=2.1em]
\item a finite singlet spectral value of \(H_H^{\rm QCD}\);
\item a source derivative of \(W_H[J]\);
\item a finite Schur/Feshbach compression and its real spectral
      projector/residue; a continued resonance pole only after the separate
      open-channel limit exported to Paper~VII-B.
\end{enumerate}
Thus the corresponding Paper-VII-B spectra, residues, distributions, thermal traces, and multi-hadron observables require no independent QCD input after
\((K_H^{\rm QCD},\mathcal O_L)\) are fixed.
\end{theorem}
\begin{proof}
The source packet deforms only the already-fixed transfer compiler.  Differentiating
\(W_H[J]\) inserts finite normal-ordered gauge-singlet operators, while spectral and
Schur extraction use the same finite Hamiltonian blocks.  Therefore all such observables
are generated by fixed algebraic operations on the native finite datum.  Any additional
normalization chosen from comparison data would change \(K_H^{\rm QCD}[J]\) or
\(\mathcal O_A^{\rm BH}\) and would therefore be an illicit new QCD input.
\end{proof}

\begin{corollary}[Finite measures and sector readouts reuse the QCD parent]
\label{p10:cor:finite-qcd-parent-readouts}
Retain the complete finite source-transfer datum
\(\mathfrak S_{\rm QCD}^{\rm BH}\), together with the actual state, thermal
datum and normalized source test where the readout consumes them.
Every covered finite hadron/nuclear spectral or Schur extraction, topological
transfer weight and current transition measure is then a licensed function
of that same packet.  It needs no additional occurrence or independent
source/measure normalization.
In particular a declared complete sector-effect family
\(E_\nu\ge0\), \(\sum_\nu E_\nu=I\), has
\(p_\nu=\omega(E_\nu)\); on the theta-transfer presentation the weights
instead use its declared normalized source-functional coefficients.
For an owned finite Hermitian current and vacuum, its connected transition
weights are the squared current matrix elements at the same energy
differences.  If the complete parent is retained as a compatible family,
retain the joint image of these readouts; only common values close without
a required member selection.
\end{corollary}
\begin{proof}
The master finite closure and source-interface theorems give deterministic
spectral, source, Schur, transition, trace and cumulant maps.
Paper~IV's finite Born evaluation fixes the effect weights; finite spectral
calculus fixes the current weights and their energy support.
Their inputs and normalizations are already retained in the same source
packet.  Paper~II-A complete-output and whole-family reuse supplies their
compatible image without another G0 premise or a new measure law.
\end{proof}
Physical sector identification, populated histories, ordered chiral or
thermodynamic limits, uniform electromagnetic IR/UV/contact bounds and
continued resonance/scattering claims remain certificates only where
consumed.  APBI is a sufficient physical-confinement route on its domain,
not an antecedent for executing every finite singlet functional.



% Source Paper VII-A, original line 1474; inherited proof bodies unchanged.
\subsubsection{Spectral readout theorem}
\label{p10:subsec:spectral-readout-theorem}

Let \(H_L=H_L(0)\) have a \emph{simple isolated} lowest eigenvalue \(E_{0,L}\), with gap \(E_{1,L}-E_{0,L}>0\), and write
\[
H_L|n,L\rangle=E_{n,L}|n,L\rangle,
\qquad
|\Omega_L\rangle:=|0,L\rangle .
\]
Let \(H_L(J)\) be a twice differentiable self-adjoint finite source family preserving the stated sector, and put \(O_A^{(L)}=\partial_AH_L(0)\), \(O_{AB}^{(L)}=\partial_A\partial_BH_L(0)\). For these specified source insertions define
\begin{equation}
\label{p10:eq:finite-residue-definition}
R_{A,n}^{(L)}:=\langle\Omega_L|O_A^{(L)}|n,L\rangle,
\qquad
M_{n,L}:=E_{n,L}-E_{0,L}.
\end{equation}
Then the connected Euclidean two-point function is
\begin{equation}
\label{p10:eq:finite-two-point-spectral-readout}
C_{AB}^{(L)}(t)
=
\sum_{n>0}R_{A,n}^{(L)}R_{B,n}^{(L)*}e^{-M_{n,L}t}.
\end{equation}
Define the static source susceptibility by \(\chi_{AB}^{(L)}:=-\partial_A\partial_BE_{0,L}(J)|_{J=0}\). Nondegenerate perturbation theory gives the \emph{complete} contact-aware identity
\begin{equation}
\label{p10:eq:finite-susceptibility-readout}
\chi_{AB}^{(L)}
=
2\sum_{n>0}
\frac{\operatorname{Re}\!\bigl(R_{A,n}^{(L)}R_{B,n}^{(L)*}\bigr)}{M_{n,L}}
-\langle\Omega_L|O_{AB}^{(L)}|\Omega_L\rangle.
\end{equation}
The spectral sum alone applies to a locally \emph{linear} Hamiltonian source family with \(O_{AB}^{(L)}=0\). A degenerate or gapless ground space requires a separate ground-subspace response, not this divided-gap formula.
Three-point matrix elements are
\begin{equation}
\label{p10:eq:finite-three-point-matrix-element}
\mathcal M_{mn}^{A,(L)}:=\langle m,L|O_A^{(L)}|n,L\rangle .
\end{equation}
Thus finite spectral masses, residues, decay constants, charges, magnetic
moments, axial couplings, scalar/tensor charges, and transition matrix elements
are all fixed by \((H_L,\mathcal O_L)\).  Their later interpretation as
asymptotic scattering amplitudes is not a finite-kernel claim.

For a finite \(P/Q\) block and
\(z\in\rho(H_L)\cap\rho(H_{QQ})\), the exact Schur operator is
\begin{equation}
\label{p10:eq:finite-schur-source-block}
H_{{\rm eff},L}(z)
=H_{PP}+H_{PQ}(z-H_{QQ})^{-1}H_{QP},
\qquad
P(z-H_L)^{-1}P=[z-H_{{\rm eff},L}(z)]^{-1}.
\end{equation}
As a meromorphic identity, every nonremovable pole of the projected full
resolvent is a real eigenvalue of the self-adjoint \(H_L\), and its residue is
a compressed orthogonal spectral projector.  Poles of the intermediate Schur
coefficient at eigenvalues of \(H_{QQ}\) may cancel and are not by themselves
full-system poles.  In particular,
the finite kernel contains neither a \(+i0\) prescription nor a complex width.
Paper~VII-B owns two distinct limits: norm-resolvent/Riesz convergence of isolated
stable poles and, only after a continuous open-channel limit and meromorphic
continuation, convergence of resonance poles and left/right matrix residues.
The latter are Schur readouts, not LSZ-stable particle inputs.

For compatibility with downstream notation, the former dressed-residue label
is retained on the finite projected-resolvent residue:
\begin{equation}
\label{p10:eq:schur-dressed-residue-normalization}
\operatorname*{Res}_{z=E_{p,L}}
P(z-H_L)^{-1}P=P P_{p,L}P .
\end{equation}


% BEGIN INLINE SOURCE: Paper7A-CompleteSchurSource-Main-v1_0.tex
% Framework v2.164: source and normalization specialization of exact Schur reuse.

% Source Paper VII-A, original line 1550; inherited proof bodies unchanged.
\subsubsection{Complete scalar source and norm of a finite Schur eigenvector}
\label{p10:sec:complete-schur-scalar-source}
\begin{theorem}[Hamiltonian Schur source with inherited residue norm]
\label{p10:thm:complete-schur-scalar-source}
Use a licensed fixed source-independent orthogonal split in the
canonical frame of a complete family $H(h)=H-hO$. Write
\[
 H(h)=\begin{pmatrix}A&X\\Y&C\end{pmatrix},\quad Y=X^\dagger,\quad
 R=(C-E)^{-1},\quad F=A-E-XRY.
\]
For real $E$ with invertible complement, $F(E,h)u=0$ reconstructs
the full eigenvector $(u,-RYu)$. Its inherited norm and scalar source are
\[
 \boxed{\mathscr Z=-\partial_EF=I+XR^2Y,\quad
 O_{\rm eff}=-\partial_hF
 =O_{PP}-O_{PQ}RY-XRO_{QP}+XRO_{QQ}RY.}
\]
Thus
\[
 \boxed{\langle O\rangle
 =\frac{u^\dagger O_{\rm eff}u}{u^\dagger\mathscr Zu}.}
\]
On a differentiable eigenbranch $E'(h)=-\langle O\rangle$.
If $E<\lambda_{\min}(C)$, $R\succ0$ and $XRY\succeq0$.
\end{theorem}
\begin{proof}
Solve the complementary block equation for the eliminated vector.
Its squared norm is $u^\dagger(I+XR^2Y)u$.
Use $\partial_ER=R^2$, $\partial_hR=RO_{QQ}R$,
$\partial_hX=-O_{PQ}$, $\partial_hY=-O_{QP}$, and
$\partial_hA=-O_{PP}$ to differentiate $F$.
Evaluation of the full source in the reconstructed vector gives
the same four terms. Differentiating $F(E(h),h)u(h)=0$ and
pairing with $u$ yields the energy derivative.
The last sign follows from the positive complementary pivot.
\end{proof}
This is the full many-body Hamiltonian source, with both mixed columns,
the eliminated-sector insertion and residue norm retained. The
positive return below the complement lowers the direct Hamiltonian
block, but does not prove a macroscopic scalar response. An indefinite
Dirac complement has no such universal sign.
Source-dependent projectors and frames require their complete jets.
Regular composition reuses Corollary~\ref{p1b:cor:complete-successive-reductions}
and Corollary~\ref{p3f:cor:complete-reduction-reuse}.
A spectral identity alone introduces no Gaussian quantum integration.


\begin{corollary}[Same-channel positive-gap pole-weight certificate]
\label{p10:cor:schur-positive-gap-weight}
On the fixed canonical Hamiltonian split of
Theorem~\ref{p10:thm:complete-schur-scalar-source}, suppose
$0<\delta I\preceq C-EI\preceq DI$ and $u^\dagger u=1$.
With $s=u^\dagger X(C-EI)^{-1}Yu$, the direct-block Rayleigh lowering
and retained population satisfy
\[
 s=u^\dagger Au-E,\qquad
 w_{\rm ret}=\frac1{u^\dagger\mathscr Zu},\qquad
 \boxed{\frac{\delta}{\delta+s}\le w_{\rm ret}
                         \le\frac{D}{D+s}.}
\]
At a simple pole its retained resolvent residue is
$w_{\rm ret}uu^\dagger$.  At a degenerate pole the nullspace matrix
formula of Theorem~\ref{p1b:thm:response-positive-gap-weight}, rather than
one chosen rank-one lift, gives the full residue.
\end{corollary}
\begin{proof}
The inherited norm is $I+X(C-EI)^{-2}Y$.  Spectral calculus gives
$(C-EI)^{-1}/D\preceq(C-EI)^{-2}
\preceq(C-EI)^{-1}/\delta$; contract with $Yu$ and invert the scalar
positive norm.  The same nullspace calculation gives the stated matrix
extension.
\end{proof}
This certificate travels with the actual split, gap and source packet
into Paper~VII-B.  It ties a same-sector lowering to its residue norm;
it supplies neither an absolute QCD energy nor a current overlap.
The complete scalar insertion still contains all four source terms.
An indefinite Dirac pivot or a continued resonance has no such universal
positive-gap estimate.

\begin{corollary}[Determinant and scalar insertion in a fixed Dirac split]
\label{p10:cor:complete-dirac-schur-source}
On the normalized Cayley support let $K=-iD_c$ be mass-independent
Hermitian and $P+Q=I$ fixed. Put $X=PKQ$, $Y=QKP$,
$E_Q=mQ+iQKQ$ and
$S_E=mP+iPKP+XE_Q^{-1}Y$. For $m>0$,
\[
 \det(mI+iK)=\det E_Q\det S_E,\qquad
 \partial_mS_E=P-XE_Q^{-2}Y,
\]
\[
 \boxed{\operatorname{Tr}(mI+iK)^{-1}
 =\operatorname{Tr}E_Q^{-1}
  +\operatorname{Tr}[S_E^{-1}(P-XE_Q^{-2}Y)].}
\]
\end{corollary}
\begin{proof}
Both diagonal blocks are accretive for $m>0$. Block elimination
uses off-diagonal factors $iX,iY$ and gives the plus Schur sign.
Differentiate the determinant factorization to obtain the trace.
\end{proof}
Keep the complementary determinant in the physical weighting.
The overlap factor $R_c=(I-V)/2$ is mass-independent at fixed
configuration and support but can depend on that configuration,
so its determinant also belongs to the weight.
Source-dependent kinetic normalization, pivots or support require
their additional derivatives and the complete sea response of
Theorem~\ref{p12:thm:v2123-complete-finite-bh-scalar-response}.

% END INLINE SOURCE: Paper7A-CompleteSchurSource-Main-v1_0.tex



% Source Paper VII-A, original line 1661; inherited proof bodies unchanged.
\subsubsection{Paper-VII-B source-interface handoff}
\label{p10:subsec:paper-xi-source-interface-handoff}

\begin{theorem}[Executable Paper-VII-A to Paper-VII-B handoff]
\label{p10:thm:source-interface-handoff}
The finite package
\begin{equation}
\label{p10:eq:paper-xi-source-interface-package}
\left(
\mathcal H_{L,\rm phys}^{\rm QCD},
H_L,
K_L,
\mathcal O_L,
\{E_{n,L},\Psi_{n,L}\},
W_H[J]
\right)
\end{equation}
fixes every downstream execution that is already typed as a singlet spectral
value, source derivative, finite matrix element, transition spectral measure,
Schur/Feshbach pole, finite trace, or multi-baryon transfer cumulant on this
carrier.  A light-front distribution additionally requires the null-action
representation certificate in
Definition~\ref{p10:def:finite-null-representation-certificate}; the six-tuple
does not acquire that cross-owner intertwiner merely because it contains a
Hamiltonian.  Paper~VII-B performs the licensed evaluations, and no result of an
evaluation may alter the package
\eqref{p10:eq:paper-xi-source-interface-package}.
\end{theorem}
\begin{proof}
The singlet basis and sparse-entry theorem fix the physical matrices, while the
source dictionary fixes all operator insertions and their normalizations.
Finite-dimensional spectral theory, source differentiation, Schur extraction,
and traces are deterministic operations on these data.  By contrast, relating
the independently typed clock--Hopf null flow, good-field carrier, QCD
Hamiltonian, Gauss projection, and color holonomy requires an explicit
intertwining diagram.  Paper~II-A's no-insertion clause forbids deriving that
diagram from simultaneous presence alone.  Reversing either dependency would
violate the readout-chronology theorem.
\end{proof}


\begin{theorem}[Derived elapsed null flow and physical null generator]
\label{p10:thm:v252-derived-null-generator-identity}
Let \(\mathscr F_L^{\rm LF}\) be a finite good-field/QCD carrier and let
\(\mathscr A_{0,L}^{\rm LF}\) be a finite equal-time test space.  Denote the
normalized spatial Hopf flow by \(\Phi_H(r)\).  Suppose there is a linear
same-carrier map
\[
 q_{0,L}:\mathscr A_{0,L}^{\rm LF}
 \longrightarrow \End(\mathscr F_L^{\rm LF})
\]
which is Hopf equivariant,
\begin{equation}
\label{p10:eq:v252-equal-time-hopf-seed}
 e^{-irP_{H,L}/\hbar}q_{0,L}(f)e^{irP_{H,L}/\hbar}
 =q_{0,L}(f\circ\Phi_H(-r)),
\end{equation}
where \(H_L=H_L^\dagger\), \(P_{H,L}=P_{H,L}^\dagger\), and
\([H_L,P_{H,L}]=0\) on this carrier.  Put
\[
 \mathscr A_L^{\rm ray}:=C_c^1(\RR;\mathscr A_{0,L}^{\rm LF})
\]
and, for \(F\in\mathscr A_L^{\rm ray}\), define the elapsed-time map
\begin{equation}
\label{p10:eq:v252-elapsed-test-map}
 Q_L(F):=\int_{\RR}dt\,
 e^{-itH_L/\hbar}q_{0,L}(F_t)e^{itH_L/\hbar}.
\end{equation}
Then the one-parameter group
\begin{equation}
\label{p10:eq:v252-derived-null-unitary}
 U_{+,L}(s):=
 \exp\!\left[-\frac{is}{\sqrt2\hbar}
 \left(\frac{H_L}{c_{\rm op}}+P_{H,L}\right)\right]
\end{equation}
implements the derived clock--Hopf null flow
\[
 N_+=\frac1{\sqrt2}\left(c_{\rm op}^{-1}u_t+e_H\right)
\]
on the retained test domain:
\begin{equation}
\label{p10:eq:v252-derived-null-equivariance}
 U_{+,L}(s)Q_L(F)U_{+,L}(s)^*
 =Q_L\!\left(F\circ\operatorname{Fl}^{N_+}_{-s}\right).
\end{equation}
Consequently its unique finite self-adjoint generator is
\begin{equation}
\label{p10:eq:v252-physical-null-generator-identity}
 \boxed{
 P_{N,L}^+=P_{\rm phys,L}^+
 =\frac1{\sqrt2}\left(\frac{H_L}{c_{\rm op}}+P_{H,L}\right)}.
\end{equation}
No separate cross-owner equality between a pre-existing geometric null
generator and a dynamical positive-polar generator is required.  The theorem
uses only the globally defined one-parameter null-ray flow and therefore does
not require a Frobenius-integrable null hypersurface.
\end{theorem}
\begin{proof}
Strong commutation gives, with
\(a=s/(\sqrt2c_{\rm op})\) and \(b=s/\sqrt2\),
\[
 U_{+,L}(s)=e^{-iaH_L/\hbar}e^{-ibP_{H,L}/\hbar}.
\]
Hence, using \eqref{p10:eq:v252-equal-time-hopf-seed},
\begin{align*}
 U_{+,L}(s)Q_L(F)U_{+,L}(s)^*
 &=\int dt\,
 e^{-i(t+a)H_L/\hbar}
 q_{0,L}(F_t\circ\Phi_H(-b))
 e^{i(t+a)H_L/\hbar}.
\end{align*}
Set \(u=t+a\).  The integrand becomes
\(q_{0,L}(F_{u-a}\circ\Phi_H(-b))\), and by the defining clock--Hopf
flow this is exactly the \(u\)-slice of
\(F\circ\operatorname{Fl}^{N_+}_{-s}\).  This proves
\eqref{p10:eq:v252-derived-null-equivariance}.  Finite-dimensional unitary
generator uniqueness then gives
\eqref{p10:eq:v252-physical-null-generator-identity}.  Equality of the
unitary groups, rather than only of their adjoint actions, also excludes an
additive central-generator ambiguity.
\end{proof}

\begin{corollary}[Infinitesimal null covariance and total conservation]
\label{p10:cor:v252-null-generator-conservation}
Under the hypotheses of
Theorem~\ref{p10:thm:v252-derived-null-generator-identity},
\begin{equation}
\label{p10:eq:v252-infinitesimal-null-covariance}
 [P_{\rm phys,L}^+,Q_L(F)]
 =-i\hbar Q_L(\mathcal L_{N_+}F)
\end{equation}
and
\begin{equation}
\label{p10:eq:v252-total-null-conservation}
 \boxed{[H_L,P_{\rm phys,L}^+]=0}.
\end{equation}
If in addition \(H_L/c_{\rm op}\ge |P_{H,L}|\), then
\(P_{\rm phys,L}^+\ge0\).
\end{corollary}
\begin{proof}
Differentiate
\eqref{p10:eq:v252-derived-null-equivariance} at \(s=0\).  The conservation
law follows from
\([H_L,P_{H,L}]=0\) and
\eqref{p10:eq:v252-physical-null-generator-identity}.  The final statement is
the positive-polar inequality.
\end{proof}

\begin{definition}[Finite QCD representation of the derived null-ray diagram]
\label{p10:def:finite-null-representation-certificate}
At level \(L\), a certificate \(\LFRep_L\) is one compatible proof of the
following clauses.
\begin{enumerate}[label=\textnormal{(R\arabic*)},leftmargin=2.8em]
\item \textbf{Native equal-time seed and derived retained-domain null equivariance.}
There are a finite good-field carrier \(\mathfrak h_L^{\rm LF}\), its physical
occupation carrier \(\mathscr F_L^{\rm LF}\), an equal-time retained test
space \(\mathscr A_{0,L}^{\rm LF}\), and a same-carrier map \(q_{0,L}\)
satisfying the Hopf-equivariance condition
\eqref{p10:eq:v252-equal-time-hopf-seed}.  The elapsed test domain and map are
then
\[
 \mathscr A_L^{\rm ray}=C_c^1(\RR;\mathscr A_{0,L}^{\rm LF}),
 \qquad
 Q_L(F)=\int dt\,e^{-itH_L/\hbar}q_{0,L}(F_t)e^{itH_L/\hbar}.
\]
With the same \(H_L\) and \(P_{H,L}\) appearing in item~\textnormal{(R2)},
Theorem~\ref{p10:thm:v252-derived-null-generator-identity} constructs
\(U_{+,L}(s)\) and proves
\[
 U_{+,L}(s)Q_L(F)U_{+,L}(s)^*
 =Q_L\!\left(F\circ\operatorname{Fl}^{N_+}_{-s}\right).
\]
Thus the R1 elapsed-null generator is not an independently supplied operator:
\[
 P_{N,L}^+=P_{\rm phys,L}^+.
\]
No covariance on an undeclared full continuum function algebra is inferred
from this retained finite test map.

\item \textbf{Same-carrier positive-polar ancestry and normalization.}
The physical Hamiltonian \(H_L\) and a self-adjoint Hopf generator
\(P_{H,L}\) act on the carrier of item~\textnormal{(R1)}, strongly commute,
the latter implements the normalized spatial Hopf action and satisfies
\[
 e^{-i(2\pi R_x)P_{H,L}/\hbar}=I,
\]
and their joint spectrum obeys
\[
 H_L/c_{\rm op}\ge |P_{H,L}|.
\]
Define
\begin{equation}
\label{p10:eq:finite-null-generator-pair}
 P_{\rm phys,L}^{\pm}:=\frac1{\sqrt2}
 \left(H_L/c_{\rm op}\pm P_{H,L}\right)\ge0 .
\end{equation}
Theorem~\ref{p10:thm:v252-derived-null-generator-identity} proves that the
state-oriented R1 generator is automatically the plus member,
\[
 P_{N,L}^+=P_{\rm phys,L}^+,
\]
so there is no additional ``same-\(P^+\)'' identification premise between R1
and R2.  Strong commutation gives
\begin{equation}
\label{p10:eq:v133-total-null-generator-conservation}
 [H_L,P_{\rm phys,L}^{\pm}]=0,
 \qquad
 e^{itH_L/\hbar}P_{\rm phys,L}^{\pm}e^{-itH_L/\hbar}
 =P_{\rm phys,L}^{\pm}.
\end{equation}
Integer weights belong to the compact Hopf generator; they do not force the
full elapsed null generator to have integer spectrum.

\item \textbf{Species resolution and physical/Fock compatibility.}
There are orthogonal species projections
\(\pi_{i,L}\) on
\(\mathfrak h_L^{\rm LF}=\bigoplus_{i\in\mathcal I}
\mathfrak h_{i,L}^{\rm LF}\), with quarks, antiquarks, and gluons distinct,
and a positive one-body generator \(p_L^+\) such that
\[
 [p_L^+,\pi_{i,L}]=0,
 \qquad \sum_i\pi_{i,L}=I,
 \qquad
 P_{i,L}^+=\mathrm d\Gamma(p_L^+\pi_{i,L}).
\]
Define the additive finite occupation generator
\begin{equation}
\label{p10:eq:v252-fock-plus-generator}
 P_{F,L}^+:=\sum_iP_{i,L}^+=\mathrm d\Gamma(p_L^+).
\end{equation}
The exact finite species clause is the nontrivial compatibility condition
\begin{equation}
\label{p10:eq:v252-r3-fock-physical-equality}
 \boxed{P_{F,L}^+=P_{\rm phys,L}^+}.
\end{equation}
This equality is not a consequence of R1--R2 and is not hidden by notation.
Total conservation of \(P_{\rm phys,L}^+\) is already derived, but it neither
constructs the species projectors or \(p_L^+\) nor implies
\([H_L,P_{i,L}^+]=0\).  In particular, interacting quark, antiquark, and gluon
contributions need not be separately conserved, and positivity of interacting
or scheme-renormalized species contributions is not inferred from positivity
of their sum.

\item \textbf{Gauss and zero-mode compatibility.}
The physical projector \(G_L\) commutes with \(P_L^+\), every
\(P_{i,L}^+\), and every species-resolved spectral counting operator used by
the readout.  The space \(\ker p_L^+\), when nonzero, is retained as explicit
\(x=0\) constraint data.

\item \textbf{Color-path compatibility.}
For every actually selected basepoint/aligned ray used by the claim, the QCD
connection and good-field realization have a stated common base or refinement
with the derived path \(\gamma_{+,z}\) of
Paper~III-E, Eq.~\eqref{p3e:eq:clock-hopf-null-path-family}.  The transporter
is
\begin{equation}
\label{p10:eq:derived-path-color-holonomy}
 W_{+,L}(s,0)
 :=\operatorname{Hol}_{A_{c,L}}
       (\gamma_{+,z}|_{[0,s]}),
\end{equation}
and its action on the finite fields intertwines the gauge-invariant bilocal
with the occupation spectral operators.  The path is derived; the color
configuration, common-base map, gauge covariance, and intertwiner are proved
in this clause rather than identified with the Hopf \(U(1)\) connection.

\item \textbf{Native-carrier comparison.}
If the continuous action is first realized on the position-side Peter--Weyl
carrier (or on its already generated \(\iota_1\)-transport), a stated
sampling/compression map to the native momentum-side Hurwitz QCD carrier
satisfies the exact or controlled approximate intertwining estimate used by
the claim.  Neither
Theorem~\ref{p10:thm:connected-flow-finite-shell-no-go} nor the existence of a
harmonic representation supplies this map.
\end{enumerate}
The certificate contains no hadron-state selection, PDF shape, fitted
normalization, light-ray renormalization, TMD soft factor, scattering state,
or continuum-limit assertion.
\end{definition}

\begin{theorem}[Internal total-null conservation and residual light-front frontier]
\label{p10:thm:v133-internal-total-null-conservation}
If clauses~\textnormal{(R1)--(R2)} of
Definition~\ref{p10:def:finite-null-representation-certificate} hold, then the
represented elapsed-null generator is not an additional cross-owner datum but
is the derived operator
\[
 \boxed{P_{N,L}^+=P_{\rm phys,L}^+
 =\frac1{\sqrt2}\left(H_L/c_{\rm op}+P_{H,L}\right)},
 \qquad
 \boxed{[H_L,P_{\rm phys,L}^+]=0.}
\]
Consequently no independent ``R1--R2 same-\(P^+\)'', ``Noether'',
``total-null-momentum conservation'', or finite-Stone premise remains in
\(\LFRep_L\).  This result does not construct the equal-time native seed of
clause~\textnormal{(R1)}, does not prove the positive-polar inequality of
clause~\textnormal{(R2)}, does not prove the physical/Fock equality of
clause~\textnormal{(R3)}, and does not supply clauses~\textnormal{(R4)--(R6)}.
\end{theorem}
\begin{proof}
Theorem~\ref{p10:thm:v252-derived-null-generator-identity} constructs the R1
unitary group directly from the same \(H_L\) and \(P_{H,L}\) and proves its
geometric null covariance.  Finite generator uniqueness yields the first
boxed equality.  Corollary~\ref{p10:cor:v252-null-generator-conservation}
gives the second.  The remaining statements are type firewalls: neither
these identities nor conservation of the total generator construct the
one-body species realization or the other representation intertwiners.
\end{proof}

\begin{theorem}[Status of the finite null-representation handoff]
\label{p10:thm:finite-null-representation-status}
The current finite compiler exports \(\LFRep_L\) only on a subdomain on which
all six clauses of
Definition~\ref{p10:def:finite-null-representation-certificate} have been
proved for the same carrier and operators.  Theorem~\ref{p10:thm:v133-internal-total-null-conservation}
removes the former R1--R2 generator-identification and separate
total-generator conservation/Noether subpremises.  A generic unitary group,
the 24-cell incidence data, an invariant harmonic carrier, or co-realization
of the geometric and QCD roles separately still does not construct the native
equal-time good-field seed, prove the positive-polar inequality on that
carrier, prove the physical/Fock equality of clause~\textnormal{(R3)}, or
discharge the Gauss/zero-mode, color-path, or native-carrier comparison
clauses.  Until those typed constructions are executed,
\(\LFRep_L\) remains an explicit theorem obligation for any claim that asks
for the stronger exact finite species/Fock simplex or POVM realization; it is
not a completed Paper-VII-A output and is not required by the source-owned
current/twist-two hierarchy.
\end{theorem}

\begin{proof}
Each remaining clause asserts a typed relation between objects with different
immediate owners.  The universal rooted co-realization theorem retains an
already licensed relation but forbids insertion of a missing cross-owner edge.
The finite-shell no-go excludes the point-permutation shortcut, and
Corollary~\ref{p10:cor:harmonic-null-action-shell-boundary} supplies only a
candidate kinematic action until its QCD equal-time seed/intertwiner is proved.
Once that seed and the same-carrier \(H_L,P_{H,L}\) data exist,
Theorem~\ref{p10:thm:v252-derived-null-generator-identity} constructs the
elapsed null action and removes the former generator-equality edge.  The
species/Fock equality and clauses~\textnormal{(R4)--(R6)} remain distinct
relations and cannot be inferred from simultaneous occurrence.
\end{proof}

\begin{corollary}[Scope of the exact finite LFRep certificate after the source-owner split]
\label{p10:cor:v264-lfrep-optional-finite-readout-scope}
The unresolved clauses of \(\LFRep_L\) remain a legitimate exact finite representation problem, but they are consumed only by claims that require the finite species/Fock simplex, its positive operator measures, or an exact occupation-space identification with the physical null generator.  Physical light-ray/twist-two source operators, the current normalization at \(n=1\), and the forward stress-tensor normalization at \(n=2\) are already owned by the QCD connection/action/source family and therefore do not wait on clause~\textnormal{(R3)}.  Any continuum PDF identification still requires its own light-ray/refinement/renormalization hypotheses.
\end{corollary}
\begin{proof}
The first statement is exactly the typing of clauses~\textnormal{(R3)--(R6)} in Definition~\ref{p10:def:finite-null-representation-certificate}.  The second is Theorem~\ref{p10:thm:v264-twist-two-current-stress-owner}.  The last statement is a limit/renormalization firewall: a finite source operator and an exact finite occupation representation are both distinct from a renormalized nonlocal continuum distribution.
\end{proof}

\begin{proposition}[Executable native-QCD protocol]
\label{p10:prop:executable-native-qcd-protocol}
A reproducible implementation records the six compiler rows in
Table~\ref{p10:tab:executable-finite-qcd-transcript-schema} before comparison.
No row may be chosen or altered to improve agreement with external data.
\end{proposition}
\begin{proof}
The six rows reconstruct the physical finite matrix problem, its normalized
source algebra, and every permitted extraction rule.  Comparison is therefore
strictly downstream.
\end{proof}

\begin{longtable}{>{\raggedright\arraybackslash}p{0.13\textwidth} >{\raggedright\arraybackslash}p{0.29\textwidth} >{\raggedright\arraybackslash}p{0.46\textwidth}}
\caption[Paper-VII-A compiler audit interface]{Paper~VII-A compiler audit interface.}\label{p10:tab:executable-finite-qcd-transcript-schema}\\
\toprule
Row & Object & Required native record \\
\midrule
\endfirsthead
\toprule
Row & Object & Required native record \\
\midrule
\endhead
R1 & basis and blocks & ordered raw basis, singlet basis, exact block labels, and dimensions \\
R2 & projectors & \(P_{\rm sing}^{(L)}\) and inherited exact charge, spin, flavor, and topology projectors \\
R3 & Hamiltonian and transfer & sparse entries in \(\eta NSCF\kappa\) form, \(H_L\), and \(K_L\) \\
R4 & source dictionary & finite source operators and their sparse entry records \\
R5 & spectral data & eigenvalues, eigenvectors, vacuum, and exact block quantum numbers \\
R6 & extraction rules & residues, matrix elements, source derivatives, susceptibilities, and Schur-pole normalization \\
\bottomrule
\end{longtable}
\TableNotes{Hadron, distribution, fragmentation, thermal, and nuclear rows are Paper-VII-B executions of this interface, not additional Paper-VII-A kernel rows.}



% BEGIN FRAMEWORK v2.323: PETERSEN--HOPF FINITE STRONG SOURCE INTERFACE
\section{Petersen--Hopf coherence criteria for the finite strong source}
\label{p10:sec:v323-petersen-strong-source}
\noindent\textit{Status: exact conditional-geometry consequences and explicit
same-action obligations; no new QCD interaction, pole, or scale is selected.}
The Petersen--Hopf packet is generated by an \emph{occurring embedded}
$2I$ and is not assumed for a bare $2T$ Hurwitz shell.  Its oriented
fifteen-edge space $W$ is not the $24$-vertex Hurwitz cell complex
$\mathcal K_H$, the level-one momentum fifteen-carrier, or the QCD color
triplet.  Hence no edge of the Petersen graph is silently promoted to a
fundamental gauge link or a constituent quark.

\begin{definition}[Mirror contrast and core occupation projectors]
\label{p10:def:v323-petersen-strong-contrast}
Write $W=\bigoplus_{m=1}^5 V_m$, where each $V_m\cong\mathbb R^3$
is the oriented Clifford-edge triple omitting core $m$.  Let $P_m$ be
its orthogonal projector and let $\Pi_3,\Pi_{3'}$ be the two rank-three
icosahedral cycle projectors derived in Paper~I-A. Define
\[
 D:=\Pi_3-\Pi_{3'},\qquad
 P_S:=\sum_{m\in S}P_m,\qquad D_S:=P_SDP_S.
\]
The sign of $D$ records which Galois-conjugate representation is called
ordinary; its physical interpretation is not fixed by this convention.
\end{definition}

\begin{theorem}[Rooted no-go and first coherent two-core detection]
\label{p10:thm:v323-first-coherent-mirror-source}
Every single core satisfies $P_mDP_m=0$.  For any two distinct cores,
$P_mDP_n$ has rank two and singular values
$\sqrt{3/20},\sqrt{3/20},0$. Therefore every state on $W$ that is
block-incoherent in the five-core decomposition,
$\rho=\sum_mP_m\rho P_m$, obeys
\[
 \operatorname{Tr}(\rho D)=0.
\]
A coherent two-core state can instead have nonzero contrast; on any
$V_m\oplus V_n$ the largest possible absolute expectation of $D$ in
a normalized pure state is $\sqrt{15}/10$.
\end{theorem}
\begin{proof}
The exact ten-axis matrices of Paper~I-A give
$P_m\Pi_3P_m=P_m\Pi_{3'}P_m=P_m/5$.  Their difference has zero
single-core blocks.  The off-diagonal $3\times3$ blocks have
$XX^{\mathsf T}$-spectrum $\{(3/20)^{(2)},0\}$, as certified
independently in Paper~VII-D.  Cyclicity of trace gives the
block-incoherent statement.  The two-core compression is
$\left(\begin{smallmatrix}0&X\\X^{\mathsf T}&0\end{smallmatrix}\right)$,
whose extreme eigenvalues are $\pm\sqrt{3/20}=\pm\sqrt{15}/10$.
\end{proof}

\begin{corollary}[When strong evolution cannot create mirror contrast]
\label{p10:cor:v323-incoherent-evolution}
If a proposed self-adjoint operator on $W$ commutes with each $P_m$,
and the initial density matrix is block-incoherent, then its unitary
evolution preserves $\operatorname{Tr}(\rho(t)D)=0$ for all times.
Consequently a nonzero contrast from such an initial state requires
an inter-core transition in the \emph{actual} evolution and a suitable
source; a nonzero geometric $D$ alone is insufficient.
\end{corollary}
\begin{proof}
The commuting Hamiltonian preserves each $V_m$, so conjugation of
$\rho$ preserves block diagonality. Apply the preceding theorem at
every time.
\end{proof}

\begin{definition}[Typed Petersen strong-source incidence test]
\label{p10:def:v323-action-incidence-test}
For the compiled Paper--VII-A Hamiltonian $H_L$ and singlet space
$\mathcal H_{L,\mathrm{sing}}$, a claimed physical Petersen term must
supply an explicit color-gauge-compatible finite map
$\iota_L:W_{\mathbb C}\rightarrow\mathcal H_{L,\mathrm{sing}}$ (or its
source-owned submodule), together with an actual source $\mathcal O_L$.
For an isometric embedding, the candidate compression is
$H_{W,L}=\iota_L^{\dagger}H_L\iota_L$.  An autonomous spectrum on $W$
requires, in addition, the reducing-space certificate
\[
 (I-\iota_L\iota_L^\dagger)H_L\iota_L=0;
\]
otherwise pole calculations must use the full $H_L$ or its exact Schur
resolvent, not the eigenvalues of $H_{W,L}$.  A nonzero pair-to-pair
matrix element $P_mH_{W,L}P_n$ is necessary for producing coherence
from a core-block-diagonal initial state.  These data are \emph{tests},
not newly asserted output of the current QCD compiler.
\end{definition}

\begin{proposition}[Color-center and three-face-fan separation]
\label{p10:prop:v323-center-firewall}
The Petersen signed-axis holonomy takes values in $\mathbb Z_2$;
fundamental $SU(3)_c$ center flux takes values in $\mathbb Z_3$.
There is no nontrivial group homomorphism
$\mathbb Z_2\rightarrow\mathbb Z_3$.
In particular, the Petersen $\pm1$ loop cannot be identified with the
$SU(3)$ baryon-junction center phase or the Wilson-loop area-law source.
The three Hurwitz faces incident to a $24$-cell edge in the established
D21 baryon construction and the three Petersen edges incident to a
Hopf-axis vertex are distinct incidence systems; equality of valency
supplies no gauge-covariant map between them.
\end{proposition}
\begin{proof}
Any image of the generator of $\mathbb Z_2$ must have order dividing
both $2$ and $3$, hence is the identity.  The actual Hurwitz three-face
fan is defined by the existing $SU(3)$ link matrices on $\mathcal K_H$;
the Petersen graph carries no such link assignment in the compiled
datum.  Therefore its holonomy cannot replace that sourced matrix word.
\end{proof}

% BEGIN FRAMEWORK v2.327: MIRROR TRANSPORT / SPIN TEST ON THE NATIVE QCD OWNER
\subsection{Two-core mirror eligibility, source leakage and conditional spin blocks}
\label{p10:sec:v327-mirror-spin-interface}
\noindent\textit{Status: exact finite operator identities and an explicitly
conditional candidate; no meson identification, interaction coefficient,
or physical pole is produced by the Petersen projection alone.}

\begin{theorem}[Exact compression--leakage Pythagorean identity]
\label{p10:thm:v327-compression-leakage-pythagoras}
Let $K=H_L-E_{0,L}I$ be the actual self-adjoint finite singlet
Hamiltonian gap operator and let $\iota:\mathcal V\to
\mathcal H_{L,\mathrm{sing}}$ be any isometry. Put
$P=\iota\iota^\dagger$, $K_{\rm eff}=\iota^\dagger K\iota$
and $B=(I-P)K\iota$. Then, exactly,
\[
 \boxed{\iota^\dagger K^2\iota=K_{\rm eff}^{\,2}+B^\dagger B},
 \qquad
 \|K\iota v\|^2=\|K_{\rm eff}v\|^2+\|Bv\|^2.
\]
The correction is positive semidefinite and vanishes precisely when
$\operatorname{ran}\iota$ is invariant under $K$ (and therefore
reducing, since $K$ is self-adjoint).
The compressed eigenvalues are physical eigenvalues of that sector
only after this reducing condition, or when the appropriate full
Schur-resolvent pole/residue calculation has been performed.
\end{theorem}
\begin{proof}
Insert $I=P+(I-P)$ between the two factors of $K$.
The $P$ term is $(\iota^\dagger K\iota)^2$ and the
complement term equals $B^\dagger B$.
The norm identity is ordinary orthogonality of the two projected
components. The vanishing criterion follows from $B^\dagger B=0$
if and only if $B=0$; self-adjointness makes an invariant subspace
reducing.
\end{proof}

\begin{definition}[Optional two-core spin source test]
\label{p10:def:v327-conditional-pair-spin-block}
Only if the actual source-owned $\mathcal V$ is shown to contain the
specified two-core factor $W_{mn,\mathbb C}$ and an actual
spin-$\tfrac12\otimes\tfrac12$ factor $\mathbb C^4$,
consider the \emph{candidate}, not a compiler output,
\[
 \mathcal T=\sum_{a=1}^3\sigma_a\otimes\sigma_a,
 \qquad \operatorname{Spec}(\mathcal T)=\{-3^{(1)},+1^{(3)}\},
\]
\[
 \mathscr M_{mn,f}^{\,2}
 =a_f I_{24}+b_f (U_{mn}\otimes I_4)
 +\kappa_f(Q_{mn}\otimes\mathcal T).
\]
The real coefficients $a_f,b_f,\kappa_f$ are \emph{not} supplied by
support counts, the Petersen spectrum or $E_{32},E_{36},E_{40}$.
The index $f$ denotes a yet-to-be-constructed physical flavor/source
sector, not an assignment of $m,n$ to constituent flavors.
The identification of this $24$-dimensional \emph{tensor product}
with the Hurwitz $24$-cell or with a color-singlet block is not assumed.
\end{definition}

\begin{theorem}[Conditional 24-dimensional mirror--spin spectrum]
\label{p10:thm:v327-conditional-mirror-spin-spectrum}
For the abstract block in the preceding definition, the exact
squared-mass-operator spectrum is
\[
 \begin{array}{c|c}
 \text{formal eigenvalue}&\text{multiplicity}\\\hline
 a_f&8\\
 a_f+b_f-3\kappa_f&2\\
 a_f-b_f-3\kappa_f&2\\
 a_f+b_f+\kappa_f&6\\
 a_f-b_f+\kappa_f&6
 \end{array}
\]
(including the usual merging of degeneracies when coefficients vanish).
In each fixed mirror parity $\pm$, the formal triplet/singlet gap is
$4\kappa_f$ and the spin-weighted average is
$a_f\pm b_f$.
Its trace and Hilbert--Schmidt second moment satisfy the exact
orthogonal decomposition
\[
 \operatorname{Tr}\mathscr M_{mn,f}^{\,2}=24a_f,
 \quad \operatorname{Tr}\!\left[(\mathscr M_{mn,f}^{\,2}
 -a_fI)^2\right]=16b_f^2+48\kappa_f^2.
\]
An arbitrary positive physical squared-mass interpretation also
requires the resulting operator to be positive and, crucially,
requires a verified reducing or source-resolvent embedding into the
actual finite QCD kernel.
\end{theorem}
\begin{proof}
On $\operatorname{ran}Q_{mn}$, $U_{mn}$ has $+1,-1$ each twice,
and the dark two-space has $U_{mn}=Q_{mn}=0$.
The singlet/triplet spin operator has eigenvalues $-3,+1$ with
multiplicities $1,3$. The displayed eigenvalues follow by tensor
products. Trace cancellation uses
$\operatorname{Tr}U_{mn}=\operatorname{Tr}\mathcal T=0$.
The two squared components are Hilbert--Schmidt orthogonal and have
$\operatorname{Tr}(U_{mn}^2\otimes I_4)=16$ and
$\operatorname{Tr}(Q_{mn}^2\otimes\mathcal T^2)=4(9+3)=48$.
\end{proof}

\begin{proposition}[Spin-blind no-go and exact source obligations]
\label{p10:prop:v327-spinblind-no-go}
Neither $D_{mn}\otimes I_4$ nor $Q_{mn}\otimes I_4$
can lift the singlet/triplet degeneracy: they are spin blind.
$Q_{mn}$ is also core block diagonal, so a Hamiltonian using only
$Q_{mn}$ cannot create core coherence from a block-diagonal state.
For a physical meson claim one must compute, from the \emph{same}
$H_L$ and source/vacuum packet: (i) the color-gauge-compatible
embedding, (ii) its leakage $B$ or the full Schur resolvent,
(iii) the actual spin and source projectors, (iv) all coefficients of
the projected operator, and (v) source residues and pole stability.
No equality of $\kappa_f$ across flavors follows from these claims.
\end{proposition}
\begin{proof}
Both spin-blind operators commute with every operator on
$\mathbb C^4$, and $Q_{mn}$ commutes with the pair core grading.
The remaining requirements are the typed spectral-owner test of
Definition~\ref{p10:def:v323-action-incidence-test} and the
preceding compression identity.
\end{proof}
% END FRAMEWORK v2.327: MIRROR TRANSPORT / SPIN TEST


% END FRAMEWORK v2.323: PETERSEN--HOPF FINITE STRONG SOURCE INTERFACE


\section{Gauge action, heat normalization, sources, and finite Ward response}
\label{p10:sec:v237-a3}
Common-heat normalization belongs to its actual heat datum. Heat time, transfer time, physical scale and dynamical evolution are separately typed; finite source derivatives and multiface interactions use the same action owner.

% Source Paper VII-A, original line 2044; inherited proof bodies unchanged.
\subsection{Color connection space and Yang--Mills normalization}
\label{p10:sec:connection-space}


% Source Paper VII-A, original line 2047; inherited proof bodies unchanged.
\subsubsection{Connection ledger and admissible color geometry}
\label{p10:subsec:connection-ledger-admissible-geometry}


\begin{table}[htbp]
\centering
\footnotesize
\begin{tabular}{>{\raggedright\arraybackslash}p{0.34\textwidth} >{\raggedright\arraybackslash}p{0.56\textwidth}}
\toprule
Claim & Discharged by / status \\
\midrule
Finite color connection space & defined from the imported \(SU(3)_c\) color carrier and the finite singlet support \\
Admissible Yang--Mills curvature channel & selected by the finite adjoint connection and no-branch closure \\
Gauge normalization & source-native branch: \(k_i g_i^2=4\pi^2\) at the common-heat matching datum; named external-scale coupling obtained by RG/scale transport.  The historical \(\Hseedzero\) member remains only a broad-branch comparison chart. \\
No external Yang--Mills coupling insertion & excluded because an inserted normalization would add a new strong-sector coefficient \\
Finite pure-glue datum prepared & exported to the transfer-kernel, gap, and confinement modules below \\
\bottomrule
\end{tabular}
\caption[Color connection and Yang--Mills normalization]{Color connection and Yang--Mills normalization.}
\label{p10:tab:connection-ym-ledger}
\TableNotes{This ledger turns the section into a table-led derivation block.  The following paragraphs and theorems discharge the rows in order before the strong clock and mass-gap modules use the resulting finite pure-glue datum.}
\end{table}

This section is one of the closed parts of the QCD construction.  The color
carrier, admissible traceless connection \emph{space}, and invariant local
Yang--Mills kinetic form are fixed.  A particular field configuration
\(A_c\), its state weight, and its holonomy on a selected path are not uniquely
selected by that structural result.

\begin{definition}[Admissible color connection]
\label{p10:def:admissible-color-connection}
Let \(C\cong\CC^3\) be the finite color triplet.  The admissible color connection is
\begin{equation}
\label{p10:eq:color-connection-space}
A_c\in\mathcal A_c:=\Omega^1(M,\mathfrak{su}(C))\cong\Omega^1(M,\mathfrak{su}(3)).
\end{equation}
We use Hermitian generators \(T_a\) with
\begin{equation}
\label{p10:eq:hermitian-generator-normalization}
\operatorname{Tr}(T_aT_b)=\frac12\delta_{ab},
\qquad
[T_a,T_b]=i f_{abc}T_c.
\end{equation}
The quark covariant derivative and curvature are
\begin{equation}
\label{p10:eq:qcd-covariant-derivative}
D_\mu q=(\partial_\mu-i g_s A_\mu^aT_a)q,
\end{equation}
\begin{equation}
\label{p10:eq:qcd-field-strength}
[D_\mu,D_\nu]q=-i g_s F_{\mu\nu}^aT_aq,
\qquad
F_{\mu\nu}^a=
\partial_\mu A_\nu^a-\partial_\nu A_\mu^a+g_s f^{abc}A_\mu^bA_\nu^c.
\end{equation}
\end{definition}

\begin{theorem}[Color connection and kinetic uniqueness]
\label{p10:thm:color-connection-kinetic-uniqueness}
The Paper~V-A color carrier fixes the interacting local color group and the CP-even two-derivative gauge kinetic term:
\begin{equation}
\label{p10:eq:ym-kinetic-output}
C\cong\CC^3
\Longrightarrow
G_c=SU(3)_c,
\qquad
\mathcal L_{\rm YM}^{(c)}=-\frac14F_{\mu\nu}^aF^{a\mu\nu}.
\end{equation}
\end{theorem}
\begin{proof}
Sector-preserving compact transformations of \(C\cong\CC^3\) give \(U(3)\).  The central \(U(1)\) is not an independent color force because the interacting color response is traceless and the abelian sectors have already been separated upstream.  Hence the interacting color algebra is \(\mathfrak{su}(3)\).  For a simple Lie algebra, every invariant quadratic form is proportional to the trace/Killing form.  With \eqref{p10:eq:hermitian-generator-normalization}, the unique CP-even local two-derivative curvature scalar is \(-\frac12\operatorname{Tr}(F_{\mu\nu}F^{\mu\nu})=-\frac14F_{\mu\nu}^aF^{a\mu\nu}\).\qedhere
\end{proof}

\begin{corollary}[Derived local color phase space and Yang--Mills Hamiltonian]
\label{p10:cor:v174-local-color-phase-space-hamiltonian}
On the native stationary-product, parity-even, pure-gauge domain of
Paper~III-F,
Theorem~\ref{p3f:thm:v174-generated-joint-gauge-phase-space-action}, the
color restriction of the generated joint connection carries
\[
 \mathcal P_{\rm YM,c}^{\rm reg},
 \qquad
 d_{A_c}\Pi_c=0,
 \qquad
 \mathcal P_{\rm YM,c}^{\rm red}
 =\mu_c^{-1}(0)\mathbin{/\mkern-6mu/}\mathcal G_{c,\Sigma},
\]
together with its positive local Yang--Mills Hamiltonian and full nonlinear
local Lorentzian action.  Thus no separate local continuum
\((\mathrm{LEG}_{\rm YM})\), \((\mathrm{LOR}_{\rm dir})\), regular
gauge-phase-space, or source-free Gauss-law premise remains for this color
branch.

This conclusion does not choose the finite regulator action representative
\(\mathcal A_L\), finite transfer/Haar measure, execution-scale running
coupling, quark-mass operator, continuum limit, boundary edge data, or an
actual color-field history.  Those remain exactly the later interfaces of
Theorem~\ref{p10:thm:finite-strong-hamiltonian-closed} and Paper~VII-C.
\end{corollary}



\begin{theorem}[Free transverse gluon octet and primitive-clock locking]
\label{p10:thm:v245-free-gluon-octet-clock-lock}
On the trivial color background of the stationary source-free parity-even
Yang--Mills branch, write \(A=\varepsilon a^aT_a\).  Since
\[
 F^a=\varepsilon\,da^a+O(\varepsilon^2),
\]
the physical transverse quadratic Hessian is
\[
 \boxed{K_{g,T}
 =\left(\partial_t^2+c_{\rm op}^2\Delta_1^{\rm co}\right)
 \otimes I_{\mathfrak{su}(3)}.}
\]
Equivalently, the eight adjoint color directions are eight copies of the same
massless coexact gauge-wave operator.  On the calibrated round branch,
\[
 \boxed{H_g^{(2)}=\Eper N_x\otimes I_8,
 \qquad E_{g,n}=n\Eper,\quad n\ge2,}
\]
and the gauge-covariant mode multiplicity is
\[
 \boxed{d_n^{(g)}=8\times2(n^2-1)=16(n^2-1).}
\]
Thus the complete free transverse gluon octet recurs after one primitive clock
turn,
\[
 \boxed{U_g^{(2)}(T_t)=I,}
\]
up to a common scalar vacuum-energy phase on a second-quantized realization.
This theorem is a quadratic/trivial-background statement.  It does not make
the full interacting QCD Hamiltonian periodic and does not identify a
covariant one-gluon mode with a physical singlet state.
\end{theorem}

\begin{proof}
The Yang--Mills curvature has
\(F^a=da^a+g_sf^{abc}a^b\wedge a^c\).  Its second variation at \(A=0\)
therefore contains only the invariant color metric \(\delta_{ab}\); the
\(f^{abc}\) terms first enter the cubic and quartic interactions.  Paper~III-F's
Gauss-reduced massless transverse Hessian gives the differential factor, and
\(\dim\mathfrak{su}(3)=8\) gives the internal identity factor.  The energy,
degeneracy and recurrence follow from the Paper-III-A massless-gauge
clock-lock theorem and the two coexact helicities.
\end{proof}



% Source Paper VII-A, original line 2195; inherited proof bodies unchanged.
\subsubsection{Coupling normalization and CP separation}
\label{p10:subsec:coupling-normalization-cp-separation}

\begin{proposition}[Strong coupling at the native and comparison scales]
\label{p10:prop:strong-coupling-normalization}
On the source-native common-heat branch,
\begin{equation}
\label{p10:eq:gs-pivot}
 \frac{g_{s,\rm heat}^2}{4\pi}=\alpha_s^{\rm heat}=\frac\pi2,
 \qquad g_{s,\rm heat}=\sqrt2\,\pi.
\end{equation}
A coupling at a named external scale \(\mu\) is a readout of the same marginal
owner after the typed RG/scale map:
\[
 \alpha_s(\mu)=\frac{g_s^2(\mu)}{4\pi}.
\]
The historical specialization \(\alphasseed=\deltageom\) is retained only as
the broad-branch \(t=1/2\) comparison member \(\Hseedzero\); it is not derived
from the color counts and is not the source-native heat normalization.
\end{proposition}
\begin{proof}
The native statement is Paper~V-C's imported common-heat result.  The
external-scale equation is the definition of the canonically normalized
running coupling after the marginal-owner transport; no new source scalar is
introduced.
\end{proof}

\begin{definition}[Strong-CP term separation]
\label{p10:def:strong-cp-term-separation}
The topological density is
\begin{equation}
\label{p10:eq:theta-density}
q_t(x)=\frac{1}{32\pi^2}F_{\mu\nu}^a\widetilde F^{a\mu\nu},
\end{equation}
and a conventional QCD action could contain \(\theta q_t\).  On the declared CP source member used for the printed numerical branch, the coefficient is specialized to:
\begin{equation}
\label{p10:eq:theta-zero-separation}
\bar\theta_{\rm QCD}=0.
\end{equation}
Thus the CP-even Yang--Mills uniqueness theorem and the CP-interface member are logically separate: the former selects the CP-even quadratic invariant, while the latter is a conditional physical strong-CP realization.  Paper~VII-A does not derive \(\bar\theta_{\rm QCD}=0\) from the finite carrier.
\end{definition}


% Source Paper VII-A, original line 2237; inherited proof bodies unchanged.
\subsubsection{Finite Hamiltonian and transfer target}
\label{p10:subsec:finite-hamiltonian-transfer-target}

\begin{definition}[Finite strong Hamiltonian and transfer target]
\label{p10:def:finite-strong-hamiltonian}
The finite regulated Born--Hopf QCD Hamiltonian is the self-adjoint finite operator
\begin{equation}
\label{p10:eq:finite-strong-hamiltonian}
\HQCDfin
:=
\Ps
\left(
H_{\YM}^{\rm fin}[A_c;g_s]
+H_q^{\rm fin}[q,A_c,M_q^{\rm phys}]
+H_{\rm int}^{\rm fin}[q,A_c;g_s]
\right)
\Ps,
\end{equation}
where the color connection, quark support, singlet projector, and finite operator calculus are inherited structurally.  A selected numerical full-QCD representative additionally requires the gauge coupling transported to the execution scale from its owning common-heat/RG branch (or an explicitly labelled comparison-scale coordinate) and a typed physical mass operator \(M_q^{\rm phys}\).  Before that member is selected, Paper~II-A whole-family reuse allows the entire generated mass-operator family to feed the full-QCD Hamiltonian family parametrically.  After subtracting the finite ground energy, the dimensionless transfer family is
\begin{equation}
\label{p10:eq:qcd-hamiltonian-transfer-target}
\widehat K_L(t)=\exp[-t(\HQCDfin-E_{0,L}\mathbf 1)],\qquad t>0.
\end{equation}
A continuum/refinement object may be defined only from a specified sequence \(L\to\infty\), a specified transfer-time scaling \(t=t_L\), and a proved convergence topology.  No preferred dimensional wall scale is inserted at this stage.
\end{definition}


\begin{theorem}[Exact finite singlet Floquet reduction for the pure-glue kernel]
\label{p10:thm:v247-singlet-floquet-qcd}
At fixed finite regulator level, decompose an admissible pure-glue Hamiltonian
as
\[
 H_{\rm YM,L}=H_{g,0,L}+V_{\rm YM,L},
\]
where \(H_{g,0,L}\) is the free transverse adjoint Hamiltonian of
Theorem~\ref{p10:thm:v245-free-gluon-octet-clock-lock} on the retained mode
carrier and \(V_{\rm YM,L}\) is the remaining time-independent self-adjoint
gauge-invariant finite interaction/regulator block.  After any common scalar
vacuum-energy offset is removed from \(H_{g,0,L}\),
\[
 U_{g,0,L}(T_t)=I.
\]
Hence the interaction-picture Hamiltonian
\[
 V_I(t)=U_{g,0,L}(t)^\dagger V_{\rm YM,L}U_{g,0,L}(t)
\]
is exactly periodic,
\[
 \boxed{V_I(t+T_t)=V_I(t).}
\]
The one-turn interaction-picture unitary is therefore
\[
 \boxed{\mathcal F_{\rm YM,L}
 =\mathcal T\exp\!\left[-\frac{i}{\hbar}
 \int_0^{T_t}V_I(t)\,dt\right].}
\]
Because both the free block and the interaction commute with the exact Haar
singlet projector \(\Ps\), the physical one-turn operator is the unitary
restriction
\[
 \boxed{\mathcal F_{\rm YM,L}^{\rm sing}
 =\Ps\mathcal F_{\rm YM,L}\Ps
 \big|_{\Ps\mathcal H_L}.}
\]
Its eigenphases determine the exact finite singlet evolution over one
primitive turn and define quasienergies modulo \(\Eper\).  Since the free
one-turn operator is the identity, this is also the full pure-glue one-turn
evolution up to the removed common phase.

This theorem is a finite spectral reorganization, not a proof of a positive
Yang--Mills mass gap, continuum confinement, or continuum recurrence.  An
absolute interacting energy requires a choice of quasienergy lift and direct
spectral comparison with \(H_{\rm YM,L}\).
\end{theorem}

\begin{proof}
The free periodicity is Theorem~\ref{p10:thm:v245-free-gluon-octet-clock-lock}.
Therefore
\(U_{g,0,L}(t+T_t)=U_{g,0,L}(t)\), which gives the interaction-picture
periodicity by conjugation.  Finite-dimensional Schr\"odinger evolution gives
\(\mathcal F_{\rm YM,L}\) without a perturbative assumption.  Gauge
invariance makes \(\Ps\) commute with both terms of the Hamiltonian, so the
singlet subspace is invariant and the displayed compression is its exact
unitary restriction.  The quasienergy period is
\(2\pi\hbar/T_t=\Eper\).
\end{proof}

\begin{theorem}[Finite regulated strong Hamiltonian: pure-glue closure and full-QCD conditional extension]
\label{p10:thm:finite-strong-hamiltonian-closed}
The finite gauge carrier and exact singlet projection define a parametric pure-glue operator family \(H_{\YM,L}[g_s;\mathcal A_L]\), where \(\mathcal A_L\) denotes an admissible gauge-invariant finite action representative.  A fixed-L numerical representative requires the coupling at that execution scale and one admissible finite action representative.  The generated upstream source owner supplies a mass-operator family \(\mathfrak M_q\) even before an actual G0 member occurs, so the full quark construction yields the family-valued descendant
\[
 \{H_{\QCD,L}[M_q]:M_q\in\mathfrak M_q\}.
\]
A selected numerical full-QCD representative additionally requires the actual same-action family assembly transported to the execution scale as \(M_q^{\rm phys}(\mu)\).  For the base phase-selected bounded-local continuum export, Paper~VII-C v2.83 does not require ultraviolet cutoff removal or regulator-universality across the admissible finite-action family: the chosen finite member remains regulator/presentation data on that certified phase branch.  A universality theorem is required only for a stronger claim of representative-independence, same-phase canonicity, or equivalence across a specified action family.  Thus finite operator existence and whole-family propagation are structural, while low-scale RG/threshold transport and any stronger action-comparison claim remain explicitly typed.
\end{theorem}
\begin{proof}
Gauge covariance and the finite carrier determine the allowed matrix structure and let one assemble \(H_{\YM,L}[g_s;\mathcal A_L]\) for every typed pair \((g_s,\mathcal A_L)\).  The representation and common-heat trace fix the native matching value of
\(g_s\), while the generic auxiliary label \(L\) does not by itself identify a
physical external scale.  The finite nonlinear action is not unique and is a regulator representative.  Paper~VII-C supplies bounded phase-cluster existence and the conditional
ordinary local-QFT reconstruction on its actual locality/full-covariance
certificate, without requiring universal identification of all representatives; universality is consumed only by the stronger comparison claims just stated.  Paper~II-A's generated exact-solution/whole-family theorem permits the quark-mass family to be propagated through the Hamiltonian constructor before an actual member occurs; numerical execution then consumes the assembled family member after G16 RG transport as \(M_q^{\rm phys}(\mu)\).  The statement is therefore a direct separation of structural compilation from occurrence and analytic transport.
\end{proof}

\begin{corollary}[Whole-family quark-mass transport into the finite QCD kernel]
\label{p10:cor:whole-family-quark-mass-transport}
Let \(\mathfrak M_q\) be a nonempty generated upstream family of typed quark-mass operators, with every member satisfying the same frozen quark-support and gauge-covariance requirements used by the finite QCD constructor.  Then for each finite level \(L\),
\[
 M_q\longmapsto H_{\QCD,L}[M_q]
\]
is a licensed whole-family descendant and the complete Hamiltonian family is generated without a section \(M_q^{\rm phys}\).  Any spectral or response readout constant on that family is therefore already common generated content; member-sensitive numerical hadron/QCD predictions consume the actual G0 quark-source member and its G16 RG/threshold transport.
\end{corollary}

\begin{proof}
This is Paper~II-A whole-family parent reuse applied to the frozen finite QCD Hamiltonian word.  The mass operator enters as a typed coefficient port, so evaluating that word on every generated family member produces the entire Hamiltonian family.  GP3 may not select a mass member merely to force a numerical readout.
\end{proof}

\begin{corollary}[G34 disposition]
\label{p10:cor:v163-g34-disposition}
Consequently G34 contains no independent quark-mass constructor or selector.
The complete finite-QCD Hamiltonian family is already generated for the entire
upstream mass-operator family.  A numerical physical full-QCD member requires
the physical flavor/source orbit member together with the claim-scoped G16
renormalized pole/RG analytic evaluation inherited from upstream.  This does
not close confinement, the continuum mass gap, hadron/nuclear source
occurrence, or stronger representative-universality claims.
\end{corollary}


% Source Paper VII-A, original line 4340; inherited proof bodies unchanged.
\subsection{Source-native gauge measure and refinement universality update}
\label{p10:sec:heat-action-universality-update}

\begin{theorem}[Canonical reference measure and common-heat regulator representative]
\label{p10:thm:canonical-haar-heat-representative}
For any finite link set $E_L$, the configuration group is
$SU(3)^{E_L}$ and its unique normalized measure invariant under independent
link translations is the product Haar measure
\[
 \boxed{d\mu_{0,L}=\prod_{e\in E_L}dU_e.}
\]
If the finite gauge regulator process is required to be conservative,
symmetric, bi-invariant, and strongly local/second order on $SU(3)$, its
generator is the quadratic Casimir/Laplace--Beltrami operator up to the single
rate/coupling owner.  The corresponding central heat kernel therefore gives a
canonical finite face-action representative.  Central convolution/Markov
naturality without strong locality is insufficient because symmetric jump
semigroups form a nonlocal counterfamily.
\end{theorem}

\begin{remark}[Three distinct heat parameters]
\label{p10:rem:three-distinct-heat-parameters}
The source-native spectral functional
$\operatorname{Tr}\exp(-D^2/\Lambda^2)$ fixes the quadratic gauge coefficient
at its matching datum.  A compact-group face density $K_{t_{f,L}}(U_f)$
additionally requires a Lie-algebra metric, cell weight/quadrature convention,
and group heat time $t_{f,L}$.  The physical Euclidean transfer
$\exp[-\tau(H_L-E_{0,L})]$ is a third object and must be constructed from the
same reflection-positive action owner when it is used for a gap claim.
Consequently the current source-native identity $g_3^2=2\pi^2$ does not by
itself imply $t_{f,L}=g_3^2$, $g_3^2/2$, or $1/g_3^2$.  Such an identification
is an explicitly normalized regulator member, not a consequence of the word
``heat.''
\end{remark}

\begin{remark}[Electric kinetic heat realization is conditional and does not select the magnetic face law]
\label{p10:rem:v2159-electric-heat-scope}
Appendix~\ref{p10:app:v2159-apbi-action-owner-diagnostics} proves the following
precise scope statement.  If a same-owner finite cylindrical holonomy--flux
realization of the generated quadratic electric Yang--Mills Hamiltonian is
supplied, then its standard left/right-invariant cotangent quantization has a
Laplace--Beltrami generator and hence a product compact-group heat kernel.
Holonomy alone does not construct the required forward cotangent/flux map, and
this kinetic result does not identify the nonlinear magnetic spatial-face
action with the canonical heat representative.  Thus it removes no APBI
action-owner premise.
\end{remark}


\begin{theorem}[Casimir-normalized heat-Hessian inversion]
\label{p10:thm:v185-heat-hessian-inversion}
Fix the $SU(3)$ Lie-algebra normalization used by the fundamental color
representation,
\[
 -2\,\operatorname{tr}_{\mathbf3}(T_aT_b)=\delta_{ab},
\]
so the corresponding quadratic Casimir has
$C_2(p,q)=(p^2+q^2+pq+3p+3q)/3$.  More generally let $G$ be compact simple,
let $q_C$ be the invariant metric used by its Laplace--Beltrami operator, and
let
\[
 \partial_tK_t=\frac12\Delta_{q_C}K_t,
 \qquad
 s_t(g):=-\log\frac{K_t(g)}{K_t(e)}.
\]
Then the identity Hessian is
\[
 D^2s_t|_e=h_G(t)q_C,
\]
where
\[
 \boxed{
 h_G(t)=\frac1{\dim G}
 \frac{\sum_\lambda d_\lambda^2C_2(\lambda)e^{-tC_2(\lambda)/2}}
      {\sum_\lambda d_\lambda^2e^{-tC_2(\lambda)/2}}.}
\]
Moreover
\[
 \boxed{
 h_G'(t)=-\frac1{2\dim G}\operatorname{Var}_t(C_2)<0,}
\]
and
\(
 h_G(t)\to\infty
\)
as $t\downarrow0$ while $h_G(t)\to0$ as $t\to\infty$.  Hence
\[
 \boxed{h_G:(0,\infty)\overset\sim\longrightarrow(0,\infty)}
\]
is a bijection.  A positive target face stiffness $\kappa$ therefore fixes a
unique heat time
\[
 \boxed{t=h_G^{-1}(\kappa).}
\]
\end{theorem}
\begin{proof}
For $SU(3)$ the Hessian and variance identities are those of
Paper~1B, Proposition~\ref{p1b:prop:heat-block}
\cite{Kugelmann2026Paper1B}. For the stated general compact simple $G$,
centrality makes the identity Hessian $\operatorname{Ad}(G)$-invariant, and
simplicity makes every invariant quadratic form proportional to $q_C$.
Peter--Weyl expansion gives
\[
 K_t(e)=\sum_\lambda d_\lambda^2e^{-tC_2(\lambda)/2}.
\]
Taking the $q_C$-Laplacian of $-\log K_t$ at the identity, where the gradient
vanishes, gives the displayed weighted Casimir mean divided by $\dim G$.
Differentiating that mean gives the negative weighted Casimir variance.  At
finite $t$ the trivial and at least one nontrivial representation both have
positive weight, so the variance is strictly positive.  The large-$t$ limit
is the trivial representation.  The local heat-kernel asymptotic
$K_t(\exp X)/K_t(e)=\exp[-q_C(X,X)/(2t)](1+O(t)+O(\|X\|^2))$ gives
$h_G(t)=t^{-1}+O(1)$ as $t\downarrow0$.  Strict monotonicity and the two
endpoint limits prove bijectivity.
\end{proof}

\begin{corollary}[Common-heat face stiffness and unique heat time on an exact finite face presentation]
\label{p10:cor:v185-common-heat-face-stiffness}
On the source-native color branch, use Paper~V-A's
$g_3^2=2\pi^2$.  Suppose an exact finite reciprocal-position face presentation of the retained
quadratic Yang--Mills action assigns a four-volume
weight $w_{f,L}$ to an independent oriented two-plane, metric face area
$A_{f,L}$, and face logarithm
\[
 X_{f,L}=A_{f,L}\,\mathcal F_{f,L}+o(A_{f,L})
\]
in the Casimir-normalized compact curvature coordinate.  Here ``face'' and
``presentation'' refer to a derived finite presentation/regulator realization
on the reciprocal-position role; they do not constitute a microscopic position
lattice or a new position carrier.  Then its quadratic face action is
\[
 S_{L}^{(2)}
 =\frac12\sum_f\kappa_{f,L}\,q_C(X_{f,L},X_{f,L}),
 \qquad
 \boxed{
 \kappa_{f,L}=\frac{w_{f,L}}{g_3^2A_{f,L}^2}.}
\]
On the canonical strongly-local heat representative the matching face time is
therefore uniquely
\[
 \boxed{t_{f,L}=h_{SU(3)}^{-1}(\kappa_{f,L}).}
\]
For an isotropic presentation cell $w_{f,L}=A_{f,L}^2$,
\[
 \kappa_{\rm iso}=\frac1{2\pi^2},
 \qquad
 t_{\rm iso}=5.7532958466\ldots .
\]
This number is a group-heat time in the stated Casimir convention; it is not
the spectral heat parameter and not the physical Euclidean transfer time.
\end{corollary}
\begin{proof}
Paper~V-A matches the common-heat color term to the canonical convention
$\frac14F^a_{\mu\nu}F^{a\mu\nu}$.  Writing the compact curvature coordinate
as $\mathcal F=g_3F$ and summing one representative for each independent
$\mu<\nu$ plane gives
\[
 S_{\rm YM}^{(2)}
 =\frac1{2g_3^2}\sum_{\mu<\nu}\int
 q_C(\mathcal F_{\mu\nu},\mathcal F_{\mu\nu})\,d\mu.
\]
Exact cell quadrature and
$\mathcal F_{f,L}=X_{f,L}/A_{f,L}$ give the displayed stiffness.  Theorem~\ref{p10:thm:v185-heat-hessian-inversion}
then fixes the heat time uniquely.  In the isotropic case the stiffness is
$1/g_3^2=1/(2\pi^2)$; direct evaluation of the absolutely convergent
$SU(3)$ character sums in the theorem gives the displayed root.  No equality
with the distinct spectral or transfer heat parameters is used.
\end{proof}

\begin{proposition}[Four-dimensional blocking is not closed in one-face heat factors]
\label{p10:prop:multiface-heat-blocking}
On an action-first cellular branch, let an eliminated link variable occur in
three face holonomies.  The Haar pushforward of a product of nonconstant
central heat kernels generically contains a joint three-face interaction and
cannot in general be written as a product of one-face heat kernels with
renormalized times.  The exact positive coarse density and its effective
action nevertheless exist; the latter belongs to the larger quasilocal
multi-face action space used by Paper~VII-C.
\end{proposition}
\begin{proof}
Expand the three factors in characters.  Integrating the common link does not
apply a pairwise convolution identity; it projects the tensor product of the
three representation factors onto its invariant subspace.  For $SU(3)$,
$\mathbf3\otimes\mathbf3\otimes\mathbf3$ contains the invariant epsilon
tensor.  On central external holonomies the resulting function therefore has
a nonzero joint triality mode $(1,1,1)$, while every single-variable triality
mode vanishes by Haar orthogonality.  A product of three nonconstant one-face
heat kernels has nonzero single-variable modes, so no such factorization is
possible.  Positivity and existence of the pushforward follow directly from
integration of the positive fine density.  Paper~VII-C's appendix gives the
corresponding explicit $U(1)$ Fourier certificate.
\end{proof}

\begin{corollary}[Correct refinement target for the common action]
\label{p10:cor:multiface-effective-action-target}
Heat-semigroup convolution exactly collapses degree-two elimination steps,
but genuine four-dimensional refinement must control the generated
multi-face activities.  Exact projectivity of measures is therefore not a
substitute for the interacting effective-action/RG stability estimate.
\end{corollary}
\begin{proof}
The first statement is the heat-kernel semigroup identity.  The second is
Proposition~\ref{p10:prop:multiface-heat-blocking} combined with the exact
pushforward construction.
\end{proof}

\begin{proposition}[Exact heat subdivision does not manufacture an area exponent]
\label{p10:prop:v2158-heat-subdivision-area-nogo}
Let one fixed face heat factor of time $t>0$ be represented exactly by a
sequence of degree-two heat-semigroup subdivisions with positive times
$t_1,\ldots,t_n$ satisfying $\sum_i t_i=t$.  In the charged triality sector,
put
\[
 u_i=e^{-2t_i/3},\qquad u=e^{-2t/3}.
\]
Then
\[
 \boxed{\prod_{i=1}^n u_i=u.}
\]
Consequently exact subdivision of one already-owned heat face does not replace
one charged suppression factor by $u^n$.  In particular, attaching a proof
fugacity $z>1$ to the number of subdivision faces produces $z^n u$, not a
uniformly decaying area weight.  A positive physical surface free energy
therefore requires genuinely extensive local action support, not merely a
semigroup refinement of one fixed heat factor.
\end{proposition}
\begin{proof}
The heat semigroup adds times under degree-two convolution.  The charged
operator norm is $e^{-2t/3}$, so
$\prod_i e^{-2t_i/3}=e^{-2(\sum_i t_i)/3}=e^{-2t/3}$.  The final statement is
then immediate for every $z>1$.
\end{proof}

\begin{proposition}[Finite nonlinear action member is regulator data; universality is claim-scoped]
\label{p10:prop:finite-action-universality-status}
Gauge invariance and the Yang--Mills quadratic Hessian do not uniquely fix the
full finite joint action: local multi-face terms beginning at fourth order in
curvature leave the same $F^2$ Hessian while changing finite correlators.
Accordingly, no privileged nonlinear finite formula is selected by the quadratic data.  Paper~VII-C v2.83 nevertheless proves that the base bounded-local continuum export may be taken phase-selectively at fixed ultraviolet momentum ceiling, so existence of that export does not require a theorem identifying every admissible regulator representative.  If one additionally claims representative-independence, canonical same-phase refinement, or equality across an action family, then the corresponding universality/comparison theorem must still be proved on that claimed domain.
\end{proposition}

\begin{corollary}[Explicit same-Hessian non-heat face-action family]
\label{p10:cor:v2159-same-hessian-nonheat-family}
Let
\[
 q(U):=1-\frac13\operatorname{Re}\operatorname{Tr}_{\mathbf3}U,
 \qquad
 s_H(U):=-\log\frac{K_{t_H}(U)}{K_{t_H}(e)}.
\]
For every $\lambda\ge0$ define the central gauge-invariant one-face action
\[
 \boxed{s_\lambda(U)=s_H(U)+\lambda q(U)^2.}
\]
Then
\[
 Ds_\lambda|_e=0,
 \qquad
 \boxed{D^2s_\lambda|_e=D^2s_H|_e}
\]
for every $\lambda$, while $s_\lambda-s_H$ is nonconstant for
$\lambda>0$.  Hence even within positive central one-face actions the exact
Yang--Mills quadratic stiffness does not select the full nonlinear heat
density.
\end{corollary}
\begin{proof}
For $U=e^X$ with $X\in\mathfrak{su}(3)$, tracelessness removes the linear
term and the exponential expansion gives
$ q(e^X)=O(\|X\|^2)$.  Therefore $q(e^X)^2=O(\|X\|^4)$, so its first and
second derivatives vanish at the identity.  For $\lambda>0$ the added term
is nonconstant and nonnegative.  Thus the normalized positive density changes
without changing the identity Hessian; two such densities cannot define the
same probability measure for all bounded observables unless their action
difference is almost everywhere constant, which it is not.
\end{proof}


% BEGIN FRAMEWORK v3.025: NONLINEAR QCD ACTION-SENSITIVITY CERTIFICATES
\begin{proposition}[Exact isolated three-face action sensitivity]
\label{p10:prop:v3025-threeface-action-sensitivity}
Let $K_t$ be the normalized central $SU(3)$ heat kernel for a supplied
$t>0$.  Set $x(U)=\frac13\operatorname{Re}\operatorname{Tr}_{\mathbf3}U$,
$q(U)=1-x(U)$, and for $\lambda\ge0$ define the normalized one-face law
\[
 d\nu_{t,\lambda}(U)=Z_{t,\lambda}^{-1}
 K_t(U)e^{-\lambda q(U)^2}\,dU,\qquad
 a_{t,\lambda}=\int x(U)\,d\nu_{t,\lambda}(U).
\]
On the \emph{isolated independent-face reference} made of three copies
of this law, evaluate the already defined polarized color determinant
$W_e$ of Proposition~\ref{p10:prop:v2220-edge-fan-baryon}.  Then
\[
 \boxed{B_{t,\lambda}:=\langle W_e\rangle_{t,\lambda}
 =a_{t,\lambda}^{3},\qquad
 \partial_\lambda B_{t,\lambda}
 =3a_{t,\lambda}^2\operatorname{Cov}_{t,\lambda}(q,q^2)>0.}
\]
This comparison law is \emph{not} the interacting 96-face Hurwitz
ensemble, and its derivative is not a physical baryon-mass derivative.
At the separately stipulated isotropic exact-face time $t=t_{\rm iso}$
from Corollary~\ref{p10:cor:v185-common-heat-face-stiffness}, the
character calculation in
Proposition~\ref{p13:prop:v3025-character-moment-certificate} gives
$B_{t_{\rm iso},0}\simeq1.0063538725\times10^{-5}$ and
$\partial_\lambda B_{t_{\rm iso},\lambda}|_{\lambda=0}
\simeq1.4817641219\times10^{-4}$.  The face time $t_{\rm iso}$ is
not identified with a primitive momentum-Hurwitz time or a temporal
transfer time without the separate matching certificate.
\end{proposition}
\begin{proof}
Centrality and Schur's lemma give
$\int U\,d\nu_{t,\lambda}=a_{t,\lambda}I_3$; inversion symmetry
makes $a_{t,\lambda}$ real.  Independence of the three reference
face variables and the two epsilon tensors in $W_e$ give $a^3$.
Differentiating the compact normalized integral yields
$\partial_\lambda a=-\operatorname{Cov}(x,q^2)
=\operatorname{Cov}(q,q^2)$.  Since $q\ge0$ is nonconstant and the heat
law has strictly positive full-support density, the covariance
$\frac12\mathbb E[(q-q')(q^2-q'^2)]$ is strictly positive for
independent copies $q,q'$.  At $\lambda=0$ the fundamental heat
character yields $a_{t,0}=e^{-2t/3}>0$; monotonicity then gives
$a_{t,\lambda}>0$.  The numerical specialization and its distinct
heat-time scope are recorded in Paper~VII-D.
\end{proof}

\begin{proposition}[Nonuniversality of the complete finite 96-face reference law]
\label{p10:prop:v3025-fullface-variance}
On the marked momentum-Hurwitz $24$-cell with its \emph{existing}
$96$ triangular faces, product-Haar link measure and strictly positive
finite heat-action weight $e^{-S_H}$, put
\[
 Q(U)=\sum_{f=1}^{96}q(U_f)^2,\qquad
 d\mu_\lambda=\frac{e^{-S_H-\lambda Q}}{Z_\lambda}
 \prod_{e}dU_e,\qquad \lambda\ge0.
\]
Then
\[
 \boxed{\partial_\lambda\langle Q\rangle_\lambda
 =-\operatorname{Var}_\lambda(Q)<0.}
\]
Thus this full finite gauge-reference system has gauge-invariant
observables sensitive to the nonlinear action member, even though the
identity Yang--Mills Hessian is unchanged.  It makes no sign claim for
the full-ensemble covariance of $Q$ with a particular baryon source.
\end{proposition}
\begin{proof}
Differentiate the finite normalized compact-group integral.
Strict positivity follows because $Q$ is continuous and nonconstant
(one may change a face holonomy around a nontrivial triangular loop),
while the strictly positive heat density has full support.
\end{proof}

\begin{remark}[Momentum-incidence and physical-action ownership]
\label{p10:rem:v3025-momentum-2i-action-owner}
The $2T$, $2I$ and $600$-cell incidence objects are rooted in
$S^3_{\rm mom}$, not a primitive position lattice.  The existing
marked-$2T$ QCD transfer belongs to its fixed momentum carrier.
A mobile $2I$ sampling/presentation is an exact presentation of the
\emph{same} physical QCD operator only after the complete spin/color
intertwiner, Gauss quotient, nonlinear action, quantum measure and
source maps are transported together.  Free/Cartan one-body
intertwiners have narrower stated domains; for general nonflat color
connections the complete Schur complement retains the nonzero return
$B^\dagger(z-C)^{-1}B$.  Independent links at all momentum samples
are not, by themselves, a second licensed physical QCD regulator.
The conditional APBI presentation of
Definition~\ref{p10:def:v365-apbi-hadron-presentation}
remains a distinct proof obligation.
\end{remark}
% END FRAMEWORK v3.025



\begin{proposition}[Central face transfer requires positive type]
\label{p10:prop:v2211-central-transfer-positive-type}
Let $w\in L^1(G)$ be a real inversion-symmetric central density on a compact
group, with coefficients $a_R=\int_Gw(g)\overline{\chi_R(g)}\,dg$.
The convolution transfer $T_wf(x)=\int w(xy^{-1})f(y)\,dy$ is positive
self-adjoint if and only if $a_R\ge0$ for every irreducible $R$.
In particular the heat coefficients $a_R=d_Re^{-tC_R/2}$ are positive.
Pointwise positivity and the identity Hessian alone are insufficient.
\end{proposition}
\begin{proof}
Peter--Weyl orthogonality makes convolution by a central density scalar
$a_R/d_R$ on the $R$ matrix-coefficient space. Positivity is therefore
equivalent to the stated inequalities. For a sharp distinction, on $U(1)$
the strictly positive density $1-\varepsilon\cos(2\theta)$,
$0<\varepsilon<1$, has a negative second Fourier coefficient.
\end{proof}


% BEGIN FRAMEWORK v3.025: TEMPORAL POSITIVE-TYPE EXCLUSION
\begin{proposition}[A same-Hessian face law failing positive temporal transfer]
\label{p10:prop:v3025-sextet-temporal-exclusion}
With $q(U)=1-\frac13\operatorname{Re}\operatorname{Tr}_{\mathbf3}U$,
consider the central density $w_{t,1/10}=K_t(U)e^{-q(U)^2/10}$.
For every $5.75\le t\le5.76$, its normalized sextet character
coefficient satisfies
\[
 \boxed{\frac16\int_{SU(3)}w_{t,1/10}(U)
 \overline{\chi_{\mathbf6}(U)}\,dU<-4.5\times10^{-5}.}
\]
Thus this law \emph{cannot} be a positive self-adjoint central
\emph{temporal convolution transfer} on that interval, by
Proposition~\ref{p10:prop:v2211-central-transfer-positive-type}.
The exact $SU(3)$ representation-ring calculation and uniform tail
certificate are given in
Proposition~\ref{p13:prop:v3025-sextet-character-certificate}.
This does not exclude $\lambda=1/10$ as a spatial magnetic potential,
does not establish full reflection positivity, and does not uniquely
select the heat magnetic action.  Identification of $t$ with the
isotropic face time uses
Corollary~\ref{p10:cor:v185-common-heat-face-stiffness}; no physical
temporal time is inferred.
\end{proposition}
% END FRAMEWORK v3.025

This is the exact test when a face law is used as a temporal convolution
transfer. A spatial magnetic potential may instead enter a supplied
self-adjoint Hamiltonian and is not required to pass this convolution test.
The positive same-Hessian $SU(3)$ family above consequently proves
nonuniqueness of density, not reflection positivity of every member. A
negative character coefficient seen by numerical Weyl quadrature is a
diagnostic until certified tail and quadrature errors make its sign rigorous;
it is not by itself an exact exclusion of a native physical member.

\begin{corollary}[Finite QCD normalization owner versus claim-scoped action universality]
\label{p10:cor:v186-g33-finite-continuum-split}
The product Haar reference measure,
Theorem~\ref{p10:thm:v185-heat-hessian-inversion}, and
Corollary~\ref{p10:cor:v185-common-heat-face-stiffness} belong to a finite
normalization/representative owner.  They remove any separate face-stiffness
or group-heat-time selector on the stated exact-cell heat branch.  They do not by themselves prove source-compatible reflected multi-face effective-action control or representative-independence across a specified action family.  However Paper~VII-C v2.83 proves the fixed-UV phase-selected bounded-local QFT export by QFT4$_E$ and weak-* cluster compactness.  Hence the native finite G33 owner is closed/retired as an independent framework gap.  What remains is only a claim-scoped action/comparison certificate (G33-COMP) when a stronger regulator-universality or canonical-representative statement is actually consumed.
\end{corollary}

\begin{proposition}[Stable topological data do not close fixed-rank bundle geometry]
\label{p10:prop:v152-unstable-fixed-rank-topology-firewall}
Vanishing positive-degree Chern classes, vanishing reduced stable
\(K^0\)-class, and vanishing complex \(K\)-homology/index pairings do not, by
themselves, imply that a fixed-rank Born--Hopf bundle is trivial or admits a
flat connection.  Paper~I-C exhibits a nontrivial rank-two sector
\(\widetilde E\to M_8\) with
\[
 c_j(\widetilde E)=0\ (j>0),\qquad
 [\widetilde E]-2=0\in K^0(M_8),\qquad
 \langle[\widetilde E]-2,x\rangle=0\quad(x\in K_0(M_8)),
\]
yet no flat connection.  Consequently an unrestricted Paper-VII-A
bundle/topological/index closure must either classify or exclude such
unstable fixed-rank sectors, and any Laplace/Dirac twist by such a sector must
supply an additional nonflat connection.
\end{proposition}
\begin{proof}
This is the direct downstream typing consequence of
Paper~I-C's theorem ``Stable blindness, nonflatness and spectral scope of the unstable sector.  Stable
\(K\)-theory and characteristic/index pairings forget the unstable fixed-rank
class, while flatness would force the relevant pullback bundle to be trivial,
contradicting that theorem.  Hence stable invariants cannot be used as a
surrogate for fixed-rank triviality or connection data.
\end{proof}

% BEGIN M8 v1.78 X_IMPORT
\begin{remark}[Finite paired algebra and the QCD regulator]
\label{p10:rem:m8-batch-qcd-inputs}
The finite Gaussian uses
Proposition~\ref{p4:prop:paired-berezin-determinant-line}
when its exact field/dual and ordering roles are licensed.
This identity does not construct the finite QCD regulator
action/measure, its refinement compatibility or the
light-front physical species intertwiner.  Those remain
the QCD and representation obligations (G33 and G35).
The topological index and induced interaction retain
their sector/weight and operator hypotheses; neither
local anomaly cancellation nor finite paired evaluation
establishes a continuum claim.  Noether applies to the
same complete regulator/action owner when that owner exists.
\end{remark}
% END M8 v1.78 X_IMPORT



\begin{remark}[Typed mode refinement]
\label{p10:rem:typed-mode-refinement}
The primitive strong carrier remains $m_p=1$.  On a typed reciprocal
mode-resolution branch, scalar, spinor, and transverse-vector comparison
carriers are obtained by intersecting their Paper-I-A homogeneous-bundle
harmonic branches with the already generated total-weight filtration.  Their
inclusions are canonical; exterior/symmetric Fock functoriality and the fixed
color action then make the Haar singlet projector commute with refinement.
This is comparison/readout structure only and does not create new compact
momentum cells or assert physical occupation of every comparison mode.
\end{remark}



% Source Paper VII-A, original line 4709; inherited proof bodies unchanged.
\subsubsection{Quartic heat vertex and exact finite interaction compiler}
\label{p10:sec:v199-heat-interactions}
Retain the Casimir convention and unique exact-cell heat time of
Theorem~\ref{p10:thm:v185-heat-hessian-inversion}.  In particular
\[
 K_t(g)=\sum_Rd_Re^{-tC_R/2}\chi_R(g),\quad
 \mathbb P_t(R)=\frac{d_R^2e^{-tC_R/2}}{K_t(e)},\quad
 h=\frac{\mathbb E_tC_R}{8},\quad h'=-\frac{\operatorname{Var}_t(C_R)}{16}.
\]
For anti-Hermitian \(X\in\mathfrak{su}(3)\), write
\(q_C(X)=-2\operatorname{tr}_{\mathbf3}X^2\).

\begin{theorem}[Exact quartic color-face vertex]
\label{p10:thm:v199-quartic-heat-vertex}
Near the group identity,
\begin{equation}
 -\log\frac{K_t(e^X)}{K_t(e)}
 =\frac{h(t)}2q_C(X)+k_4(t)q_C(X)^2+O(\|X\|^6),\qquad
 k_4(t)=\frac{h(t)^2+h'(t)+h(t)/4}{40}.
 \label{p10:eq:v199-quartic-heat-vertex}
\end{equation}
\end{theorem}
\begin{proof}
Apply Paper~1B, Proposition~\ref{p1b:prop:heat-block}
\cite{Kugelmann2026Paper1B}, with the present Casimir normalization
$C_{\rm fund}=4/3$, $C_{\rm adj}=3$ and the heat time fixed by
Theorem~\ref{p10:thm:v185-heat-hessian-inversion}.
Its quadratic form is the present $q_C$, so its fourth coefficient is
exactly \eqref{p10:eq:v199-quartic-heat-vertex}.
\end{proof}

On the existing isotropic branch \(h(t_{\rm iso})=1/(2\pi^2)\), the
controlled representation sums give
\[
 \begin{gathered}
 t_{\rm iso}=5.753295846633\ldots,\\
 h'(t_{\rm iso})=-0.0276143453622\ldots,\quad
 k_4(t_{\rm iso})=-0.000309567546080\ldots.
 \end{gathered}
\]
The negative quartic is compatible with the positive exact heat kernel and
quadratic Hessian.  Its Taylor polynomial is local, and these finite heat
coefficients are not running couplings at a selected physical energy.

\begin{theorem}[Haar projectors compile exact finite blocked interactions]
\label{p10:thm:v199-haar-interaction-compiler}
On a finite oriented gauge complex with positive face times, character
expansion and normalized Haar elimination of each interior edge give
\begin{equation}
 P_e=\int\bigotimes_{f\supset e}\rho_{R_f}^{\epsilon_{fe}}(U)\,dU
 =P_{\operatorname{Inv}(\bigotimes_{f\supset e}V_{R_f}^{\epsilon_{fe}})}
 =\sum_\alpha|\iota_{e,\alpha}\rangle\langle\iota_{e,\alpha}|.
 \label{p10:eq:v199-haar-projector}
\end{equation}
Here negative incidence denotes the dual representation after arranging
indices, and \(\iota_{e,\alpha}\) is any orthonormal invariant basis.
The exact boundary weight of a finite block is
\begin{equation}
 F_{\mathcal B}(U_\partial)=\sum_{\{R_f\}}
 \left(\prod_f d_{R_f}e^{-t_fC_{R_f}/2}\right)
 \operatorname{Contr}_{\mathcal B}
 \left(\bigotimes_{e\ {
m interior}}P_e;
          \{\rho_{R_f}(U_\partial)\}\right).
 \label{p10:eq:v199-blocked-tensor-weight}
\end{equation}
The contractions follow the existing face-edge incidence.  This tensor class
is closed under exact finite blocking and gluing.
\end{theorem}
\begin{proof}
The Haar average of a unitary representation is a self-adjoint idempotent
with image exactly its invariant subspace.  Apply it to every interior
edge of the character-expanded face product.  At fixed positive times,
\(|\chi_R|\le d_R\) and \(\sum_Rd_R^2e^{-tC_R/2}<\infty\) justify
interchanging the sums and finite Haar integrals.  Gluing is contraction of
shared boundary indices.  This is the compact-group duality construction
\cite{OecklPfeiffer2001Dual} specialized to the retained native incidence.
\end{proof}

One explicit coupled channel uses the invariant
\(\iota_{ijk}=\varepsilon_{ijk}/\sqrt6\) in \(\mathbf3^{\otimes3}\).
For external color matrices \(A_c,B_c,C_c\),
\begin{align}
 &\int\chi_{\mathbf3}(A_cU)\chi_{\mathbf3}(B_cU)
                  \chi_{\mathbf3}(C_cU)\,dU\nonumber\\
 &=\frac16\bigl[\operatorname{tr}A_c\operatorname{tr}B_c\operatorname{tr}C_c
 -\operatorname{tr}A_c\operatorname{tr}(B_cC_c)
 -\operatorname{tr}B_c\operatorname{tr}(A_cC_c)\nonumber\\
 &\hspace{1.8em}-\operatorname{tr}C_c\operatorname{tr}(A_cB_c)
 +\operatorname{tr}(A_cB_cC_c)+\operatorname{tr}(A_cC_cB_c)\bigr].
 \label{p10:eq:v199-three-face-channel}
\end{align}
Expand the two epsilon tensors to prove this identity; setting all three
matrices equal gives their determinant.  In the unnormalized three-face
heat product its coefficient is \(27e^{-2(t_1+t_2+t_3)/3}\), with the
conjugate channel contributing the complex conjugate.  This is one exact
channel, not a dominance assumption.

The full \(F_{\mathcal B}\) is positive and its action is
\(-\log F_{\mathcal B}\), generally coupling several boundary faces.
Individual channel amplitudes need not be positive.  Appendix~
\ref{p10:app:v199-heat-tails} supplies explicit uniform truncation bounds for
each fixed positive-time finite block.  Tensor closure alone provides no
small interaction norm, source-uniform RG estimate or continuum state;
Paper~VII-C retains these as separate complete-action obligations.

% BEGIN V202 POSITIVE NATIVE HEAT BLOCKS

% Source Paper VII-A, original line 4814; inherited proof bodies unchanged.
\subsubsection{Explicit positive bounds from the native heat stiffness}
\label{p10:sec:v202-positive-heat-blocks}
\begin{theorem}[Native heat minimum and finite-block relative error]
\label{p10:thm:v202-native-heat-bounds}
Use the admitted $SU(3)$ heat representative and normalized Haar measure
of Theorem~\ref{p10:thm:v185-heat-hessian-inversion}. At
$h(t_{\rm iso})=1/(2\pi^2)$,
\begin{equation}
 a:=\frac{\pi^2-6}{\pi^2-3}\le K_{t_{\rm iso}}(g)
 \le b_H:=\frac{\pi^2}{\pi^2-3},\qquad
 a>\frac{14}{25},\quad b_H<\frac{36}{25}.
 \label{p10:eq:v202-native-heat-minimum}
\end{equation}
Consequently the full Haar-integrated boundary weight of a block of $N$
such faces satisfies $a^N\le F_{\mathcal B}\le b_H^N$.
Moreover the exact native root and total-degree-17 truncation obey
\begin{equation}
\begin{gathered}
 t_{\rm iso}>\frac{14383}{2500},\\
 \|K_{t_{\rm iso}}-K_{t_{\rm iso}}^{(17)}\|_\infty<4\,10^{-117},\\
 \left\|\frac{F_{\mathcal B}^{(17)}}{F_{\mathcal B}}-1\right\|_\infty
 <\varepsilon_N:=\frac{25N}{14}(4\,10^{-117})
                  \left(\frac{18}{7}\right)^{N-1}.
\end{gathered}
 \label{p10:eq:v202-native-heat-relative-error}
\end{equation}
For $N\le96$, $\varepsilon_N<6.35\,10^{-76}$. The face-count condition
does not identify 96 with the full four-dimensional three-layer complex.
\end{theorem}
\begin{proof}
Set $\kappa=1/(2\pi^2)<1/12$ in Paper~1B,
Proposition~\ref{p1b:prop:heat-block} \cite{Kugelmann2026Paper1B}.
Its pointwise and positive-block bounds give
\eqref{p10:eq:v202-native-heat-minimum}, including the displayed
rational comparisons. The rational certificate in
Appendix~\ref{p10:app:v202-native-heat-certificate} proves the native
root and face-tail margins. Telescoping the face product, using
$|K_t^{(17)}|\le K_t(e)\le b_H$, then dividing by $a^N$, gives
the relative error and its displayed rational estimate. This argument
uses the actual $N$ faces and does not factor the blocked weight.
\end{proof}

The same proof allows face-dependent $0<\kappa_f<1/12$, with
$a_f=(1-12\kappa_f)/(1-6\kappa_f)$ and
$b_f=1/(1-6\kappa_f)$. For the convention $K_t/K_t(e)$ divide each
face bound by $K_t(e)$; relative errors and normalized fixed-background
readouts agree, while source- or metric-dependent constants remain in the
complete action. Paper~VII-C now derives fixed bounded-source derivatives
from these estimates in Theorem~\ref{p12:thm:v202-source-errors}.
% END V202 POSITIVE NATIVE HEAT BLOCKS


% Source Paper VII-A, original line 4864; inherited proof bodies unchanged.
\subsection{Refinement, area, and finite-index closure update}
\label{p10:sec:g7-g8-g9-update}

The Paper-II-C evaluation-spectrum isomorphism
$\iota_1:S^3_{\rm mom}\to S^3_x$ transports the fixed Hurwitz regulator
complex canonically:
\[
 \mathcal K_H^x:=\iota_1(\mathcal K_H),
 \qquad
 \iota_{1\#}:C_q(\mathcal K_H)\overset\sim\longrightarrow C_q(\mathcal K_H^x),
\]
with
\[
 \partial\iota_{1\#}=\iota_{1\#}\partial.
\]
Thus fixed-base face/chain ancestry is derived.  A macroscopic/refined area
current remains separate because increasing spectral resolution does not
select a unique mesh.

\begin{proposition}[Center flux does not force quadratic area extensivity]
\label{p10:prop:center-fractionalization-no-area}
Let
\[
 H_c=\operatorname{diag}(1,1,-2),
 \qquad
 \exp\!\left(i\frac{2\pi}{3}H_c\right)=\omega I.
\]
On a spanning disk of $N$ commuting faces set
\[
 U_f=\exp\!\left(i\frac{2\pi}{3N}H_c\right).
\]
Then the boundary holonomy is the nontrivial center element $\omega I$, while
for the integrated face coordinates
$F_f=(2\pi/3N)H_c$,
\[
 \boxed{\sum_{f=1}^N\|F_f\|_{\rm HS}^2=\frac{8\pi^2}{3N}\to0.}
\]
Hence center charge plus the geometric Hodge lower bound cannot imply
$\|F\|^2\ge cA$.  A quantized center-defect/no-fractionalization theorem or a
direct positive surface-free-energy theorem is required for area extensivity.
\end{proposition}

\begin{proposition}[Finite exact-chirality index obstruction]
\label{p10:prop:finite-exact-chiral-index-zero}
Let $D_+:H_+\to H_-$ be the chiral block of an exact odd finite Dirac operator
with $\dim H_+=\dim H_-<\infty$.  Then
\[
 \boxed{\operatorname{ind}D_+=0.}
\]
\end{proposition}

\begin{proof}
Rank-nullity gives
$\dim\ker D_+=N-\operatorname{rank}D_+$ and
$\dim\ker D_+^*=N-\operatorname{rank}D_+^*$; the ranks agree.
\end{proof}

Therefore nonzero finite topology cannot come from simply compressing the G1
exact-chiral symbol.  The finite-regulator theorem below supplies the required
GW/overlap family from the already generated Hurwitz spatial incidence,
transfer ribbons and protected four-dimensional Clifford grading.  Integer
sector weights and topological susceptibility remain outputs of the physical
action/state/transfer law, not of the index classification itself.

\begin{theorem}[Generated four-dimensional Hurwitz--transfer GW regulator]
\label{p10:thm:v190-generated-4d-gw-regulator}
Let $C_{N_\tau}$ be any finite cyclic representative obtained by composing the
one-step transfer complex $\mathcal K_H^\tau$, and let
\[
 \mathscr H_{4,N_\tau}
 =\ell^2(C_{N_\tau}\times V_H)\otimes\mathbb C^4_{\rm spin}
 \otimes\mathbb C^3_{\rm color}.
\]
Assign spatial links $U^{(j)}_{qq'}\in SU(3)$ and temporal ribbon links
$V_{j,q}\in SU(3)$.  Using the first-symbol components of the existing
Hurwitz stencil and the centered temporal covariant difference, define
anti-Hermitian covariant derivatives $\nabla_1,\ldots,\nabla_4$ and
\[
 \boxed{D_4^{(0)}=\sum_{\mu=1}^4\Gamma_\mu\nabla_\mu.}
\]
Then
\[
 (D_4^{(0)})^\dagger=-D_4^{(0)},\qquad
 \{\Gamma_5,D_4^{(0)}\}=0.
\]
The same covariant incidence gives positive spin-even Laplacians
$\Delta_H^U,\Delta_\tau^V\ge0$.  For regulator coordinates
$r_H,r_\tau\ge0$ and $\rho\in\mathbb R$, put
\[
 W=\frac{r_H\ell_H}{2}\Delta_H^U
   +\frac{r_\tau a_\tau}{2}\Delta_\tau^V,
 \qquad K_{\rho,r}=D_4^{(0)}+W-\rho I,
\]
\[
 \boxed{H_{\rho,r}=\Gamma_5K_{\rho,r}.}
\]
Then $H_{\rho,r}$ is Hermitian.  On every connected gapped component
\[
 \mathcal G=\{(U,V;r,\rho):0\notin\operatorname{Spec}H_{\rho,r}\},
\]
finite functional calculus gives $\epsilon_H=\operatorname{sign}H_{\rho,r}$
and
\[
 \boxed{D_{\rm ov}=\frac1{a_R}(I+\Gamma_5\epsilon_H),}
\]
which satisfies exactly
\[
 \boxed{\Gamma_5D_{\rm ov}+D_{\rm ov}\Gamma_5
 =a_RD_{\rm ov}\Gamma_5D_{\rm ov}.}
\]
On the balanced protected spin carrier its finite index is
\[
 \boxed{\nu_{\rm ov}=-\frac12\operatorname{Tr}\epsilon_H}
\]
up to the fixed global chirality/orientation convention.  Motion inside one
connected gapped component is regulator-presentation data and does not select
a physical Wilson mass or coupling.
\end{theorem}

\begin{proof}
The existing Hurwitz stencil and cyclic temporal ribbons supply the
product graph and link data of Paper~1B,
Theorem~\ref{p1b:thm:overlap} \cite{Kugelmann2026Paper1B}, with rank
$r=3$, spatial spacing $\ell_H$ and the present temporal spacing
$a_\tau$. Reverse edges have opposite tangents and adjoint links;
the corresponding Laplacians are $d_U^\dagger d_U$ and spin scalar.
The theorem therefore applies to this $D_4^{(0)},W,H_{\rho,r}$ and
gives the displayed adjoint, chiral and GW identities and index formula
on each stated gapped component. The trace uses the balanced protected
spin carrier $\operatorname{Tr}\Gamma_5=0$.
\end{proof}

\begin{definition}[Native finite clutching target and remaining index match]
\label{p10:def:v190-native-clutching-target}
Cut a finite cyclic transfer representative at one layer.  If its closing
slice differs from its initial slice by a gauge transformation
\[
 g:|\mathcal K_H|\simeq S^3\longrightarrow SU(3),
\]
its native topological target is the clutching class
\[
 [g]\in\pi_3(SU(3))\cong K^1(S^3)\cong\mathbb Z.
\]
The canonical unit candidate is the embedded quaternion
\[
 \boxed{g_1(q)=\operatorname{diag}(q,1).}
\]
For claims that export the native overlap charge to this clutching presentation,
the comparison theorem must supply a source-level admissibility/interpolation
and homotopy certificate on the claimed admitted family and prove
\[
 \boxed{\nu_{\rm ov}=\nu_{\rm clutch}}
\]
on that whole family. Theorem~\ref{p10:thm:v199-native-clock-covers} proves the explicit clock-cover families; Theorem~\ref{p10:thm:m8td-all-link-neighborhood} extends them to arbitrary nearby links on the same supplied clutched bundle. Neither asserts the unrestricted comparison theorem.  Existence of the gapped GW operator already supplies native finite topology; existence of a clutching candidate does not by itself prove equality with that comparison presentation.
\end{definition}

\begin{corollary}[Native finite regulator and topology are generated]
\label{p10:cor:v190-g20-finite-regulator-generated}
The generated gapped Hurwitz--transfer family supplies the 4D massless operator,
positive degree-two regulator, Hermitian overlap kernel, sign calculus, GW
operator and native finite index as same-owner descendants.  Native finite
topology therefore does not wait for an unrestricted clutching comparison.
Definition~\ref{p10:def:v190-native-clutching-target} is retained only for
claims that identify this native charge with that geometric/continuum
presentation.  Sector probabilities and susceptibility remain dynamical
outputs rather than selector data.
\end{corollary}


\begin{remark}[Scalar spin-twist and four-dimensional regulator firewall]
\label{p10:rem:v190-scalar-spin-g20-firewall}
A flat temporal spin sign may enter the cyclic transfer boundary condition, but
$S_E^1$ is not a fifth spacetime coordinate of the four-dimensional QCD
regulator.  The scalar spin twists carry no gauge curvature and therefore do
not supply the local anomaly density or the native integer gauge topological
class.  Once the native index-match theorem is proved, the regulated gauge
index must be independent of such flat scalar spin twists on the same closed
index domain.  Global determinant magnitudes or boundary spectra may still
depend on the selected global lift.
\end{remark}

\begin{definition}[Historical G20 layer split and v1.61 retyping]
\label{p10:def:v186-g20-layer-split}
The older G20 split is retained only as provenance.  In the active v1.61
classification:
\begin{enumerate}[label=\textnormal{(\Alph*)},leftmargin=3.8em]
\item the finite regulator owner and native overlap index are generated;
\item positive primitive \(\nu=\pm1\) sectors are generated on open native
      neighborhoods by Theorem~\ref{p10:thm:v161-positive-primitive-native-sectors};
\item the strong sector character is fixed by the generated action and the
      quark determinant orientation by the rooted source geometry;
\item unrestricted overlap/clutching comparison and any global
      gauge-field-dependent chiral determinant line are continuum/chiral
      certification tasks exported to Paper~VII-C.
\end{enumerate}
Thus this historical label no longer denotes an independent headline gap.
No clause weakens Proposition~\ref{p10:prop:finite-exact-chiral-index-zero}.
\end{definition}




% Source Paper VII-A, original line 7549; inherited proof bodies unchanged.
\subsection{same-action QCD source certificates}
\label{p10:sec:v2210-lit-source-certificates}
The following results act on the complete finite source member of
Definition~\ref{p10:def:finite-qcd-compilation-datum} and
Corollary~\ref{p10:cor:finite-qcd-parent-readouts}. They do not select that
member or replace its physical clock, quark/sea allocation, metric jets or
comparison map. Their purpose is to turn the applicable literature routes
into explicit finite tests. Finite massive spectroscopy and current
responses do not wait on an exact ordered massless condensate, an optional
continuous Yang--Mills gap or the stronger exact finite $\LFRep_L$.


% Source Paper VII-A, original line 7560; inherited proof bodies unchanged.
\subsubsection{Exact blocked source jets and the multiface obligation}
\begin{theorem}[Complete source derivatives after same-action Haar blocking]
\label{p10:thm:v2210-lit-blocked-source-jets}
On a licensed positive-weight source chart let $U$ be the eliminated
compact gauge variables with source-independent normalized Haar measure
$d\mu(U)$. At fixed retained variables $y$, put
\[
 Z_y(J)=\int e^{-S(y,U;J)}d\mu(U)>0,
 \qquad A(y;J)=-\log Z_y(J).
\]
Assume $S$ is real and twice source differentiable, with the derivatives
dominated on the declared chart. Determinant, sea and normalization factors,
when present, are retained in the same positive weight; if represented by
$S$, their logarithms and jets are included. With brackets denoting this
normalized conditional weight,
\[
 \boxed{A_a=\langle S_a\rangle_y,\qquad
 A_{ab}=\langle S_{ab}\rangle_y
       -\langle(S_a-\langle S_a\rangle_y)
                (S_b-\langle S_b\rangle_y)\rangle_y.}
\]
For an actual nonlinear source pullback $J=J(\eta)$, the complete Hessian is
\[
 \boxed{\partial_i\partial_j A(J(\eta))
 =J^a_{,i}A_{ab}J^b_{,j}+A_aJ^a_{,ij}.}
\]
If a licensed gauge/source transformation preserves the original action and
Haar integral, the blocked functional obeys the same source transformation
identity. This conclusion applies to the full effective functional, including
the generated multiface terms.
\end{theorem}
\begin{proof}
Differentiating $Z_y$ gives $Z_{y,a}=-Z_y\langle S_a\rangle_y$.
For a differentiable insertion $F$,
\[
 \partial_b\langle F\rangle_y
 =\langle F_b\rangle_y-
   \langle(F-\langle F\rangle_y)(S_b-\langle S_b\rangle_y)\rangle_y.
\]
Apply this with $F=S_a$ and use the ordinary second-order chain rule.
For the last statement change integration variables by the actual licensed
transformation; normalized Haar invariance and action covariance give the
identity before any source derivative is taken.
\end{proof}
The negative covariance term is not an independent coupling. It records the
eliminated variables' response. Deleting it, a determinant factor or a source
chart contact changes the same-action Hessian. The valence-three interaction
and Hessian already executed in
Sections~\ref{p10:sec:v2201-derived-d21}--\ref{p10:sec:v2201-derived-d22}
are restrictions of this calculation; independent face-time renormalization
does not remove their mixed terms.

\begin{proposition}[A sufficient domain for analytic multiface source blocking]
\label{p10:prop:v2210-lit-multiface-polymer-domain}
Suppose the actual blocked weight admits an exact polymer representation
with activities $z_\gamma(J)$, its actual incompatibility relation, and
nonnegative functions $a(\gamma)$. Assume the activities extend holomorphically
to a declared source polydisc and obey there, uniformly in the consumed
finite presentations,
\[
 \sum_{\gamma'\not\sim\gamma}
 |z_{\gamma'}(J)|e^{a(\gamma')}\le a(\gamma).
\]
Then the Koteck\'{y}--Preiss theorem gives the absolutely convergent connected
expansion of the log of that polymer partition function. On smaller source
polydiscs its source jets are obtained by differentiating the full connected
expansion. Any cross-level/local observable limit additionally retains
convergence of its actual activities and uniform tail bounds.
\end{proposition}
\begin{proof}
Apply the abstract polymer convergence theorem to these activities and this
incompatibility relation~\cite{KoteckyPreiss1986}. The uniform holomorphic
bound gives local normal convergence; Cauchy integral differentiation on a
smaller polydisc commutes with the connected sum. The last sentence follows
by splitting that sum into a finite part and a uniformly bounded tail.
\end{proof}
This is a test of the actual multiface action. It constructs neither its
activities nor their bounds from a product-face ansatz. Analytic local-source
blocking alone is not an area-law theorem. APBI still needs its same-owner
extensive action, physical-area and flux estimates. The 2026 decorated-loop
heat-kernel result is an additional technique on large-$N$ $U(N)$ lattice
theory, not a finite-$SU(3)$ physical bridge~\cite{LitQCDLemoine2026HeatMaster}.


% Source Paper VII-A, original line 7643; inherited proof bodies unchanged.
\subsubsection{Finite Ward defects and a consistent axial reduction}
\begin{theorem}[Finite same-source Ward defect and response bound]
\label{p10:thm:v2210-lit-ward-defect}
Let $H(J)$ be the actual finite Hermitian QCD member on a twice differentiable
source chart, $\beta>0$, and
$\rho_J=e^{-\beta H(J)}/\operatorname{Tr}e^{-\beta H(J)}$.
Let $v^a(J)$ be a declared real source transformation and $Q(J)$ its actual
Hermitian generator. Define
\[
 \mathcal R(J)=v^a(J)H_a(J)-i[Q(J),H(J)],\qquad
 \mathcal W(J)=v^a(J)\partial_a\log Z(J).
\]
Then
\[
 \boxed{\mathcal W=-\beta\langle\mathcal R\rangle_J,\qquad
 |\mathcal W|\le\beta\|\mathcal R\|,}
\]
and, retaining all derivatives of $v,Q,H$,
\[
 \boxed{\partial_b\mathcal W
 =-\beta\langle\mathcal R_b\rangle_J
   +\beta^2\mathcal C_{\rho_J}(\mathcal R,H_b),\qquad
 |\partial_b\mathcal W|
 \le\beta\|\mathcal R_b\|+\beta^2\|\mathcal R\|\|H_b\|,}
\]
where
\[
 \mathcal C_\rho(A,B)=\int_0^1
  \operatorname{Tr}(\rho^s\delta A\rho^{1-s}\delta B)\,ds,
 \qquad\delta A=A-\operatorname{Tr}(\rho A)I.
\]
Thus exact vanishing of the actual finite defect closes this Ward identity
and its differentiated identity on this chart. It does not assert that the
physical axial or anomaly defect vanishes.
\end{theorem}
\begin{proof}
Duhamel differentiation gives $\partial_a\log Z=-\beta\langle H_a\rangle$.
Since $[\rho_J,H]=0$, the trace of $\rho_J[Q,H]$ vanishes. This proves the
first identity. Differentiating the expectation of $\mathcal R$ gives the
second. In a density eigenbasis the Duhamel diagonal quadratic form uses
positive logarithmic means. They are bounded by arithmetic means, hence
\[
 0\le\mathcal C_\rho(A,A)
 \le\operatorname{Tr}(\rho A^2)-\langle A\rangle^2\le\|A\|^2.
\]
Cauchy--Schwarz for this positive form bounds its mixed value by
$\|A\|\|B\|$; the operator norm bounds the remaining expectations.
\end{proof}
The contacts $\mathcal R_b$ include $v^a_{,b}H_a$, $v^aH_{ab}$ and the
derivative of $Q$. They must not be replaced by a current charge count.
A zero-temperature or physical axial-pole conclusion additionally consumes
its actual vacuum/source limit and pole overlap.

\begin{corollary}[Same-functional kernel differentiation is an axial test]
\label{p10:cor:v2210-lit-axial-functional-test}
If a licensed differentiable effective-action branch supplies a dressed
quark propagator $S_q$ and self-energy map $\Sigma_q[S_q;J]$, and the actual
infinitesimal axial variation on that branch obeys
\[
 \Sigma_q[S_q+\varepsilon\delta_A S_q;J+\varepsilon\delta_AJ]
 =\Sigma_q[S_q;J]+\varepsilon\delta_A\Sigma_q+o(\varepsilon),
\]
then its derivative satisfies
\[
 D_S\Sigma_q[\delta_A S_q]+D_J\Sigma_q[\delta_AJ]
 =\delta_A\Sigma_q.
\]
In the convention $K=-D_S\Sigma_q$, replacing this derivative by a separately
supplied two-body kernel $\widetilde K$ changes the covariance equation's
left side by $-(\widetilde K-K)[\delta_A S_q]$, bounded in norm by
$\|\widetilde K-K\|\,\|\delta_A S_q\|$.
\end{corollary}
\begin{proof}
Differentiate the printed same-functional covariance identity. Subtract the
kernel derivative identity and apply the induced operator norm.
\end{proof}
This is the consistency route underlying
\cite{LitQCDMunczek1995Axial,LitQCDMarisRobertsTandy1998Pion}.
It assumes the actual effective functional and its axial covariance; it
creates no independent two-particle-irreducible owner or fitted interaction.
The native Hamiltonian Schur reduction instead already retains all four
source terms and its residue norm in
Theorem~\ref{p10:thm:complete-schur-scalar-source}. Both tests prohibit
dressing a propagator and choosing an unrelated pion kernel. A Ward test is
not a proof of nonzero ordered $\Sigma$, $F_\pi$ or Goldstone residue.


% Source Paper VII-A, original line 7729; inherited proof bodies unchanged.
\subsubsection{Observable scheme transport without scale insertion}
\begin{proposition}[Observable step scaling on an actual one-scale branch]
\label{p10:prop:v2210-lit-observable-step-scaling}
Retain the actual member, physical scale map and all other physical marks.
On a declared interval suppose an owned dimensionless readout $u(\mu)$ is
continuously differentiable and locally one-to-one. For scale factors $s,r$
whose arguments remain on this branch define
\[
 \sigma_s(u)=u(s\mu(u)).
\]
Then
\[
 \boxed{\sigma_s\circ\sigma_r=\sigma_{sr},\qquad
 \left.\frac{d}{d\ell}\sigma_{e^\ell}(u)\right|_{\ell=0}
 =\mu\frac{du}{d\mu}.}
\]
If a second observable scheme is related on this same branch by a
differentiable invertible map $v=f(u)$, its step scaling is
$\sigma_s^{(v)}=f\circ\sigma_s^{(u)}\circ f^{-1}$ and its beta function is
$\beta_v=f'(u)\beta_u$. Actual mass/source transport maps likewise compose
only when their retained source, state and scale domains match.
\end{proposition}
\begin{proof}
Use the actual inverse scale coordinate twice:
$\mu(u(r\mu))=r\mu$. Differentiation at $s=e^\ell=1$ gives the logarithmic
scale derivative. The scheme statement is the chain rule. Composition of
mass/source maps is their ordinary composition on the matched domains.
\end{proof}
Observable gradient-flow/SF schemes supply a possible comparison technology,
with their actual boundary and zero-mode premises~\cite{LitQCDFritzschRamos2013GF}.
The native common-heat value does not supply $\mu(u)$, face heat time or a
physical transfer clock. Its strong boundary cannot be evolved by a
weak-coupling series without a controlled bridge. The applicable mass
normalization/threshold map remains the existing G16 calculation; the
nonperturbative step-scaling literature supplies a method rather than a
new quark member~\cite{LitQCDALPHA2018MassRunning}.


% Source Paper VII-A, original line 7766; inherited proof bodies unchanged.
\subsubsection{Topological response and its exact positivity scope}
\begin{theorem}[Native sector cumulants and the unsigned higher response]
\label{p10:thm:v2210-lit-sector-cumulants}
Retain the complete native sector effects/weights of
Corollary~\ref{p10:cor:finite-qcd-parent-readouts}. For a finite integer-charge
support with normalized probabilities $p_\nu\ge0$, define the auxiliary
real tilt $M(h)=\sum_\nu p_\nu e^{h\nu}$ and the Fourier source
$\Phi(\theta)=\sum_\nu p_\nu e^{i\theta\nu}$. Then
\[
 (\log M)''(0)=\operatorname{Var}(\nu)\ge0,
 \quad -(\log\Phi)''(0)=\operatorname{Var}(\nu),
\]
\[
 (\log M)^{(4)}(0)
 =\langle(\nu-\langle\nu\rangle)^4\rangle
       -3\operatorname{Var}(\nu)^2.
\]
The fourth cumulant has no universal sign. CP-symmetric weights imply
vanishing odd cumulants. Neither tilt changes the already fixed native
 character $\chi_{\rm top}=1$ or the rooted finite invariant
$\bar\theta=0\pmod{2\pi}$ of Paper~V-C,
Corollary~\ref{p7v2:cor:v259-native-strong-cp}.
\end{theorem}
\begin{proof}
Differentiate the finite characteristic/moment-generating sums and their
logarithms. For $p_1=p_{-1}=1/2$ the fourth cumulant is $-2$.
For $p_1=p_{-1}=\epsilon/2$, $p_0=1-\epsilon$ it is
$\epsilon-3\epsilon^2>0$ when $0<\epsilon<1/3$.
Reflection of the charge makes $\Phi$ even, giving the odd-cumulant claim.
An auxiliary source derivative does not insert a physical sector character.
\end{proof}
Index/sector existence supplies none of the $p_\nu$ values. Matching these
native cumulants to a regulated susceptibility or anomalous-Ward definition
retains its volume/contact/operator comparison; gradient-flow topology gives
a comparison route on its own domain~\cite{LitQCDCe2015Topological}.
The finite CP character is closed; regulated physical transport is retained
only by claims that consume it. The existing three-flavor anomaly
contradiction in Appendix~\ref{p10:app:conditional-chiral-anomaly} remains
conditional on its actual local infrared/anomaly content, as in the verified
v2 proof of 2 March 2025~\cite{CiambrielloEtAl2024ChiralStrategy}. It does not compute a native
bilinear condensate.
% END VERIFIED QCD LITERATURE TRANSFERS v2.210


\begingroup
\small
\RaggedRight
\setlength{\emergencystretch}{10em}
\sloppy
% Bibliography consolidated at the end of this volume
\endgroup

% BEGIN CHARGED-SOURCE-IMPORT v1.68

% Source Paper VII-A, original line 8536; inherited proof bodies unchanged.
\subsection{Geometry-first consequences and physical scope}
\label{p10:sec:v2201-geometry-first-consequences}
\paragraph{Geometry ancestry.}
Rooted color representation and licensed incidence $\to$ positive Haar/heat operator $\to$ exact elimination and complete mixed source Hessian.
\paragraph{Status.}
These are consequences on the domains stated below.
A mathematical certificate or an executed restriction retains every
physical source, state, clock and readout hypothesis of its inputs.


% Source Paper VII-A, original line 8545; inherited proof bodies unchanged.
\subsubsection{D21: Explicit $SU(3)$ three-face interaction after Haar elimination}
\label{p10:sec:v2201-derived-d21}
\noindent\textit{Status: proved result on stated domain.}

On a licensed valence-three heat block let
\[
 I_t(A,B,C)=\int_{SU(3)}K_t(AU)K_t(BU)K_t(CU)\,dU,
 \quad r=e^{-2t/3}.
\]
The exact fundamental triple invariant is
\begin{align*}
 \mathcal T(A,B,C)&=\int\operatorname{tr}(AU)\operatorname{tr}(BU)
                              \operatorname{tr}(CU)\,dU\\
 &=\frac16\{\operatorname{tr}A\operatorname{tr}B\operatorname{tr}C
  +\operatorname{tr}(ABC)+\operatorname{tr}(ACB)\\
 &\hspace{10mm}-\operatorname{tr}A\operatorname{tr}(BC)
  -\operatorname{tr}B\operatorname{tr}(AC)
  -\operatorname{tr}C\operatorname{tr}(AB)\}.
\end{align*}
\begin{proof}
Haar integration of three fundamental matrix entries is the
rank-one invariant projector $\epsilon\otimes\epsilon/6$. Contract its
two epsilon tensors with $A,B,C$; the six permutation terms give the
displayed polarized determinant. The code checks this against the direct
epsilon contraction with Gaussian-integer-compatible complex matrices.
\end{proof}
The first nonconstant character contributions are
\[
 I_t=1+3r^2\sum_{i\ne j}\operatorname{tr}(A_iA_j^\dagger)
       +27r^3\{\mathcal T+\overline{\mathcal T}\}
       +O(e^{-17t/6}).
\]
Therefore $-\log I_t$ contains the joint baryonic triality term at order
$r^3$, and also the square of the pair term at order $r^4$. Renormalizing
three separate face times does not reproduce this action. This executes
the nonabelian block behind the source's valence-three obstruction.
\par\noindent\textit{Inherited inputs:} \ref{p10:thm:canonical-haar-heat-representative}, \ref{p10:prop:multiface-heat-blocking}, \ref{p12:app:valence-three-heat-projectivity-obstruction}.


% Source Paper VII-A, original line 8583; inherited proof bodies unchanged.
\subsubsection{D22: The complete three-face source Hessian is now evaluated}
\label{p10:sec:v2201-derived-d22}
\noindent\textit{Status: executed scoped calculation.}

Put $Z_3(t)=\int K_t(U)^3dU$. For Hermitian traceless infinitesimal
sources $X_i$ at $A_i=I$, the exact Hessian of $-\log I_t$ is
\[
 \boxed{\mathcal H_{ij}(X,Y)=
 k(t)(3\delta_{ij}-1)\operatorname{tr}_{\mathbf3}(XY),\qquad
 k(t)=-\frac1{12}\frac{d}{dt}\log Z_3(t)>0.}
\]
\begin{proof}
A common right source shift is absorbed into $U$, so the
three-by-three Hessian has zero row sums. Face permutation and adjoint
symmetry leave one scalar. Schur orthogonality makes the second derivative
on one face $Z_3'/6$ times $\operatorname{tr}(X^2)$ before taking
the logarithm: $\operatorname{tr}_R(X^2)
=d_RC_2(R)\operatorname{tr}_{\mathbf3}(X^2)/4$ and
$Z_3'=3\int(\partial_tK_t)K_t^2$. This gives the formula.
\end{proof}
Moreover, the character expansion is a positive sum
\[
 Z_3(t)=\sum_{R,S,T}d_Rd_Sd_T\,dim\operatorname{Inv}(R\otimes S\otimes T)
 e^{-t(C_R+C_S+C_T)/2}.
\]
Thus $k=\langle C_R+C_S+C_T\rangle/24$ and
$k'=-\operatorname{Var}(C_R+C_S+C_T)/48<0$.
The common eight-dimensional gauge direction is the kernel; the Hessian
has rank sixteen on the relative source space. This rank is unrelated to
the cosmological K16 label.

\begin{center}\small\begin{tabular}{@{}rrr@{}}\toprule Group heat time $t$ & $Z_3(t)$ & $k(t)$\\\midrule
2.000000000 & 10.600683275211 & 0.191261055113\\
3.000000000 & 2.308243464644 & 0.077910519060\\
4.420998862 & 1.159159223915 & 0.015934397826\\
5.753295847 & 1.025768004130 & 0.002827495980\\
6.000000000 & 1.018473138968 & 0.002036816511\\
\bottomrule\end{tabular}\end{center}
The last two named heat times are the existing optional $D_4$ and isotropic
face presentations. They are not selected physical transfer times.
The numerical computation retains complete $(p,q)$ representations through
$p,q\le12$, uses exact zero-mode Weyl-torus quadrature for that character
polynomial, and bounds the omitted compact character sum by dimensions.
Different cutoffs/grids reproduce $k$ to $6\times10^{-17}$ absolute.
The analytic representation-tail bound is reported separately from
floating-point error; machine rounding is not interval certified.
\par\noindent\textit{Inherited inputs:} \ref{p10:thm:v185-heat-hessian-inversion}, \ref{p10:cor:v185-common-heat-face-stiffness}, \ref{p12:prop:v2158-d4-voronoi-constant-curvature}.

% END GEOMETRY-FIRST CONSEQUENCES v2.201

% BEGIN FRAMEWORK v2.220: EXACT HURWITZ BARYON SPIN--FLUX OPERATOR


\section{Strong clock, spectral gap, confinement, and area-law obligations}
\label{p10:sec:v237-a4}
Finite spectral and Schur results are separated from a physical uniform Yang--Mills mass gap, strict string tension or generated dimensional normalization. These stronger claims require their own phase, state and limit certificates.

% Source Paper VII-A, original line 2364; inherited proof bodies unchanged.
% Former logical part: Strong Clock, Mass Gap, and Confinement Closure; integrated into the new thematic argument.
\label{p10:part:strong-clock-gap-confinement}

\noindent\textbf{Part orientation.}
This part defines the strong-scale, gap, flux, and confinement computation targets of the finite kernel.  Exact geometric and kinematic statements are separated from physical spectral/area-law conclusions, which require the explicit source, finite-action, pairing, and uniform-limit hypotheses printed below.


% Source Paper VII-A, original line 2370; inherited proof bodies unchanged.
\subsection{Strong clock and dimensional transmutation}
\label{p10:sec:strong-clock}

The perturbative clock is retained as a conditional comparison calculation.  It is not used to manufacture a native nonperturbative energy from unresolved mass coordinates.

\begin{definition}[Strong inverse clock]
\label{p10:def:strong-inverse-clock}
On a fixed \(n_f\)-flavor perturbative interval,
\begin{equation}
\label{p10:eq:strong-inverse-clock}
\tau_{\rm QCD}(\mu):=\alpha_s(\mu)^{-1},\qquad
\frac{d\tau_{\rm QCD}}{d\log\mu}=\frac{33-2n_f}{6\pi}+O(\alpha_s).
\end{equation}
\end{definition}

\begin{theorem}[Born--Hopf asymptotic-freedom comparison]
\label{p10:thm:born-hopf-asymptotic-freedom}
For \(SU(3)_c\) with \(n_f\) Dirac fundamentals,
\begin{equation}
\label{p10:eq:one-loop-alpha-beta-derived}
\mu\frac{d\alpha_s}{d\mu}=-\frac{33-2n_f}{6\pi}\alpha_s^2+O(\alpha_s^3).
\end{equation}
Hence the one-loop comparison branch is asymptotically free for \(n_f\le6\).  This standard coefficient is fixed by \(C_A=3\) and \(T_F=1/2\); its numerical integration remains conditional on the chosen seed and physical threshold map.
\end{theorem}
\begin{proof}
The standard one-loop coefficient is \(\beta_0=(11/3)C_A-(4/3)T_F n_f=11-2n_f/3\), giving the displayed equation.
\end{proof}

\begin{lemma}[Strong-running dependency boundary]
\label{p10:lem:strong-running-import-chain}
The finite color carrier fixes the representation coefficient of the beta function, while the source-native common-heat branch fixes the boundary gauge normalization at its own matching datum.  Numerical threshold running to a named external scale remains conditional on the typed scale/RG map and on \(M_q^{\rm phys}\); the historical Paper--V-C exponent \(t\) is only a broad-branch calibration chart and is not a second source owner.
\end{lemma}
\begin{proof}
This follows directly from the upstream status theorems summarized in Section~\ref{p10:sec:imported-interface}.
\end{proof}

\begin{lemma}[\(E_{36}\) alternating adjoint block is geometric, not a clock identity]
\label{p10:lem:e36-alternating-adjoint-block}
The quotient \(2O/2T\) has symmetric and alternating one-dimensional quotient lines, and tensoring the alternating line with \(C\otimes C^*\) contains an adjoint color block.  This representation statement does not identify \(\sqrt2E_{36}^{\rm mix}\) with a zero of a QCD inverse clock.
\end{lemma}
\begin{proof}
The quotient-line decomposition and \(\mathbf1\oplus\mathbf8\) color decomposition are algebraic.  A clock zero additionally depends on a running functional and threshold spectrum; those data are not determined by the quotient representation.
\end{proof}

\begin{corollary}[Defect provenance of the alternating quotient line]
\label{p10:cor:v188-alternating-defect-provenance}
The alternating line in Lemma~\ref{p10:lem:e36-alternating-adjoint-block} is
the nontrivial real character of
\[
 2O/[2O,2O]=2O/2T\cong C_2.
\]
Thus its sign/parity has parent provenance in the restricted $M_8$ order-swap
defect.  This strengthens the representation owner only; it supplies neither
a QCD inverse-clock zero nor a threshold/RG map.
\end{corollary}

\begin{lemma}[Projection separation of strong and photon running]
\label{p10:lem:color-singlet-running-no-double-counting}
The singlet and adjoint projectors in \(C\otimes C^*=\mathbf1\oplus\mathbf8\) are orthogonal.  Consequently a color-adjoint running block and a color-singlet electromagnetic current block are distinct operator projections and may not be identified term by term.
\end{lemma}
\begin{proof}
Orthogonal irreducible projectors satisfy \(P_{\bf1}P_{\bf8}=0\).
\end{proof}

\begin{lemma}[Shell determinant versus Ward-current determinant is a choice of observable]
\label{p10:lem:native-shell-vs-ward-current-determinant}
A smooth threshold primitive such as \(\log(Q^2+m^2)\) and a Ward-subtracted current--current determinant are different functions because they encode different observables.  The finite carrier alone does not select one as the unique physical strong-running kernel; that choice must come from the QRG/observable map.
\end{lemma}
\begin{proof}
The two functions are derivatives of different source deformations: one differentiates a color-curvature shell response, the other a conserved abelian current correlator subject to a Ward subtraction.  Gauge representation and current conservation do not identify these source functionals.  Therefore equality can hold only after an additional operator identity is proved.
\end{proof}

\begin{definition}[Legacy smooth mixed-clock diagnostic, typed only as a coordinate test]
\label{p10:def:legacy-mixed-clock-diagnostic}
To test the v1.6 numerical identity without promoting the Paper--V-B coordinates to physical masses, define
\begin{equation}
\label{p10:eq:legacy-mixed-clock-coordinate-ledger}
\begin{gathered}
M_Z^{\rm coord}=91.1872058\ \mathrm{GeV},\\
(m_u,m_d,m_s,m_c,m_b)_{\rm coord}\\
=(0.002156185,0.004700044,0.093377822,1.264775,4.207890)\ \mathrm{GeV},\\
E_{36}^{\rm mix}=0.219578897269190\ \mathrm{GeV},\quad
\deltageom=\frac1{6\sqrt2}.
\end{gathered}
\end{equation}
For $0<Q<m_b^{\rm coord}$ set
\begin{equation}
\label{p10:eq:legacy-mixed-clock-explicit}
\begin{split}
\tau_{32}^{\rm diag}(Q)
={}&\deltageom^{-1}+\frac{11}{2\pi}\log\frac{Q}{M_Z^{\rm coord}}
-\frac1{6\pi}\sum_{q=u,d,s,c,b}
\log\frac{Q^2+(m_q^{\rm coord})^2}{(M_Z^{\rm coord})^2+(m_q^{\rm coord})^2},\\
\mathcal R_{36}^{\rm diag}(Q)
={}&\frac3{4\pi}\log\frac{Q^2+(E_{36}^{\rm mix})^2}{(m_b^{\rm coord})^2+(E_{36}^{\rm mix})^2},\\
\tau_c^{\rm diag}(Q)={}&\tau_{32}^{\rm diag}(Q)+\mathcal R_{36}^{\rm diag}(Q).
\end{split}
\end{equation}
This definition reproduces the old smooth function solely for falsification/replay.  It is not the physical QRG clock because the coordinate masses in \eqref{p10:eq:legacy-mixed-clock-coordinate-ledger} are not yet eigenvalues of $M_q^{\rm phys}$.
\end{definition}

\begin{theorem}[Conditional mixed-clock diagnostic and corrected root]
\label{p10:thm:native-strong-shell-transport-e36-wall}
For the explicit diagnostic of Definition~\ref{p10:def:legacy-mixed-clock-diagnostic}, \(\tau_c^{\rm diag}(Q)\) is strictly increasing on \(0<Q<m_b^{\rm coord}\) and has the unique roots
\begin{equation}
\label{p10:eq:corrected-mixed-clock-root}
\begin{gathered}
\tau_c^{\rm diag}(Q_0)=0:\quad
Q_0^{\rm diag}=0.3158981634\ldots\ \mathrm{GeV},\\
\tau_c^{\rm diag}(Q_1)=1:\quad
Q_1^{\rm diag}=0.5510277102\ldots\ \mathrm{GeV}.
\end{gathered}
\end{equation}
At the former claimed values \(\sqrt2E_{36}^{\rm mix}=0.3105314545\ldots\,\mathrm{GeV}\) and \(\sqrt6E_{36}^{\rm mix}\), the same clock does not satisfy the corresponding zero/one equations; in particular \(\tau_c^{\rm diag}(\sqrt2E_{36}^{\rm mix})\simeq-0.03006\).  Since the inserted masses are not yet a physical mass operator, both roots are coordinate diagnostics, not native QCD scales.
\end{theorem}
\begin{proof}
Differentiating \eqref{p10:eq:legacy-mixed-clock-explicit} with respect to $\log Q$ gives
\[
\frac{d\tau_c^{\rm diag}}{d\log Q}
=\frac{11}{2\pi}
-\frac1{3\pi}\sum_q\frac{Q^2}{Q^2+(m_q^{\rm coord})^2}
+\frac3{2\pi}\frac{Q^2}{Q^2+(E_{36}^{\rm mix})^2}.
\]
Because every fraction lies in $[0,1]$, the derivative is bounded below by $11/(2\pi)-5/(3\pi)=23/(6\pi)>0$.  Hence each level equation has at most one root.  Direct interval evaluation gives sign changes around the quoted decimals, and bisection/Newton evaluation yields the displayed roots.  The typing conclusion follows from the unresolved operator map \eqref{p10:eq:qrg-quark-mass-operator-needed}.
\end{proof}

\begin{theorem}[Comparison one-loop strong handoff]
\label{p10:thm:derived-perturbative-strong-handoff}
On the explicit seed/coordinate branch of Definition~\ref{p10:def:legacy-mixed-clock-diagnostic}, the sharp threshold-matched one-loop comparison clock gives
\begin{align}
\tau_c^{\rm sharp}
&=\deltageom^{-1}
+\frac{23}{6\pi}\log\frac{m_b^{\rm coord}}{M_Z^{\rm coord}}
+\frac{25}{6\pi}\log\frac{m_c^{\rm coord}}{m_b^{\rm coord}},\\
\label{p10:eq:sharp-one-loop-handoff}
\muQCD
&=m_c^{\rm coord}\exp\!\left[-\frac{2\pi}{9}\bigl(\tau_c^{\rm sharp}-1\bigr)\right]
=0.284353696978\ldots\ \mathrm{GeV}.
\end{align}
It is retained only as a comparison coordinate: the formula uses the conditional \(t=1/2\) seed member and the pre-operator mass-coordinate ledger.
\end{theorem}
\begin{proof}
Integrate Definition~\ref{p10:def:strong-inverse-clock} with \(n_f=5\) from \(M_Z^{\rm coord}\) to \(m_b^{\rm coord}\), then with \(n_f=4\) to \(m_c^{\rm coord}\).  Below charm the one-loop coefficient is \(9/(2\pi)\); solving \(\tau(\muQCD)=1\) gives \eqref{p10:eq:sharp-one-loop-handoff}.  Substitution of the displayed coordinate ledger yields the stated decimal.
\end{proof}

\begin{theorem}[Heavy-quark threshold decoupling for a supplied mass operator]
\label{p10:thm:heavy-quark-threshold-decoupling}
Once the actual G0 quark-source member has been transported to the relevant scale as \(M_q^{\rm phys}(\mu)\), ordinary perturbative decoupling may be applied at its heavy eigenvalues with the standard matching conditions.  No Paper--V-B coordinate is used as a physical threshold in place of that inherited operator.
\end{theorem}
\begin{proof}
Once a self-adjoint physical mass operator is available, its spectral projectors give the heavy and light subspaces.  In any smooth determinant/resolvent threshold function, a mode of mass $m$ contributes powers of $Q^2/m^2$ for $Q\ll m$ and approaches the massless coefficient for $Q\gg m$.  The statement is therefore a consequence of the supplied operator spectrum, not of the pre-operator coordinate packet.
\end{proof}

\begin{theorem}[Handoff-scale scheme separation]
\label{p10:thm:handoff-scale-scheme-separation}
The quantities \(\muQCD\), \(Q_0^{\rm diag}\), \(E_{36}^{\rm mix}\), and any eventual nonperturbative spectral scale are logically distinct: the first two are conditional running diagnostics, the third is a geometric energy coordinate, and the last must be computed from the physical finite Hamiltonian or a proved continuum limit.  No equality among them is asserted without an explicit theorem.
\end{theorem}
\begin{proof}
Each displayed scale is defined by a different equation or source map.  Equality of two such scales would require that both defining equations have the same solution.  Since no theorem supplies that identity, and the mixed-clock replay explicitly separates two of them, they must remain distinct typed quantities.
\end{proof}

\begin{proposition}[Dimensional-transmutation readout]
\label{p10:prop:dimensional-transmutation-readout}
For a fully specified renormalization scheme, common-heat matching-scale identification, and threshold operator, the usual RG invariant scale is a deterministic functional of those inputs.  Before the scale map and threshold operator are typed, dimensional transmutation is a computation target rather than a parameter-free Paper--VII-A prediction.
\end{proposition}
\begin{proof}
For a specified beta function $d\alpha/d\log\mu=\beta(\alpha)$ and a specified boundary value, separation of variables gives $\log(\mu/\mu_0)=\int_{\alpha_0}^{\alpha(\mu)}d a/\beta(a)$.  A conventional transmutation scale is the integration constant of this equation.  Changing the scheme, thresholds, or boundary member changes that constant, proving the stated conditionality.
\end{proof}

\begin{theorem}[No count-derived absolute nonperturbative scale]
\label{p10:thm:native-nonperturbative-scale-ratio}
The finite 24-cell counts and Hodge eigenvalues determine dimensionless geometric invariants but do not fix an absolute QCD energy.  An absolute nonperturbative scale must arise from the physically typed Hamiltonian/source map or from a proved dimensional-transmutation branch; the former \(24/15\)-based scale is therefore withdrawn.
\end{theorem}
\begin{proof}
The ratio \(24/15\) is dimensionless and Paper~I-A supplies no amplitude map for it.  Multiplying a candidate dimensionful Hamiltonian normalization by a positive constant leaves those counts unchanged, so the counts alone cannot determine the scale.
\end{proof}

\begin{corollary}[Exact geometric scalar-shell spectrum and moments]
\label{p10:cor:exact-minimal-scalar-shell-moments}
All purely geometric scalar-shell/Hodge eigenvalues and multiplicities remain exact finite outputs.  Their conversion to an energy spectrum is performed only after a physical Hamiltonian normalization has been supplied.
\end{corollary}
\begin{proof}
The Hodge operator is a fixed finite Hermitian matrix.  Its eigenvalues, multiplicities, traces, and all polynomial spectral moments are therefore obtained by finite spectral calculus and are independent of any dimensional Hamiltonian normalization.
\end{proof}

\begin{theorem}[Positive nonperturbative scale requires a positive spectral computation]
\label{p10:thm:positive-native-nonperturbative-scale-floor}
A strictly positive finite QCD scale or gap follows exactly when the corresponding generalized eigenvalue/Rayleigh infimum of the compiled Hamiltonian is strictly positive.  Positivity is therefore a finite matrix certificate, not a consequence of hierarchy-rung absence or a count ratio.
\end{theorem}
\begin{proof}
For a finite self-adjoint Hamiltonian the smallest excitation on a specified invariant subspace is the Rayleigh minimum of $H-E_0$.  It is strictly positive exactly when that minimum is positive.  Geometry can specify the subspace but cannot determine this dimensional Rayleigh quotient without the Hamiltonian normalization.
\end{proof}

\begin{corollary}[Scale, gap, and confinement are distinct spectral functionals]
\label{p10:cor:single-shell-origin-scale-gap-confinement}
A single geometric shell may enter all three calculations, but their numerical values are respectively an energy normalization, a spectral difference, and a flux/area response.  Equality or a universal common coefficient must be proved from the compiled operator and is not assumed.
\end{corollary}
\begin{proof}
The three quantities are defined by different maps: a dimensional normalization/readout, a spectral infimum, and an extended-flux/area limit.  Equality among them is not an algebraic consequence of sharing a geometric band.  Hence each must be computed or bounded separately.
\end{proof}


% Source Paper VII-A, original line 2568; inherited proof bodies unchanged.
\subsection{Finite mass-gap, color projection, and confinement targets}
\label{p10:sec:finite-mass-gap-floor}

\begin{theorem}[Hierarchy-rung absence does not bound a dynamical infrared spectrum]
\label{p10:thm:desert-shell-infrared-cascade-block}
The absence of an additional protected hierarchy rung does not imply a positive QCD spectral gap.  On a fixed finite carrier the family \(K_\varepsilon=P_0+(1-\varepsilon)(1-P_0)\), \(\varepsilon\downarrow0\), preserves all carrier labels while its excitation gap tends to zero.
\end{theorem}
\begin{proof}Direct diagonalization gives transfer eigenvalues \(1\) and \(1-\varepsilon\).  No new hierarchy label is introduced.\end{proof}

\begin{theorem}[Exact finite color-singlet projector]
\label{p10:thm:exact-finite-singlet-projector}
For any finite unitary \(SU(3)_c\) representation \(U_L\),
\begin{equation}
\Ps^{(L)}=\int_{SU(3)}U_L(g)\,dg
\end{equation}
is the orthogonal projector onto the exact color-singlet subspace.
\end{theorem}
\begin{proof}Haar averaging is self-adjoint, idempotent, and fixes precisely the invariant vectors.\end{proof}

\begin{corollary}[Nonzero triality is removed by the singlet readout]
\label{p10:cor:nonzero-triality-no-asymptotic-ray}
A state with nonzero \(\ZZ_3\) triality is annihilated by \(\Ps^{(L)}\).  Thus it is absent from the asymptotic state space if physical asymptotic readout is defined by the exact singlet projector.
\end{corollary}
\begin{proof}
For a center element $z=e^{2\pi i/3}$, a triality-$r$ state transforms by $z^r$.  Haar averaging over $SU(3)$ includes averaging over the center, and $1+z^r+z^{2r}=0$ for $r\not\equiv0\pmod3$.  Thus the singlet projector annihilates nonzero-triality states.
\end{proof}

\begin{theorem}[Kinematic color confinement as a readout obstruction]
\label{p10:thm:kinematic-color-confinement}
The exact singlet projector gives a rigorous kinematic obstruction to every
isolated color representation with no trivial subrepresentation.  This
representation-theoretic fact is distinct from a dynamical area law or mass
gap.
\end{theorem}
\begin{proof}
For any finite unitary representation \(R\), Haar averaging is the orthogonal
projector onto \(R^{SU(3)}\).  Hence \(P_{\rm sing}|_R=0\) exactly when
\(R\) contains no trivial summand.  Nonzero triality is one sufficient test,
by Corollary~\ref{p10:cor:nonzero-triality-no-asymptotic-ray}, but it is not a
necessary one: the adjoint \(\mathbf8\) has zero triality and still contains no
singlet.  The theorem is only a readout obstruction and makes no statement
about a dynamical area law.
\end{proof}

\begin{corollary}[An isolated gluon wave is not an isolated physical singlet ray]
\label{p10:cor:v245-isolated-gluon-no-singlet-ray}
For one gluon, \(R=\mathbf8\) is a nontrivial irreducible adjoint
representation, so
\[
 \boxed{P_{\rm sing}|_{\mathbf8}=0.}
\]
Thus the covariant free gluon modes of
Theorem~\ref{p10:thm:v245-free-gluon-octet-clock-lock} exist as harmonic
connection modes, while no isolated one-gluon state lies in the physical
singlet readout space.
\end{corollary}

\begin{corollary}[Lowest free projected pure-glue threshold]
\label{p10:cor:v245-free-two-gluon-singlet-threshold}
On the free quadratic branch,
\[
 \mathbf8\otimes\mathbf8
 =\mathbf1\oplus\mathbf8_s\oplus\mathbf8_a
  \oplus\mathbf{10}\oplus\overline{\mathbf{10}}\oplus\mathbf{27}.
\]
The lowest one-gluon coexact level has \(n=2\) and energy \(2\Eper\).  The
symmetric invariant tensor \(\delta_{ab}\) allows two bosons in the same
lowest spatial/helicity mode to form the singlet.  Therefore the lowest free
singlet two-gluon threshold is
\[
 \boxed{E_{gg,{\rm min}}^{(0),{\rm sing}}=4\Eper.}
\]
This is neither the interacting Yang--Mills mass gap nor a glueball mass; both
belong to the full projected interacting Hamiltonian.
\end{corollary}

\begin{theorem}[Center-triality superselection]
\label{p10:thm:center-triality-superselection}
The center \(Z(SU(3))\cong\ZZ_3\) grades tensor products by triality; singlets have total triality zero.
\end{theorem}
\begin{proof}
The center of $SU(3)$ acts as $z$ on $\mathbf3$, $z^{-1}$ on $\bar{\mathbf3}$, and trivially on the adjoint.  Tensor-product center characters multiply, so triality is additive modulo three.  Invariants must carry the trivial center character.
\end{proof}

\begin{corollary}[Integer asymptotic baryon number]
\label{p10:cor:integer-baryon-number}
On quark/antiquark singlet sectors, \(n_q-n_{\bar q}\equiv0\pmod3\), so \(B=(n_q-n_{\bar q})/3\in\ZZ\).
\end{corollary}
\begin{proof}
Singletness implies $n_q-n_{\bar q}\equiv0\pmod3$ by the center-triality theorem.  Dividing this integer by three therefore gives an integer baryon number.
\end{proof}

\begin{theorem}[Integer hadronic electric charge]
\label{p10:thm:integer-hadronic-electric-charge}
With the standard quark electric charges, every color-singlet quark composite of integer baryon number has integer total electric charge after the usual flavor assignments are imposed.
\end{theorem}
\begin{proof}
All quark charges are congruent to $-1/3$ modulo integers, while antiquark charges have the opposite congruence.  Hence the total charge modulo one is $-(n_q-n_{\bar q})/3$.  Singlet triality makes the numerator divisible by three, so the total charge is integral.
\end{proof}

\begin{theorem}[Vector flavor-charge conservation]
\label{p10:thm:vector-flavor-charge-conservation}
On the vectorlike QCD branch with diagonal mass/source operators, the corresponding vector flavor charges commute with the strong Hamiltonian.
\end{theorem}
\begin{proof}
With diagonal vectorlike mass and source terms, independent global phase rotations of each conserved vector flavor preserve the finite QCD Hamiltonian.  Their finite one-parameter unitary groups have unique self-adjoint generators by Paper~III-A; differentiating the Hamiltonian-invariance relation gives $[H,Q_f]=0$ as the internal finite Noether-I/commutator consequence.  No separate conservation premise is inserted.  Axial charges are excluded from this statement.
\end{proof}

\begin{theorem}[Pure-glue center obstruction to screening]
\label{p10:thm:pure-glue-center-obstruction-screening}
Adjoint gluonic excitations have zero center triality and cannot screen a source of nonzero \(\ZZ_3\) triality to a singlet using gluons alone.
\end{theorem}
\begin{proof}
The adjoint representation is center-trivial.  Tensoring any number of adjoint factors therefore leaves the center character of an external source unchanged.  A source of nonzero triality cannot be screened to the trivial representation by gluons alone.
\end{proof}


% BEGIN INLINE SOURCE: Paper7A-WilsonSectorScope-Main-v1_0.tex
% Framework v2.187: essential sector scope for the positive Wilson-area target.
\begin{definition}[Sector-indexed Wilson-area and confinement readouts]
\label{p10:def:v2187-wilson-sector-scope}
The finite singlet/triality restrictions and finite Wilson-source values are
readouts of the complete finite QCD owner.  A separate claim of strictly
positive asymptotic fundamental string tension names a pure-glue or otherwise
certified unscreened static-source sector, its complete action/measure, and
its ordered Wilson limits.  The heat-action APBI theorem is a sufficient
route on that stated sector; APBC is its square-face specialization, and a
direct same-owner multiface/flux estimate is an alternative.

The pure-glue center obstruction does not extend to a source sector containing
dynamical fundamental quarks merely by retaining the same color count.
Its complete static-source ground-state problem can include static-light
singlet channels.  Consequently positive asymptotic fundamental tension is
not a universal antecedent of finite full-QCD singlet spectra or a definition
of every native confinement readout.  Appendix~\ref{p10:app:v2187-readout-domain-comparisons}
records the conditional variational screening test; no actual screening or
vanishing-tension result is imported from it.
\end{definition}

% END INLINE SOURCE: Paper7A-WilsonSectorScope-Main-v1_0.tex


\begin{theorem}[Polyakov-loop center criterion]
\label{p10:thm:polyakov-center-criterion}
In a thermal/infinite-volume regime in which center symmetry is well defined, a vanishing fundamental Polyakov-loop expectation is an order-parameter criterion for unbroken center symmetry.  This is a conditional continuum/thermodynamic criterion, not a finite-carrier proof of confinement.
\end{theorem}
\begin{proof}
In a center-symmetric thermal measure, the Polyakov loop in the fundamental transforms by a nontrivial center phase.  Invariance of the measure then forces its expectation value to equal that phase times itself, hence to vanish.  The usual free-energy interpretation requires the stated thermodynamic setting.
\end{proof}

\begin{theorem}[Finite spectral discreteness and exact finite gap functional]
\label{p10:thm:finite-spectral-discreteness}
Every finite self-adjoint \(H_L\) has a finite real spectrum.  Order its eigenvalues with multiplicity as \(E_{0,L}\le E_{1,L}\le\cdots\) and define
\begin{equation}
\label{p10:eq:finite-gap-functional}
\gamma_L:=E_{1,L}-E_{0,L}\ge0
\end{equation}
when the block has dimension at least two.  Thus \(\gamma_L>0\) is exactly the statement that the finite ground level is simple and separated from the first excitation; a continuum mass-gap claim additionally requires a uniformly positive, correctly scaled refinement limit.
\end{theorem}
\begin{proof}
A finite self-adjoint matrix is unitarily diagonalizable with finitely many real eigenvalues.  Ordering them gives the displayed difference.  It is positive exactly when the lowest eigenvalue has multiplicity one and the next level lies strictly above it.  Finite positivity alone gives no uniform-in-refinement lower bound.
\end{proof}

\begin{theorem}[Positive finite Euclidean transfer family]
\label{p10:thm:reflection-positive-transfer-representative}
For every finite self-adjoint \(H_L\) bounded below and every \(t>0\), \(\widehat K_L(t)=e^{-t(H_L-E_0)}\) is positive self-adjoint with operator norm one.
\end{theorem}
\begin{proof}
Functional calculus gives eigenvalues $e^{-t(E_n-E_0)}\in(0,1]$ for $t>0$.  Therefore the exponential is positive self-adjoint, its largest eigenvalue is one, and its operator norm is one.
\end{proof}

\begin{theorem}[Dynamical transfer components are spectral/algebraic outputs]
\label{p10:thm:finite-heat-kernel-gap}
The invariant components of \(\widehat K_L(t)\) are obtained from its reducing projectors or, more strongly, from the Wedderburn decomposition of \(\mathfrak A_{L,{\rm str}}\).  Finite positivity does not force a single connected component.
\end{theorem}
\begin{proof}
Reducing projectors are precisely orthogonal projectors commuting with the transfer and the chosen generated algebra.  The finite Wedderburn decomposition enumerates them algebraically.  Positivity of an operator places no restriction preventing several reducing blocks, as direct sums of positive operators show.
\end{proof}

\begin{corollary}[Correct Perron--Frobenius gap criterion]
\label{p10:cor:finite-pure-glue-pf-gap-derived}
If \(\widehat K_L(t)\) is positivity improving on a closed generating cone of a finite block, then its spectral radius is a simple eigenvalue with a strictly positive eigenvector and the next eigenvalue has smaller modulus.  Without positivity improvement/irreducibility, positivity and graph connectivity alone do not imply this conclusion.
\end{corollary}
\begin{proof}
A positivity-improving finite operator maps every nonzero vector in the positive cone to its interior.  The finite Perron--Frobenius theorem then makes the spectral radius simple with a strictly positive eigenvector; all other eigenvalues have smaller modulus.  The counterexamples in the audit show why mere positive semidefiniteness is insufficient.
\end{proof}

\begin{theorem}[Finite plaquette--Hodge curvature Hessian]
\label{p10:thm:finite-plaquette-hodge-curvature-hessian}
For the oriented 24-cell boundary complex, the incidence matrices satisfy \(\partial^2=0\).  The independently reconstructed face Hodge Hessian has the exact spectrum
\begin{equation}
\begin{split}
1^{(6)}\oplus2^{(16)}\oplus\left(\frac{9-\sqrt{17}}2\right)^{(9)}\oplus4^{(16)}\oplus5^{(16)}\\
\oplus\left(\frac{9+\sqrt{17}}2\right)^{(9)}\oplus8^{(13)}\oplus9^{(1)}\oplus10^{(8)}\oplus12^{(2)}.
\end{split}
\end{equation}
In particular its smallest eigenvalue is exactly one.
\end{theorem}
\begin{proof}
Direct construction of the oriented vertex-edge, edge-face, and face-cell incidence matrices gives the displayed finite spectrum.  The multiplicities sum to 96, the number of triangular faces.
\end{proof}

% BEGIN S3 LITERATURE INTEGRATION v3.040 [VII-A: GEOMETRIC EDGE-CURL]
\begin{theorem}[Exact native Hurwitz edge--curvature spectrum]
\label{p10:thm:v340-s3lit-hurwitz-edge-spectrum}
On the oriented unit-incidence boundary complex of the native Hurwitz
$24$-cell, let $d_0:C^0\to C^1$ and $d_1:C^1\to C^2$ be the
$24\to96\to96$ coboundaries, with $d_1d_0=0$.  The geometric edge
curl--curl operator $M_1=d_1^*d_1$ has the exact spectrum
\[
\begin{split}
 \operatorname{Spec}(M_1)={}&0^{(23)}\oplus1^{(6)}\oplus2^{(16)}
 \oplus\Bigl(\frac{9-\sqrt{17}}2\Bigr)^{(9)}\\
 &\oplus4^{(12)}\oplus5^{(16)}
 \oplus\Bigl(\frac{9+\sqrt{17}}2\Bigr)^{(9)}
 \oplus8^{(4)}\oplus9^{(1)}.
\end{split}
\]
Consequently $\ker M_1=\operatorname{im}d_0$,
$\operatorname{rank}M_1=73$, $\lambda_{\min}^+(M_1)=1$, and
\[
 \operatorname{pdet}M_1=2^{88}3^2 5^{16}.
\]
This is the unit-incidence \emph{spatial geometric} operator and is not
the compiled full electromagnetic heat Hessian or a count of photon species.
\end{theorem}
\begin{proof}
For the explicitly oriented integer incidence matrix one checks the
exact annihilating polynomial
\[
 P(M_1)=0,\qquad
 P(t)=t(t-1)(t-2)(t-4)(t-5)(t-8)(t-9)(t^2-9t+16).
\]
Its integer traces $\operatorname{Tr}M_1^k$, $k=0,\ldots,7$, are
\[
 (96,288,1440,8352,53280,363168,2600352,19337760).
\]
As $M_1$ is real symmetric, the distinct real roots of $P$ contain its
spectrum; the traces and equal multiplicities of the two conjugate
quadratic roots determine the stated multiplicities.  The native
$24$-cell boundary is $S^3$, with $H^1(S^3)=0$, so
$\ker d_1=\operatorname{im}d_0$ and $\dim\operatorname{im}d_0=23$.
Multiply all nonzero eigenvalues, including the conjugate pair whose
product is $16$, for the displayed pseudodeterminant.
\end{proof}

\begin{corollary}[Positive-face-weight spectral bridge]
\label{p10:cor:v340-s3lit-weighted-face-bound}
On the same complex, assume a \emph{separately action-selected} local
Abelian spatial face-quadratic form is represented as
$M_W=d_1^*Wd_1$ for a positive face matrix satisfying
$0<w_-I\preceq W\preceq w_+I$.  Then
\[
 \ker M_W=\operatorname{im}d_0,\qquad
 w_-\lambda_i^+(M_1)\le\lambda_i^+(M_W)
 \le w_+\lambda_i^+(M_1),\quad 1\le i\le73,
\]
and
\[
 w_-^{73}\operatorname{pdet}M_1\le
 \operatorname{pdet}M_W\le
 w_+^{73}\operatorname{pdet}M_1.
\]
No positivity or local face-weight factorization of the \emph{full}
clock/contact/nonlocal photon Hessian is inferred from geometry.
\end{corollary}
\begin{proof}
Quadratic forms give $w_-M_1\preceq M_W\preceq w_+M_1$.
Strict positivity of $W$ makes their kernels equal; apply min--max
on the common orthogonal complement and multiply the $73$ positive
spectral bounds.
\end{proof}
% END S3 LITERATURE INTEGRATION v3.040 [VII-A]


\begin{theorem}[Exact geometric Hodge coercivity and physical stiffness functional]
\label{p10:thm:bilinear-curvature-heat-time}
Let \(\Delta_{2,H}\) be the positive face-Hodge operator of Theorem~\ref{p10:thm:finite-plaquette-hodge-curvature-hessian}.  Then
\begin{equation}
\langle F,\Delta_{2,H}F\rangle\ge\|F\|^2
\end{equation}
for every face field \(F\).  For any compiled Yang--Mills quadratic form \(Q_{\YM,L}\), define its exact generalized stiffness relative to the geometric form by
\begin{equation}
\kappa_{\YM,L}:=\inf_{F\ne0}\frac{Q_{\YM,L}[F]}{\langle F,\Delta_{2,H}F\rangle}\ge0.
\end{equation}
Then \(Q_{\YM,L}[F]\ge\kappa_{\YM,L}\|F\|^2\).  The number \(\kappa_{\YM,L}\) is computed from the physical finite matrix; geometry alone does not set it to a prescribed coupling-dependent value.
\end{theorem}
\begin{proof}
The spectrum in Theorem~\ref{p10:thm:finite-plaquette-hodge-curvature-hessian} lies in $[1,12]$, so the spectral theorem gives $\Delta_{2,H}\ge I$.  The generalized Rayleigh infimum defining $\kappa_{\YM,L}$ then gives $Q_{\YM,L}[F]\ge\kappa_{\YM,L}\langle F,\Delta_{2,H}F\rangle\ge\kappa_{\YM,L}\|F\|^2$.
\end{proof}

\begin{theorem}[Scalar shell as a specified adjoint-curvature spectral subspace]
\label{p10:thm:scalar-shell-primitive-curvature-norm}
Any scalar adjoint-curvature shell used in Paper~VII-A is defined by an explicit spectral projector \(P_{{\rm sh},L}=1_I(\Delta_{2,H})\) onto a specified Hodge eigenspace/band.  This makes its norm and multiplicity exact geometric data while leaving its physical Hamiltonian response to \(\gamma_{{\rm sh},L}\).
\end{theorem}
\begin{proof}
Because $\Delta_{2,H}$ is finite self-adjoint, the spectral theorem defines $1_I(\Delta_{2,H})$ uniquely for every specified spectral set $I$.  Its rank and induced norm are geometric spectral data.  The Hamiltonian response requires the separate quadratic form and is therefore not fixed by the projector.
\end{proof}

\begin{theorem}[Finite flat-connection quadratic support modulo gauge]
\label{p10:thm:finite-flat-connection-vacuum-support}
The 24-cell boundary is a cell decomposition of \(S^3\), so \(H^1(S^3)=H^2(S^3)=0\).  Consequently the linearized flat-connection sector has no non-gauge harmonic one-form obstruction: a closed infinitesimal connection is exact, and a zero face curvature is pure gauge at the cochain level.  This is a geometric quadratic statement and does not prove uniqueness of the interacting vacuum.
\end{theorem}
\begin{proof}
The cellular cohomology of the 24-cell boundary equals that of $S^3$, so $H^1=H^2=0$.  At the linearized level, zero curvature means the connection cochain is closed; vanishing $H^1$ makes it exact, i.e. an infinitesimal gauge mode.  Nonlinear vacuum uniqueness does not follow from this cohomological statement.
\end{proof}

\begin{lemma}[Scalar-shell Dirichlet projector]
\label{p10:lem:scalar-shell-dirichlet-kernel}
For every isolated finite Hodge eigenvalue or spectral interval \(I\), the projector \(P_{{\rm sh},L}=1_I(\Delta_{2,H})\) is uniquely defined by spectral calculus.  No \(24/15\) amplitude normalization enters this projector.
\end{lemma}
\begin{proof}
This is direct finite spectral calculus: an isolated eigenspace is the range of its Riesz/spectral projector, and a union of isolated eigenvalues is projected by summing those orthogonal projectors.  No dimension ratio enters the construction.
\end{proof}

\begin{theorem}[Hodge spectral completeness, not scalar-shell saturation]
\label{p10:thm:scalar-shell-support-saturation}
The face space is the orthogonal direct sum of all Hodge spectral projectors.  A chosen scalar shell saturates only its own spectral band; it does not exhaust the interacting face space unless its projector is the identity.
\end{theorem}
\begin{proof}
The finite spectral theorem resolves the identity as the orthogonal sum of all spectral projectors of $\Delta_{2,H}$.  The range of one proper projector is a proper orthogonal summand unless that projector equals the identity.
\end{proof}

\begin{lemma}[Computed scalar-shell domination coefficient]
\label{p10:lem:scalar-shell-transfer-domination}
Define
\begin{equation}
\kappa_{{\rm sh},L}:=\inf_{\substack{\psi\in P_{{\rm sh},L}\mathcal H_L\\\|\psi\|=1}}
\langle\psi,(H_L-E_0)\psi\rangle.
\end{equation}
Then \(H_L-E_0\ge\kappa_{{\rm sh},L}P_{{\rm sh},L}\) as a quadratic-form inequality on that shell.  This is exact by definition; \(\kappa_{{\rm sh},L}\) may be zero and is not set by \(\deltageom\).
\end{lemma}
\begin{proof}
For every unit vector in the shell range, the Rayleigh quotient is at least its infimum $\kappa_{{\rm sh},L}$.  Multiplying by the squared norm gives the quadratic-form inequality on that range.  No positivity beyond the definition is asserted when the infimum is zero.
\end{proof}

\begin{theorem}[Uniform scalar-shell gap criterion]
\label{p10:thm:uniform-scalar-shell-transfer-gap}
For a refinement family, a uniform positive shell gap exists exactly when
\begin{equation}
\kappa_{\rm sh,*}:=\inf_L\kappa_{{\rm sh},L}>0.
\end{equation}
Under this checkable condition the shell excitation energies are bounded below by \(\kappa_{\rm sh,*}\); without it no uniform floor follows.
\end{theorem}
\begin{proof}
If $\inf_L\kappa_{{\rm sh},L}=c>0$, every level has shell energy at least $c$.  Conversely, any level-independent lower bound $c>0$ forces each $\kappa_{{\rm sh},L}\ge c$, hence their infimum is positive.  This proves the equivalence.
\end{proof}

\begin{theorem}[Pure Yang--Mills gap is not the full-QCD chiral gap]
\label{p10:thm:pure-ym-gap-not-full-qcd-chiral-gap}
A positive pure-glue spectral gap, if established, does not imply a positive full-QCD gap in the massless chiral limit, where Goldstone modes may appear after the thermodynamic/refinement limit.
\end{theorem}
\begin{proof}
Adding massless quark degrees of freedom changes both the Hilbert space and Hamiltonian.  In a chirally broken thermodynamic limit Goldstone's theorem supplies arbitrarily soft non-singlet modes even if the pure-glue sector remains gapped.  Therefore a pure-Yang--Mills gap cannot imply a full-QCD chiral gap.
\end{proof}

\begin{theorem}[Strong-coupling area-law comparison lemma]
\label{p10:thm:strong-coupling-area-law-import}
Under the hypotheses of a strong-coupling character/cluster expansion, Wilson-loop expectations obey an area-law bound.  Paper~VII-A uses this only as a hypothesis-matched comparison theorem; it does not infer the hypotheses from the 24-cell counts.
\end{theorem}
\begin{proof}
This is an external hypothesis-matched theorem: when the character/cluster expansion converges in its strong-coupling domain, the leading nonzero surface term is proportional to the minimal spanning area and the remainder is controlled.  Paper~VII-A invokes the conclusion only under those imported hypotheses and supplies no converse.
\end{proof}

\begin{theorem}[Separated color readout requires a Wilson connector]
\label{p10:thm:separated-color-readout-wilson-connector}
A gauge-invariant operator containing separated fundamental color sources requires a path-ordered Wilson connector between them.
\end{theorem}
\begin{proof}
Under local gauge transformations, fields at distinct points acquire independent group factors.  A path-ordered parallel transporter transforms with inverse endpoint factors, so inserting it is required to contract the endpoints into a gauge-invariant operator.
\end{proof}

\begin{lemma}[Nonzero center flux is non-flat face support]
\label{p10:lem:nonzero-center-flux-nonflat-face}
A finite loop with nontrivial center holonomy cannot be represented by a globally flat connection on the spanning face complex.
\end{lemma}
\begin{proof}
A globally flat connection on a simply connected spanning complex has trivial holonomy around contractible loops up to gauge.  A loop with nontrivial center holonomy is therefore incompatible with global flatness on its spanning face complex.
\end{proof}

\begin{lemma}[Flux-face spectral activation]
\label{p10:lem:flux-face-scalar-shell-activation}
If a nonzero-center-flux face field has nonzero projection into a chosen Hodge band \(P_I\), then that band is activated.  The statement fixes support, not a universal energy coefficient.
\end{lemma}
\begin{proof}
Orthogonal projection satisfies $P_I F\ne0$ exactly when $F$ has a component in the chosen spectral band.  Thus the band is present in the field decomposition, while its energy remains controlled by the separate Hamiltonian quadratic form.
\end{proof}

\begin{theorem}[Computed facewise flux coercivity]
\label{p10:thm:scalar-shell-contraction-nonzero-nality}
On a specified nonzero-\(N\)-ality flux subspace \(\mathcal F_{N,L}\), define
\begin{equation}
\kappa_{{\rm flux},L}:=
\inf_{\substack{F\in\mathcal F_{N,L}\\F\ne0}}
\frac{Q_{\YM,L}[F]}{\|F\|^2}\ge0.
\end{equation}
This is the exact finite facewise contraction/coercivity coefficient.  A positive bound is obtained by computing/proving \(\kappa_{{\rm flux},L}>0\), not by inserting a shell-count amplitude.
\end{theorem}
\begin{proof}
The quotient defining $\kappa_{{\rm flux},L}$ is the Rayleigh infimum of the positive quadratic form on the specified flux subspace.  Therefore $Q_{\YM,L}[F]\ge\kappa_{{\rm flux},L}\|F\|^2$ for every such $F$.  Strict positivity is an additional computed property, not a consequence of the definition.
\end{proof}

\begin{theorem}[Generated position area and the remaining string-tension calculation]
\label{p10:thm:hurwitz-reciprocal-area-string-tension}
The retained Hurwitz embedding transports chains to the reciprocal-position
complex, where the inherited metric supplies the area
\(A_{g_x}(\Sigma)=\int_\Sigma dA_{g_x}\) and the same position connection
supplies Wilson holonomy.  As made explicit in
Theorem~\ref{p10:thm:v164-position-area-factorization}, no monolithic
\(\mathrm{PAIR}_{xp}^{\rm op}\) intertwiner is a live prerequisite.
On a specified pure-glue or otherwise certified unscreened static-source
sector, positive asymptotic fundamental string tension still needs the
appropriate same-owner nonlinear extensive/ordered-limit bound and
dimensional/RG realization; finite face counts or center flux alone do not
supply it.  The domain is Definition~\ref{p10:def:v2187-wilson-sector-scope}.
\end{theorem}
\begin{proof}
The fixed embedding intertwines the chain boundary operators.  Position
area and Wilson holonomy are then defined from the already inherited
position metric and connection on those transported chains.  This is the
factorization, not a universal linear momentum-to-position operator map.
The remaining positive tension is an extensive energy/area statement and
is not implied by the existence of those observables.
\end{proof}

\begin{lemma}[Flux-sector form lower bound]
\label{p10:lem:flux-sector-form-lower-bound}
For every finite flux sector, \(Q_{\YM,L}[F]\ge\kappa_{{\rm flux},L}\|F\|^2\) by the definition of \(\kappa_{{\rm flux},L}\).  The inequality is nontrivial precisely when the computed coefficient is positive.
\end{lemma}
\begin{proof}
For nonzero $F$ in the flux sector, the definition of $\kappa_{{\rm flux},L}$ gives $Q_{\YM,L}[F]/\|F\|^2\ge\kappa_{{\rm flux},L}$.  Multiplication by $\|F\|^2$ yields the displayed inequality.
\end{proof}

\begin{theorem}[Static-potential floor from generated area, flux coercivity, and extensivity]
\label{p10:thm:hurwitz-static-potential-floor}
On a specified pure-glue or certified unscreened static-source sector, let \(C_L\) be a sequence of Wilson contours and assume the full admissible energy domain is covered by the stated flux bounds.  Use (i) the generated position area \(A_x(C_L)\) in the stated dimensional realization; suppose (ii) \(\kappa_{\rm flux,*}:=\inf_L\kappa_{{\rm flux},L}>0\); and (iii) every admissible flux representative spanning \(C_L\) obeys an extensive support bound \(\|F_L\|^2\ge c_A A_x(C_L)\) with \(c_A>0\) independent of \(L\).  Then
\begin{equation}
E_L(C_L)\ge \kappa_{\rm flux,*}c_A A_x(C_L),
\end{equation}
so any convergent static-potential limit has string-tension lower bound \(\sigma\ge\kappa_{\rm flux,*}c_A>0\).  Without all three ingredients Paper~VII-A makes no physical static-potential floor claim.
\end{theorem}
\begin{proof}
Assumption (ii) and the flux form bound give $E_L(C_L)\ge\kappa_{\rm flux,*}\|F_L\|^2$.  Assumption (iii) then gives $E_L(C_L)\ge\kappa_{\rm flux,*}c_A A_x(C_L)$.  Dividing by area and taking any limit permitted by (i) proves the stated lower bound.
\end{proof}

\begin{corollary}[No confinement theorem from support alone]
\label{p10:cor:no-strong-coupling-assumption-native-confinement}
Exact singlet/triality projection gives kinematic color confinement.  A dynamical unscreened area-law claim uses the generated position-area/Wilson structure and still requires a same-owner positive flux coercivity and an extensive spanning-flux/cluster estimate.  These claim-scoped obligations are not antecedents of every finite full-QCD singlet readout.
\end{corollary}
\begin{proof}
The singlet projector and center character establish only which states can be physical.  Theorem~\ref{p10:thm:hurwitz-static-potential-floor} exhibits the additional quantitative ingredients required for an area-law lower bound.  Omitting any of them leaves a counterfamily with the same support data and no uniform tension.
\end{proof}

\begin{theorem}[Finite mass-gap/confinement computation map and continuum handoff]
\label{p10:thm:mass-gap-confinement-closure-map}
Paper~VII-A exports the exact finite Hamiltonians, their commutant decomposition, the Hodge operators, and the computable quantities \(\gamma_L\), \(\kappa_{\YM,L}\), and \(\kappa_{{\rm flux},L}\).  If an enlarged continuum branch claims a Yang--Mills mass gap, it requires a refinement theorem with \(\inf_L\gamma_L>0\); that claim is not the framework's native carrier gap and is not an active Paper-VII-A prediction.  A claimed positive asymptotic fundamental string tension on a pure-glue or certified unscreened static-source sector uses the generated position-area/Wilson structure and additionally requires a same-owner positive flux coercivity and an area-extensive spanning-flux/cluster limit.  These are claim-specific handoff conditions, not already closed results.
\end{theorem}
\begin{proof}
All listed finite quantities are functions of the compiled matrices and are therefore executable.  A continuum gap would be a uniform-in-$L$ lower-bound statement, while a string tension additionally refers to a typed position-space area and extensive flux.  Neither limit property follows from finite spectral existence; the gap condition is retained only for a branch that independently makes that claim.
\end{proof}


% Source Paper VII-A, original line 6183; inherited proof bodies unchanged.
\subsection{finite vacuum, physical gap, and factorized confinement}
\label{p10:sec:v164-vacuum-gap-confinement}

\begin{theorem}[Generated symmetric finite vacuum]
\label{p10:thm:v164-generated-finite-vacuum}
For every compiled finite Hamiltonian \(H_L\), let \(P_{0,L}\) be its ground
projector and \(d_{0,L}=\operatorname{rank}P_{0,L}\).  On the declared
symmetric vacuum branch,
\[
 \rho_{\beta,L}
 =\frac{e^{-\beta H_L}}{\operatorname{Tr}e^{-\beta H_L}}
 \longrightarrow
 \rho_{0,L}=\frac{P_{0,L}}{d_{0,L}}.
\]
If the positive transfer is positivity improving, \(d_{0,L}=1\).  Thus no
independent finite vacuum selector remains on this branch.  Selection of one
pure broken phase would instead require its actual boundary/source limit.
\end{theorem}

\begin{theorem}[Physical finite transfer gap]
\label{p10:thm:v164-physical-transfer-gap}
If \(a_L\) is the physical duration of one transfer step and \(q_L\) is the
largest vacuum-orthogonal transfer eigenvalue, then
\[
 \boxed{\Delta_L^{\rm phys}=-\frac1{a_L}\log q_L.}
\]
Neither positivity of every dimensionless finite gap nor \(q_L\to1\)
decides a continuum mass gap.  On an enlarged branch that makes that claim,
the required antecedent is \(\inf_L\Delta_L^{\rm phys}>0\), followed by
state/OS/cyclicity transfer.  It is not an unfinished part of the native
carrier-gap theorem.
\end{theorem}

\begin{theorem}[Position-area factorization and center-flux no-go]
\label{p10:thm:v164-position-area-factorization}
The fixed Hurwitz embedding already transports chains by
\[
 \iota_{1\#}:C_q(\mathcal K_H)\cong C_q(\mathcal K_H^x),
 \qquad \partial\iota_{1\#}=\iota_{1\#}\partial.
\]
Once a surface is on the reciprocal-position side, its area is
\(A_{g_x}(\Sigma)=\int_\Sigma dA_{g_x}\), and the same position connection
defines its Wilson holonomy.  No universal global linear \(x\to p\)
intertwiner is required or generated; the old monolithic
\(\mathrm{PAIR}_{xp}^{\rm op}\) premise is retired.

Center flux and quadratic Hodge coercivity still do not imply an area law.
Distributing the fundamental center flux over \(N\) commuting faces gives
\[
 \sum_{f=1}^N\|F_f\|_{\rm HS}^2=\frac{8\pi^2}{3N}\longrightarrow0.
\]
Strict area extensivity must come from the nonlinear dynamics.
\end{theorem}

\begin{theorem}[Ordered Wilson limits]
\label{p10:thm:v164-ordered-wilson-limits}
For every fixed compact rectangle, Wilson convergence is the bounded-cylinder
instance of QFT4.  The physical order is
\[
\begin{gathered}
 \omega(W(C_{R,T}))=\lim_{L\to\infty}\omega_L(W_L(C_{R,T})),\\
 V(R)=-\lim_{T\to\infty}\frac1T\log\omega(W(C_{R,T})),\quad
 \sigma=\lim_{R\to\infty}\frac{V(R)}R.
\end{gathered}
\]
The local-flat chart exhaustion supplies cofinal rectangles with diverging
metric area, so no independent `WCONV' or `SURF' premise remains.  Positive
transfer and spatial reflection inequalities give existence of \(V(R)\) and
of the asymptotic slope \(\sigma\) on the reconstructed static sector.  The
strict-positivity problem is reduced below to an action-preserving bounded-
incidence presentation theorem; no additional momentum--position pairing or
area selector is required.  Numerical evaluation and unit/RG matching remain
separate readouts.
\end{theorem}

\begin{theorem}[Triality-resolved heat-surface majorant on a bounded-incidence presentation]
\label{p10:thm:v2155-triality-branch-core-majorant}
Work on the canonical isotropic \(SU(3)\) face-heat member fixed by
Corollary~\ref{p10:cor:v185-common-heat-face-stiffness}, with
\[
 t_H=h_{SU(3)}^{-1}\!\left(\frac1{2\pi^2}\right)
 =5.7532958466\ldots .
\]
Suppose a same-action reciprocal-position face presentation of the finite heat
measure has square faces, at most six incident faces at every interior edge,
and at most three unused boundary edges on a child square after attachment
through one parent edge.  Let \(\omega=e^{2\pi i/3}\) and define the exact
triality projectors
\[
 K_t^{[r]}(g):=\frac13\sum_{j=0}^{2}\omega^{-rj}K_t(\omega^jg),
 \qquad r\in\mathbb Z_3.
\]
Then the complete nonzero-triality heat semigroup satisfies
\[
 \bigl\|T_{t_H}^{[\pm1]}\bigr\|
 =u:=e^{-2t_H/3}=0.02158988061\ldots,
\]
while the nontrivial neutral sector satisfies
\[
 \bigl\|T_{t_H}^{[0]\ominus\mathbf1}\bigr\|
 =v:=e^{-3t_H/2}=1.7867469\times10^{-4}.
\]
For the area fugacity \(z=11/10\), the neutral-decoration recursion admits
\(N_*=1/500\) as a strict supersolution.  Writing \(y=1+N_*\), the complete
triality-compatible rooted charged recursion
\[
 x=uz(y+B)^3,
 \qquad
 \Phi_5(x;y)=\frac{(y+2x)^5-(y-x)^5}{3}
\]
admits \(B_*=9/20\) as a strict supersolution:
\[
 \Phi_5\!\left(uz(y+B_*)^3;y\right)
 =0.4318926829\ldots <\frac9{20}.
\]
Consequently every connected center-triality surface polymer generated by the
full Peter--Weyl heat expansion, including neutral decorations, branchings,
closed representation loops and handles, is absolutely majorized by
\[
 C_\partial^{P(C)}\,z^{-A_{\rm cell}(C)}
\]
for a finite perimeter constant \(C_\partial\).  In particular the cell-area
surface free energy is strictly positive and obeys
\[
 \sigma_{\rm cell}\ge\log(11/10)=0.0953101798\ldots .
\]
\end{theorem}
\begin{proof}
The center average projects the Peter--Weyl expansion onto the three exact
triality sectors.  Since the smallest charged Casimir is \(C_F=4/3\), the
operator norm of the complete charged heat sector is exactly \(e^{-2t/3}\);
no truncation to the fundamental representation is used.  The smallest
nontrivial neutral Casimir is \(C_A=3\), giving the displayed neutral norm.
Normalized Haar integration is the orthogonal projector onto the invariant
tensor subspace and hence has norm one.  Degree-one vacuum branches vanish,
while degree-two corridors compose by the heat semigroup and therefore cost
\(u^\ell\) for corridor length \(\ell\).

At an entered square face, the five remaining incident-face slots carry
triality \(0,\pm1\).  The root-of-unity filter gives exactly
\(\Phi_5(x;y)=[(y+2x)^5-(y-x)^5]/3\).  The stated values of \(N_*\) and \(B_*\)
are strict supersolutions of the neutral and charged monotone recursions.
Closing a charged or neutral corridor replaces an operator-norm factor by a
trace-class contraction.  The existing heat-tail certificate gives
\(q_{\pm}(1)<0.197\) and \(q_0^*(1)<0.012\), so Schatten--H\"older bounds show
that cutting a cycle or handle open can only increase the majorant.  The
rooted-tree recursion therefore dominates the complete connected surface
sum.  The extra factor \(z^{A_{\rm cell}}\) in the convergent majorant yields
the displayed strictly positive area exponent.
\end{proof}


\begin{theorem}[Finite-type triality branch-core criterion]
\label{p10:thm:v2158-finite-type-triality-branch-core}
Consider a same-action finite face-heat presentation for which the rooted
surface exploration has only finitely many local incidence states.  For every
face type $\alpha$ let its heat time be $t_\alpha>0$ and define
\[
 u_\alpha=e^{-2t_\alpha/3},\qquad
 v_\alpha=e^{-3t_\alpha/2}.
\]
At each rooted incidence state, assign triality $0,\pm1$ to the finitely many
available incident-face slots.  Weight a charged face of type $\alpha$ by
$u_\alpha z$ times the charged/neutral majorants on its unused boundary
states, weight a nontrivial neutral face by $v_\alpha z$ times the neutral
majorants on those states, and impose the exact $\mathbb Z_3$ root-of-unity
projection at the entered edge.  This produces a finite-dimensional monotone
polynomial map
\[
 \mathcal F_z:\mathbb R_+^{2r}\longrightarrow\mathbb R_+^{2r},
 \qquad X=(N_1,\ldots,N_r,B_1,\ldots,B_r),
\]
where the $N_i$ bound rooted neutral decorations and the $B_i$ bound rooted
charged branches.

Assume that for some $z>1$ there is a componentwise strict supersolution
$X_*>0$,
\[
 \boxed{\mathcal F_z(X_*)<X_*},
\]
and that closing every charged or nontrivial-neutral corridor which can occur
in the finite type set is a trace-class contraction of norm strictly below
one.  Then every connected nonzero-triality surface polymer generated by the
complete Peter--Weyl heat expansion is absolutely majorized by
\[
 C_\partial^{P(C)}z^{-A_{\rm cell}(C)}
\]
for a finite perimeter constant $C_\partial$.  Hence
\[
 \boxed{\sigma_{\rm cell}\ge\log z>0.}
\]
The square-face theorem
Theorem~\ref{p10:thm:v2155-triality-branch-core-majorant} is the one-state
specialization with five remaining incident-face slots and three unused child
edges.
\end{theorem}
\begin{proof}
The center average decomposes every face heat kernel into the exact three
triality sectors.  Charged and neutral nontrivial sectors are bounded by the
stated $u_\alpha$ and $v_\alpha$, while normalized Haar integration at an edge
is an orthogonal projector of norm one.  Because the rooted incidence set is
finite, the local surface exploration is dominated by the monotone recursion
$X\mapsto\mathcal F_z(X)$.  Starting from zero, monotonicity and
$\mathcal F_z(X_*)<X_*$ keep every rooted tree contribution below $X_*$.  A
closed channel or handle is obtained by contracting an open corridor; the
trace-class hypothesis makes that contraction no larger than the open
majorant.  Thus the rooted-tree bound dominates the complete connected
surface expansion.  Removing the proof fugacity gives one factor $z^{-1}$ per
surface face and hence the claimed cell-area exponent.
\end{proof}

\begin{corollary}[Native Hurwitz triangle/ribbon branch certificate]
\label{p10:cor:v2158-hurwitz-two-type-triality}
For the native one-step Hurwitz transfer complex, an interior spatial edge has
three spatial triangular faces and two temporal ribbons, while an interior
temporal edge has eight temporal ribbons.  A triangle entered through a
spatial parent edge leaves two spatial child edges.  A ribbon entered through
a spatial parent leaves one spatial and two temporal child edges, while a
ribbon entered through a temporal parent leaves two spatial and one temporal
child edge.

Let $t_s,t_\tau>0$ be the spatial-triangle and temporal-ribbon heat times,
respectively, and put
\[
 u_s=e^{-2t_s/3},\quad v_s=e^{-3t_s/2},\qquad
 u_\tau=e^{-2t_\tau/3},\quad v_\tau=e^{-3t_\tau/2}.
\]
For $y_s=1+N_s$, $y_t=1+N_t$, define
\[
 n_T=v_szy_s^2,\qquad
 n_{R,s}=v_\tau z y_s y_t^2,\qquad
 n_{R,t}=v_\tau z y_s^2y_t,
\]
\[
 \mathcal N_s=
 \max\!\left\{
 (1+n_T)^2(1+n_{R,s})^2-1,
 (1+n_T)^3(1+n_{R,s})-1
 \right\},
 \qquad
 \mathcal N_t=(1+n_{R,t})^7-1.
\]
For the charged branches put
\[
 x_T=u_s z(y_s+B_s)^2,
\]
\[
 x_{R,s}=u_\tau z(y_s+B_s)(y_t+B_t)^2,
 \qquad
 x_{R,t}=u_\tau z(y_s+B_s)^2(y_t+B_t),
\]
and
\[
 \Psi_{p,q}(x_1,x_2;y)
 =\frac{(y+2x_1)^p(y+2x_2)^q-(y-x_1)^p(y-x_2)^q}{3}.
\]
Then
\[
 \mathcal B_s=\max\{\Psi_{2,2}(x_T,x_{R,s};y_s),
                       \Psi_{3,1}(x_T,x_{R,s};y_s)\},
\]
\[
 \mathcal B_t=
 \frac{(y_t+2x_{R,t})^7-(y_t-x_{R,t})^7}{3}.
\]
On the explicitly labelled equal-heat comparison specialization
$t_s=t_\tau=t_H$ and $z=1.1$, the point
\[
 (N_s,N_t,B_s,B_t)=(0.001,0.002,0.20,0.41)
\]
is a strict supersolution:
\[
 \mathcal N_s=0.0007891572042\ldots<0.001,
 \qquad
 \mathcal N_t=0.001382123174\ldots<0.002,
\]
\[
 \mathcal B_s=0.1960993513\ldots<0.20,
 \qquad
 \mathcal B_t=0.4052908673\ldots<0.41.
\]
Thus the native Hurwitz triangle/ribbon incidence creates no local branching
obstruction to the heat-surface mechanism.  This corollary does not identify
$t_s$ with $t_\tau$ on the physical branch and does not turn the fixed compact
Hurwitz regulator into a macroscopic position-area net.
\end{corollary}
\begin{proof}
The incidence counts are those of the regular $24$-cell transfer complex.
The two alternatives in $\mathcal N_s$ and $\mathcal B_s$ correspond to
entering a spatial edge through a triangle or through a ribbon.  The formulas
then follow by the same neutral product and $\mathbb Z_3$ root-of-unity filter
used in Theorem~\ref{p10:thm:v2155-triality-branch-core-majorant}.  Direct
evaluation at the displayed specialization gives the strict inequalities.
\end{proof}

\begin{corollary}[APBI reduction of positive asymptotic string tension]
\label{p10:cor:v2158-apbi-string-tension-reduction}
On a specified pure-glue or certified unscreened fundamental Wilson sector,
assume for its complete measure on one selected heat-action branch an
\emph{action-preserving bounded-incidence presentation} (APBI): along a cofinal fixed-UV local-flat
family the same nonlinear heat action admits finite face presentations with a
uniform finite set of rooted incidence types, physical face areas bounded by
$A_+<\infty$, and one level-independent strict supersolution of
Theorem~\ref{p10:thm:v2158-finite-type-triality-branch-core} at some $z>1$.
Then
\[
 |\omega(W(C))|
 \le C_0\exp\!\left(c_\partial P_{g_x}(C)
 -\frac{\log z}{A_+}A_{g_x}(\Sigma_C)\right),
\]
and
\[
 \boxed{
 V(R)\ge \frac{\log z}{A_+}R-O(1),\qquad
 \sigma\ge\frac{\log z}{A_+}>0.}
\]
APBC is the square-face/valence-six sufficient specialization with $z=1.1$.
\end{corollary}
\begin{proof}
The finite-type theorem supplies $z^{-A_{\rm cell}}$.  Since every face has
physical area at most $A_+$,
$A_{g_x}(\Sigma)\le A_+A_{\rm cell}(\Sigma)$.  The ordered Wilson limits then
give the stated static-potential and asymptotic-tension inequalities.
\end{proof}

\begin{corollary}[APBC reduction of positive asymptotic string tension]
\label{p10:cor:v2155-apbc-string-tension-reduction}
On the same pure-glue or certified unscreened Wilson sector and complete
measure, assume on the canonical heat branch an \emph{action-preserving bounded-
incidence cubulation} (APBC): the generated reciprocal-position action admits a
cofinal same-action local-flat square-face presentation satisfying the
incidence hypotheses of
Theorem~\ref{p10:thm:v2155-triality-branch-core-majorant}, with physical face
areas uniformly bounded above by \(A_+<\infty\) on the fixed-UV exhaustion.
Then
\[
 |\omega(W(C))|
 \le C_0\exp\!\left(c_\partial P_{g_x}(C)
 -\frac{\log(11/10)}{A_+}A_{g_x}(\Sigma_C)\right),
\]
and the ordered static-potential limit of
Theorem~\ref{p10:thm:v164-ordered-wilson-limits} obeys
\[
 V(R)\ge \sigma_*R-O(1),
 \qquad
 \sigma\ge\sigma_*:=\frac{\log(11/10)}{A_+}>0.
\]
Thus APBC closes the square-face specialization.  The broader physical
frontier is the APBI condition of
Corollary~\ref{p10:cor:v2158-apbi-string-tension-reduction}; no additional
color, Wilson, area, center-flux or representation-theory premise remains.
\end{corollary}
\begin{proof}
Theorem~\ref{p10:thm:v2155-triality-branch-core-majorant} gives the cell-area
bound.  Since \(A_{g_x}(\Sigma)\le A_+A_{\rm cell}(\Sigma)\), the cell exponent
implies the displayed physical-area exponent.  Taking the regulator limit at
fixed \((R,T)\), then \(T\to\infty\) and finally \(R\to\infty\), gives the
static-potential inequality and positive asymptotic slope.
\end{proof}

\begin{remark}[APBI is a same-action geometric bridge, not a lattice postulate]
\label{p10:rem:v2155-apbc-firewall}
\label{p10:rem:v2158-apbi-firewall}
Theorems~\ref{p10:thm:v2155-triality-branch-core-majorant} and
\ref{p10:thm:v2158-finite-type-triality-branch-core} are theorems on
same-action face presentations satisfying their stated incidence and heat
bounds.  They do not infer those bounds merely from the existence of a
local-flat chart.  In particular
Proposition~\ref{p10:prop:multiface-heat-blocking} continues to forbid
replacing exact four-dimensional blocking by a product of one-face heat
factors, while Proposition~\ref{p10:prop:v2158-heat-subdivision-area-nogo}
forbids manufacturing an extensive area fugacity by semigroup subdivision of
one fixed heat face.  APBI must therefore be proved on a genuinely extensive
fine same-owner action family.  APBC remains one sufficient square-face
specialization.
\end{remark}

\begin{theorem}[Native mass-gap import and continuum boundary]
\label{p10:thm:v165-native-gap-continuum-boundary}
The framework-native mass gap is the Paper-V-B carrier-support theorem
\[
 \mathbf1_{(E_{56},E_{40})}(H_{\rm car})=0,
 \qquad E_{48}=\sqrt{E_{40}E_{56}},
 \qquad P_{48}^{\rm act}=0.
\]
It is derived from the GP1--GP3 active carrier and does not depend on the QCD
refinement sequence.  By contrast,
\(\Delta_L^{\rm phys}=-a_L^{-1}\log q_L\) is a finite QCD diagnostic.  The
implication
\[
 \inf_L\Delta_L^{\rm phys}>0
 \Longrightarrow \text{positive reconstructed continuum gap}
\]
is retained only as a conditional handoff.  The active framework makes no
positive continuous Yang--Mills mass-gap claim, so this antecedent is not a
mandatory unfinished prediction of Paper~VII-A.
\end{theorem}



\section{Chiral, CP and topological source response}
\label{p10:sec:v237-a5}
Finite chiral source jets and charge weights are computed on their stated domains. The ordered thermodynamic/refinement massless limit, source-normalized correlators and physical susceptibility are not inferred from support counts.

% Source Paper VII-A, original line 2956; inherited proof bodies unchanged.
% Former logical part: Chiral and Anomaly Readout Closure; integrated into the new thematic argument.
\label{p10:part:native-nonperturbative-readouts}


% Source Paper VII-A, original line 2959; inherited proof bodies unchanged.
\subsection{Chiral symmetry, the required order of limits, and source readouts}
\label{p10:sec:chiral-symmetry-breaking}

\begin{proposition}[Minimal scalar singlet order parameter]
\label{p10:prop:minimal-scalar-singlet-order-parameter}
For three massless flavors the standard scalar bifundamental order parameter is
\(\Phi^i{}_j\sim\bar q_{R,j}q_L^i\), transforming as \((\mathbf3,\bar{\mathbf3})\) under \(SU(3)_L\times SU(3)_R\).  A vector-preserving condensate is proportional to \(I_3\).
\end{proposition}
\begin{proof}
A left-handed triplet and right-handed triplet transform independently under $SU(3)_L\times SU(3)_R$.  The bilinear $\bar q_R q_L$ therefore transforms as $(\mathbf3,\bar{\mathbf3})$ and is a Lorentz scalar/color singlet.  Its expectation value proportional to the identity leaves the diagonal $SU(3)_V$ subgroup invariant.
\end{proof}

\begin{theorem}[Finite unique-vacuum symmetry theorem]
\label{p10:thm:vector-flavor-vacuum-preservation}
At finite \(L\), if the exact massless Hamiltonian commutes with the connected group \(SU(3)_L\times SU(3)_R\) and has a unique vacuum, that vacuum is invariant under the full group.  Therefore exact spontaneous chiral breaking cannot occur in the same finite unique-vacuum problem.
\end{theorem}
\begin{proof}
A unique vacuum transforms by a one-dimensional continuous character.  A connected semisimple group has only the trivial continuous character, so the vacuum is invariant.
\end{proof}

\begin{theorem}[Nonabelian Goldstone package in the correct order of limits]
\label{p10:thm:conditional-goldstone-package}
Introduce a common positive scalar mass source \(h>0\), coupled to \(\sum_f\bar q_fq_f\), and a refinement/thermodynamic sequence \(L\to\infty\).  If the per-flavor condensate limit
\begin{equation}
\Sigma:=-\lim_{h\downarrow0}\lim_{L\to\infty}\frac{1}{N_fV_L}\,\frac{\partial}{\partial h}\log Z_L(h),\qquad N_f=3
\end{equation}
exists and is nonzero with unbroken subgroup \(SU(3)_V\), then the standard Goldstone theorem gives eight non-singlet Goldstone modes.  Reversing the limits at finite volume gives no spontaneous breaking.
\end{theorem}
\begin{proof}
At every finite level and $h=0$, a unique vacuum cannot break the connected chiral group, as shown by Theorem~\ref{p10:thm:vector-flavor-vacuum-preservation}.  Taking $L\to\infty$ at fixed $h>0$ first allows distinct limiting phases; if the subsequent $h\downarrow0$ limit has stabilizer $SU(3)_V$, the broken-generator count is $\dim SU(3)_L+\dim SU(3)_R-\dim SU(3)_V=8$, and Goldstone's theorem gives eight massless modes.
\end{proof}

\begin{theorem}[Axial-singlet exclusion from the Goldstone count]
\label{p10:thm:axial-singlet-goldstone-exclusion}
When the \(U(1)_A\) anomaly is nonzero in the limiting theory, the singlet axial direction is not an additional exact Goldstone generator; the non-singlet coset has dimension eight.
\end{theorem}
\begin{proof}
The anomalous divergence of the singlet axial current means $U(1)_A$ is not an exact continuous symmetry of the quantum theory on this branch.  Goldstone's theorem counts generators of exact spontaneously broken symmetries, so the anomalous singlet generator is excluded from the octet count.
\end{proof}

\begin{theorem}[Ward normalization fixes charges, not the Goldstone residue]
\label{p10:thm:no-residual-goldstone-current-renormalization}
A Ward identity fixes the normalization of the conserved/partially conserved axial current once the source convention is fixed.  It does not determine the vacuum-to-Goldstone matrix element \(F_\pi\), which is a dynamical residue.
\end{theorem}
\begin{proof}
A Ward identity fixes the normalization of the conserved charge by its action on fields.  The vacuum-to-one-particle residue also depends on the normalized vacuum and eigenstate and therefore is not determined by the charge algebra.  This separates current normalization from $F_\pi$.
\end{proof}

\begin{lemma}[Axial matrix-element definition of \(F_\pi\)]
\label{p10:lem:finite-color-singlet-axial-matrix-element}
In the symmetry-broken limit, define
\begin{equation}
\label{p10:eq:axial-fpi-definition}
\langle0|A_\mu^a(0)|\pi^b(p)\rangle=i\delta^{ab}F_\pi p_\mu .
\end{equation}
The finite approximants to \(F_\pi\) are extracted from the axial two-point residue of the compiled source operator; \(C_F=4/3\) and Ward normalization do not fix its value.
\end{lemma}
\begin{proof}
Lorentz covariance and the pseudoscalar quantum numbers make the one-pion axial-current matrix element proportional to $p_\mu\delta^{ab}$; the proportionality coefficient defines $F_\pi$.  Once the current and states are normalized, the coefficient is uniquely extracted from the corresponding two-point residue.
\end{proof}

\begin{theorem}[Goldstone-tangent axial residue as a computation target]
\label{p10:thm:goldstone-tangent-axial-saturation}
Provided the double limit of Theorem~\ref{p10:thm:conditional-goldstone-package} exists, \(F_\pi\) is the pole residue defined by \eqref{p10:eq:axial-fpi-definition}.  It is a native correlator readout of the limiting QCD source functional, not an independently adjustable low-energy parameter and not a count-derived number.
\end{theorem}
\begin{proof}
Insert a complete set of energy eigenstates in the axial two-point function.  In the broken phase the Goldstone state contributes a pole whose residue is the square of the matrix element in Lemma~\ref{p10:lem:finite-color-singlet-axial-matrix-element}.  Thus $F_\pi$ is a spectral computation target, not an independent normalization.
\end{proof}

\begin{theorem}[GMOR chiral bridge]
\label{p10:thm:gmor-chiral-bridge}
Under the standard chiral-limit hypotheses,
\begin{equation}
F_\pi^2m_\pi^2=-(m_u+m_d)\langle\bar u u+\bar d d\rangle/2+O(m_q^2)
\end{equation}
and equivalently \(m_\pi^2=B_0(m_u+m_d)+O(m_q^2)\) with \(B_0=\Sigma/F_\pi^2\) for the corresponding condensate convention.
\end{theorem}
\begin{proof}
The axial Ward identity relates the divergence of the nonsinglet axial current to the quark-mass pseudoscalar density.  Saturating the Ward identity by the Goldstone pole and expanding to first order in the light masses yields the displayed GMOR relation.  The statement is restricted to the order-of-limits and analyticity hypotheses under which this standard derivation is valid.
\end{proof}

\begin{lemma}[Finite Ward-dual inverse response]
\label{p10:lem:primitive-source-response-duality}
For an invertible finite source Hessian \(D_{J,L}\) on a specified tangent block, the response operator \(R_{J,L}:=D_{J,L}^{-1}\) is uniquely defined.  This is an algebraic inverse, not a prediction of its scalar eigenvalue.
\end{lemma}
\begin{proof}
An invertible finite matrix has a unique inverse.  Defining the response as that inverse fixes the operator algebraically; it does not constrain the eigenvalues of the original Hessian beyond invertibility.
\end{proof}

\begin{lemma}[Symmetry reduction of the scalar chiral susceptibility]
\label{p10:lem:inverse-defect-scalar-spurion-selector}
If the chiral tangent block is an irreducible \(SU(3)_V\) octet and the source Hessian is \(SU(3)_V\)-equivariant, Schur's lemma gives
\begin{equation}
D_{\chi,L}=d_{\chi,L}I_8
\end{equation}
for a real scalar \(d_{\chi,L}\).  Symmetry fixes the matrix shape but does not fix \(d_{\chi,L}=\deltageom\); the scalar is computed from the source two-point function.
\end{lemma}
\begin{proof}
On an irreducible complex $SU(3)_V$ octet, any equivariant endomorphism commutes with the representation.  Schur's lemma makes it a scalar multiple of the identity.  Self-adjointness makes the scalar real, but symmetry provides no equation fixing its numerical value.
\end{proof}

\begin{theorem}[Dynamical chiral scalar coefficient]
\label{p10:thm:native-gmor-scalar-coefficient}
In the chiral limit, define
\begin{equation}
B_0:=\frac{\Sigma}{F_\pi^2}
\end{equation}
when \(F_\pi\ne0\).  Equivalently it may be extracted from the appropriate scalar-source derivative of the low-energy generating functional.  No formula \(B_0=Q/\deltageom\) is implied by finite color symmetry.
\end{theorem}
\begin{proof}
The leading chiral mass term may be written with coefficient $B_0$ or with the condensate $\Sigma$; differentiating the vacuum energy with respect to the mass source identifies $\Sigma=F_\pi^2B_0$ in the stated convention.  Solving gives $B_0=\Sigma/F_\pi^2$.
\end{proof}

% BEGIN OVERLAP-SCALAR-SOURCE-NORMALIZATION v1.65

% Source Paper VII-A, original line 3073; inherited proof bodies unchanged.
\subsubsection{Native overlap scalar-source normalization and Cayley response}
\label{p10:sec:v165-overlap-scalar-source-normalization}

\begin{corollary}[Mass derivative of the massive overlap branch]
\label{p10:cor:v165-overlap-scalar-source-derivative}
On the declared massive vectorlike overlap branch
\[
 D_m=mI+\left(1-\frac{a_Rm}{2}\right)D_{\rm ov},
 \qquad
 D_{\rm ov}=\frac1{a_R}(I+\Gamma_5\epsilon_H),
\]
the native scalar-source derivative is
\[
 \boxed{
 \mathcal O_m^{\rm BH}:=\frac{\partial D_m}{\partial m}
 =I-\frac{a_R}{2}D_{\rm ov}
 =\frac12(I-\Gamma_5\epsilon_H).}
\]
In particular the arbitrary overall overlap normalization coordinate $a_R$
cancels from the native scalar-source insertion.
\end{corollary}

\begin{theorem}[Exact overlap Cayley reduction of the native scalar response]
\label{p10:thm:v2158-overlap-cayley-scalar-response}
On the declared massive vectorlike overlap branch set
\[
 V:=\Gamma_5\epsilon_H,\qquad
 R:=I-\frac{a_R}{2}D_{\rm ov}=\frac12(I-V).
\]
On the finite subspace $\mathcal H_c=(\ker R)^\perp$, define
\[
 D_c:=D_{\rm ov}R^{-1}.
\]
Then
\[
 \boxed{D_m=R(D_c+m)},
 \qquad
 \boxed{D_m^{-1}\frac{\partial D_m}{\partial m}=(D_c+m)^{-1}}
\]
on $\mathcal H_c$.  On $\ker R$ the native scalar insertion vanishes.  If
$V\psi=e^{i\theta}\psi$ with $\theta\not\equiv0\pmod{2\pi}$, then
\[
 \boxed{D_c\psi=i\lambda(\theta)\psi,\qquad
 \lambda(\theta)=\frac{2}{a_R}\cot\frac\theta2.}
\]
The physical Dirac zero therefore corresponds exactly to $V=-I$,
$\theta=\pi$.  The $\Gamma_5$-Hermiticity of $V$ pairs nonreal phases
$e^{\pm i\theta}$, and each pair contributes to the scalar trace
\[
 \boxed{\frac{2m}{\lambda^2+m^2}.}
\]
Thus the current overlap source normalization reduces exactly to the standard
Banks--Casher resolvent kernel before any thermodynamic/chiral limit is taken.
\end{theorem}
\begin{proof}
The identities
$D_m=D_{\rm ov}+m(I-a_RD_{\rm ov}/2)=D_{\rm ov}+mR$ and
$[R,D_{\rm ov}]=0$ give $D_m=R(D_c+m)$ on $\mathcal H_c$.  Since
$\partial_mD_m=R$, multiplication by the inverse gives the second boxed
identity.  For an eigenphase of $V$,
\[
 D_c=\frac{2}{a_R}\frac{1+e^{i\theta}}{1-e^{i\theta}}
 =i\frac{2}{a_R}\cot\frac\theta2.
\]
The remaining statements follow from
$\Gamma_5V\Gamma_5=V^\dagger$ and the elementary conjugate-pair resolvent
sum.
\end{proof}

\begin{corollary}[Isolated-zero numerator and normalized-ensemble scope]
\label{p10:cor:v2158-finite-volume-zero-mode-firewall}
At fixed finite volume let $n_0$ be the actual number of exact overlap
zero modes and let $N_f$ positive degenerate masses be consumed.
The determinant contains $m^{N_fn_0}$ and their explicit scalar
trace is $n_0/m$, so the isolated-zero numerator is
$O(m^{N_fn_0-1})$. For the unit-index otherwise nondegenerate
branch, $n_0=|\nu|=1$ and $N_f=3$, this is $O(m^2)$.
Suppression in a normalized full ensemble additionally requires a
nonzero massless partition contribution and the corresponding
controlled remaining weight. Conditioning on a sector whose
partition function carries the same determinant power can cancel
that power and leave $n_0/m$. Neither finite-sector singularity
nor finite zero support constitutes a thermodynamic density.
Zero-mode-saturating anomaly insertions remain outside this scalar
numerator statement.
\end{corollary}
\begin{proof}
Each exact zero supplies mass eigenvalue $m$ and unit scalar
derivative, hence one determinant factor per flavor and one
resolvent factor $1/m$. Normalize only after retaining the full
partition function. A restricted sector can cancel its common
mass power; additional paired zeros require the actual $n_0$.
\end{proof}

\begin{remark}[Scope]
\label{p10:rem:v165-overlap-scalar-source-scope}
The corollary fixes the native source derivative on the declared overlap
branch.  It does not identify that finite presentation with
$\overline{\rm MS}$ or any other named continuum subtraction scheme, and it
does not select the physical fermion mass member.
\end{remark}
% END OVERLAP-SCALAR-SOURCE-NORMALIZATION v1.65


% BEGIN INLINE SOURCE: Paper7A-ChiralSourceCertificates-Main-v1_0.tex
% Framework v2.164: finite prefix and conditional chiral response certificates.

% Source Paper VII-A, original line 3179; inherited proof bodies unchanged.
\subsubsection{Native chiral source certificates and their limiting domains}
\label{p10:sec:chiral-source-certificates}
Use the complete packet of
Definition~\ref{p10:def:complete-finite-qcd-source-transfer-datum}.
All primitive spatial structure remains on the fixed compact momentum
carrier, $m_p=1$, as in Definition~\ref{p10:def:fixed-mp-large-mx}.
The extent and clock normalizations below are inherited reciprocal
readouts. No primitive position lattice, additional flavor, fitted
condensate or new interaction coefficient is introduced.

On the stationary massless claim branch in a fixed canonical kinetic
frame put $H_L(h)=H_L-hO_L$, $\mathcal N_L=N_fV_L$, $N_f=3$,
and $p_L(h)=[E_{0,L}(0)-E_{0,L}(h)]/\mathcal N_L$.
Here $V_L$ is the inherited spatial normalization of the ground-energy
claim and $O_L$ the already normalized common scalar source.
Write $\mathcal C_L(h)=p'_{L,+}(h)$ with the chosen positive-probe
orientation. The signed condensate retains its source convention.
Identifying this with a Euclidean overlap response consumes the licensed
clock/transfer normalization and stationary ground-state passage.

\begin{theorem}[Ordered chiral energy/source certificate]
\label{p10:thm:chiral-energy-source-certificate}
For any physical singlet density matrices define
$q_L=\operatorname{Tr}(\rho_LO_L)/\mathcal N_L$ and
$\delta_L=\operatorname{Tr}\{\rho_L[H_L-E_{0,L}(0)]\}/\mathcal N_L$.
Then for $h>0$,
\[
 \mathcal C_L(h)\ge q_L-\delta_L/h.
\]
If $\mathcal N_L\to\infty$, $\liminf_Lq_L\ge q_*>0$ and
$\delta_L\to0$, then
\[
 \boxed{\liminf_{h\downarrow0}\liminf_{L\to\infty}
             \mathcal C_L(h)\ge q_*.}
\]
Existence of an ordinary matched condensate limit remains part of that claim.
\end{theorem}
\begin{proof}
Apply Theorem~\ref{p4:thm:finite-energy-source-certificate}.
For each fixed positive $h$ take the extent first, then lower $h$.
\end{proof}

\begin{theorem}[Positive pressure slope and subextensive excess energy]
\label{p10:thm:chiral-pressure-energy-equivalence}
Assume $\|O_L\|/\mathcal N_L\le B$, a zero-source physical ground with
vanishing scalar expectation, and a full compatible pressure limit
$p(h)=\lim_Lp_L(h)$ on an interval around zero. Then $p$ is convex
and Lipschitz, $p(0)=0$, and
\[
 \Sigma_{\rm mag}:=p'_+(0)=\lim_{h\downarrow0}p(h)/h>0
\]
if and only if physical states on the same sequence have
$\liminf_Lq_L>0$ and $\delta_L\to0$.
The finite constrained energy of
Corollary~\ref{p4:cor:finite-constrained-energy-duality} obeys
$p_L(h)=\sup_q[hq-I_L(q)]$; convergence of $I_L$ is not a premise.
\end{theorem}
\begin{proof}
The uniform source bound gives $B$-Lipschitz pressures. Convexity
passes to the full limit. At a differentiability point $h$, bounding
left and right secants squeeze every finite ground-branch derivative
to $p'(h)$. Choose differentiability points $h_j\downarrow0$;
their limiting expectations approach $p'_+(0)$. In their sourced
ground states,
$\delta_L=h_jq_L-p_L(h_j)$.
The zero-source ground with zero expectation gives $p_L(h_j)\ge0$,
so $0\le\delta_L\le Bh_j$. A diagonal choice with $j$ increasing
sufficiently slowly as $L$ grows yields the required physical states.
The reverse implication follows from the preceding variational bound
and $p(h)\ge hq_*$.
\end{proof}
A bounded source supplies subsequential pressures, not the unique full
compatible limit assumed here. Where the ordinary ordered condensate
exists on the matched branch, it agrees with this slope convention.

\begin{theorem}[Action-weighted shrinking-window lower bound]
\label{p10:thm:weighted-chiral-mode-window}
On the positive vectorlike overlap/Cayley branch retain the actual
mass-dependent action, determinant, state and scalar insertion. Define
the per-flavor measure and positive response by
\[
 \mu_{L,m}(A)=\mathcal V_L^{-1}
       \langle\operatorname{Tr}1_A(-iD_{c,L})\rangle_{L,m},\qquad
 C_L(m)=\int\frac{m}{\lambda^2+m^2}\,d\mu_{L,m}(\lambda).
\]
The spectral normalization $\mathcal V_L$ includes the clock factor
where that presentation consumes it. For every $\alpha,m>0$,
\[
 \boxed{C_L(m)\ge
 \frac{\mu_{L,m}([-\alpha m,\alpha m])}{m(1+\alpha^2)}.}
\]
If $b_\alpha=\liminf_{m\downarrow0}\liminf_L
\mu_{L,m}([-\alpha m,\alpha m])/m>0$, the ordered lower response
limit is at least $b_\alpha/(1+\alpha^2)$.
\end{theorem}
\begin{proof}
The exact scalar insertion in
Theorem~\ref{p10:thm:v2158-overlap-cayley-scalar-response}
gives the positive resolvent kernel. Bound that kernel on the
displayed interval and retain the remaining nonnegative integral.
\end{proof}
Both signs of each paired eigenvalue and exact zeros are counted;
flavor degeneracy is not counted twice. The theorem proves
nonvanishing, with ordinary-limit existence and finiteness still required.
Fixed-$\lambda$ convergence alone does not control the mass-dependent
boundary layer $|\lambda|\sim m$.

\begin{theorem}[Native source-moment sufficient certificate]
\label{p10:thm:chiral-source-moment-certificate}
Let $\Omega_L$ be the unique invariant zero-source ground, so
$\langle O_L\rangle_0=0$. Define
\[
 M_{0,L}=\langle O_L^2\rangle_0,\quad
 M_{1,L}=\langle O_L(H_L-E_{0,L})O_L\rangle_0
       =\tfrac12\langle[O_L,[H_L,O_L]]\rangle_0.
\]
For $M_{0,L}>0$, set
$A_L=\sqrt{M_{0,L}}$, $B_L=\langle O_L^3\rangle_0/M_{0,L}$,
$o_{+,L}=(B_L+\sqrt{B_L^2+4A_L^2})/2$ and
$\Delta_L=M_{1,L}/M_{0,L}$. Then
\[
 \chi_L(0)\ge\frac{2M_{0,L}^2}{\mathcal N_LM_{1,L}},\qquad
 \mathcal C_L(h)\ge\frac{o_{+,L}}{\mathcal N_L}
                 -\frac{\Delta_L}{h\mathcal N_L}.
\]
If $\liminf_LM_{0,L}/\mathcal N_L^2\ge a^2>0$,
$\|O_L\|/\mathcal N_L\le B$, and
$M_{1,L}/(\mathcal N_LM_{0,L})\to0$, the ordered lower response
limit is at least $(\sqrt{B^2+4a^2}-B)/2>0$.
\end{theorem}
\begin{proof}
Spectral expansion gives the double commutator and nonnegative
transition moments. Cauchy--Schwarz applied to
Eq.~\eqref{p10:eq:finite-susceptibility-readout}, divided by
$\mathcal N_L$, gives the susceptibility bound. The physical state
$\psi_L=O_L\Omega_L/\sqrt{M_{0,L}}$ is orthogonal to $\Omega_L$
and has energy excess $\Delta_L$. On their span the source matrix is
$\left(\begin{smallmatrix}0&A_L\\A_L&B_L\end{smallmatrix}\right)$.
Its positive eigenvector has source $o_{+,L}$ and energy excess
between zero and $\Delta_L$. Apply the full-Hamiltonian variational
certificate. Finally $B_L/\mathcal N_L\ge-B$ and
$A_L/\mathcal N_L$ has lower limit at least $a$, while the energy
penalty vanishes at each fixed $h>0$.
\end{proof}
Keep the third moment: three flavors do not imply an assumed scalar
sign-flip symmetry. The stronger sufficient estimates
$M_0\gtrsim\mathcal N^2$, $M_1\lesssim\mathcal N$ give
$\Delta=O(\mathcal N^{-1})$; they have not been established here for
the actual interacting QCD sequence.

For $O_L=\sum_aO_{a,L}$, generated co-consumption retains
\[
 M_0=\sum_{a,b}\langle O_aO_b\rangle_c,\qquad
 M_1=\tfrac12\sum_{a,b}\langle[O_a,[H,O_b]]\rangle.
\]
Theorems~\ref{p2a:thm:v3101-generated-coconsumption} and
\ref{p2a:thm:v3102-preterminal-relation-preservation} preserve
these mixed terms; they do not replace a state expectation by a count.
On the fixed numerical CAR source branch,
Theorem~\ref{p4:thm:car-source-variance} gives $M_0=W+\Xi$.
Under its bounded mode-density and source hypotheses, the positive
quadratic variance premise is exactly a positive quadratic connected
four-point return, by Corollary~\ref{p4:cor:bounded-source-connected-return}.
This is one sufficient route, not a prerequisite of the energy criterion.

The Cartan-twisted defining-doublet kernel, intrinsic momentum Dirac,
multiplication symbol, full overlap Dirac and many-body Hamiltonian
have different roles. The canonical SM anchor has zero color component
and nonzero weak component; its complete affine family has color
fluctuations, by Theorem~\ref{p5:thm:canonical-full-joint-connection}.
Reciprocity with momentum inversion is not automatically axial
$\Gamma_5$. A uniformly protected fixed-rank geometric kernel cannot
supply an extensive mode count. Appendix~\ref{p10:app:chiral-diagnostics}
records that restriction and the conditional trial, free, clock and
anomaly comparisons. They supply no physical QCD state or amplitude.

The least finite energy/source receiver reuses
Corollary~\ref{p2a:cor:least-dynamic-test-algebra} with
$I,H_L,O_L$ and all consumed jets. Its Wedderburn blocks retain the
minimum sourced ground energy, including central energy and actual
physical sectors. This finite construction does not supply the limiting
pressure. The fixed noncritical chart of
Theorem~\ref{p12:thm:v283-local-1pi-transfer} does not certify the
growing critical source limit. A07 remains an open evaluation and
ordered-limit claim on the existing action.

% END INLINE SOURCE: Paper7A-ChiralSourceCertificates-Main-v1_0.tex


% BEGIN INLINE SOURCE: Paper7A-ClaimScopedChiralHierarchy-Main-v1_1.tex
% Framework v2.189: constructor occurrence and ordered source consumption are distinct.
% Stable labels preserve the existing finite/ordered hierarchy references.

% Source Paper VII-A, original line 3371; inherited proof bodies unchanged.
\subsubsection{Finite massive readouts and the ordered chiral claim}
\label{p10:sec:v2186-chiral-claim-hierarchy}

\begin{corollary}[Covered finite massive source readouts precede the ordered chiral limit]
\label{p10:cor:v2186-finite-massive-before-ssb}
Retain the complete finite source-transfer datum of
Definition~\ref{p10:def:complete-finite-qcd-source-transfer-datum}, the actual
mass/source member and the state or transfer prescription consumed by the
readout.  Its finite scalar and pseudoscalar source responses, axial transition
matrix elements, singlet spectral values and Schur residues are licensed
readouts of that same datum.  Their finite evaluation does not consume the
nonzero ordered massless condensate of
Theorem~\ref{p10:thm:conditional-goldstone-package}, an extensive fermion
family, APBI, or a continuum uniqueness theorem.
An interpretation using the exact Goldstone limit or the GMOR chiral bridge
does consume the additional hypotheses printed in those results.
\end{corollary}
\begin{proof}
The master finite closure and source-interface theorems express each covered
quantity as a finite spectral, source-derivative, transition or Schur map.
Corollary~\ref{p10:cor:finite-qcd-parent-readouts} retains the same parent,
member and state through those maps and attaches limiting certificates only
where consumed.  Neither finite diagonalization nor finite source
differentiation invokes a thermodynamic or massless limit.  The additional
limits enter the separately stated Goldstone and GMOR conclusions.
\end{proof}

\begin{definition}[Extensive fermion construction and pressure comparison are separate]
\label{p10:def:v2186-xfe-constructor-pressure-split}
For a claim that actually consumes a cofinal reciprocal-position fermion
family, denote by \texttt{XFE-EXT-CAR} its same-owner constructor component:
finite quark mode spaces $\mathcal K_L^F$, inherited fermionic CAR/Fock
realizations, covariant Dirac operators and normalized scalar, pseudoscalar
and axial source maps. On the existing licensed quark-action/statistics
branch, Theorem~\ref{p10:thm:v2189-native-orbital-quark-constructor} and
Propositions~\ref{p10:prop:v2189-spatial-car-density} and
\ref{p10:prop:v2189-spatial-quark-source-word} construct this component
using the physical rank-$72$ quark fibre, the native complete component
bands, reciprocal homothety and the actual field/dual restriction.
They retain $m_p=1$, prove the stated mode density/free-symbol comparison,
and give exact covariance between the transported embedding fibres with
the inherited adjoint and source jets. They introduce neither a new
physical cutoff nor a spatial lattice. The scalar $\Lambda_F$ here is the
existing frontier comparison parameter; its occurrence is not a selection
of a same-numerical-cutoff spinor frontier.
The density refers to $\dim\mathcal K_L^F/V_L$, not the exponentially
larger dimension of its exterior Fock space.

The temporal Euclidean CAR and positive-mass source components of
Propositions~\ref{p10:prop:v2189-native-clock-car}--
\ref{p10:prop:v2189-native-clock-bloch-limit} are distinct from those
physical equal-time spatial modes. Identification of either constructed
word with the particular finite or ordered VII-A QCD readout must retain
the actual field/dual, mass/source member, gauge/constraint reduction,
central/sea allocation, extent and clock normalization. This comparison
is a physical consumer of the constructed family; it does not reopen the
native bundle, charged source occurrence or CAR construction.
Denote by \texttt{XFE-PRESSURE-COMP} the downstream analytic certificate
for the full compatible massive pressure/source limit of the constructed
same-action family, with the actual extent/clock normalization, state or
measure, source chart, and cross-level or boundary estimates.  For example,
on the stationary source branch used above it certifies the specified limit
$p(h)=\lim_L p_L(h)$ on its claimed source domain.  Finite CAR identities,
finite determinant positivity and bounded subsequential compactness alone
are not that certificate.  Pressure convergence is not an additional quark,
statistics or carrier-occurrence constructor.  This pressure certificate is
consumed by a claim using the pressure route; a direct matched ordered
source-limit theorem can instead supply the limit it actually requires.
In particular, the lower-limit certificates above do not assume a unique full
pressure merely to establish a nonzero lower response.
\end{definition}

For the exact ordered spontaneous-breaking claim, the constructed spatial
owner is available on its stated licensed branch. Its identification with
the consumed QCD source/sea/clock word, matched pressure or direct ordered
source limit, and nonzero action-weighted critical response retain their
own hypotheses. The existing \texttt{A07-RPSUSC} ordered
residual-susceptibility target remains downstream of that matched family.
This ordering leaves every already covered finite massive readout in the
preceding corollary available now. The constructor results do not assert a
full physical pressure limit, condensate value or $\Sigma>0$.

Appendix~\ref{p10:app:v2189-native-spatial-quark-comparison} proves that
fixed-band multiplier covariance and leading density alone cannot identify
the consumed QCD source/sea word. The conditional covariant-spectral-band
comparison of Appendix~\ref{p10:app:v2186-covariant-spectral-band} remains
a mathematical alternative with its printed global/domain/physical-origin
premises. The positive spatial construction above uses the already licensed
component field restriction, rather than treating that conditional
comparison as a carrier-occurrence theorem.

% END INLINE SOURCE: Paper7A-ClaimScopedChiralHierarchy-Main-v1_1.tex


% BEGIN INLINE SOURCE: Paper7A-NativeClockSourceComponent-Main-v1_0.tex
% Framework v2.189: native temporal Euclidean CAR and massive-source components.
% The distinct equal-time spatial quark construction is in the spatial module.

% Source Paper VII-A, original line 3468; inherited proof bodies unchanged.
\subsubsection{Canonical global bundle and native temporal source components}
\label{p10:sec:v2189-native-clock-source-component}

Every physical internal structure remains rooted at $m_p=1$.
The following finite family uses the already generated clock pullbacks;
it introduces no spatial lattice, physical cutoff, coupling or state.
Its CAR role is the exact algebra over the Euclidean quark carrier,
not an identification with physical equal-time many-body quark modes.

\begin{proposition}[Canonical bundle gluing has no independent topology root]
\label{p10:prop:v2189-canonical-global-bundle}
Use the global canonical factor restriction $j_x:K_x\to K_{\rm rel}$
and the generated principal bundle $X_7^{\rm cmp}\to K_{\rm rel}$ of
Corollary~\ref{p3b:cor:canonical-factor-restrictions} and
Theorem~\ref{p3f:thm:canonical-local-gred-pullback}. Then
\[
 P_{{\rm B},\rm glob}:=j_x^*X_7^{\rm cmp}
\]
restricts to the printed $P_{{\rm B},U}$ on every admitted product chart.
Its generated connection, structure-group extensions and associated
fibre-spin bundles restrict compatibly and need no new topology or overlap
datum. This does not identify fibre spin with tangent spin or select a
global scalar-circle spin lift.
\end{proposition}
\begin{proof}
For inclusion $\iota_U$, functoriality gives
$\iota_U^*j_x^*X_7^{\rm cmp}=(j_x\circ\iota_U)^*X_7^{\rm cmp}$,
exactly the local pullback already printed. The same identity on overlaps
proves the gluing cocycle and connection compatibility. Extension and
associated-bundle transport commute with pullback. Their origin is the
already licensed word in
Theorem~\ref{p3f:thm:generated-fibre-spin-joint-operator}.
The tangent-spin factor in its Dirac anchor is a distinct factor; the
round spatial specialization already has its unique spin structure and
self-adjoint elliptic domains, while the scalar-circle lifts keep
Corollary~\ref{p3b:cor:v190-scalar-spin-twists}.
\end{proof}

\begin{proposition}[Compatible native clock pullbacks and exact CAR algebra]
\label{p10:prop:v2189-native-clock-car}
Use the actual repeated three-layer clutched family of
Theorem~\ref{p10:thm:v199-native-clock-covers}, on
\[
 \mathscr H_k=\ell^2(C_{3k}\times V_H)
       \otimes\mathbb C^4_{\rm spin}\otimes\mathbb C^3_{\rm color},
 \qquad \dim\mathscr H_k=864k.
\]
For $k\mid l=nk$, retain compatible spin signs $s_l=s_k^n$ and the
actual bundle pullback. The normalized pullback
\[
 I_{k,l}=n^{-1/2}\pi_{l,k}^*
\]
is an isometry with $H_lI_{k,l}=I_{k,l}H_k$ and
$I_{l,m}I_{k,l}=I_{k,m}$. The same intertwining holds for chirality,
sign, overlap Dirac, the massive scalar-source operator, and every
actually repeated compatible source jet. The maps
$c_k(f)\mapsto c_l(I_{k,l}f)$ give exact injective CAR algebra maps,
and the exterior functor gives compatible Fock isometries.
Periodic spin uses every divisor cover. Fixed antiperiodic spin uses
odd divisor covers; $(2j+1)!!$ is a cofinal tower in their divisor order.
\end{proposition}
\begin{proof}
Each base vertex has exactly $n$ preimages, proving $I^*I=1$ and the
composition rule. Every native clutched stencil/link is the actual pullback
of its base word, so local assembly commutes with $\pi^*$, including the
closing transition. Compatibility of the spin signs is exactly its
holonomy pullback. Hermiticity makes the isometry image reducing;
functional calculus therefore intertwines sign and overlap. The
same calculation applies to the already normalized repeated source word,
including its actual derivatives. Pull back gauge maps and connection
together for exact gauge covariance.
Let $\mathcal J_{k,l}:=\bigoplus_r\bigwedge^r I_{k,l}$, with vacuum sent
to vacuum. Exterior multiplication and the wedge-removal formula give,
on the entire base Fock space,
\[
 c_l^\dagger(I_{k,l}f)\mathcal J_{k,l}
 =\mathcal J_{k,l}c_k^\dagger(f),\qquad
 c_l(I_{k,l}f)\mathcal J_{k,l}
 =\mathcal J_{k,l}c_k(f).
\]
Indeed $\mathscr H_l=I_{k,l}\mathscr H_k\oplus(I_{k,l}\mathscr H_k)^\perp$,
and the Fock image has complementary modes in their vacuum; the inner
products in every annihilator contraction are preserved by $I_{k,l}$.
Thus this is an intertwining of annihilators, not just of creators or a
compression of an unrelated Fock word. The universal finite CAR algebra
map is unital and injective. For any actually compatible one-body source
$A_l I_{k,l}=I_{k,l}A_k$, also
$d\Gamma(A_l)\mathcal J_{k,l}=\mathcal J_{k,l}d\Gamma(A_k)$;
no unproved many-body interaction intertwining is asserted.
Naturality of exterior powers proves composition. Odd cover multiplicities preserve the
antiperiodic sign. Every odd positive integer divides a sufficiently large
odd double factorial, proving its cofinality. No independently replicated
physical spatial volume or additional fermion species has been inserted.
\end{proof}

\begin{proposition}[Uniform massive scalar-source tube on the native covers]
\label{p10:prop:v2189-native-clock-massive-tube}
On the declared positive vectorlike overlap branch retain
\[
 V_k=\Gamma_5\operatorname{sign}H_k,\quad
 R_k=\tfrac12(1-V_k),\quad
 D_{k,m}=a_R^{-1}(1+V_k)+mR_k.
\]
For $m\in[m_0,m_1]\Subset(0,\infty)$, set
$b_0=\min(m_0,2/a_R)>0$. On the complex scalar tube
$|z-m|<b_0/2$,
\[
 \|D_{k,z}^{-1}\|\le2/b_0,\qquad
 \left|k^{-1}\partial_z^q\log\det D_{k,z}\right|
 \le864(q-1)!(2/b_0)^q\quad(q\ge1).
\]
For real positive $m$, the determinant is positive and the sharper first
derivative bound is $864/\min(m,2/a_R)$. All constants are independent
of cover degree. This is a scalar-source component certificate, not
positivity of every local, axial, chiral or phase-crossing source chart.
\end{proposition}
\begin{proof}
Proposition~\ref{p10:prop:v202-vectorlike-positive} gives
$\sigma_{\min}(D_{k,m})\ge\min(m,2/a_R)$, and $\|R_k\|\le1$.
The same native mass derivative $R_k$ is used at every level.
The resolvent identity and Neumann series give the complex-tube bound.
Finite determinant differentiation gives
\[
 \partial_z^q\log\det D_{k,z}
 =(-1)^{q-1}(q-1)!\operatorname{Tr}
                 [(D_{k,z}^{-1}R_k)^q].
\]
Bound the trace by $864k$ times the operator norm. The real branch retains
the actual positive determinant and all actual action/sea normalizations.
No moving spectral projector is frozen and no physical mass is selected.
\end{proof}

\begin{proposition}[Periodic temporal determinant and source limit]
\label{p10:prop:v2189-native-clock-bloch-limit}
On the actually repeated three-layer connection/source word, let
$\exp(i\theta)=s_k$. The native Fourier decomposition gives
\[
 H_k\simeq\bigoplus_{j=0}^{k-1}H(\phi_j),\qquad
 \phi_j=(\theta+2\pi j)/k.
\]
On every compact positive-mass source tube of the preceding proposition,
\[
 \lim_{k\to\infty}k^{-1}\log\det D_{k,z}
 =\frac1{2\pi}\int_0^{2\pi}\log\det D_z(\phi)\,d\phi,
\]
locally uniformly in $z$, together with every fixed scalar derivative.
The two temporal spin signs give the same limit. With the inherited
cover normalization $\mathcal V_{4,k}=k\mathcal V_{4,1}$, the finite
Euclidean matrix density is $864/\mathcal V_{4,1}$.
\end{proposition}
\begin{proof}
The Fourier blocks are those of the actual bundle-cover operator, not
independent replicas. Functional calculus and repeated source assembly
preserve the blocks. The all-phase sign-kernel gap in
Theorem~\ref{p10:thm:v199-native-clock-covers} makes
$H(\phi)$ and its sign continuous on the phase circle. The positive
real determinant fixes its logarithm; the uniform massive tube supplies
the analytic continuation and uniform derivative bounds. Ordinary Riemann
sums converge uniformly on compact source subdomains, and Cauchy
differentiation gives all fixed derivatives. The shifted phase sampling
does not change the integral. Covering degree multiplies the inherited
four-volume exactly. This is not a full interacting gauge pressure or a
stationary ground-energy pressure theorem.
\end{proof}

\paragraph{Scoped outcome.}
The native global bundle, compatible temporal Euclidean CAR algebra,
uniform positive-mass scalar-source tube and periodic temporal determinant
limit are proved above. Physical equal-time quark modes are supplied by the
distinct spatial construction of
Theorem~\ref{p10:thm:v2189-native-orbital-quark-constructor}; temporal
layers are not independent spatial modes. The actual ordered condensate,
residual susceptibility, Goldstone and GMOR obligations retain their stated
hypotheses. Appendix~\ref{p10:app:v2189-native-clock-comparison} proves
that the fixed clock-sector index density produces a divergent massless
scalar response, rather than a finite condensate, and records the exact
source/domain distinctions needed when applying these components.

% END INLINE SOURCE: Paper7A-NativeClockSourceComponent-Main-v1_0.tex


% BEGIN INLINE SOURCE: Paper7A-NativeSpatialQuarkConstructor-Main-v1_0.tex
% Framework v2.189: native spatial quark constructor main module.

% Source Paper VII-A, original line 3651; inherited proof bodies unchanged.
\subsubsection{Native orbital quark-field restriction and its extensive CAR family}
\label{p10:sec:v2189-native-spatial-quark-constructor}

This construction uses the existing licensed clocked quark-action and
fermionic branch. It inherits the one-particle/statistics assignment of
Theorem~\ref{p4:thm:finite-second-quantization} and the claim domain of
Theorem~\ref{p5:thm:two-residual-sm-constitution}; it does not infer physical
statistics from spinor shell dimensions alone. Its source member is the
generated charged packet of
Theorem~\ref{p5:thm:v397-charged-preterminal-srcval-realization}, transported
by Theorem~\ref{p3f:thm:v177-local-clocked-fermion-jet} and
Corollary~\ref{p3f:cor:v1122-root-generated-fermion-member-handoff}.
The actual field, dual field, interaction, constraint and central/sea
allocation are retained through
Theorem~\ref{p3f:thm:v1148-finite-metric-constructor}.

\begin{theorem}[Generated physical orbital quark restriction]
\label{p10:thm:v2189-native-orbital-quark-constructor}
On that branch, use the generated global invariant spin frame on the
reciprocal round slice $S_R^3$. The global connection and associated bundles
are the pullback restrictions of
Theorems~\ref{p3f:thm:canonical-local-gred-pullback} and
\ref{p3f:thm:generated-fibre-spin-joint-operator}; the compact elliptic
domain is the one in Section~\ref{p3f:sec:m8td-spatial-operators}.
Let
\[
 \mathcal Q_+
 =(S_L\otimes W_L\otimes C\otimes G)
 \oplus
 (S_R\otimes(W_{R+}\oplus W_{R-})\otimes C\otimes G).
\]
This is the quark restriction of
\eqref{p5:eq:physical-one-particle-factorization}, so
$\operatorname{rank}\mathcal Q_+=2\cdot18+2\cdot18=72$.
For the existing complete orbital frontier of
Theorem~\ref{p2c:thm:exact-reciprocal-scalar-bands}, define
\[
 \mathcal K_m^q(R)=\mathcal P_x(m)\otimes\mathbb C^{72},\qquad
 (J_{m,R}c)(Rq)=R^{-3/2}\sum_{A,\alpha}
        h_A(q)e_\alpha(q)c_{A\alpha},
\]
where $h_A$ is a complete orthonormal orbital band in $d\mu_1$,
$\operatorname{vol}(S_1^3)=V_\gamma=2\pi^2$, and $e_\alpha$ is the generated
physical quark frame. Equivalently $h_A=f_A/\sqrt{V_\gamma}$ for an
orthonormal band $f_A$ in the probability Haar measure $d\mu_1/V_\gamma$.
Then $J_{m,R}$ is an actual isometric field embedding with
\[
 \boxed{\dim\mathcal K_m^q(R)
 =72D_m,
 \qquad D_m=\frac{(m+1)(m+2)(2m+3)}6.}
\]
The pulled-back external and dual-field word is exactly
\[
 S_{F,m}[\bar c,c;J]
 =S_{F,\rm parent}[\bar cJ_{m,R}^{\dagger},J_{m,R}c;J].
\]
The full same-action finite matrix and every consumed interaction/source
slot are obtained by the finite-metric constructor, rather than by replacing
this line with a determinant or by choosing a mass matrix after truncation.
No independent spinor threshold, spatial lattice, physical mirror sector or
increase of $m_p=1$ occurs.
\end{theorem}
\begin{proof}
The reciprocal spin transport of
Corollary~\ref{p2c:cor:v244-reciprocal-spin-dirac-transport} carries the
already generated spin geometry. The unique round spin structure and the
invariant frame are the existing harmonic package of
Theorem~\ref{p1:thm:representation-harmonic-package}.
The quark frame is the color summand of the generated matter factorization.
The equal-time principal symbol $c(n_{\rm clk})c(\xi)$ is even in external
Clifford degree; the mass term has an odd clock factor and the inherited odd
internal source. Both preserve
$P_+=\tfrac12(1+\gamma_x\otimes\Gamma_W)$.
Thus the rank-$144$ analytic quark heat completion of
Theorem~\ref{p5:thm:least-self-adjoint-heat-carrier} is not a second physical
equal-time carrier.

Orthonormality and $d\mu_R=R^3d\mu_1$ give $J^\dagger J=I$.
Each component is smooth and lies in the existing Dirac domain. The finite
field/dual substitution is the licensed evaluation word of
Theorem~\ref{p4:thm:external-fermion-line-generated-closure}.
Its kinetic matrix on the static uninserted chart is $G_m=c_FI$;
canonical normalization removes this same $c_F$ from both field legs and
the inherited finite source. Moving frames use the actual positive-root
and embedding jets of the finite-metric theorem.
This supplies the required constructor in the permitted product case of
Corollary~\ref{p3f:cor:v2180-scalar-frontier-type-firewall}: the quark
sector and its action have already been independently generated. It does
not infer a spinor spectral projector from a scalar projector.
\end{proof}

\begin{proposition}[Exact orbital Dirac blocks and native free-symbol comparison]
\label{p10:prop:v2189-orbital-dirac-symbol}
In the generated invariant frame, the rank-two spatial spinor component
block at orbital $j=n/2$ is
\[
 R\mathcal D_j=\tfrac32I+2\sum_a\sigma_a\otimes J_a^{(j)}.
\]
The complete component-orbital band through $n=m$ has positive spin-Dirac
shells $k=0,\ldots,m$ and negative shells $k=0,\ldots,m-1$, with
\[
 \lambda_{n,+}=(n+3/2)/R,
 \quad d_{n,+}=(n+1)(n+2),\qquad
 \lambda_{n,-}=-(n+1/2)/R,
 \quad d_{n,-}=n(n+1).
\]
It is reducing for the round Dirac and for the generated homogeneous
Cartan operator
$R\mathcal D_f=4S\cdot J+4fS\cdot I+3/2$.
The physical left/right blocks and their finite source are retained in
$\mathcal Q_+$. The principal symbol remains the actual clocked Clifford
symbol; the homogeneous curvature/connection corrections are $O(R^{-1})$.
\end{proposition}
\begin{proof}
The first expression is already the component-orbit formula of
Corollary~\ref{p3f:cor:level-one-quadrupole-incidence}.
Couple $S=1/2$ to $j=n/2$ and use
$4S\cdot J=2[t(t+1)-j(j+1)-3/4]$.
The branches $t=j\pm1/2$, including the spectator multiplicity $n+1$,
give the stated eigenvalues/ranks. Reindex the negative branch by $k=n-1$.
Summing gives $2D_m$ for one rank-two spatial spinor component.
Invariant-frame derivatives preserve every orbital Casimir block; the
finite spin/internal term of Theorem~\ref{p3f:thm:m8td-hodge-dirac} does
likewise. The exact symbol comes from the generated joint Dirac, not from
counting. General nonconstant fields use the actual action restriction,
without an unsupported reducing-subspace claim.
\end{proof}

\begin{proposition}[Spatial CAR maps and finite extensive density]
\label{p10:prop:v2189-spatial-car-density}
For admitted reciprocal representatives with $m\le n$, normalized constant
homothety gives
\[
 (U_{R'R}\psi)(R'q)=(R/R')^{3/2}\psi(Rq),\qquad
 i_{nm}=U_{R_nR_m}|_{\mathcal K_m^q}:
 \mathcal K_m^q(R_m)\hookrightarrow\mathcal K_n^q(R_n).
\]
These are isometries with exact composition. Their exterior lifts
$\mathcal I_{nm}=\bigoplus_r\Lambda^ri_{nm}$ satisfy
\[
 c_n^\dagger(i_{nm}f)\mathcal I_{nm}
 =\mathcal I_{nm}c_m^\dagger(f),\qquad
 c_n(i_{nm}f)\mathcal I_{nm}
 =\mathcal I_{nm}c_m(f).
\]
Consequently they give exact injective CAR-algebra maps on the inherited
fermionic branch. On the existing scalar-frontier comparison
$m/R_m\to\Lambda_F$,
\[
 \boxed{\frac{\dim\mathcal K_m^q(R_m)}{2\pi^2R_m^3}
       \longrightarrow\frac{12\Lambda_F^3}{\pi^2}.}
\]
The scale here is the existing scalar-frontier comparison parameter, not a
selected same-numerical-cutoff spinor frontier. This is external mode density;
the complete internal root remains at $m_p=1$.
\end{proposition}
\begin{proof}
The factor $R^{-3/2}$ is exactly the three-dimensional $L^2$ normalization.
Complete orbital bands have canonical inclusions and exhaust the smooth
comparison field domain. The complement of $i\mathcal K_m^q$ is in its
vacuum on the Fock image; wedge multiplication and contraction therefore
give both displayed identities, not just creation compatibility.
The mode count and the inherited frontier asymptotic of
Corollary~\ref{p3f:cor:complete-frontier-mode-density} give the limit.
For the free geometric term the actual interlevel identity is
$\mathcal D_{R_n}i=(R_m/R_n)i\mathcal D_{R_m}$.
This does not replace physical homothety by an unscaled dynamics
intertwiner or assert a cosmological history reaches every level.
\end{proof}

\begin{proposition}[Exact transported-fibre covariance and inherited source jets]
\label{p10:prop:v2189-spatial-quark-source-word}
Let $(a,\Phi,J_m)$ be the actual parent field/embedding word. Under an
admitted generated gauge map $M_g$, retain
\[
 (a,\Phi,J_m)\longmapsto(a^g,\Phi^g,M_gJ_m).
\]
Then the finite action, Dirac matrices and all compatible sources are exactly
covariant between these retained embedding fibres; their exterior lifts
intertwine the corresponding CAR generators. This is not the compressed
multiplier $P_mM_gP_m$ on one fixed band.
For an actual bounded numerical Hermitian bilinear symbol $S_a(q)$ on a
fixed normalized chart,
\[
 S_{a,m}=J_m^\dagger S_aJ_m,
 \qquad\|S_{a,m}\|\le\|S_a\|_\infty.
\]
The scalar, pseudoscalar and axial field/dual ports inherit their actual
Clifford and internal source slots before physical singlet reduction.
Source or metric derivatives retain $J_a$, the kinetic normalization,
interaction slots, projectors and central/sea contacts of the complete
finite-metric theorem.
\end{proposition}
\begin{proof}
Use $D_{a^g,\Phi^g}=M_gD_{a,\Phi}M_g^{-1}$ together with the transported
embedding and dual embedding. Adjacent gauge maps cancel exactly in the
pulled-back action and compose on its fibres. No multiplication closure of
a fixed orbital band is used. Isometry gives the norm bound and preserves
adjoints. Differentiation of $J^\dagger D J$ gives all three terms
$J_a^\dagger DJ+J^\dagger D_aJ+J^\dagger DJ_a$;
the full kinetic and interaction jets are those already compiled in
Theorem~\ref{p3f:thm:v1148-finite-metric-constructor}.
On static mass/current charts $J_a=0$ where the actual embedding is fixed;
moving metric/frontier charts retain their real fixed-rank domain instead.
\end{proof}

\paragraph{Constructor and ordered-readout handoff.}
This proves a native spatially extensive CAR/action/source restriction family
on the existing licensed quark branch. It supplies the constructor component
and its density, symbol, covariance and source structure rather than merely
multiplying a scalar count by an unassigned bundle rank.
Identification with the particular finite/ordered VII-A QCD source, physical
gauge/sea reduction and clock normalization retains its actual comparison
word. The pressure, action-weighted critical response, residual susceptibility,
Goldstone and GMOR limits remain downstream. Appendix~\ref{p10:app:v2189-native-spatial-quark-comparison} shows why
leading density and exact transported-fibre covariance do not identify a
particular QCD source/sea word or its ordered response. No nonzero condensate
or physical susceptibility value is asserted.

% END INLINE SOURCE: Paper7A-NativeSpatialQuarkConstructor-Main-v1_0.tex


\begin{theorem}[Banks--Casher bridge]
\label{p10:thm:banks-casher-bridge}
Under the infinite-volume/chiral order of limits for which the Dirac spectral density exists,
\begin{equation}
\Sigma=\pi\rho_D(0)
\end{equation}
(up to the stated sign/convention).  A finite matrix by itself has only a discrete spectrum and does not supply a nonzero density at the origin.
\end{theorem}
\begin{proof}
This is the Banks--Casher spectral theorem under the stated
infinite-volume/\allowbreak chiral order of limits.
The finite-level spectrum is a sum of delta functions; only after the
volume-normalized spectral measure converges can its value at zero define
$\rho_D(0)$ and hence $\Sigma$.
\end{proof}

\begin{theorem}[Near-zero Dirac universality class]
\label{p10:thm:near-zero-dirac-universality}
Random-matrix/epsilon-regime universality may be imported only after the symmetry class, topology, volume scaling, and chiral limit are matched.  It is a conditional universality theorem, not a finite-carrier identity.
\end{theorem}
\begin{proof}
Universality identifies microscopic spectral correlation functions only after the antiunitary symmetry class, index sector, scaling variable, and volume/chiral regime match those of the universal ensemble.  Since those hypotheses are independent data, the result is correctly conditional rather than a finite-count identity.
\end{proof}

\begin{theorem}[Unit topological representative: conditional regulated construction]
\label{p10:thm:unit-topological-representative-finite-carrier}
If an admissible Ginsparg--Wilson/overlap-type Dirac operator is constructed on a refinement and has index \(\nu=\pm1\), then that refinement contains a unit topological representative.  The 24-cell topology makes such sectors geometrically admissible but does not by itself construct the gauge field/operator pair.
\end{theorem}
\begin{proof}
For an admissible Ginsparg--Wilson operator, the index is an integer invariant and an index $\pm1$ implies a configuration in the corresponding unit sector.  The implication is conditional because neither the admissible gauge field nor the overlap/Ginsparg--Wilson operator follows from the 24-cell incidence data alone.
\end{proof}

\begin{theorem}[Topological extensivity criterion]
\label{p10:thm:primitive-topological-cell-extensivity}
For local topological contributions \(\nu_x\),
\begin{equation}
\operatorname{Var}\!\left(\sum_x\nu_x\right)=\sum_{x,y}\operatorname{Cov}(\nu_x,\nu_y).
\end{equation}
A nonzero finite susceptibility density follows if the covariance is summable with a positive extensive limit.  Existence of a single \(\nu=\pm1\) configuration does not imply that condition.
\end{theorem}
\begin{proof}
Expanding the variance of a sum gives the double covariance sum exactly.  Dividing by volume has a nonzero finite limit only when the integrated covariance has the corresponding extensive limit.  A single admissible charged configuration gives no information about its statistical weight or correlations.
\end{proof}

\begin{theorem}[Topological quenching is a dynamical question]
\label{p10:thm:no-topological-quench}
Topology is unquenched at finite level exactly when the \(\theta=0\) transfer measure gives nonzero variance to \(\nu_L\).  This is computed from the transfer measure; it is not excluded by the absence of a commuting projector.
\end{theorem}
\begin{proof}
By Theorem~\ref{p10:thm:native-topological-susceptibility-transfer-curvature}, the finite topological response is proportional to $\operatorname{Var}(\nu_L)$.  It vanishes exactly when the charge is almost surely constant under the transfer measure.  This is a property of the measure, not of the abstract existence of a topological sector.
\end{proof}

\begin{theorem}[Index forces zero modes, not arbitrary nonzero matrix elements]
\label{p10:thm:no-zero-mode-decoupling}
Under the regulated index theorem, \(\nu_L\ne0\) forces chiral zero modes with net chirality equal to the index.  It does not force every symmetry-allowed zero-mode mixing or source matrix element to be nonzero.
\end{theorem}
\begin{proof}
The index theorem fixes the difference between right- and left-handed zero-mode multiplicities.  It contains no statement about matrix elements of an unrelated source between those modes and other states.  Such a matrix element can vanish by dynamics even when the zero modes are required.
\end{proof}

\begin{theorem}[Topological zero-mode support without a density claim]
\label{p10:thm:topological-cell-density-near-zero-support}
A nonzero regulated index supplies exact zero-mode support in that topological sector.  A nonzero near-zero spectral density \(\rho_D(0)\) still requires the thermodynamic/refinement limit and is not inferred from finite support alone.
\end{theorem}
\begin{proof}
A nonzero finite index gives exact kernel vectors of the regulated Dirac operator.  A density $\rho_D(0)$ is instead a volume-normalized limit of the number of eigenvalues in shrinking intervals around zero.  The former finite existence statement therefore does not imply the latter limit.
\end{proof}

\begin{theorem}[Chiral breaking and topology are compatible but logically distinct]
\label{p10:thm:native-topological-chiral-condensation}
A QCD refinement may simultaneously have nontrivial topology and a nonzero chiral condensate, but neither property follows from the other without the corresponding transfer/index and order-of-limits hypotheses.  Paper~VII-A therefore treats them as coupled computation targets rather than a single forced activation theorem.
\end{theorem}
\begin{proof}
Topology is controlled by the gauge-sector index/measure, whereas chiral breaking is controlled by the thermodynamic order parameter or near-zero spectral density.  Neither defining condition logically contains the other.  They can coexist in one solution only after both independent conditions are met.
\end{proof}

\begin{theorem}[Chiral vacuum alignment]
\label{p10:thm:chiral-vacuum-alignment}
Once the symmetry-broken chiral limit and a small mass spurion \(M\) are supplied, leading-order vacuum alignment is obtained by minimizing \(-\Sigma\operatorname{ReTr}(MU^\dagger)\) over \(U\in SU(3)\).  This is a low-energy consequence, not a finite unique-vacuum theorem.
\end{theorem}
\begin{proof}
The leading mass-dependent chiral potential is $V(U)=-\Sigma\operatorname{ReTr}(MU^\dagger)$.  Once $\Sigma$ and $M$ are supplied, minimizing this explicit finite-dimensional function over $SU(3)$ determines the leading vacuum orientation.  No finite unique-vacuum claim is used.
\end{proof}

\begin{theorem}[Conditional chiral coset action]
\label{p10:thm:native-chiral-coset-action}
Under the SSB hypotheses, the leading \(SU(3)\) chiral action is
\begin{equation}
\mathcal L_2=\frac{F_\pi^2}{4}\operatorname{Tr}(D_\mu U D^\mu U^\dagger)+\frac{F_\pi^2B_0}{2}\operatorname{Tr}(MU^\dagger+UM^\dagger),
\end{equation}
with \(F_\pi\) and \(B_0\) supplied by the native correlator/source readouts above.
\end{theorem}
\begin{proof}
The coset $[SU(3)_L\times SU(3)_R]/SU(3)_V$ is represented by $U\in SU(3)$.  At two derivatives and first order in the mass spurion, chiral covariance and parity allow the two displayed invariants; their coefficients are conventionally denoted $F_\pi^2/4$ and $F_\pi^2B_0/2$.  The preceding correlator/source definitions determine those coefficients.
\end{proof}

\begin{corollary}[Leading pseudoscalar mass pattern]
\label{p10:cor:leading-pseudoscalar-mass-pattern}
The standard leading pseudoscalar mass relations follow algebraically from Theorem~\ref{p10:thm:native-chiral-coset-action} once the physical light-quark spurion is supplied.
\end{corollary}
\begin{proof}
Write $U=\exp(i\phi/F_\pi)$ and expand the mass term of Theorem~\ref{p10:thm:native-chiral-coset-action} to quadratic order in the octet field.  Reading off the quadratic form gives the standard leading pseudoscalar mass relations.
\end{proof}



\begin{theorem}[Physical light-flavor QCD spurion inherits the microscopic mass/source operator]
\label{p10:thm:visible-shell-light-chiral-spurion}
The chiral mass spurion is
\begin{equation}
M_\chi=\operatorname{diag}(m_u^{\rm phys},m_d^{\rm phys},m_s^{\rm phys})
\end{equation}
(or its source-equivalent operator) obtained from the derived charged preterminal source member after the typed G16 RG/threshold transport.  The harmonic number \(H_8\) is not used: current Paper~III-D proves that \(H_8\) requires an added probability measure over rank alternatives and is not a native physical constant.
\end{theorem}
\begin{proof}
The chiral Lagrangian couples to the same scalar mass/source operator that appears in the underlying quark Hamiltonian.  Therefore its eigenvalues are inherited from the derived charged preterminal source member after the certified G16 RG/threshold transport; no second chiral mass selector is introduced.  Paper~III-D explicitly leaves $H_8$ dependent on an added rank measure, so replacing the mass operator by $H_8$ would reintroduce an unowned selector.
\end{proof}

\begin{corollary}[Light-flavor ratios and Leutwyler package are conditional algebraic readouts]
\label{p10:cor:native-light-quark-ratio-package}
Once the inherited \(M_\chi\) is available at the chiral scale, its mass ratios and the standard leading chiral ratio combinations are algebraic readouts.  No native numerical ratio is claimed before the actual source member and required RG transport exist.
\end{corollary}
\begin{proof}
Once the three eigenvalues of $M_\chi$ are supplied, any ratio or standard leading chiral combination is obtained by ordinary algebra.  Because no new parameter is introduced by forming a ratio, its logical status is exactly that of the supplied masses.
\end{proof}

\begin{theorem}[Pseudoscalar source-matrix export to Paper VII-B]
\label{p10:thm:pseudoscalar-source-matrix-export}
Paper~VII-A exports \((M_\chi,F_\pi,B_0,L_i,\ldots)\) with their computed/conditional statuses to Paper~VII-B.  Paper~VII-B may evaluate hadronic observables but may not use them to retroactively choose these kernel/source objects.
\end{theorem}
\begin{proof}
Each listed object is either a typed source operator or a coefficient extracted from the Paper-VII-A generating functional.  Passing this tuple forward does not alter it; using a Paper-VII-B observable to choose it would reverse the declared dependency arrow.
\end{proof}

\begin{theorem}[WZW anomaly normalization]
\label{p10:thm:wzw-anomaly-normalization}
Once the chiral EFT exists, the quantized Wess--Zumino--Witten coefficient is fixed by \(N_c=3\) and the anomaly-matching convention.  Physical decay amplitudes still depend on dynamical residues such as \(F_\pi\).
\end{theorem}
\begin{proof}
Anomaly matching fixes the coefficient of the Wess--Zumino--Witten functional to the integer color multiplicity $N_c$ in the adopted normalization.  Matrix elements of that functional also contain the normalization of the Goldstone fields, hence dynamical residues such as $F_\pi$.
\end{proof}

\begin{theorem}[Anomaly-source normalization and execution boundary]
\label{p10:thm:anomaly-source-paper-xi-boundary}
The anomaly coefficient/current normalization is fixed by the UV fermion content; amplitudes obtained by sandwiching the anomaly operator between hadron states are Paper-VII-B spectral/source readouts and inherit the computed \(F_\pi\) and mixing data.
\end{theorem}
\begin{proof}
The ultraviolet triangle anomaly is fixed by fermion representations and charges.  A hadronic amplitude is obtained only after evaluating the corresponding operator between nonperturbative states.  Therefore anomaly normalization is upstream of, and logically distinct from, the Paper-VII-B spectral matrix element.
\end{proof}

\begin{theorem}[Native chiral EFT hierarchy]
\label{p10:thm:native-chiral-eft-hierarchy}
If a broken chiral phase with scale separation exists, the derivative/mass expansion defines the usual \(p^2,p^4,\ldots\) hierarchy.  Its low-energy constants are source-cumulant targets of the QCD kernel, not fitted inputs in Paper~VII-A.
\end{theorem}
\begin{proof}
When external momenta and quark-mass insertions are small compared with the non-Goldstone scale, locality and symmetry organize the effective functional by derivative/mass order.  Coefficients at each order are Taylor/source-cumulant coefficients of the same generating functional; scale separation is the stated hypothesis that controls the expansion.
\end{proof}

\begin{corollary}[NLO chiral constants are native computation targets]
\label{p10:cor:nlo-chiral-constants-native-targets}
The coefficients \(L_1,\ldots,L_{10},H_1,H_2\) are obtained, when the chiral limit exists, by projecting the low-momentum/source expansion of the connected generating functional onto the standard \(p^4\) invariant basis.
\end{corollary}
\begin{proof}
The standard $p^4$ invariants form a finite basis modulo identities and integrations by parts.  Expanding the connected source functional to fourth chiral order and projecting onto that basis uniquely extracts the coefficients whenever the expansion exists.
\end{proof}

\begin{theorem}[Leading pseudoscalar axial-source normalization]
\label{p10:thm:native-leading-pseudoscalar-weak-decays}
At leading chiral order the octet axial-source coupling is normalized by the dynamically computed \(F_\pi\) (and \(F_K,F_{\eta_8}\) when symmetry breaking is included).  There is no count/Casimir formula fixing these decay constants.  Decay widths remain Paper-VII-B execution problems.
\end{theorem}
\begin{proof}
The axial source couples through the normalized axial current.  In the Goldstone sector its one-particle matrix element is exactly the residue defining the corresponding decay constant.  Hence the leading source coefficient is the computed $F$-type residue and cannot be replaced by a color Casimir or count ratio.
\end{proof}

\begin{theorem}[Native electromagnetic-current interface]
\label{p10:thm:native-electromagnetic-form-factor-bridge}
The electromagnetic current is
\begin{equation}
J_{\rm em}^{\mu}=\sum_f Q_f\bar q_f\gamma^\mu q_f.
\end{equation}
Its exact charge normalization is fixed by the quark charges, while nonzero-momentum form factors and radii are dynamical matrix elements exported to Paper~VII-B.
\end{theorem}
\begin{proof}
On the declared quark-source branch, the electromagnetic phase rotation is a symmetry of the same typed quark/source owner.  The displayed current is therefore its native same-owner Noether current rather than an independently normalized current.  Its integrated charge acts with eigenvalues $Q_f$, fixing zero-momentum normalization.  This statement does not promote the still-open Lorentzian fermionic real-form problem elsewhere in the series: Paper~VII-A uses the already declared finite quark/source branch.  Momentum dependence probes matrix elements of the interacting eigenstates and is therefore a downstream spectral readout.
\end{proof}


% Source Paper VII-A, original line 4055; inherited proof bodies unchanged.
\subsection{Strong CP and topological transfer sectors}
\label{p10:sec:strong-cp-topology}

\begin{theorem}[Conditional finite topology and regulated index representative]
\label{p10:thm:admissible-finite-topology-index}
If the finite gauge regulator supplies an admissible Ginsparg--Wilson/overlap Dirac operator \(D_L\), define
\begin{equation}
\nu_L:=\operatorname{index}D_L\in\ZZ.
\end{equation}
The integer is stable under continuous deformations that remain inside the same admissible component.  The existence and admissibility bound of such a regulator are explicit construction obligations; they are not inferred from the 24-cell count.
\end{theorem}
\begin{proof}
The Ginsparg--Wilson/overlap index is integer-valued and locally constant on an admissible connected component because zero can be crossed only when the operator loses the spectral gap required by admissibility.  The premise explicitly requires constructing such an operator and admissible domain.
\end{proof}

\begin{corollary}[Admissibility threshold is regulator data]
\label{p10:cor:no-free-admissibility-constant}
Topology fixes the integer index once an admissible regulator is given, but it does not determine a universal numerical admissibility cutoff.  Any cutoff must be generated by the chosen regulator theorem and shown sufficient for index stability.
\end{corollary}
\begin{proof}
Index stability is a qualitative topological consequence once admissibility holds.  Rescaling or changing the regulator can change a sufficient plaquette/spectral cutoff without changing the integer index, so topology alone cannot fix a universal numerical threshold.
\end{proof}

\begin{theorem}[Trivial native sector character leaves topology intact]
\label{p10:thm:topological-sectors-survive}
The native strong action assigns the trivial topological character
\(\chi_{\rm top}(\nu)=1\).  This does not eliminate configurations with
nonzero overlap index: topological sectors are components of the admitted
finite gauge family, not states created by a CP-odd coupling.
\end{theorem}
\begin{proof}
The overlap index is defined and locally constant on every gapped component.
Multiplying components by a character changes their action weight but not the
existence of the components.  The native action simply uses unit character.
\end{proof}

\begin{theorem}[Topological susceptibility as exact transfer curvature and charge variance]
\label{p10:thm:native-topological-susceptibility-transfer-curvature}
For a finite Euclidean volume \(V_L\) and a well-defined integer charge \(\nu_L\), let
\begin{equation}
Z_L(\theta)=\sum_{\nu\in\ZZ} Z_{L,\nu}e^{i\nu\theta}.
\end{equation}
Then
\begin{equation}
\label{p10:eq:topological-susceptibility-variance}
\chi_{t,L}:=-\frac1{V_L}\left.\frac{d^2}{d\theta^2}\log Z_L(\theta)\right|_{\theta=0}
=\frac{\operatorname{Var}_{0}(\nu_L)}{V_L}\ge0.
\end{equation}
If CP holds at \(\theta=0\), \(\langle\nu_L\rangle_0=0\).
\end{theorem}
\begin{proof}
Differentiate the characteristic function of the integer-valued charge distribution.  The first derivative gives \(i\langle\nu\rangle\), and the second logarithmic derivative gives minus the variance.
\end{proof}

\begin{remark}[The angle in the Fourier transform is an auxiliary source]
\label{p10:rem:v161-theta-fourier-source}
In \(Z_L(\theta)=\sum_\nu Z_{L,\nu}e^{i\nu\theta}\), the variable \(\theta\)
is the characteristic/Fourier source for the integer-valued overlap charge.
Its derivatives at zero generate cumulants.  It is a physical strong-action
coupling only if a separate licensed constructor generates the sector
character \(e^{i\theta\nu}\); the current native strong action does not.
\end{remark}

\begin{theorem}[Exact CP-even Fourier form and nearest-winding truncation]
\label{p10:thm:primitive-unit-winding-theta-transfer-form}
Normalize \(p_{L,\nu}:=Z_{L,\nu}/Z_L(0)\).  CP implies \(p_{L,\nu}=p_{L,-\nu}\), so
\begin{equation}
\frac{Z_L(\theta)}{Z_L(0)}=p_{L,0}+2\sum_{\nu\ge1}p_{L,\nu}\cos(\nu\theta).
\end{equation}
If a nearest-winding truncation \(\nu=0,\pm1\) is independently justified, then
\begin{equation}
\frac{Z_L(\theta)}{Z_L(0)}=1-r_L(1-\cos\theta),\qquad r_L=2p_{L,1}\in[0,1].
\end{equation}
Thus the cosine form is derived, but its activity \(r_L\) is a transfer-measure output and is not fixed by a shell coefficient.
\end{theorem}
\begin{proof}
Group the exact Fourier series for the integer charge by the pairs $\pm\nu$.  CP makes their weights equal, turning each pair into $2p_{L,\nu}\cos(\nu\theta)$.  If all weights with $|\nu|>1$ vanish, normalization gives $p_0+2p_1=1$, which rearranges to the displayed form with $r_L=2p_1$.
\end{proof}

\begin{theorem}[No dimensional-analysis normalization of susceptibility]
\label{p10:thm:minimal-native-topological-susceptibility}
Equation~\eqref{p10:eq:topological-susceptibility-variance} is the native finite susceptibility.  Writing \(\chi_t=c\Lambda^4\) is only a dimensional parametrization; neither the dimensionless \(c\) nor a cell four-volume is fixed by dimensional analysis.  Both must be computed from the transfer measure and the physical volume map.
\end{theorem}
\begin{proof}
Susceptibility has mass dimension four, so one may always write it as $c\Lambda^4$ after choosing a nonzero reference scale.  Dimensional analysis does not constrain the dimensionless number $c$: multiplying the transfer weights or changing correlations changes $c$ without changing dimensions.  Equation~\eqref{p10:eq:topological-susceptibility-variance} therefore contains the actual normalization.
\end{proof}

\begin{corollary}[Topological enhancement is a finite covariance readout]
\label{p10:cor:topological-enhancement-readout}
For two specified transfer measures or channel packets, an enhancement factor is the ratio of their computed charge variances/susceptibilities.  It is not fixed by a representation dimension alone.
\end{corollary}
\begin{proof}
By definition an enhancement compares two susceptibilities or variances computed with two specified measures.  The ratio depends on their sector weights and covariances; a representation dimension supplies neither quantity.
\end{proof}

\begin{theorem}[Mixed-shell winding variance is a covariance-matrix quadratic form]
\label{p10:thm:orthogonal-hurwitz-adjoint-winding-variance}
For channel charges \(\nu_i\) with real weights \(w_i\),
\begin{equation}
\label{p10:eq:mixed-shell-topological-covariance}
\operatorname{Var}\!\left(\sum_iw_i\nu_i\right)=w^TCw,
\qquad C_{ij}=\operatorname{Cov}(\nu_i,\nu_j).
\end{equation}
A factor such as \(36\eta^2\) follows only under additional conditions (uniform weights, specified carrier overlap, and a covariance proportional to the identity).  Current Paper~V-B admits different overlap values on different carriers, so no representation-independent \(24/25\) factor is used.
\end{theorem}
\begin{proof}
Let $X=\sum_iw_i\nu_i$. Expanding the variance gives
\[
\operatorname{Var}X
=\langle(X-\langle X\rangle)^2\rangle
=\sum_{ij}w_iw_j\operatorname{Cov}(\nu_i,\nu_j)
=w^TCw.
\]
A count-only factor follows only after imposing the additional covariance
and weight assumptions stated in the theorem.
\end{proof}

\begin{corollary}[Full primitive topological winding block]
\label{p10:cor:full-primitive-topological-winding-block}
The Paper-VII-A topological export is the finite charge operator \(\nu_L\), sector weights \(p_{L,\nu}\), and covariance matrix \(C_L\) (or the equivalent \(\theta\)-dependent partition function).  This package is sufficient to compute the susceptibility and higher cumulants without a count-based variance ansatz.
\end{corollary}
\begin{proof}
The characteristic function $Z_L(\theta)/Z_L(0)$ is determined by the sector weights $p_{L,\nu}$, and all cumulants are its logarithmic derivatives.  Equivalently, for a channel decomposition the same second and higher moments are encoded by the covariance/cumulant tensors.  Thus the exported tuple is sufficient.
\end{proof}

\begin{theorem}[Screened full-QCD topology in the chiral regime]
\label{p10:thm:light-quark-screened-full-qcd-topology}
Once a physical light-quark mass spurion and a nonzero chiral condensate have been established, the standard anomalous chiral effective theory determines the leading small-mass screening of \(\chi_t\).  For three light flavors at leading order,
\begin{equation}
\chi_t=\frac{\Sigma}{m_u^{-1}+m_d^{-1}+m_s^{-1}}+O(m_q^2)
\end{equation}
within the usual chiral-domain conventions.  This formula is conditional on the physical mass operator and condensate; it is not a native finite-count number.
\end{theorem}
\begin{proof}
At leading chiral order the vacuum energy is obtained by minimizing the anomalous/mass chiral potential subject to the $\theta$ phase.  Expanding the minimum to second order in $\theta$ gives $\chi_t=\Sigma/(m_u^{-1}+m_d^{-1}+m_s^{-1})$ for three flavors.  The derivation requires the stated condensate, masses, and chiral regime.
\end{proof}

\begin{corollary}[Auxiliary-source curvature and qualified vacuum-energy shape]
\label{p10:cor:native-full-qcd-theta-shape}
Let $E(\theta)$ be a twice differentiable, CP-even source-free-energy limit
whose curvature identification $E''(0)=\chi_t$ is justified on the stated
comparison domain.  Then
\begin{equation}
E(\theta)-E(0)=\frac12\chi_t\theta^2+o(\theta^2).
\end{equation}
If $E$ is $C^4$ near zero, the remainder may be strengthened to
$O(\theta^4)$.  Higher even cumulants are identified only when the required
higher derivatives and their source/limit transport exist.  Interpreting
$E$ as a physical strong-CP vacuum energy additionally consumes a licensed
physical sector-character action; the native auxiliary Fourier source alone
does not constitute it.
\end{corollary}
\begin{proof}
CP evenness gives $E'(0)=0$.  The second-order Taylor formula with Peano
remainder gives the displayed $o(\theta^2)$ term.  Under the additional $C^4$
hypothesis the third derivative at zero vanishes by evenness and the locally
bounded fourth derivative gives $O(\theta^4)$.  Differentiation of the same
normalized source functional gives each higher cumulant only on its admitted
derivative/limit domain.  Appendix~\ref{p10:rem:v2187-c2-source-remainder-comparison}
records the $C^2$ counterexample to a stronger remainder without that premise.
\end{proof}

\begin{theorem}[Topological index bridge]
\label{p10:thm:topological-index-bridge}
Under the hypotheses of an admissible Ginsparg--Wilson/continuum index theorem, the regulated integer charge equals the chiral Dirac index.  Paper~VII-A invokes this only after those hypotheses are met.
\end{theorem}
\begin{proof}
This is the hypothesis-matched lattice/continuum index theorem: admissibility and the Ginsparg--Wilson relation identify the regulated trace/index with the topological charge.  Paper~VII-A uses the equality only on regulator sequences satisfying those hypotheses.
\end{proof}

\begin{theorem}[Topology-induced \('t\) Hooft interaction]
\label{p10:thm:topology-induced-thooft-interaction}
In a continuum/admissible sector with nonzero index and the corresponding fermion zero modes, integrating over the zero modes generates the standard determinant \('t\) Hooft vertex.  Its coefficient is a dynamical sector-weight integral, not fixed by topology alone.
\end{theorem}
\begin{proof}
In a sector of index $\nu$, the fermion determinant has the corresponding chiral zero modes.  Correlation functions are nonzero only when external fermion insertions saturate those Grassmann zero modes, producing the determinant flavor vertex.  The remaining nonzero-mode determinant and sector measure set its coefficient.
\end{proof}

\begin{theorem}[Topological singlet-pseudoscalar bridge]
\label{p10:thm:topological-singlet-pseudoscalar-bridge}
Under the large-\(N_c\)/Witten--Veneziano hypotheses, pure-Yang--Mills susceptibility contributes to the anomalous singlet pseudoscalar mass parameter.
\end{theorem}
\begin{proof}
The Witten--Veneziano large-$N_c$ argument relates the anomalous singlet curvature of the chiral effective action to the pure-Yang--Mills topological susceptibility.  Since this derivation invokes the large-$N_c$ and chiral assumptions, Paper~VII-A records it only as a conditional bridge.
\end{proof}

\begin{theorem}[Anomalous singlet pseudoscalar mass is a correlator/topology readout]
\label{p10:thm:native-anomalous-singlet-pseudoscalar-mass}
On the Witten--Veneziano branch,
\begin{equation}
 m_0^2=\frac{2N_f\,\chi_t^{\YM}}{F_0^2}+O(1/N_c)
\end{equation}
with \(\chi_t^{\YM}\) and \(F_0\) computed from their respective transfer/correlator problems.  No numerical mass is fixed before those quantities are evaluated.
\end{theorem}
\begin{proof}
Substituting the computed pure-Yang--Mills susceptibility and singlet decay constant into the Witten--Veneziano relation gives the displayed mass parameter.  The formula contains no additional count-derived coefficient, and its numerical value is undefined until both inputs are evaluated.
\end{proof}

\begin{corollary}[No universal anomalous singlet mass floor from counts]
\label{p10:cor:primitive-u1a-singlet-pseudoscalar-mass-floor}
A positive lower bound on \(m_0\) requires positive quantitative lower bounds on \(\chi_t^{\YM}\) and an upper control of \(F_0^2\).  Finite channel counts alone provide neither; the former count-derived mass floor is withdrawn.
\end{corollary}
\begin{proof}
From $m_0^2\propto\chi_t^{\YM}/F_0^2$, a positive lower bound requires the numerator to be bounded away from zero and the denominator to be controlled from above.  Representation counts alone provide neither analytic inequality, so no universal floor follows.
\end{proof}

\begin{theorem}[Geometric Hopf sectors of root \(SU(2)\) reductions]
\label{p10:thm:native-hopf-sectors}
For every root embedding \(SU(2)\subset SU(3)\), the quotient \(SU(2)/U(1)\simeq S^2\) admits maps \(S^3\to S^2\) classified by \(\pi_3(S^2)=\ZZ\).  This is a mathematical topological subchannel.  Promotion of a Hopf class to a physical finite pure-glue state requires an explicit map from the gauge configuration/transfer algebra to the \(S^2\)-valued field and a nonzero transfer weight.
\end{theorem}
\begin{proof}
For a root $SU(2)$ subgroup, quotienting by its maximal torus gives
\[
SU(2)/U(1)\simeq S^2.
\]
Homotopy classes of continuous maps from the carrier $S^3$ to this
sphere are $\pi_3(S^2)\cong\mathbb Z$.

A physical state additionally requires a specified gauge-to-sphere map
and nonzero transfer weight, which are not consequences of the homotopy
classification.
\end{proof}



\section{Cartan clock, physical light-front null generator and finite transfer}
\label{p10:sec:v237-a6}
The native equal-time Hopf source and finite Hamiltonian produce the physical null-generator result under its stated seed and positivity hypotheses. A species/Fock light-front realization is a separately certified readout.

% Source Paper VII-A, original line 5065; inherited proof bodies unchanged.
\subsection{Exact Cartan links and native clock-cover index matching}
\label{p10:sec:v199-native-geometric-index}
Use the already generated Hurwitz--transfer regulator of
Theorem~\ref{p10:thm:v190-generated-4d-gw-regulator} and the unit clutch
\(g(q)=\operatorname{diag}(q,1)\) of
Definition~\ref{p10:def:v190-native-clutching-target}.  The spatial complex
is transported to reciprocal position by the existing \(\iota_1\).
The connection family below is the momentum restriction of
Theorem~\ref{p3b:thm:v199-cartan-deformation} composed with that transport.
It does not identify the abstract relative sphere with position by itself.

\begin{theorem}[Exact native Cartan links and face curvature]
\label{p10:thm:v199-exact-cartan-links}
In the right Maurer--Cartan frame let
\(\mathscr A_s=-s\,dg\,g^{-1}\), \(\mathscr A_\tau=0\).
On each shorter round Hurwitz edge, its inverse parallel transport is exactly
\begin{equation}
 U_{qq'}(s)=\exp\{s\operatorname{Log}[g(q)g(q')^{-1}]\}.
 \label{p10:eq:v199-native-link}
\end{equation}
Every edge has relative quaternion angle \(\theta=\pi/3\).  Each spatial
triangular face holonomy is conjugate in \(SU(3)\) to
\(\operatorname{diag}(e^{i\varphi(s)},e^{-i\varphi(s)},1)\), where
\begin{equation}
 w(s)=\cos(\theta s)+\frac4{3\sqrt3}\sin^3(\theta s),\qquad
 \varphi(s)=\arccos w(s).
 \label{p10:eq:v199-face-holonomy}
\end{equation}
For \(0\le s\le1\), \(w(s)=w(1-s)\), and
\begin{equation}
 \max_s\|I-P_s\|_{\rm op}=\sqrt{2-\frac{10}{3\sqrt3}}.
 \label{p10:eq:v199-spatial-curvature-bound}
\end{equation}
An interior temporal ribbon with identity temporal links has
\(\|I-P_{\tau,j}\|_{\rm op}
=2|\sin[\theta(s_{j+1}-s_j)/2]|\).
\end{theorem}
\begin{proof}
This is Paper~1B, Theorem~\ref{p1b:thm:cartan-links}
\cite{Kugelmann2026Paper1B}, on the native Hurwitz edges with
$g(q)=\operatorname{diag}(q,1)$ and $\theta=\pi/3$.
Both papers use the right-frame connection $-s\,dg\,g^{-1}$ and
inverse parallel transport; consequently the edge order and face
orientation agree in \eqref{p10:eq:v199-native-link}--
\eqref{p10:eq:v199-spatial-curvature-bound}. Identity temporal links
give the stated ribbon specialization.
\end{proof}

The gauge convention \(A\mapsto gAg^{-1}-dg\,g^{-1}\) sends the left
representative \(f g^{-1}dg\) to \(-(1-f)dg\,g^{-1}\), so \(s=1-f\).
This parameter reversal must be retained when comparing oriented histories.
The native chirality/orientation convention in the following index theorem
assigns the clutch \(g\) the integer \(+1\).

\begin{proposition}[Exact heat-action dynamics on the Cartan family]
\label{p10:prop:v199-restricted-heat-dynamics}
On the product face-heat branch, with positive times constant within the
spatial and ribbon face orbits, define
\[
 \Psi_s(s)=-\log K_{t_s}(\operatorname{diag}(e^{i\varphi(s)},e^{-i\varphi(s)},1)),
 \quad
 \Psi_\tau(d)=-\log K_{t_\tau}(\operatorname{diag}(e^{i\theta d},e^{-i\theta d},1)).
\]
With \(s_0=0\), \(s_{N-1}=1\), identity interior temporal links and closing
link \(g(q)\), the action restricted to this family is exactly
\begin{equation}
 S_{\rm heat}=96\sum_{j=0}^{N-1}\Psi_s(s_j)
 +96\sum_{j=0}^{N-2}\Psi_\tau(s_{j+1}-s_j)+\text{constant}.
 \label{p10:eq:v199-finite-heat-action}
\end{equation}
Its interior equations are
\[
 \Psi_s'(s_j)+\Psi_\tau'(s_j-s_{j-1})
 -\Psi_\tau'(s_{j+1}-s_j)=0.
\]
For \(N=3\), the profile \((0,1/2,1)\) is stationary within this family
for every positive \((t_s,t_\tau)\).
\end{proposition}
\begin{proof}
Specialize the restricted-function calculation following Paper~1B,
Theorem~\ref{p1b:thm:cartan-links} \cite{Kugelmann2026Paper1B}, to
$K_s=K_{t_s}$ and $K_\tau=K_{t_\tau}$. The native complex has its
96 triangles per layer and 96 ribbons per interior step, while the
clutched closing ribbons are identity. Thus that calculation is exactly
the restriction of the present face-heat action, including its interior
equations and the reflection-stationary three-layer profile.
\end{proof}
The heat times are fixed by the exact-cell stiffness theorem where its
hypotheses apply.  This is an action restriction, not a Haar pushforward to
a profile measure or stationarity under every gauge variation.  Stability
and occurrence of the profile require their own action/state calculation.

\begin{theorem}[Certified native unit component]
\label{p10:thm:v199-native-unit-component}
Take exactly \(N_\tau=3\),
\(\ell_H=a_\tau=r_H=r_\tau=\rho=1\), \((s_0,s_1,s_2)=(0,1/2,1)\),
\(q\mapsto q_0I-iq_a\sigma_a\), and
\[
 \Gamma_a=\sigma_1\otimes\sigma_a\ (a=1,2,3),\quad
 \Gamma_4=\sigma_2\otimes I,\quad
 \Gamma_5=\Gamma_1\Gamma_2\Gamma_3\Gamma_4.
\]
The exact integer certificate in
Appendix~\ref{p10:app:v256-exact-inertia-certificate} proves that the exact
\(864\times864\) algebraic Hermitian kernel \(H_*\) obeys
\begin{equation}
 n_-(H_*)=433,\quad n_+(H_*)=431,\quad
 \nu_{\rm ov}=\nu_{\rm clutch}=1,\quad
 \min|\operatorname{Spec}H_*|>0.256.
 \label{p10:eq:v199-unit-inertia}
\end{equation}
For the same endpoint clutch and temporal links, the match persists whenever
\begin{equation}
 2\pi|s_1-1/2|+|\rho-1|<0.256.
 \label{p10:eq:v199-unit-profile-tube}
\end{equation}
\end{theorem}
\begin{proof}
The exact kernel is \((H_a+\sqrt3H_b)/16\), with Gaussian-integer Hermitian
coefficients reconstructed from the native stencil by the supplied exact
certificate. The appendix gives integer-verified Gram/congruence inequalities
and rational radical enclosures proving the inertia and gap, without a
floating-product error premise;
the smooth unit clutch fixes its topological target.  For the neighborhood,
the eight-neighbor spatial stencil, with Dirac and Wilson edge coefficients
\(3/(8\ell_H)\) and \(3r_H/(8\ell_H)\), gives
\[
 \|H(s,\rho)-H(t,\rho_*)\|_{\rm op}
 \le\frac{3(1+r_H)}{\ell_H}\max_e\|U_e(s)-U_e(t)\|_{\rm op}
     +|\rho-\rho_*|.
\]
Each Clifford tangent has norm one; the block row and column sums prove the
bound.  Since the edge difference is at most \((\pi/3)|s-t|\), it becomes
the stated \(2\pi\) bound at this representative.  The straight profile
path stays invertible, so its inertia and index cannot change.
\end{proof}

\begin{theorem}[All finite clock covers and their stable native index]
\label{p10:thm:v199-native-clock-covers}
Let \(H(\phi)\) be the same three-layer kernel with closing fermion boundary
transporter \(e^{i\phi}g(q)\).  The phase is an auxiliary flat-line/Fourier
parameter, not a new physical spin structure or color gauge field. The exact
integer mesh certificate of Theorem~\ref{p10:thm:v256-exact-native-inertia} gives
\begin{equation}
 \operatorname{gap}H(\phi)>417/7000>0.059,\qquad
 \nu_{\rm ov}(H(\phi))=1\quad(\phi\in\mathbb R).
 \label{p10:eq:v199-uniform-twist-gap}
\end{equation}
For every positive integer \(k\), repeat the three spatial profiles on
\(3k\) layers with the native clutch on every third temporal edge.  This is
the degree-\(k\) clock pullback of the unit bundle.  Both physical temporal
spin choices give
\begin{equation}
 \dim H_{3k}=864k,\qquad \operatorname{gap}H_{3k}>0.059,\qquad
 \nu_{\rm ov}(H_{3k})=\nu_{\rm clutch}=k.
 \label{p10:eq:v199-cover-index}
\end{equation}
Allowing each block's midpoint to vary independently preserves the match if
\begin{equation}
 2\pi\max_\ell|s_{\ell,1}-1/2|+|\rho-1|<0.059.
 \label{p10:eq:v199-cover-profile-tube}
\end{equation}
Time-orientation reversal gives matched charges \(-k\) with the same gaps.
\end{theorem}
\begin{proof}
The appendix certifies the 16 exact phases $j\pi/8$, each with inertia
$(433,431)$ and gap greater than $0.256$. Only the closing blocks
vary; their spin-color norm is one, giving
$\|H(\phi)-H(\psi)\|\le|e^{i\phi}-e^{i\psi}|\le|\phi-\psi|$.
Every phase is within $\pi/16<11/56$ of a sample, so the gap exceeds
$0.256-11/56=417/7000$. Continuity fixes the index.

These native estimates supply the uniform phase-gap premise of
Paper~1B, Proposition~\ref{p1b:prop:cover-transfer}
\cite{Kugelmann2026Paper1B}. Its Fourier decomposition, for both
temporal spin signs, gives \eqref{p10:eq:v199-cover-index}; pullback
multiplies the supplied clutch charge by $k$. The preceding native
profile row bound is independent of $k$ and gives
\eqref{p10:eq:v199-cover-profile-tube}.
For reversed gauge data, the time permutation $R_t$ and unitary
$J_s=\Gamma_5\Gamma_4$ satisfy
\[
 H[U^{\rm rev}]=-J_sR_tH[U]R_t^{-1}J_s^{-1}.
\]
$J_s$ preserves the three spatial Clifford generators and reverses
$\Gamma_4,\Gamma_5$, so both the sign trace and clutch orientation reverse.
\end{proof}

% BEGIN M8 PREPARED ALL-LINK NEIGHBORHOODS
\begin{theorem}[Native index neighborhoods in every link direction]
\label{p10:thm:m8td-all-link-neighborhood}
Let $H_k^*$ be a certified $3k$-layer clock-cover center of
Theorem~\ref{p10:thm:v199-native-clock-covers}, at
$\ell_H=a_\tau=r_H=r_\tau=\rho=1$. Its exact integer certificate gives
$\operatorname{gap}H_k^*>g_*:=417/7000$ and its
overlap index is $k$. Let arbitrary $SU(3)$ spatial and temporal links
$U_e,V_e$ have reverse links given by adjoints. Put
\[
 \epsilon_s=\max_{e\ {\rm spatial}}\|U_e-U_e^*\|,\quad
 \epsilon_t=\max_{e\ {\rm temporal}}\|V_e-V_e^*\|,\quad
 \epsilon_\rho=|\rho-1|,\quad b=6\epsilon_s+2\epsilon_t+\epsilon_\rho.
\]
Here the superscript $*$ denotes the fixed center assignment, not an
adjoint; adjoints of links are denoted by $\dagger$. Then
\begin{equation}
 \|H_k(U,V,\rho)-H_k^*\|\le b,
 \qquad b<g_*\ \Longrightarrow\
 \operatorname{gap}H_k>g_*-b,\quad\nu_{\rm ov}(H_k)=k.
 \label{p10:eq:m8td-all-link-bound}
\end{equation}
The bound is independent of $k$, allows noncommuting color perturbations
on every edge, and requires no residual clock periodicity. If the links
are furnished as holonomies on the retained degree-$k$ clutched bundle,
the overlap index equals that bundle's clutch index throughout this set.
For bare link arrays the conclusion is the overlap index only.
\end{theorem}
\begin{proof}
At the stated unit parameters each of the eight spatial neighbors has
spin block
\[
 B_e=\Gamma_5\frac{\sqrt3\sum_{a=1}^3s_a\Gamma_a-3I}{8},
 \qquad s_a\in\{\pm1\},\qquad \|B_e\|=3/4,
\]
using $(\sum s_a\Gamma_a)^2=3I$; the two temporal blocks have norm one.
Thus these native blocks meet the perturbation estimate of Paper~1B,
Proposition~\ref{p1b:prop:cover-transfer} \cite{Kugelmann2026Paper1B},
with the present $\epsilon_s,\epsilon_t,\epsilon_\rho$ and the gap
already certified by Theorem~\ref{p10:thm:v199-native-clock-covers}.
It gives \eqref{p10:eq:m8td-all-link-bound} for arbitrary link directions.
A supplied connection on the same clutched bundle has unchanged clutch
charge; the finite matrix conclusion itself is inertia stability.
\end{proof}

\begin{corollary}[Gauge-orbit neighborhoods and overlap continuity]
\label{p10:cor:m8td-overlap-continuity}
The preceding bound may be evaluated after any vertex gauge transformation;
thus it defines a neighborhood of the center's gauge orbit. For two
Hermitian kernels with gaps at least $g>0$,
\[
 \|\operatorname{sign}H-\operatorname{sign}K\|\le\|H-K\|/g.
\]
Consequently, for the member in \eqref{p10:eq:m8td-all-link-bound},
\begin{equation}
 \|D_{\rm ov}(H_k)-D_{\rm ov}(H_k^*)\|
 \le\frac{b}{a_R(g_*-b)}.
 \label{p10:eq:m8td-overlap-continuity}
\end{equation}
\end{corollary}
\begin{proof}
Gauge covariance is unitary conjugacy. Using the sign-resolvent
calculus of Paper~1B, Theorem~\ref{p1b:thm:locality}
\cite{Kugelmann2026Paper1B}, subtract the resolvent integrals
for the two sign functions and use
\[
 \operatorname{sign}H-\operatorname{sign}K
 =\frac1\pi\int_{\mathbb R}
       [(H-it)^{-1}-(K-it)^{-1}]\,dt.
\]
The difference integral converges in norm, with integrand bounded by
$\|H-K\|/(g^2+t^2)$. Integration and the overlap definition prove the
claim. The index match remains local to the supplied bundle/interpolation
domain; this estimate is neither a spatial refinement theorem nor a
physical fermion mass gap.
\end{proof}
\begin{theorem}[Positive primitive native overlap sectors on the pure-heat magnitude branch]
\label{p10:thm:v161-positive-primitive-native-sectors}
Let \(\mathcal U_1\) be any nonempty all-link neighborhood supplied by
Theorem~\ref{p10:thm:m8td-all-link-neighborhood} with \(\nu_{\rm ov}=1\), and
let \(\mathcal U_{-1}\) be its orientation-reversed partner.  Then normalized
product Haar measure satisfies
\[
 \mu_{\rm Haar}(\mathcal U_{\pm1})>0.
\]
On the native isotropic pure-product heat magnitude branch,
Theorem~\ref{p10:thm:v202-native-heat-bounds} gives a pointwise lower bound
\(F_{\mathcal B}\ge a^N>0\).  Therefore its pure-heat magnitude
sector coefficients obey
\[
 \boxed{Z^{(+)}_{1}>0,\qquad Z^{(+)}_{-1}>0,}
\]
and the populated native charge lattice is primitive:
\[
 \boxed{\gcd\{|\nu|:Z^{(+)}_\nu>0\}=1.}
\]
The statement uses only the native overlap index.  Unrestricted equality with
a clutching presentation is not required.
\end{theorem}
\begin{proof}
Every nonempty open subset of a compact group has positive normalized Haar
measure: otherwise finitely many left translates would cover the group and
force total measure zero.  Multiply this positive measure by the strictly
positive pure-heat lower bound and integrate.  Orientation reversal gives
the \(-1\) sector.  Since both unit charges are populated, the gcd is one.
\end{proof}

% END M8 PREPARED ALL-LINK NEIGHBORHOODS

% BEGIN FINITE TOPOLOGICAL WEIGHT EVALUATION v2.189

% BEGIN INLINE SOURCE: Paper7A-NativeHeatInterval-Main-v1_0.tex
% Framework v2.189: explicit finite pure-heat susceptibility interval; no physical theta selected.

% Source Paper VII-A, original line 5364; inherited proof bodies unchanged.
\subsubsection{An explicit finite native pure-heat susceptibility interval}
\label{p10:sec:next187-native-heat-interval}

Retain the native three-layer Hurwitz regulator with its actual spin--colour
carrier, $\rho=1$, and the normalized product Haar measure on all independent
spatial and temporal $SU(3)$ links. The consumed VII-A owners are
\ref{p10:thm:v202-native-heat-bounds},
\ref{p10:thm:v199-native-unit-component},
\ref{p10:thm:v199-native-clock-covers},
\ref{p10:thm:m8td-all-link-neighborhood}, and
\ref{p10:thm:v161-positive-primitive-native-sectors}.
The kernel-gap assertion below now uses the exact certificate of
Theorem~\ref{p10:thm:v256-exact-native-inertia}.
This section derives a quantitative bound on the pure, positive, product heat
magnitude readout. An additional determinant or other weight must satisfy its
own certified bounds before this interval is applied to that readout.
Theorem~\ref{p10:thm:v2189-whole-vectorlike-determinant-bounds} supplies
them for the declared positive real scalar-mass vectorlike branch.

\begin{lemma}[A normalized Haar ball bound]
\label{p10:lem:next187-su3-haar-ball}
For $0<\varepsilon\le1/256$, $\delta=\varepsilon/4$, and normalized
$SU(3)$ Haar measure $\mu$, set
\[
 B_\varepsilon=\{U:\|U-I\|_{\rm op}<\varepsilon\},\qquad
 q(\varepsilon)=\frac{1225}{3}\left(\frac7{22}\right)^8
                   \left(\frac\varepsilon4\right)^8.
\]
Then $\mu(B_\varepsilon)>q(\varepsilon)$. The same estimate holds for every
left translate of this ball. In particular,
\[
 q_*:=q(1/256)
 =\frac{7061881225}{199022581174732616540136492402475008}.
\]
\end{lemma}
\begin{proof}
Diagonalize a regular unitary with eigenangles
$(\theta_1,\theta_2,-\theta_1-\theta_2)$. The conjugation map from the maximal
torus times its flag quotient has six sheets. On each off-diagonal tangent
pair its Jacobian is $|e^{i\theta_i}-e^{i\theta_j}|^2$; thus for a class
function the normalized density on the two-dimensional torus is
\[
 \frac1{6(2\pi)^2}\prod_{i<j}|e^{i\theta_i}-e^{i\theta_j}|^2.
\]
Its normalization follows independently by expanding the Vandermonde
determinant. Of the $6^2$ products of permutation terms, torus integration
retains precisely the six identical permutation pairs: an exponent difference
has sum zero, and vanishing after $\theta_3=-\theta_1-\theta_2$ forces all
three differences equal, hence zero. Singular conjugacy classes have zero
measure and do not change the formula.

Use the rectangle $3\delta\le\theta_1\le4\delta$ and
$-\delta\le\theta_2\le-\delta/2$. All eigenangles have magnitude at most
$4\delta=\varepsilon$, and their chord lengths from $1$ are strictly smaller
than $\varepsilon$. The three positive angle differences are at least
$7\delta/2$, $5\delta$, and $\delta$ respectively, and at most
$15\delta/2<\pi$. Concavity of sine on $[0,\pi/2]$ gives
$|e^{i\alpha}-e^{i\beta}|\ge2|\alpha-\beta|/\pi$. Consequently the squared
Vandermonde is at least $19600\delta^6/\pi^6$. The rectangle has area
$\delta^2/2$, whence
\[
 \mu(B_\varepsilon)\ge\frac{1225\delta^8}{3\pi^8}>q(\varepsilon).
\]
For the last strict inequality one may use the elementary positive-integral
identity $22/7-\pi=\int_0^1 x^4(1-x)^4/(1+x^2)\,dx>0$.
Translation invariance proves the final statement.
\end{proof}

\begin{theorem}[Explicit susceptibility interval]
\label{p10:thm:next187-native-heat-interval}
Let $\Omega$ be an admitted gapped link domain containing the two full
product balls in which each of the $360$ independent links differs in
operator norm by less than $1/256$ from its link in the native index $+1$
connection or its reversed index $-1$ connection. On the actual
three-layer product use the positive
heat weight $w$ with $576$ faces and the current native heat representative,
so that
\[
 a^{576}\le w\le b^{576},\qquad
 a=\frac{\pi^2-6}{\pi^2-3}>\frac{14}{25},\quad
 b=\frac{\pi^2}{\pi^2-3}<\frac{36}{25}.
\]
Normalize $w\,d\mu$ on $\Omega$, let $\nu=-\tfrac12\operatorname{Tr}
\operatorname{sign}H$ be its existing index, and use the same-parent curvature
$\chi V=\operatorname{Var}(\nu)$. Then
\[
 10^{-18796}<4\left(\frac7{18}\right)^{1152}q_*^{720}
       <\chi V\le186624.
\]
The leftmost inequality is an exact rational comparison, with no floating
logarithm needed for its certificate.
\end{theorem}
\begin{proof}
The Hurwitz graph has $24$ vertices, degree $8$, $96$ unoriented edges and
$96$ triangular faces. The three-layer product has $3\cdot96+3\cdot24=360$
independent links and $3\cdot96+3\cdot96=576$ heat faces. The last $96$
ribbons are the closing ribbons; they remain counted when a center holonomy
is the identity. The temporal factor is the clock one-complex, so it supplies
no extra filled temporal triangle.

Perturb every link by operator norm less than $\varepsilon=1/256$. The
printed all-link bound gives $\|\Delta H\|\le8\varepsilon=1/32$, strictly
below the inherited lower gap $417/7000$. Each entire product ball therefore
has fixed index $+1$ or $-1$. These two balls are disjoint, and by the lemma
each has Haar mass greater than $q_*^{360}$. Since $\mu(\Omega)\le1$, its
partition integral is at most $b^{576}$. Therefore the normalized
probabilities satisfy
\[
 p_+,p_->P_*:=\left(\frac7{18}\right)^{576}q_*^{360}.
\]
For independent copies $\nu,\nu'$, the identity
$\operatorname{Var}\nu=\tfrac12\mathbb E(\nu-\nu')^2$ retains the two
oppositely ordered unit-sector events and gives
$\operatorname{Var}\nu\ge4p_+p_->4P_*^2$.
The carrier dimension is $864$, so $|\nu|\le432$ and
$\operatorname{Var}\nu\le432^2=186624$.

Writing $q_*=n/d$ as above, the rational decimal comparison is exactly
\[
 4\,7^{1152}n^{720}10^{18796}>18^{1152}d^{720}.
\]
The accompanying integer-only certificate verifies this inequality, together
with the graph and face counts. Its floating logarithm is diagnostic only.
\end{proof}

\begin{corollary}[Primitive auxiliary return]
\label{p10:cor:next187-primitive-auxiliary-return}
For the existing characteristic readout $\varphi(\vartheta)=
\mathbb E e^{i\vartheta\nu}$,
\[
 1-\operatorname{Re}\varphi(\vartheta)
    \ge(p_++p_-)(1-\cos\vartheta)
    >2P_*(1-\cos\vartheta)
\]
whenever $\vartheta\notin2\pi\mathbb Z$. Hence
$\varphi(\vartheta)=1$ precisely at $2\pi\mathbb Z$, and its least positive
period is $2\pi$. This is an auxiliary readout, not a selected physical
strong-CP coupling. The conclusion does not assert
$|\varphi(\vartheta)|<1$ everywhere else.
\end{corollary}

The finite heat interval is now an evaluated native mathematical output;
it is not a phenomenological susceptibility value, an infinite-volume lower
bound, or an estimate for an unbounded or vanishing determinant weight.
Appendix~\ref{p10:app:v2189-topological-weight-diagnostics} records the degree-$17$ heat-tail scope and the characteristic
modulus diagnostic; those guards prevent promotion outside this theorem's
consumed branch.

% END INLINE SOURCE: Paper7A-NativeHeatInterval-Main-v1_0.tex


% BEGIN INLINE SOURCE: Paper7A-NativeWeightedTopology-Main-v1_0.tex
% Framework v2.189: finite real vectorlike weight bounds and source jets.

% Source Paper VII-A, original line 5516; inherited proof bodies unchanged.
\subsubsection{Finite vectorlike weights, their massless endpoint, and native source jets}
\label{p10:sec:v2189-native-weighted-topology}

This section evaluates the declared real vectorlike branch of
Proposition~\ref{p10:prop:v202-vectorlike-positive}; it does not select that
branch or its physical masses. Retain the generated three-layer regulator,
the selected temporal spin lift, and unit stencil coordinates
$\ell_H=a_\tau=r_H=r_\tau=\rho=1$. Its spin--color space has dimension
$d=864$, its independent link group is $SU(3)^{360}$, and its full product
face-heat presentation has $576$ faces, including the closing ribbons.
Write $F_H=\prod_fK_{t_{\rm iso}}(U_f)$, so the exact source bounds give
$a^{576}\le F_H\le b_H^{576}$ and $a/b_H>7/18$.
The integration domain $\Omega$ below is a fixed measurable subset of the
gapped link domain, independent of the physical mass sources. The complete
gapped domain $\Omega_{\rm gap}=\{U:0\notin\operatorname{Spec}H(U)\}$ is
one such explicitly defined finite domain. We declare it only as a finite
representative for these bounded integration readouts, by reusing the generated
gapped family and its normalized Haar measure. No occurring physical transfer
is asserted to select this representative. Restrictions or additional state
weights retain their own support conditions.

\begin{theorem}[Uniform whole determinant bounds on the retained massive branch]
\label{p10:thm:v2189-whole-vectorlike-determinant-bounds}
Let the actually retained scalar flavor masses satisfy $m_f>0$, and put
\[
 c=2/a_R,\qquad x_f=m_f/c,\qquad
 \widehat D_{x_f}=D_{m_f}/c
 =\tfrac12(I+V)+\tfrac{x_f}{2}(I-V),\qquad
 V=\Gamma_5\operatorname{sign}H.
\]
On the whole gapped family, $V$ is unitary and
$\Gamma_5V\Gamma_5=V^\dagger$. With
$\ell_f=\min(1,x_f)$, $u_f=\max(1,x_f)$, and
$r_f=\ell_f/u_f=\min(x_f,x_f^{-1})$, one has
\[
 0<\ell_f^{864}\le\det\widehat D_{x_f}\le u_f^{864},
 \qquad
 \min(m_f,c)^{864}\le\det D_{m_f}\le\max(m_f,c)^{864}.
\]
Thus the full retained weight $F_H\prod_f\det D_{m_f}$ has explicit
positive lower and finite upper bounds. Configuration-independent positive
normalization factors cancel from its normalized sector law; the ratios
$x_f$ remain. Changing $a_R$ while keeping $m_f$ fixed is not that cancellation.
\end{theorem}
\begin{proof}
For an eigenphase $e^{i\vartheta}$ of $V$,
\[
 \left|\tfrac{1+e^{i\vartheta}}2
       +x_f\tfrac{1-e^{i\vartheta}}2\right|^2
 =\cos^2(\vartheta/2)+x_f^2\sin^2(\vartheta/2).
\]
Its singular value therefore lies between $\ell_f$ and $u_f$.
Nonreal eigenphases occur in conjugate pairs; their determinant products
are positive squared moduli. The real eigenphases $+1$ and $-1$ give the
positive eigenvalues $1$ and $x_f$. Hence the determinant is positive and
the product of the $864$ singular values. Multiplication by $c^{864}$
gives the unnormalized statement. For finitely many retained flavors these
configuration-independent factors multiply the numerator and partition
function by the same constant and cancel.
\end{proof}

\begin{corollary}[Explicit massive sector and susceptibility certificate]
\label{p10:cor:v2189-massive-topological-interval}
Suppose $\Omega$ contains both full $360$-link product balls of per-link
operator radius $\epsilon=1/256$ around the certified unit-index center
and its orientation reverse. Their gaps and indices are now certified
exactly by Theorem~\ref{p10:thm:v256-exact-native-inertia}. Set
\[
 q(\epsilon)=\frac{1225}{3}\left(\frac7{22}\right)^8
                     \left(\frac\epsilon4\right)^8,\qquad
 P_H=(7/18)^{576}q(1/256)^{360},\qquad R_m=\prod_f r_f^{864}.
\]
For the normalized full weight
\[
 d\mathbb P_m=Z_m^{-1}\,\mathbf1_\Omega F_H
                    \prod_f\det D_{m_f}\,d\mu_{\rm Haar},
\]
the actual sector probabilities and auxiliary-source curvature obey
\[
 p_{m,+1}>P_HR_m,\qquad p_{m,-1}>P_HR_m,
 \qquad
 4P_H^2R_m^2<\chi_{t,m}V_L=\operatorname{Var}_m(\nu)\le432^2.
\]
In particular $\chi_{t,m}V_L>10^{-18796}R_m^2$.
The populated charge lattice is primitive, and the characteristic function
$\phi_m(\theta)=\mathbb E_me^{i\theta\nu}$ returns to $1$ exactly at
$\theta\in2\pi\mathbb Z$. Its least positive period is $2\pi$; no strict
modulus bound is asserted.
\end{corollary}
\begin{proof}
The actual all-link row bound is $6\epsilon_s+2\epsilon_t$, so the two
balls retain their indices since $8/256<417/7000$. Lemma~\ref{p10:lem:next187-su3-haar-ball} gives mass greater than $q(1/256)^{360}$ for each.
Use the preceding whole-weight lower bound on each ball and the upper
bound $Z_m\le b_H^{576}\prod_f\max(m_f,c)^{864}$ on the full domain.
Their ratio gives the two probability bounds without any face-independence
assumption. For two independent copies of the charge,
$\operatorname{Var}\nu=\frac12\mathbb E(\nu-\nu')^2\ge4p_{+1}p_{-1}$.
The index formula gives $|\nu|\le864/2=432$. The exact integer comparison $4P_H^2>10^{-18796}$ is the
certificate of Theorem~\ref{p10:thm:next187-native-heat-interval}. Finally
$1-\Re\phi_m(\theta)\ge(p_{+1}+p_{-1})(1-\cos\theta)$, proving the
return and least-period assertions. This Fourier angle remains an
auxiliary charge source, not a newly supplied CP action.
\end{proof}

\begin{theorem}[An explicit open native zero-index, invertible-overlap domain]
\label{p10:thm:v2189-zero-index-overlap-witness}
Set every spatial link to $I_3$, every interior temporal link to $I_3$,
and the closing temporal links to $\zeta I_3$, where
$\zeta=e^{2\pi i/3}$. These are actual $SU(3)$ link variables.
For either retained temporal spin sign, this assignment $U^{(0)}$ obeys
\[
 \operatorname{gap}H(U^{(0)})>2/9,
 \qquad \sigma_{\min}(\widehat D_0(U^{(0)}))>1/99,
 \qquad \nu(U^{(0)})=0.
\]
The full product ball $\mathcal B_0$ with per-link operator radius
$\epsilon_0=1/4096$ around this assignment has
\[
 \begin{gathered}
 \operatorname{gap}H>1015/4608,\quad
 \sigma_{\min}(\widehat D_0)>\eta:=1139/200970,\\
 \nu=0,\quad\mu_{\rm Haar}(\mathcal B_0)>q(1/4096)^{360}.
 \end{gathered}
\]
No unit-center residual calculation is used for this separate analytic witness.
\end{theorem}
\begin{proof}
The temporal shift satisfies $T^3=\sigma\zeta I$, where $\sigma=\pm1$
is the retained spin sign. Its Fourier eigenangles are
$\omega_k=(2\pi k+2\pi/3+\varphi_\sigma)/3$ with
$\varphi_+=0$, $\varphi_-=\pi$. In both cases
$\min_k|\sin\omega_k|=\sin(\pi/9)>2/9$.
In each block the exact source stencil has
\[
 K=D_s+B+i s\Gamma_4,\qquad
 s=\sin\omega_k,\qquad B=W_s-\cos\omega_k I.
\]
Here $D_s^\dagger=-D_s$, $\{\Gamma_4,D_s\}=0$, while
$B^\dagger=B$ and $[\Gamma_4,B]=0$. Therefore
$A=\Gamma_4(D_s+B)$ is Hermitian and
$K=\Gamma_4(A+i sI)$. This gives
$K^\dagger K=A^2+s^2I$, hence
$\sigma_{\min}(H)=\sigma_{\min}(K)\ge|s|$.
The source's unit-stencil bound gives $\|K\|=\|H\|\le11$.
Its polar unitary is
\[
 V=K(K^\dagger K)^{-1/2}
   =\Gamma_4Q,\qquad
 Q=(A+i sI)(A^2+s^2I)^{-1/2}.
\]
For a unit vector $\psi$,
\[
 \|(I+V)\psi\|=\|(Q+\Gamma_4)\psi\|
 \ge\left|\Im\langle\psi,Q\psi\rangle\right|
 =|s|\langle\psi,(A^2+s^2I)^{-1/2}\psi\rangle
 \ge|s|/11.
\]
Thus $\widehat D_0=(I+V)/2$ is invertible with the stated lower bound.
The GW index is the net chirality of its zero modes, so it is zero.
This proof does not assume $[D_s,W_s]=0$.

For arbitrary link perturbations the same stencil block row and column
bound gives $b\le8\epsilon_0=1/512$. The kernel gap is greater than
$2/9-1/512=1015/4608$. The current sign-continuity estimate gives
\[
 \|\widehat D_0(U)-\widehat D_0(U^{(0)})\|
 \le\frac{b}{2(2/9-b)}\le\frac9{2030}.
\]
Weyl's singular-value bound leaves
$1/99-9/2030=1139/200970$. Invertibility proves index zero throughout
the ball, and Lemma~\ref{p10:lem:next187-su3-haar-ball} supplies its positive mass.
\end{proof}

\begin{theorem}[Exact finite unsaturated massless endpoint]
\label{p10:thm:v2189-massless-topological-endpoint}
Retain at least one exactly massless unsaturated vectorlike determinant,
any other fixed positive scalar-mass determinants, and the same heat
weight and mass-independent domain. If $\Omega$ contains $\mathcal B_0$
(in particular if $\Omega=\Omega_{\rm gap}$), then its massless partition
function is strictly positive and finite. Its normalized charge law obeys
\[
 \mathbb P_0(\nu\ne0)=0,\qquad
 \phi_0(\theta)=1,\qquad \chi_{t,0}V_L=0.
\]
The nonzero native index components still exist; this is their vanishing
unsaturated determinant weight, not their removal from the regulator.
\end{theorem}
\begin{proof}
At every configuration $n_0=\dim\ker D_{\rm ov}\ge|\nu|$.
An unsaturated massless determinant therefore vanishes on every
$\nu\ne0$ component. On $\mathcal B_0$ the massless normalized
determinant exceeds $\eta^{864}$; each other retained positive scalar
mass has the preceding uniform lower bound. The positive heat and Haar
mass thus give $Z_0>0$. The finite upper bounds give integrability.
Only the zero charge has nonzero normalized weight, proving all three
identities. If a restricted domain omits all invertible-overlap support,
its massless partition function can instead be zero; normalization is
then undefined and this conclusion is not applicable.
\end{proof}

% Framework v2.217 helper integration.
\begin{theorem}[Finite temporal GW source, contact and strict positivity]
\label{p10:thm:v2217-temporal-gw-source}
For an admitted invertible massless GW operator
$D=(I+V)/a_R$, put $R_D=I-a_RD/2$, $G=D^{-1}$, $Q=GR_D$.
Let $E_t$, $t=0,1,2$, be the complete cyclic Euclidean clock projectors,
commuting with $\Gamma_5$, of rank $288$. In the unit off-diagonal charged
flavor convention define
\[
 C_{PP}^{\rm GW}(t,s)=\operatorname{Tr}
       (GE_t\Gamma_5R_DGE_s\Gamma_5R_D).
\]
Then, configuration by configuration,
\[
 C_{PP}^{\rm GW}(t,s)=\|E_tQE_s\|_{\rm HS}^2\ge0,\qquad
 C_{PP}^{\rm bare}(t,s)=C_{PP}^{\rm GW}(t,s)
                         +\delta_{ts}72a_R^2,
\]
where $C_{PP}^{\rm bare}=\|E_tGE_s\|_{\rm HS}^2$.
At the zero-index center of
Theorem~\ref{p10:thm:v2189-zero-index-overlap-witness}, the reducing
spatially uniform sector contributes, for every $t\ne s$,
\[
 C_{PP,\rm unif}^{\rm GW}(t,s)=
 \frac{12a_R^2}{|1-\sigma\zeta|^2}
 =\begin{cases}4a_R^2,&\sigma=+1,\\12a_R^2,&\sigma=-1,\end{cases}
 \quad\zeta=e^{2\pi i/3}.
\]
On a mass-independent admitted domain containing the certified zero-index
ball, the full three-degenerate-massless-flavor, positive-heat-weighted
charged source has strictly positive off-clock expectation. The two source
legs belong to two distinct flavors; the third determinant is a weight,
not an extra copy of this source.
Finally the finite massless axial-divergence insertions
$A_t=[D,E_t]\widehat\Gamma_5$, $\widehat\Gamma_5=\Gamma_5(I-a_RD)$,
obey
\[
 \sum_t\operatorname{Tr}(GA_tGE_s\Gamma_5R_D)=0.
\]
\end{theorem}
\begin{proof}
GW hermiticity gives $Q^*=\Gamma_5Q\Gamma_5$; $G$ and $R_D$ commute.
Cyclicity therefore turns the source trace into
$\operatorname{Tr}(Q^*E_tQE_s)$, the stated squared norm.
The GW relation also gives $G+G^*=a_RI$, so
$Q=G-a_RI/2$ is anti-Hermitian. For unequal clocks the identity contact
vanishes. For equal clocks its cross term has zero real trace, leaving
$(a_R^2/4)\operatorname{rank}E_t=72a_R^2$.

At the stated center the normalized uniform spatial sector is reducing:
the balanced spatial first-order stencil and its Wilson Laplacian both
annihilate constants. Its temporal shift satisfies $T^3=wI$,
$w=\sigma\zeta$. On this sector
$K=-P_+T^{-1}-P_-T$ is unitary, with $P_\pm=(I\pm\Gamma_4)/2$, so
$V=K$ and $Q=(a_R/2)(I+K)^{-1}(I-K)$.
On either $\Gamma_4$ eigenspace put $A=T^{\mp1}$; then
\[
 (I+A)(I-A)^{-1}
   =\frac{(1+w^{\mp1})I+2A+2A^2}{1-w^{\mp1}}.
\]
Each off-clock entry has modulus $a_R/|1-w|$. Four spin and three color
directions give the displayed uniform norm. Orthogonal reducing sectors
add squared norms, so this is a lower bound on the full source at the center.

The inverse and sign are continuous on its certified open zero-index ball.
The source stays positive on a smaller positive-Haar ball, where all three
determinants and all 576 face heat factors are strictly positive.
The full weighted source is nonnegative elsewhere on the invertible set.
At a zero of $D$ its determinant-weighted extension is integrable: two
inverse factors cancel at most two of the three determinant factors,
equivalently use the finite adjugate formula. The bounded heat factor on
the compact finite link space supplies a finite upper bound. The massless
partition function is finite and positive by the preceding theorem.
Its normalized source integral is therefore strictly positive. Finally
$\sum_tA_t=[D,I]\widehat\Gamma_5=0$, proving the Ward sum.
\end{proof}
This is a finite cyclic Euclidean source calculation. It supplies neither
a physical positive-time transfer identification nor a pion pole, axial
residue, condensate or continuum confinement theorem.


\begin{corollary}[Quantitative approach to the finite massless endpoint]
\label{p10:cor:v2189-finite-chiral-topological-suppression}
Let $F\ge1$ retained flavors have $0\le x_f\le1$ and tend toward zero;
other retained scalar flavors have fixed positive ratios $y_j$.
Assume $\Omega$ contains $\mathcal B_0$, and set
\[
 C_F=432^2(18/7)^{576}q(1/4096)^{-360}\eta^{-864F}
          \prod_j\min(y_j,y_j^{-1})^{-864}.
\]
Then for the full normalized retained measure
\[
 0\le\chi_tV_L\le\mathbb E\nu^2
 \le C_F\prod_{f=1}^F x_f,
 \qquad
 \mathbb P(\nu\ne0)\le(C_F/432^2)\prod_{f=1}^F x_f.
\]
This is a fixed finite-carrier result, with a fixed domain and all
normalizations retained; it invokes no thermodynamic or SSB assumption.
\end{corollary}
\begin{proof}
For $x\in[0,1]$ every normalized singular value is at most one, and
every exact zero mode contributes $x$. Hence
$0\le\det\widehat D_x\le x^{n_0}\le x^{|\nu|}$.
On $\nu\ne0$ the product is at most $\prod_fx_f$.
On $\mathcal B_0$, the squared singular values
$\cos^2(\vartheta/2)+x^2\sin^2(\vartheta/2)$ are at least their
$x=0$ values, so each determinant is at least $\eta^{864}$ uniformly
down to zero. The other positive flavors have upper and lower bounds
$\max(1,y_j)^{864}$ and $\min(1,y_j)^{864}$.
Use these bounds on the numerator and the ball's lower bound on the full
partition function, then $|\nu|\le432$. The coarse heat ratio is $18/7$.
\end{proof}

\begin{theorem}[Native finite scalar-source jet hierarchy]
\label{p10:thm:v2189-topological-source-jet-hierarchy}
Keep the preceding fixed domain and a mass-independent, integrable remaining
weight that is strictly positive almost everywhere on the consumed domains. For $F$ independently sourced flavors, each sector coefficient
$Z_\nu(x_1,\ldots,x_F)$ is a finite polynomial of degree at most $864$
in each $x_f$, with nonnegative real coefficients. A nonzero coefficient
of $\prod_fx_f^{k_f}$ requires, on contributing configurations,
\[
 k_f\ge n_0\ge|\nu|,
 \qquad k_f\equiv n_0\equiv\nu\pmod2
 \quad\text{for every }f.
\]
More precisely, at a fixed configuration its determinant has the exact form
\[
 \det\widehat D_x
 =x^{n_0}\prod_{j=1}^{J}
       \bigl(\cos^2(\vartheta_j/2)+x^2\sin^2(\vartheta_j/2)\bigr),
 \qquad \vartheta_j\not\equiv0,\pi.
\]
Thus the first nonzero derivative there has order $n_0$, with coefficient
$\prod_j\cos^2(\vartheta_j/2)>0$. The index alone supplies a lower
order bound, not $n_0=|\nu|$ or a numerical anomaly coefficient.
For an open populated sector of positive remaining weight, some such
finite source jet is nonzero. These licensed mass derivatives do not by
themselves constitute an arbitrary external Grassmann insertion or a
positive probability measure made from an anomaly vertex.

On a common source $x_1=\cdots=x_F=x$, with $\mathcal B_0\subset\Omega$,
the full partition function has $Z(0)>0$. If $F\ge2$, then
$Z'(0)=0$ and $\partial_x\log Z(x)=O(x)$ as $x\downarrow0$.
For one vanishing flavor, a linear anomaly term is allowed.
None of these fixed finite analytic statements proves an ordered
thermodynamic condensate.
\end{theorem}
\begin{proof}
The eigenphase pairing in the first theorem gives the displayed determinant
polynomial; $+1$ eigenphases contribute one, $-1$ eigenphases contribute
$x$, and the remaining pairs give positive quadratic factors. Exact GW
zero modes have definite chirality, so $n_0-|\nu|$ is even.
Finite products and integration against the fixed positive remaining
weight preserve nonnegative coefficients and the printed degree, parity
and minimum-order restrictions. All coefficients are integrable by the
uniform finite determinant bound; no interchange of an infinite expansion
or a moving integration domain is needed. A sector with positive weight
at positive masses cannot have the zero polynomial.
For the common source, a zero-index configuration has even $n_0$, hence
no linear term; a nonzero-index configuration has minimum degree
$F|\nu|\ge2$ when $F\ge2$. Thus $Z'(0)=0$.
The positive constant term supplied by $\mathcal B_0$ makes the logarithm
analytic near zero, proving the $O(x)$ response.
\end{proof}

All statements concern this retained product-heat, real vectorlike
source branch and the consumed domain. A further mass-independent factor
$G(U)$ with known $0<g_-\le G\le g_+<\infty$ multiplies lower
probability bounds by $g_-/g_+$ and the upper suppression constant by
$g_+/g_-$. Positivity alone is not that quantitative bound. A determinant
phase, a gauge-dependent chiral measure, a source-dependent state factor,
or an effective multi-face action must be retained and bounded on its
own licensed branch. Physical mass transport, a physical volume map,
reflection positivity, the continuum susceptibility and ordered chiral
limits are not supplied by this finite computation.

Appendix~\ref{p10:app:v2189-topological-weight-diagnostics} preserves the
normalization, sign-continuity and source-jet diagnostics and the
excluded weight/domain cases, with every actual factor retained.

% END INLINE SOURCE: Paper7A-NativeWeightedTopology-Main-v1_0.tex

\begin{remark}[Primitive activity belongs to the retained weight]
\label{p10:rem:v2189-weight-scoped-primitive-activity}
The coefficient $Z^{(+)}_\nu$ in
Theorem~\ref{p10:thm:v161-positive-primitive-native-sectors} is the
pure-heat magnitude coefficient. The positive real massive vectorlike
weight retains primitive activity with the whole-weight bounds of
Corollary~\ref{p10:cor:v2189-massive-topological-interval}. An
unsaturated massless determinant instead annihilates nonzero-index
coefficients; on the retained domain containing the explicit zero-index
ball, Theorem~\ref{p10:thm:v2189-massless-topological-endpoint} gives
a positive partition function and zero charge variance. In all three
cases the native index components remain components of the same regulator.
Removing a phase alone does not supply a lower bound for an additional
magnitude or state factor.
\end{remark}
% END FINITE TOPOLOGICAL WEIGHT EVALUATION v2.189

% BEGIN V202 NATIVE OVERLAP LOCALITY

% Source Paper VII-A, original line 5916; inherited proof bodies unchanged.
\subsubsection{Locality and chiral transport on the certified native domain}
\label{p10:sec:v202-overlap-locality}
The gap in Theorem~\ref{p10:thm:m8td-all-link-neighborhood} also controls
the nonlocal tail of the finite sign operator. The following result uses
the actual Hurwitz/clock graph distance, not a spatial continuum limit.

\begin{theorem}[Native overlap locality on every certified cover neighborhood]
\label{p10:thm:v202-native-overlap-locality}
At $\ell_H=a_\tau=r_H=r_\tau=1$ and arbitrary unitary links,
\begin{equation}
 \|H\|\le M_\rho:=|4-\rho|+8.
 \label{p10:eq:v202-native-norm}
\end{equation}
For a range-one Hermitian kernel with $0<\gamma\le|H|\le M$, $\gamma<M$,
put $q=(M-\gamma)/(M+\gamma)$. For distinct graph vertices,
\begin{equation}
 \|D_{\rm ov}(x,y)\|
 \le\frac{M}{a_R\gamma}q^{\lfloor d(x,y)/2\rfloor}.
 \label{p10:eq:v202-native-locality}
\end{equation}
For the all-link neighborhood, take $M=11+\epsilon_\rho$ and
$\gamma=417/7000-b$, with $b$ as in
\eqref{p10:eq:m8td-all-link-bound}. These constants are uniform in clock
winding on every subfamily $b\le b_*<417/7000$ with bounded
$\epsilon_\rho$. In particular, at every central cover,
\[
 \|D_{\rm ov}(x,y)\|
 \le\frac{11000}{59a_R}
       \left(\frac{10941}{11059}\right)^{\lfloor d(x,y)/2\rfloor}.
\]
The central gap is supplied by the exact integer certificate of
Theorem~\ref{p10:thm:v256-exact-native-inertia}.
\end{theorem}
\begin{proof}
Each of the eight spatial blocks has norm at most $3/4$, each of the
two temporal blocks norm at most one, and the diagonal is $(4-\rho)I$
before the unitary factor $\Gamma_5$. Block row and column sums give
\eqref{p10:eq:v202-native-norm}. The stencil has range one, so
Paper~1B, Theorem~\ref{p1b:thm:locality}
\cite{Kugelmann2026Paper1B}, gives
\eqref{p10:eq:v202-native-locality}. Substitution of the native gap
$\gamma=417/7000-b$ and $M=11+\epsilon_\rho$ gives the stated
uniform constants; at the center one may weaken the gap to $59/1000$
to obtain the displayed numerical bound.
\end{proof}

\begin{theorem}[Smooth native chiral projectors and their induced connection]
\label{p10:thm:v202-chiral-projector-transport}
On a smooth gapped parameter family in a fixed finite Hilbert-space
trivialization, $\epsilon=\operatorname{sign}H$ is smooth and
\[
 \|D\epsilon_H[E]\|\le\frac{\|E\|}{\gamma},\qquad
 \|D^n\epsilon_H[E_1,\ldots,E_n]\|
 \le\frac{n!\,\Gamma(n/2)}{\sqrt\pi\,\Gamma((n+1)/2)}
       \gamma^{-n}\prod_i\|E_i\|\quad(n\ge1).
\]
Here $\Gamma(\cdot)$ is the Euler gamma function. The overlap projector
$P=\widehat P_-=(I+\epsilon)/2$ defines the connection and curvature
\begin{equation}
 \nabla=P\,d,\qquad \mathcal F=P(dP\wedge dP)P,
 \qquad
 \frac{|\operatorname{Tr}\mathcal F(\xi,\zeta)|}{\dim H}
 \le\frac{\|H_\xi\|\|H_\zeta\|}{2\gamma^2}.
 \label{p10:eq:v202-projector-connection}
\end{equation}
Its rank is $431k$ on the positive-winding neighborhood. Along a specified
smooth path, $U'=[P',P]U$, $U(0)=I$, has a unique unitary solution with
$P(s)U(s)=U(s)P(0)$.
\end{theorem}
\begin{proof}
Paper~1B, Theorem~\ref{p1b:thm:locality}
\cite{Kugelmann2026Paper1B}, gives the smooth sign derivatives,
projector connection and unitary path transport on this gapped family.
Here $\widehat\Gamma_5=\Gamma_5(I-a_RD_{\rm ov})=-\epsilon$,
so $P=(I+\epsilon)/2$. The native inertia counts give
$\operatorname{rank}P=431k$. Finally
$\|dP\|\le\|dH\|/(2\gamma)$ and
$\|[P_\xi,P_\zeta]\|\le2\|P_\xi\|\|P_\zeta\|$ imply
\eqref{p10:eq:v202-projector-connection} by
$|\operatorname{Tr}A|\le(\dim H)\|A\|$.
\end{proof}


\begin{theorem}[Native chiral charts and determinant-line overlap data]
\label{p10:thm:v203-chiral-norm-ball-chart}
Fix an invertible finite Hermitian centre $H_0$ with gap $\gamma_0>0$,
put $P_0=(I+\operatorname{sign}H_0)/2$, and suppose
$\|H-H_0\|\le\delta<\gamma_0$. Write
$P=(I+\operatorname{sign}H)/2$, $\gamma=\gamma_0-\delta$ and
$\|H\|\le M$. On $\operatorname{Ran}P_0$,
\[
 G=P_0PP_0\ge\frac{2\gamma}{M+\gamma}I,\qquad
 T(H,H_0)=PP_0G^{-1/2}.
\]
Extend $T$ by zero on $\ker P_0$. Then $T$ is smooth and
\[
 T^*T=P_0,\qquad TT^*=P,\qquad PT=TP_0=T.
\]
It is natural under simultaneous unitary conjugation of both centres and
kernels. Its top exterior power is a relative determinant-line
trivialization. If centre frames give local orthonormal frames
$V_\alpha(H)$, the overlap functions obey
\[
 g_{\alpha\beta}=\det(V_\alpha^*V_\beta)\in U(1),\qquad
 g_{\alpha\beta}g_{\beta\gamma}=g_{\alpha\gamma}.
\]
\end{theorem}
\begin{proof}
Apply Paper~1B, Proposition~\ref{p1b:result:explicit-chiral-bundle-charts}
\cite{Kugelmann2026Paper1B}, to this base kernel $H_0$ and norm ball.
Its positive-projector bound gives the displayed $G$ and its polar map
is the present $T$. Extending by zero on $\ker P_0$ gives
$PT=TP_0=T$. The proposition's smoothness, covariance and determinant
cocycle therefore hold for the chosen native parameter family and frames.
\end{proof}
The existing norm bounds for each certified finite clock-cover
neighborhood supply these inputs. The centre matrix path is used only to
prove a chart in the finite ambient Hermitian space; it is not asserted to
be an interpolation by native gauge configurations. A centre line phase
is not selected. This algebraic transport need not be parallel for the
induced connection and does not set its curvature to zero. Gauge maps may
move the centre: covariance with both centres transformed does not itself
give a gauge-compatible quotient measure. Global gluing and determinant-line integrability for a gauge-field-dependent
chiral presentation, together with unrestricted clutching comparison where
claimed, are exported continuum-certification obligations.  Native sector
weights/susceptibility remain dynamical readouts, not a G20 selector.


\begin{corollary}[The existing all-link certificate supplies the chart hypotheses]
\label{p10:cor:v204-native-chart-discharge}
On the certified family of
Theorem~\ref{p10:thm:m8td-all-link-neighborhood}, take
$H_0=H_k^*$, $g_*=417/7000$, $b<g_*$ and $\gamma=g_*-b$.
The existing norm bound of
Theorem~\ref{p10:thm:v202-native-overlap-locality} permits
$M=11+\epsilon_\rho$. The chart of
Theorem~\ref{p10:thm:v203-chiral-norm-ball-chart} then satisfies
\[
 P_0P(H_k)P_0\ge\frac{2(g_*-b)}{M+g_*-b}I
 \quad\hbox{on }\operatorname{Ran}P_0.
\]
Thus chart existence and the relative determinant transition functions
require no additional gap or local-triviality assumption on this family.
\end{corollary}
\begin{proof}
The native certificate already proves
$\operatorname{gap}H_k^*>g_*$ and $\|H_k-H_k^*\|\le b$;
Theorem~\ref{p10:thm:v202-native-overlap-locality} gives
$M=11+\epsilon_\rho$.  Paper~1B,
Corollary~\ref{p1b:cor:certified-chart}
\cite{Kugelmann2026Paper1B}, records exactly this inherited finite-link
certificate specialization; equivalently, substitute these estimates into
Proposition~\ref{p1b:result:explicit-chiral-bundle-charts}, as in the
preceding local theorem.
All central constants now follow from the exact integer certificate of
Theorem~\ref{p10:thm:v256-exact-native-inertia}.
\end{proof}
The outstanding global measure problem concerns compatibility of these
already constructed local charts with the complete gauge/action data and
their extension to the required topology domain.

Locality and connection existence are now derived on this entire certified
domain. A gauge-compatible physical chiral measure still requires its
integrability and global phase data; an induced connection does not select
them~\cite{Luscher2001ChiralRevisited}. The optional positive real
vectorlike determinant in Appendix~\ref{p10:app:v202-vectorlike} has a
different, explicitly declared mass branch. None of these statements is a
physical mass-gap or unrestricted native index theorem.
% END V202 NATIVE OVERLAP LOCALITY



In the nontrivial color doublet the spatial antipode \(q\mapsto-q\)
also exactly intertwines the two temporal spin signs: spatial relative
links and tangents are unchanged, while the closing doublet link changes
sign.  The spectator color line is a separate free operator; this is not an
isospectrality claim for arbitrary full \(SU(3)\) gauge fields.

The overlap index and additivity are established on these displayed
\(\pm k\) clock-cover families, profile tubes and same-bundle all-link
neighborhoods of Theorem~\ref{p10:thm:m8td-all-link-neighborhood};
Theorem~\ref{p10:thm:v161-positive-primitive-native-sectors} also proves that
unit-charge native sectors have positive pure-heat magnitude weight.
The massive and unsaturated massless weights have the separate outcomes
of Remark~\ref{p10:rem:v2189-weight-scoped-primitive-activity}. Hence native
finite topology no longer waits for an unrestricted overlap/clutching theorem.
The unrestricted intrinsic interpolation and equality
\(
u_{
m ov}=
u_{
m clutch}\) remain only where one exports the native
finite index to that geometric/continuum clutching presentation.  They are
therefore routed to the continuum/comparison certificate.  The uniformity
here is over clock covers, not spatial/ultraviolet refinement or a physical
fermion mass gap.  Appendix~\ref{p10:app:v199-topology-comparison} retains a
sufficient topology-resolving comparison criterion.  On the additional fixed
Euclidean pure-Yang--Mills branch,
Appendix~\ref{p10:app:v199-cartan-yang-mills} derives the transgression,
classical kink action and homogeneous fluctuation operator; those results do
not assert a selected history or a stationary full effective action.



\section{Hurwitz baryon spin--flux and same-action source geometry}
\label{p10:sec:v237-a7}
The three-face fan, signed contractions and native baryon spin--flux construction remain in the primary kernel. Their exact finite geometric moments are not automatically hadron mass or magnetic-moment matrix elements.

% Source Paper VII-A, original line 8634; inherited proof bodies unchanged.
\subsection{Hurwitz three-face baryon spin--flux geometry}
\label{p10:sec:v2220-hurwitz-baryon-flux}
\noindent\textit{Status: exact finite geometric, conditional kinetic and
pure-gauge reference results; no physical baryon pole or hyperfine mass
is claimed.}

The inherited physical spin doublet and its angular-momentum addition,
finite color singlet, Hurwitz-shell Dirac stencil, and the licensed
three-face heat block are used without introducing another spin field,
color coupling, or kinetic coefficient.  Throughout this section the
Hurwitz vertices belong to the compact \emph{momentum} carrier; a face
area is not an asserted spatial hadron radius.  The full interacting
Hamiltonian and its spectral readouts remain those of
Theorem~\ref{p10:thm:finite-qcd-matrix-compiler}.


% Source Paper VII-A, original line 8649; inherited proof bodies unchanged.
\subsubsection{Canonical edge fan and the existing D21 baryon channel}
\label{p10:subsec:v2220-edge-fan}
Let $\mathscr H$ be the unit Hurwitz $24$-cell, with edge and triangular
face sets $E_H$ and $F_H$.  Its verified incidence data are
\[
 |E_H|=|F_H|=96,\qquad
 |\{f\in F_H:e\subset f\}|=3\quad(e\in E_H).
\]
Choose an oriented edge $e=(p,q)$ and denote the third vertices of its
three incident faces by $r_i$, $i=1,2,3$.  For the already declared
oriented $SU(3)$ link variables $U_{ab}$, $U_{ba}=U_{ab}^{-1}$, put
\begin{equation}
 \label{p10:eq:v2220-edge-fan-paths}
 A_i(e)=U_{qr_i}U_{r_ip},\qquad
 U_{f_i}=U_{pq}A_i(e).
\end{equation}
The first expression is the unique two-edge return path in the selected
face, with no extra link variable or path choice.  Under
$U_{ab}\mapsto g_aU_{ab}g_b^{-1}$, all three $A_i$ transform as
$A_i\mapsto g_qA_i g_p^{-1}$.

\begin{proposition}[Existing baryon determinant on a Hurwitz edge fan]
\label{p10:prop:v2220-edge-fan-baryon}
The polarized color determinant
\begin{equation}
\label{p10:eq:v2220-baryon-wilson}
 W_e=\frac16\epsilon_{abc}\epsilon_{a'b'c'}
 (A_1)_{aa'}(A_2)_{bb'}(A_3)_{cc'}
\end{equation}
is gauge invariant and independent of the enumeration of the three
faces.  It agrees with the same polarized determinant evaluated on the
three face holonomies $U_{f_i}$.  Moreover it is precisely the
fundamental three-character Haar contraction in D21:
\[
 \int_{SU(3)}\prod_{i=1}^3\operatorname{tr}(A_iU)\,dU=W_e.
\]
Thus the corresponding \emph{local three-face} heat expansion has the
already fixed contribution $54e^{-2t}\operatorname{Re}W_e$.
\end{proposition}
\begin{proof}
The epsilon tensor transforms with the determinant of each common left
or right matrix; these determinants are one for $SU(3)$.  The same
argument removes the common factor $U_{pq}$ from the three $U_{f_i}$.
The Haar identity is the unique invariant projector
$\int U^{\otimes3}dU=|\epsilon\rangle\langle\epsilon|$, with
$|\epsilon\rangle=\epsilon/\sqrt6$; the heat coefficient is exactly the
existing D21 fundamental-character coefficient in
Section~\ref{p10:sec:v2201-derived-d21}.
\end{proof}
The equality identifies an operator \emph{channel} in the licensed D21
block.  Integrating a link in the \emph{complete} interacting measure also
encounters fermion, ribbon, and other gauge terms and is not in general
the isolated D21 integral.


% Source Paper VII-A, original line 8703; inherited proof bodies unchanged.
\subsubsection{Conditional all-orders spin exchange and the first flux term}
\label{p10:subsec:v2220-spin-exchange}
Let $X_H[U]$ be the inherited Hermitian Weyl/Levi--Civita finite
spin--color \emph{comparison} stencil, with dimensionless comparison
parameter $x$.  The one-quark fixed-link block is
$g_{pq}(x;U)=(e^{-xX_H[U]})_{pq}$, with two spin and three color indices.
On the three-spin tensor product define the color-determinant contraction
\begin{equation}
\label{p10:eq:v2220-singlet-transfer}
 G_{pq}(x;U)_{s_1s_2s_3,t_1t_2t_3}
 =\frac16\epsilon_{abc}\epsilon_{def}
 \prod_{j=1}^{3}(g_{pq})_{s_ja_j,t_jd_j}.
\end{equation}
All three factors in this \emph{comparison sector} use the same
one-body block.  This is not the physical finite Hamiltonian transfer
with different light and strange mass sources.

Let $\mathbf J=\tfrac12\sum_{j=1}^3\boldsymbol\sigma_j$ (in units of
$\hbar$),
$\Omega=\sum_{i<j}\boldsymbol\sigma_i\cdot\boldsymbol\sigma_j
=2(\mathbf J^2-\tfrac94I)$, and
\[
 P_{3/2}=\tfrac12(I+\Omega/3),\qquad
 P_{1/2}=\tfrac12(I-\Omega/3).
\]
The latter projects onto both spin-$\tfrac12$ copies in the eight-spin
space.  Put
\[
 S_J(x;U)=\frac1{4}\operatorname{Tr}_{(\mathbb C^2)^{\otimes3}}
   \bigl(P_JG_{pq}^{\dagger}G_{pq}\bigr)
 \quad(J=\tfrac12,\tfrac32).
\]

\begin{theorem}[Exact conditional spin comparison and its first flux coefficient]
\label{p10:thm:v2220-conditional-spinflux}
For every common one-body spin--color block in
\eqref{p10:eq:v2220-singlet-transfer},
\begin{equation}
\label{p10:eq:v2220-spin-trace}
 S_{3/2}-S_{1/2}
 =\frac1{12}\operatorname{Tr}_{(\mathbb C^2)^{\otimes3}}
    (\Omega G_{pq}^{\dagger}G_{pq}).
\end{equation}
For the licensed Hurwitz comparison stencil and neighboring $p,q$,
flat color links give $G_{pq}^{\dagger}G_{pq}=s(x)^3I_8$ with
$s(x)\ge0$ and hence $S_{3/2}=S_{1/2}$ at every $x$.  For arbitrary
fixed link configurations, the \emph{complete} conditional propagator,
with all intermediate paths included, has
\begin{equation}
\label{p10:eq:v2220-first-spinflux}
 S_{3/2}(x;U)-S_{1/2}(x;U)
 =x^8D_e(U)+O(x^9),\quad
 D_e(U)=\sum_{f\supset e}
 \left(\left|\frac{\operatorname{tr}U_f}{3}\right|^2
       -\operatorname{Re}\frac{\operatorname{tr}U_f}{3}\right).
\end{equation}
The coefficient is gauge invariant and is zero for flat links; its sign
is not universally fixed.
\end{theorem}
\begin{proof}
The color determinants at input and output produce the same sign under
simultaneous leg permutation.  Since $\Omega$ has eigenvalue $+3$ on
the four-dimensional spin-$\tfrac32$ sector and $-3$ on the two
spin-$\tfrac12$ copies of total dimension four, taking normalized
sector traces gives \eqref{p10:eq:v2220-spin-trace}.
For flat links the inherited quaternionic free-shell propagator obeys
$g_{pq}^{\dagger}g_{pq}=s(x)I_2$, hence its threefold singlet
contraction is spin-isometric.  In the nonflat case expand the full one-body exponential,
retaining every intermediate Hurwitz path.  The $x^8$ coefficient of the
squared three-quark transfer depends on amplitude orders three, four and
five, not on the sixth-order partitions $(1,1,4)$, $(1,2,3)$ and
$(2,2,2)$.  The analytic all-link reduction in
Theorem~\ref{p10:thm:v230-all-link-spinflux} gives the exact three
face traces in \eqref{p10:eq:v2220-first-spinflux} and cancels the
lower spin contrasts.  The finite-matrix regression checks remain valid.
\end{proof}

% BEGIN FRAMEWORK v2.230: ALL-LINK HURWITZ SPIN-FLUX ANALYTIC PROOF
\begin{theorem}[Analytic three-face spin--flux coefficient for every fixed link field]
\label{p10:thm:v230-all-link-spinflux}
For the common-one-body finite comparison stencil of
Theorem~\ref{p10:thm:v2220-conditional-spinflux}, the leading
spin contrast in \eqref{p10:eq:v2220-first-spinflux} holds for
\emph{every} fixed assignment of oriented $SU(3)$ edge links.
It follows from the cubic, quartic and quintic coefficients of the
three-quark amplitude alone.  Neither an independent distribution of links
nor its sixth-order amplitude term is required.
\end{theorem}
\begin{proof}
Define the polarized singlet contraction
\[
 {\cal E}(A,B,C)_{s_1s_2s_3,t_1t_2t_3}
  =\frac16\epsilon_{abc}\epsilon_{def}
       A_{s_1a,t_1d}B_{s_2b,t_2e}C_{s_3c,t_3f},
\]
with $6\times6$ spin--color endpoint matrices.  For an oriented edge
$e=(p,q)$, set $A=(X_H)_{pq}$, $B=(X_H^2)_{pq}$ and
$C=(X_H^3)_{pq}$.  Expanding the \emph{full} one-body exponential first,
\[
 g_{pq}(x)=-xA+\frac{x^2}{2}B-\frac{x^3}{6}C+O(x^4),
\]
gives
\[
 G_{pq}(x)=-x^3C_3+x^4C_4-x^5C_5+O(x^6),
\quad
\begin{aligned}
 C_3&={\cal E}(A,A,A),\\
 C_4&=\tfrac12\sum_{\rm legs}{\cal E}(B,A,A),\\
 C_5&=\tfrac14\sum_{\rm pairs}{\cal E}(B,B,A)
       +\tfrac16\sum_{\rm legs}{\cal E}(C,A,A).
\end{aligned}
\]
Here the sums run over all distinct placements of the indicated blocks.
Write $T(W)=\operatorname{Tr}(\Omega W)/12$.  The $x^8$
coefficient in $T(G_{pq}^*G_{pq})$ is exactly
\begin{equation}\label{p10:eq:v230-order-eight-input}
 T(C_4^*C_4+C_3^*C_5+C_5^*C_3).
\end{equation}
There is no $C_6$ contribution at that order.

The direct edge spin factor and link are unitary.  Removing them by a
common endpoint-frame change puts the three common-neighbor two-hop
blocks in the form
\[
 A=I_2\otimes I_3,\qquad
 B=\sum_{i=1}^3R_i\otimes V_i,\qquad V_i\in SU(3),
\]
where $V_i$ are the oriented triangular face holonomies, up to inversion.
The three edge-adapted quaternionic matrices satisfy
\[
 R_i^*=-R_i,\qquad
 R_iR_j+R_jR_i=-2\delta_{ij}I_2.
\]
The Hurwitz $24$-cell has exactly three common neighbors per edge, and
its canonical Levi--Civita edge stencil satisfies these relations on
all $96$ edges.  A simultaneous spin-frame rotation may thus use
$R_i=-i\sigma_i$ for the Pauli matrices.  Put
$z_i=\operatorname{tr}(V_i)/3$.
The one-color epsilon contraction
${\cal E}_{\rm color}(V,I,I)=\operatorname{tr}(V)/3$
yields
\[
 C_3=I_8,\qquad
 C_4=\frac12\sum_{i=1}^3z_i\sum_{\ell=1}^3R_i^{(\ell)}.
\]
One-leg operators have zero $T$-trace, while for distinct legs
\[
 T(R_i^{(\ell)}R_j^{(m)})=-\frac23\delta_{ij}
 \qquad(\ell\ne m).
\]
Consequently, the six ordered two-leg contractions give
\begin{equation}\label{p10:eq:v230-flux-square}
 T(C_4^*C_4)=\sum_{i=1}^3|z_i|^2.
\end{equation}

The $(1,1,3)$ placements of $C_5$ contain only one-leg spin operators
when paired with $C_3$, so their $T$-traces vanish.  For the
$(1,2,2)$ placements, the two-color epsilon contraction is
\[
 \kappa_{ij}={\operatorname{tr}(V_i)\operatorname{tr}(V_j)
                    -\operatorname{tr}(V_iV_j)\over6},\qquad
 C_5^{(1,2,2)}=\frac14\sum_{\ell<m}\sum_{i,j=1}^3
            \kappa_{ij}R_i^{(\ell)}R_j^{(m)}.
\]
The spin trace removes all $i\ne j$.  The $SU(3)$ determinant-one
Cayley--Hamilton identity gives
\[
 \kappa_{ii}={\bigl(\operatorname{tr}V_i\bigr)^2
                       -\operatorname{tr}(V_i^2)\over6}
       ={\operatorname{tr}(V_i^{-1})\over3}=\overline{z_i}.
\]
Using the three unordered leg pairs and adding the Hermitian conjugate,
\begin{equation}\label{p10:eq:v230-flux-cross}
 T(C_3^*C_5+C_5^*C_3)=-\sum_{i=1}^3\operatorname{Re}z_i.
\end{equation}
The $x^6$ and $x^7$ spin contrasts vanish by the same one-leg trace
identities.  Equations~\eqref{p10:eq:v230-order-eight-input}--
\eqref{p10:eq:v230-flux-cross} yield
\[
 S_{3/2}-S_{1/2}
 =x^8\sum_{i=1}^3\bigl(|z_i|^2-\operatorname{Re}z_i\bigr)
   +O(x^9)=x^8D_e(U)+O(x^9).
\]
Inverting a face orientation leaves the real class function unchanged.
The calculation is an exact fixed-background comparison \emph{amplitude}
identity; no interacting baryon pole, hyperfine energy or QCD vacuum
expectation is derived.
\end{proof}
% END FRAMEWORK v2.230: ALL-LINK HURWITZ SPIN-FLUX ANALYTIC PROOF

In particular, comparing $x^6$ coefficients \emph{after squaring} is
not equivalent to the coefficient of the squared full propagator.
Neither this squared comparison transfer nor its sign is a baryon pole.


% Source Paper VII-A, original line 8897; inherited proof bodies unchanged.
\subsubsection{Sharp bounds, center sensitivity and whole-shell reference moments}
\label{p10:subsec:v2220-global-flux}
Define on each actual triangular face
\begin{equation}
\label{p10:eq:v2220-face-functional}
 F(U_f)=\left|\frac{\operatorname{tr}U_f}{3}\right|^2
       -\operatorname{Re}\frac{\operatorname{tr}U_f}{3}.
\end{equation}
This is a bounded, real, orientation-independent gauge class function.

\begin{theorem}[Sharp local spin--flux bounds and an invariant global observable]
\label{p10:thm:v2220-sharp-flux-bounds}
For every $U\in SU(3)$,
\begin{equation}
\label{p10:eq:v2220-local-bounds}
 -\frac14\le F(U)\le\frac32,\qquad
 -\frac34\le D_e(U)\le\frac92.
\end{equation}
All four local endpoints are attained by admissible link
configurations.  For a central face holonomy $U_f=\omega^kI_3$,
$\omega^3=1$, the contributions are $F(I_3)=0$ and
$F(\omega I_3)=F(\omega^2I_3)=3/2$.
The root-independent multiplication observable
\begin{equation}
\label{p10:eq:v2220-global-observable}
 \mathcal F_H(U)=\frac1{96}\sum_{f\in F_H}F(U_f)
 =\frac1{288}\sum_{e\in E_H}D_e(U)
\end{equation}
is invariant under local $SU(3)$ gauge transformations, face
orientation reversal, and Hurwitz-cell automorphisms.  It is bounded in
$[-1/4,3/2]$; these global bounds are not asserted to be jointly sharp.
\end{theorem}
\begin{proof}
For $z=\operatorname{tr}U/3$,
$F(U)=|z-1/2|^2-1/4$, proving the lower bound, attained for
$U=\operatorname{diag}(e^{i\theta},e^{-i\theta},1)$ with
$\cos\theta=1/4$.  The eigenphases of an $SU(3)$ matrix give
$|z|\le1$ and $\operatorname{Re}z\ge-1/2$; the latter follows by
minimizing $\cos\alpha+\cos\beta+\cos(\alpha+\beta)$.  Hence
$F\le3/2$, attained by either nontrivial center element.
Each edge fan has three faces, whose holonomies can be assigned
independently using their exclusive detour links, giving sharp local
bounds for $D_e$.  Finally every face contains three edges, so
$\sum_{e\in E_H}D_e=3\sum_{f\in F_H}F(U_f)$.  Each summand is a real
class function invariant under inversion, proving the stated invariances.
\end{proof}

\begin{theorem}[Strict finite pure-heat fluctuations and exact Haar benchmark]
\label{p10:thm:v2220-haar-flux}
For any strictly positive normalized pure-gauge link density on the
finite full Hurwitz shell, with respect to product Haar measure,
\begin{equation}
\label{p10:eq:v2220-positive-variance}
  \operatorname{Var}(\mathcal F_H)>0.
\end{equation}
For the specific reference law with all $96$ unoriented edge links
independently Haar distributed, the exact whole-shell values are
\begin{equation}
\label{p10:eq:v2220-haar-reference}
 \mathbb E_{\rm Haar}\mathcal F_H=\frac19,\qquad
 \operatorname{Var}_{\rm Haar}(\mathcal F_H)=\frac{11}{15552}.
\end{equation}
\end{theorem}
\begin{proof}
Changing only one Hurwitz link to the nontrivial central value
$\omega I_3$ gives $\mathcal F_H=3/64$, while changing it to
$\operatorname{diag}(e^{i\theta},e^{-i\theta},1)$ with
$\cos\theta=1/4$ gives $-1/128$.  The other links remain identity.
Thus the continuous observable is nonconstant.  A strictly positive
density gives positive measure to neighborhoods of both configurations,
so its variance is strictly positive.
For the independent Haar reference, the fundamental character identity
$\chi_3\chi_{\bar3}=1+\chi_8$ yields
\[
 F(U)=\frac{1+\chi_8(U)}9
     -\frac{\chi_3(U)+\chi_{\bar3}(U)}6,
\qquad
 \mathbb E F=\frac19,\quad\operatorname{Var}(F)=\frac{11}{162}.
\]
Distinct triangular faces share at most one edge.  Conditioning on
that common edge still leaves an exclusive independent Haar factor in
each face holonomy, so distinct face holonomies are pairwise
independent (although not necessarily jointly independent).  Variances
of the $96$ summands therefore add, giving \eqref{p10:eq:v2220-haar-reference}.
\end{proof}
The positive-heat statement assumes a \emph{strictly positive full
finite pure-gauge} density; a signed fermion determinant or an incomplete
state-preparation measure cannot inherit it automatically.  The
independent-Haar values are reference data, not interacting QCD-vacuum
expectations, baryon mass fractions, or physical scale settings.


% BEGIN FRAMEWORK v2.235: SIGNED HURWITZ FAN AND RELATIVE-SOURCE PROVIDERS

% Source Paper VII-A, original line 8990; inherited proof bodies unchanged.
\subsubsection{Common/relative fan incidence and signed gauge observables}
\label{p10:sec:v235-signed-hurwitz-fan}
The finite link statements below are made on the same \(96\)
Hurwitz edge links and the existing two-hop stencil of
Theorem~\ref{p10:thm:v230-all-link-spinflux}.  They do not add a
particle carrier or a term to the undeformed strong Hamiltonian.

\begin{proposition}[Triality relative-source matrix of the three-face fan]
\label{p10:prop:v235-triality-relative-source}
Label the three adjacent octahedra by the dual triality classes of
Corollary~\ref{p1:cor:v235-dual-cell-triality}, and label the three
incident faces by the class of the octahedron \emph{omitting} each
face.  With the compatible cyclic orientation, the signed relative
incidence operator is, up to simultaneous relabeling,
\[
 M_e=\begin{pmatrix}0&-1&1\\1&0&-1\\-1&1&0\end{pmatrix},
 \qquad M_e^{\mathsf T}M_e=3I_3-\mathbf1\mathbf1^{\mathsf T}.
\]
The already generated D21/D22 three-face heat-action Hessian in
Section~\ref{p10:sec:v2201-derived-d22} consequently equals
\[
 \mathcal H_{ij}(X,Y)=
 k(t)(M_e^{\mathsf T}M_e)_{ij}\operatorname{tr}_{\mathbf3}(XY)
\]
on its actual three Lie-algebra source labels.  This is a match of
the \emph{relative label metric}, not an identification of the
scalar flux functions with physical fermion currents.
\end{proposition}
\begin{proof}
The three octahedra each contain two of the three faces.  Their
alternating incidence gives the displayed directed triangle matrix;
its Gram matrix is the projector onto the sum-zero plane multiplied
by three.  Substitution into the already evaluated heat Hessian
\(k(t)(3\delta_{ij}-1)\operatorname{tr}_{\mathbf3}(XY)\)
proves the same-block identity.
\end{proof}

\begin{corollary}[Lossless common/relative face analysis]
\label{p10:cor:v235-common-relative-frame}
Let \(B=\partial_2\) be the positive-oriented \(96\times96\)
face-to-edge incidence matrix of
Theorem~\ref{p1:thm:v235-hurwitz-chain-complex}, and let \(R\)
collect the three-component outputs \(M_ef_e\) of all \(96\)
edge fans from the formal face vector \(f\).  Then
\[
 B^{\mathsf T}B+R^{\mathsf T}R=9I_{96},\qquad
 f=\frac19\bigl(B^{\mathsf T}Bf+R^{\mathsf T}Rf\bigr).
\]
Equivalently, with \(D=Bf\), \(W=Rf\),
\(9\|f\|^2=\|D\|^2+\|W\|^2\).
The \(23\)-dimensional face-boundary kernel invisible to \(D\) is
recoverable when the relative mathematical channel \(W\) is also
retained; this is not an independent physical measurement claim.
\end{corollary}
\begin{proof}
Every face appears in exactly three edge fans.  In one fan,
\(M_e^{\mathsf T}M_e\) has diagonal entries two and off-diagonal
entries minus one; the common row has three unit entries.
The three diagonal contributions sum to nine, and each pair of
faces shares at most one edge fan, in which its common and relative
cross entries cancel.  Hence the displayed tight-frame identity.
\end{proof}

% BEGIN FRAMEWORK v3.026: stable common/relative inverse and typed geometric matching.
\begin{corollary}[Stable reconstruction from the formal common/relative fan frame]
\label{p10:cor:v3026-common-relative-stability}
For the \emph{mathematical} face coefficients $f\in\mathbb C^{96}$
(or $\mathbb R^{96}$), define $\mathcal T f=(Bf,Rf)$ as in
Corollary~\ref{p10:cor:v235-common-relative-frame}.  From
$\mathcal T^\dagger\mathcal T=9I_{96}$, noisy readouts
$\widetilde D=Bf+\eta_D$, $\widetilde W=Rf+\eta_W$ have
canonical least-squares reconstruction
\[
\boxed{\widehat f=\tfrac19
(B^\dagger\widetilde D+R^\dagger\widetilde W),\qquad
\|\widehat f-f\|\le\tfrac13
\sqrt{\|\eta_D\|^2+\|\eta_W\|^2}.}
\]
On $f\in\ker B$, $\|Rf\|^2=9\|f\|^2$, including the
23-dimensional face-boundary kernel. These are identities for two
mathematically available channels, not claims of two independent
laboratory detectors or new physical currents.
\end{corollary}
\begin{proof}
$\widehat f-f=\mathcal T^\dagger(\eta_D,\eta_W)/9$;
all nonzero singular values of $\mathcal T$ equal three, giving
$\|\mathcal T^\dagger\|/9=1/3$.
The last identity follows by evaluating
$B^\dagger B+R^\dagger R=9I$ on $\ker B$.
\end{proof}

\begin{proposition}[Five-core transport is a conditional operator test]
\label{p10:prop:v3026-five-core-physical-matching-gate}
Let $C_m:\mathcal V_{\rm loc}\to\mathcal V_{\rm rec}$ denote the
\emph{real geometric} maps of Paper~I-B,
Corollary~\ref{p1:cor:v3026-five-core-hs-gram}, where
$\dim_{\mathbb R}\mathcal V_{\rm loc}=14$ and
$\dim_{\mathbb R}\mathcal V_{\rm rec}=6$.
For any \emph{independently supplied} same-action physical Hilbert
sectors $\mathcal H_{\rm loc},\mathcal H_{\rm rec}$, suppose there
are specified complex-linear isometric embeddings
$U_{\rm loc}:\mathcal V_{\rm loc}\otimes\mathbb C\to
\mathcal H_{\rm loc}$ and
$U_{\rm rec}:\mathcal V_{\rm rec}\otimes\mathbb C\to
\mathcal H_{\rm rec}$ that satisfy the needed gauge, Gauss, spin,
charge, source and clock intertwiners. Then the transported matrices
$\widetilde C_m=U_{\rm rec}C_mU_{\rm loc}^\dagger$ preserve the
geometric Gram and the four-dimensional span.
The existence of such embeddings and of an \emph{actual action-derived}
block $K:\mathcal H_{\rm loc}\to\mathcal H_{\rm rec}$ is a
\textbf{separate unproved physical matching obligation}; no five-core
transition rate, decay arrow, resonance pole or width is obtained by
projecting a freely chosen comparison matrix.
\end{proposition}
\begin{proof}
Isometry gives
$U_{\rm loc}^\dagger U_{\rm loc}=I$ and
$U_{\rm rec}^\dagger U_{\rm rec}=I$. Hence
$\operatorname{Tr}(\widetilde C_m^\dagger\widetilde C_n)
=\operatorname{Tr}(C_m^\dagger C_n)$.
The asserted occurrence restrictions are not consequences of this
algebraic identity; they remain premises of any physical application.
\end{proof}
% END FRAMEWORK v3.026: reconstruction and physical matching guard.

% BEGIN FRAMEWORK v3.069 -- source-level decay matching test.
\begin{definition}[Five-core baryon decay action-matching test]
\label{p10:def:v369-five-core-baryon-decay-action-test}
Let $J_{\rm geom}$ be the \emph{operator-space} isometry of
Corollary~\ref{p1b:cor:v369-five-core-operator-cubic} and let
$W_{56}$ be the \emph{different-type} spin--flavor comparison isometry
of Theorem~\ref{p11:thm:v369-five-core-baryon-chart}.
Neither alone defines a map into the actual Gauss-reduced physical
Born/QCD Hilbert carrier.  For a proposed source, require the
\emph{same} action to supply one field-dependent, gauge-covariant
source/auxiliary pair $B(\Phi)$ and $H_Q(\Phi)$ on specified physical
Hilbert sectors, with correctly matched color, spin, flavor charge,
clock and daughter-channel projectors.  Wherever the resolvent exists,
put
\[
 R(z,\Phi)=(z-H_Q(\Phi))^{-1},\qquad
 \Sigma(z,\Phi)=B(\Phi)^\dagger R(z,\Phi)B(\Phi).
\]
For any field $\Phi_i$, write $B_i=\partial_iB$ and
$H_i=\partial_iH_Q$.  Its exact field-response identity is
\[
 \boxed{\partial_i\Sigma=B_i^\dagger RB+B^\dagger RB_i
 +B^\dagger RH_iRB.}
\]
The conditional physical decay-vertex test is whether the
\emph{on-shell, pole-residue-normalized, Ward-compatible} physical
current obtained from this \emph{same} action and its complete
field variation reproduces the baryon transitions, without
independently prescribing $G^{ia}$, channel-specific weights or an
arbitrary orientation phase.  A finite auxiliary $H_{\rm ori}$ alone
is not an absorptive retarded channel.
\end{definition}

% BEGIN v3.074 -- action-owned source-cyclic reduction and Schur response jets.
\begin{theorem}[Least source-visible realization without an independently chosen embedding]
\label{p10:thm:v374-source-cyclic-observable-realization}
Let $H=H^\dagger$ act on an \emph{already specified} finite
Gauss-reduced same-action Hilbert space $\mathcal H$ of dimension $d$.
Let $V:\mathbb C^m\to\mathcal H$ be a source matrix built from the
\emph{actual} prepared source vectors (for instance, action-defined
$O_a|\Omega\rangle$). Then
\[
 \boxed{\mathcal K_{H,V}=\sum_{n=0}^{d-1}H^n\operatorname{Ran}V}
\]
is the unique smallest $H$-reducing subspace containing those sources.
For every $z\notin\operatorname{Spec}(H)$ and $n\ge0$,
\[
 \boxed{V^\dagger(z-H)^{-1}V
 =V^\dagger(z-H|_{\mathcal K_{H,V}})^{-1}V,\qquad
 V^\dagger H^nV=V^\dagger(H|_{\mathcal K_{H,V}})^nV.}
\]
Writing $M_n=V^\dagger H^nV$, its dimension is directly readable
from the block moment Hankel matrix:
\[
 \boxed{\dim\mathcal K_{H,V}
 =\operatorname{rank}(M_{p+q})_{0\le p,q<d}.}
\]
Two source/Hamiltonian pairs with identical \emph{all} matrix moments
$M_n$ have unitarily equivalent minimal source-visible realizations.
The conclusion requires neither a freely selected comparison
$56\to240$ isometry nor an independent Schur connector. If a Schur
presentation is desired, the source projector is the orthogonal
projector onto $\operatorname{Ran}V$, its auxiliary is the orthogonal
complement \emph{inside} $\mathcal K_{H,V}$, and both $B=QHP$ and
$H_Q=QHQ$ are thereby fixed by the same $H,V$.
For several occurring Hermitian field insertions $T_i$, the analogous
iterated closure under $H,T_i$ and their adjoints is the least joint
reducing carrier and preserves all mixed resolvent/field-response
words formed from those operators and the given sources.
\end{theorem}
\begin{proof}
By Cayley--Hamilton, $H\mathcal K_{H,V}\subseteq\mathcal K_{H,V}$.
Hermiticity makes the orthogonal complement invariant, so this space
is reducing. Every invariant space containing $\operatorname{Ran}V$
contains every $H^nV$ and hence contains $\mathcal K_{H,V}$.
Resolvents and moments then restrict exactly. The block Krylov
matrix $\mathscr C=[V,HV,\ldots,H^{d-1}V]$ has Gram
$\mathscr C^\dagger\mathscr C=(M_{p+q})$, proving the rank formula.
When all moments agree between two models, the assignment
$H^nVx\mapsto (H')^nV'x$ preserves every inner product of finite
linear combinations; it therefore extends to a unitary intertwiner
between the minimal cyclic carriers. The joint claim follows by
iterating the spans under the finite $*$-closed generator set.
\end{proof}

\begin{proposition}[Full same-parent two-field Schur jet and spectral positivity]
\label{p10:prop:v374-schur-twojet-and-spectral-measure}
In Definition~\ref{p10:def:v369-five-core-baryon-decay-action-test},
let $B(\Phi),H_Q(\Phi)$ be twice differentiable at fixed spectral
argument $z$, with $H_Q(0)=H_Q(0)^\dagger$ and
$R=(z-H_Q(0))^{-1}$. For field derivatives denoted by subscripts,
\[
\begin{aligned}
\partial_{ij}\Sigma={}& B_{ij}^\dagger RB+B^\dagger RB_{ij}
 +B_i^\dagger RB_j+B_j^\dagger RB_i\\
&+B_i^\dagger RH_jRB+B_j^\dagger RH_iRB
 +B^\dagger RH_iRB_j+B^\dagger RH_jRB_i\\
&+B^\dagger R(H_{ij}+H_iRH_j+H_jRH_i)RB.
\end{aligned}
\]
No one-field current determines these additional contact terms
without information about the \emph{same action's} $B_{ij},H_i,H_{ij}$.
For the zero-field fixed self-adjoint $H_Q$ and bounded $B$, the
matrix-valued spectral measure $d\mu(\lambda)=
B^\dagger dE_{H_Q}(\lambda)B$ is positive and
\[
 \boxed{\Sigma(z)=\int\frac{d\mu(\lambda)}{z-\lambda},\quad
 \int d\mu=B^\dagger B,\quad
 \int\lambda\,d\mu=B^\dagger H_QB.}
\]
For $z=E+i\eta$ with $\eta>0$,
$-\operatorname{Im}\Sigma(z)=
\eta B^\dagger((E-H_Q)^2+\eta^2)^{-1}B\succeq0$.
On a real resolvent gap,
$\partial_E\Sigma=-B^\dagger(E-H_Q)^{-2}B\preceq0$.
If $\lambda$ is a simple full-Hamiltonian source-visible eigenvalue
outside $\operatorname{Spec}H_Q$, and
$u\in\ker[\lambda-H_P-\Sigma(\lambda)]$ is nonzero, then
\[
 \boxed{\operatorname{Res}_{z=\lambda}
  [z-H_P-\Sigma(z)]^{-1}
 =\frac{uu^\dagger}
 {\langle u,(I+B^\dagger(\lambda-H_Q)^{-2}B)u\rangle}.}
\]
These are exact pole/residue and positive spectral-measure identities,
not a computed open-channel width. In finite dimension the measure is
atomic and cannot alone generate an irreversible exponential lifetime.
\end{proposition}
\begin{proof}
Differentiate $(z-H_Q)R=I$ to obtain
$R_i=RH_iR$ and
$R_{ij}=R(H_{ij}+H_iRH_j+H_jRH_i)R$; the product rule gives the
stated two-field derivative. The spectral theorem gives the positive
measure, moments and imaginary-part identity. The real derivative is
$-B^\dagger R^2B$. For the residue, the full $P\oplus Q$
eigenvector is $(u,(\lambda-H_Q)^{-1}Bu)$ and has squared norm
$\langle u,(I+B^\dagger(\lambda-H_Q)^{-2}B)u\rangle$.
The displayed rank-one residue is its normalized source projection.
\end{proof}
% END v3.074 action-owned source-cyclic reduction and Schur response jets.


% BEGIN FULL-CARRIER SAME-ACTION SOURCE MATCHING v1.1.
\begin{definition}[Typed full-carrier current action-comparison map]
\label{p10:def:fullcarrier-current-action-comparison}
Retain the native source operator
$\mathcal O_A^{\rm BH}=\partial_{J_A}H_H[J]|_{J=0}$
with its previously specified finite-source and transfer domain.
Choose \emph{one} comparison branch
$r\in\{28,\mathrm{finite\ exact\ current\ dilation},
\mathrm{native\ other}\}$ and an explicitly type-matched isometry
$\mathcal I_r$ from its flavor-charge comparison carrier into a
specified locally Gauss-invariant Born/QCD carrier.  If the actions
and domains coincide, define the first-current mismatch
\[
 \boxed{\Delta_A^{(1)}(r)=
 \mathcal I_r^\dagger\mathcal O_A^{\rm BH}\mathcal I_r
 -\widehat Q_r(F_A).}
\]
The equation $\Delta_A^{(1)}(r)=0$ is a \emph{physical-action test},
not an inference from the existence of a representation or the
$A_5$ graph.  In particular a spinless $28$-state charge must not be
identified by type with the spatial axial current on the symmetric
$56$-state spin--flavor chart; in that rest-spin chart the conventional
source convention is $2G^{iA}$.
\end{definition}

\begin{definition}[Same-action Gauss and second source-jet gate]
\label{p10:def:fullcarrier-source-jet-and-gauss-gate}
The comparison current is physically licensed only if it arises
from the \emph{same} member $H_H[J]$ whose source contacts, Duhamel
responses, poles and daughter channels define the observable.  On
the actually admitted gauge-link and matter truncation require
\[
 P_{\rm Gauss}\mathcal I_r=\mathcal I_r,\qquad
 [P_{\rm Gauss},H_H[J]]=0
\]
on the joint source domain, plus the original color, baryon number,
flavor, spin, charge, clock and CAR/Wilson-source compatibilities.
The operator contact
$\partial_{J_A}\partial_{J_B}H_H[J]|_0$ is \emph{not}
identified with $\partial_{J_A}\partial_{J_B}\log\operatorname{Tr}
 e^{-tH_H[J]}|_0$: the latter contains Duhamel-ordered
current--current insertions and disconnected subtraction.
Use Theorem~\ref{p10:thm:v350-dynamical-source-hessian} and
Definition~\ref{p10:def:v369-five-core-baryon-decay-action-test}
for the source/return and pole requirements.  Nonzero current mismatch
or incompatible source Hessians invalidate the proposed physical
identification, not the underlying finite geometric existence theorem.
\end{definition}

\begin{proposition}[Exact conditional comparison-action families]
\label{p10:prop:fullcarrier-comparison-action-families}
On either exact geometric representation of
Theorems~\ref{p1b:thm:fullcarrier-minimal-su3-completion} and
\ref{p1b:thm:fullcarrier-finite-exact-current-dilation}, select a
self-adjoint free comparison operator $H_{0,r}$ and put
\[
 H_{\rm cmp,r}[J]=H_{0,r}+\sum_AJ_A\widehat Q_r(F_A).
\]
The first source derivatives are $\widehat Q_r(F_A)$ and the
\emph{Hamiltonian} second source-contact derivatives vanish for this
specified source-linear action.  This is a mathematical comparison,
not an action generated by G0/GP1--GP3 or the actual finite QCD
Hamiltonian.  In particular the high-dimensional extension of
$H_{0,r}$ outside the compressed graph image is not unique.
The earlier nonlinear spin-link action shares the selected triplet
first derivative but has different second contacts and must not be
silently conflated with this source-linear family.
\end{proposition}
\begin{proof}
Differentiate the explicitly source-linear self-adjoint comparison
family.  Its log-partition Hessian need not vanish because of
Duhamel-ordered correlations.
\end{proof}
% END FULL-CARRIER SAME-ACTION SOURCE MATCHING v1.1.

\begin{proposition}[Cubic five-core projection does not select a decay current]
\label{p10:prop:v369-cubic-five-core-projection-firewall}
Under the \emph{conditional} same-action embeddings of
Proposition~\ref{p10:prop:v3026-five-core-physical-matching-gate},
let $\widetilde C_m=U_{\rm rec}C_mU_{\rm loc}^\dagger$ and
$K:\mathcal H_{\rm loc}\to\mathcal H_{\rm rec}$ be a supplied
physical comparison block.  Its orthogonal projection onto the
five-core operator span is exactly
\[
 \boxed{K_{\rm geom}={40\over69}\sum_{m=1}^{5}
 \operatorname{Tr}(\widetilde C_m^\dagger K)\widetilde C_m,\qquad
 \|K\|_{\rm HS}^2=\|K_{\rm geom}\|_{\rm HS}^2
 +\|K-K_{\rm geom}\|_{\rm HS}^2.}
\]
The cubic map of Corollary~\ref{p1b:cor:v369-five-core-operator-cubic}
supplies a canonical unitary coordinate chart on this rank-four
span, but does \emph{not} determine whether $K$ occurs as the
field derivative of the actual QCD action or whether its
octet--decuplet matrix elements agree with the physical axial
current.  The bare orientation Hamiltonian of
Proposition~\ref{p11:prop:v369-oriented-core-auxiliary} has zero
direct pair-octet/decuplet blocks, so it cannot be substituted
for the physical decay vertex.
\end{proposition}
\begin{proof}
The five-core Gram is $(69/40)P_\perp$ and
$\sum_m\widetilde C_m=0$.  Its Moore--Penrose inverse is
$(40/69)P_\perp$, so the displayed finite-rank formula is the
Hilbert--Schmidt orthogonal projection.  Pythagoras gives its norm
identity.  The field/action occurrence does not follow from an
orthogonal projection, and the direct-vertex exclusion is the
representation-level kernel theorem in Paper~VII-B.
\end{proof}
% END FRAMEWORK v3.069 -- source-level decay matching test.

\begin{definition}[Two enumeration-independent signed fan functions]
\label{p10:def:v235-signed-fan-operators}
For the three rooted face holonomies \(V_i=U_{f_i}\) and their
quaternionic spin factors \(R_i\), \(i=0,1,2\), define
\[
 \begin{aligned}
 \chi_e&=\tfrac13\operatorname{Im}\operatorname{tr}
 \bigl(ABA^{-1}B^{-1}\bigr),&
 A&=V_0V_2^{-1},\quad B=V_1V_2^{-1},\\
 \psi_e&=\tfrac13\operatorname{Re}\operatorname{tr}
 (V_0V_1V_2-V_0V_2V_1),&
 \sigma_e&=\tfrac12\operatorname{tr}(R_0R_1R_2)\in\{\pm1\}.
 \end{aligned}
\]
The signed functions \(\widehat\chi_e=\sigma_e\chi_e\) and
\(\widehat\psi_e=\sigma_e\psi_e\) are gauge invariant and independent
of the enumeration of the three faces.  Under the ordinary
color-conjugation operation \(V_i\mapsto\overline{V_i}\), \(\chi_e\)
is odd and \(\psi_e\) is even.  This color parity is not an
identification with physical CP.
\end{definition}

\begin{theorem}[Exact cubic rooted Weyl contraction]
\label{p10:thm:v235-rooted-weyl-cubic}
The already appearing two-hop Weyl/Levi--Civita block in the
rooted direct-edge frame,
\(\mathcal B_e=\sum_{i=0}^2R_i\otimes V_i\), obeys
\[
 \boxed{\operatorname{Re}\operatorname{Tr}_{\rm spin\otimes color}
 (\mathcal B_e^3)=18\widehat\psi_e.}
\]
Thus the color-C-even signed function is generated as a cubic
polynomial contraction of the existing geometric comparison operator.
The identity does not, without an action-incidence certificate,
establish the occurrence of the cubic in the actual physical
fermionic determinant or Hamiltonian.
\end{theorem}
\begin{proof}
The Pauli/quaternionic relations of
Theorem~\ref{p10:thm:v230-all-link-spinflux} make the spin trace
of every cubic word with a repeated index vanish.  For distinct
indices it equals \(2\sigma_e\operatorname{sgn}(\pi)\).
The three cyclic permutations of each orientation have the same
color trace, so the cubic spin-color trace equals
\(6\sigma_e\operatorname{tr}(V_0V_1V_2-V_0V_2V_1)\).
Taking the real part proves the formula.
\end{proof}

\begin{proposition}[Exact signed Haar modes on the full link shell]
\label{p10:prop:v235-signed-haar-norms}
On \(\mathcal H_0=L^2(SU(3)^{96},d\mu_{\rm Haar}^{\otimes96})\),
with the physical gauge-invariant subspace understood, the \(192\)
real multiplication functions \(\{\widehat\chi_e,\widehat\psi_e:
e\in E_H\}\) obey
\[
 \begin{aligned}
 \langle\widehat\chi_e,\widehat\chi_f\rangle_0
   &=\tfrac5{144}\delta_{ef},&
 \langle\widehat\psi_e,\widehat\psi_f\rangle_0
   &=\tfrac8{81}\delta_{ef},&
 \langle\widehat\chi_e,\widehat\psi_f\rangle_0&=0.
 \end{aligned}
\]
The corresponding orthogonal projections \(P_\chi,P_\psi\)
onto the two \(96\)-dimensional spans are uniquely defined on
this \emph{untruncated product-Haar} Hilbert space.
\end{proposition}
\begin{proof}
The three rooted face holonomies in a fan are independent Haar
variables.  The group commutator \(W=ABA^{-1}B^{-1}\) for independent
Haar \(A,B\) satisfies
\(\mathbb E\operatorname{tr}W=1/3\),
\(\mathbb E|\operatorname{tr}W|^2=9/8\), and
\(\mathbb E(\operatorname{tr}W)^2=1/2\), by
\(3\otimes\bar3=1\oplus8\),
\(3\otimes3=6\oplus\bar3\), and character orthogonality.
Thus \(\mathbb E\chi_e^2=5/144\).
For \(\psi_e\), conditioning on any one \(V_i\) and using
matrix-entry orthogonality gives
\(\mathbb E[\psi_e^2\mid V_i]
=(8-\chi_8(V_i))/81\), hence \(\mathbb E\psi_e^2=8/81\).
The six detour links of distinct edge fans differ on at least four
exclusive links in each direction; the \(\psi\) modes use these six
links and the central edge.  Haar integration on an exclusive link
annihilates cross-edge overlaps.  Opposite color-conjugation parity
annihilates every mixed \(\chi\)--\(\psi\) overlap.
\end{proof}

\begin{corollary}[Conditional finite electric Casimir and source scope]
\label{p10:cor:v235-signed-electric-casimir}
Under the \emph{same-owner cylindrical cotangent realization} of
Theorem~\ref{p10:thm:v2159-conditional-electric-heat},
with its normalization \(\operatorname{tr}(T_aT_b)=\delta_{ab}/2\),
the signed modes obey
\[
 -\sum_{\ell\in E_H}\Delta_\ell\widehat\chi_e
 =18\widehat\chi_e,\qquad
 -\sum_{\ell\in E_H}\Delta_\ell\widehat\psi_e
 =11\widehat\psi_e.
\]
For common positive link coefficient \(\alpha\), their
\emph{pure-electric} eigenvalues are respectively
\(9\hbar^2\alpha\) and \(11\hbar^2\alpha/2\), with zero mixed
compression.  These eigenvalues are not hadron masses or interacting
QCD spectral gaps.
\end{corollary}
\begin{proof}
The commutator-trace function is independent of the central edge;
each of its six detour links carries an adjoint matrix coefficient
of Casimir three.  For the antisymmetric cubic trace the central
edge carries an adjoint coefficient of Casimir three and each of
its six detours a fundamental or antifundamental coefficient of
Casimir \(4/3\).  Bi-invariance of the Haar Laplacians yields
\(6\cdot3=18\) and \(3+6\cdot(4/3)=11\), respectively.
The conditional electric operator of the cited theorem then gives
the displayed eigenvalues.
\end{proof}
% END FRAMEWORK v2.235: SIGNED HURWITZ FAN AND RELATIVE-SOURCE PROVIDERS

\paragraph{Physical handoff.}
The operator $\mathcal F_H$ is already a gauge-invariant bounded
multiplication observable.  Its expectation in an actual baryon state
and its contribution to a mass-source derivative are distinct,
uncomputed descendants of the complete positive singlet transfer; see
Paper~VII-B.  The dimensionless factor $8/9$, the charged exclusion
factor $24/25$, and the $E_{36}$ rung do not acquire a universal hadron
mass interpretation from these geometric theorems.
% END FRAMEWORK v2.220: EXACT HURWITZ BARYON SPIN--FLUX OPERATOR

% BEGIN FRAMEWORK v2.226: KINEMATIC CLOSURES


\section{Finite QCD theorem exports and physical limits}
\label{p10:sec:v237-a8}
All completed finite constructs and their typed open physical/continuum obligations are collected here. The long auxiliary calculations reside in Paper~VII-D; Paper~VII-B evaluates readouts and Paper~VII-C handles bounded-continuum certification.

% Source Paper VII-A, original line 4261; inherited proof bodies unchanged.
% Former logical part: Paper-VII-A Completion and Paper-VII-B Export; integrated into the new thematic argument.
\label{p10:part:completion-paper-xi-export}


% Source Paper VII-A, original line 4264; inherited proof bodies unchanged.
\subsection{Finite QCD core completion}
\label{p10:sec:finite-qcd-core-completion}

\begin{definition}[Repaired Paper-VII-A finite QCD core datum]
\label{p10:def:finite-qcd-core-datum}
The rigorous Paper-VII-A export is
\begin{equation}
\label{p10:eq:finite-qcd-core-datum}
\mathfrak D_{7A}^{\rm QCD,core}:=
\left(\begin{gathered}
C,SU(3)_c,\Ps,\mathcal K_H,\Delta_{2,H},\\
\{H_L,\widehat K_L(t),\mathfrak A_{L,{\rm str}},
\mathfrak A'_{L,{\rm str}},\mathcal O_L,W_L[J]\}_L,\\
\mathfrak T_{\rm chiral},\mathfrak T_{\rm top}
\end{gathered}\right),
\end{equation}
where \(\mathfrak T_{\rm chiral}\) retains the precisely defined ordered-limit targets \((\Sigma,F_\pi,B_0,L_i,\ldots)\) in addition to the covered finite source responses.  The topological packet \(\mathfrak T_{\rm top}\) reuses the native finite overlap/index on its proved domain, the auxiliary Fourier functional \(Z_L(\theta)\), and charge covariance/cumulants; comparison limits are retained only where consumed.  The rooted charged finite mass/source member is inherited as an upstream generated output, while \(M_q^{\rm phys}(\mu)\) consumes the applicable execution-scale G16 transport.  A comparison-branch \(\alphasseed\) retains its named-scale matching status, and a physical \(\bar\theta\) action requires a separately licensed sector character.  On the source-native common-heat branch the gauge boundary itself is theorem-level and only its applicable scale/RG transport remains.
\end{definition}

\begin{corollary}[Finite QCD core is a generated structural descendant]
\label{p10:cor:v152-finite-qcd-generated-descendant}
Once the upstream color owner and the licensed Paper-VII-A constructors have been
generated, the finite carrier, gauge-covariant compiler, singlet algebra,
finite Yang--Mills Hamiltonian/transfer family, source algebra, Hodge data, and
finite spectral/source target maps collected in
\(\mathfrak D_{7A}^{\rm QCD,core}\) are structural descendants in the sense
of Paper~II-A.  They require no second premise asserting that ``QCD exists''
as a finite structural sector.  This statement introduces no new charged quark-mass member selection: that
finite source is already inherited from its upstream generated constructor.
It does not supply a particular populated state, hadron preparation, positive
unscreened string tension, scattering history, physical scale comparison, or
cross-action continuum universality claim.
\end{corollary}
\begin{proof}
Every listed object is produced by a licensed constructor from the generated
finite strong owner or is an explicitly typed target functional on that
owner.  Paper~II-A therefore co-realizes the resulting structural operation
graph without a duplicate occurrence root.  State/member selection,
numerical calibration, and analytic continuum control are different sorts and
are not outputs of that structural induction.
\end{proof}

\begin{theorem}[Paper-VII-A kernel-definition completion certificate]
\label{p10:thm:completion-certificate}
Paper~VII-A is closed as a finite kernel-definition and computation-target layer: the 24-cell geometry, exact singlet projection, finite compiler, source algebra, commutant decomposition, Hodge operators, and rigorous definitions of gap/chiral/topological readouts are fixed without downstream hadronic calibration.  Paper~VII-A does \emph{not} claim that the Yang--Mills mass gap, physical string tension, chiral low-energy constants, or dimensionful physical topological susceptibility have already been numerically solved.
The pure-heat and retained vectorlike finite charge-variance bounds, and
the exact unsaturated massless endpoint, are proved below on their
stated weight and integration domains. Each claim retains the hypotheses and domain printed in its own defining result.
\end{theorem}
\begin{proof}
Every object in the repaired export is either an exact finite construction, a finite spectral/source functional, or an explicitly marked conditional/limit obligation.  The withdrawn scale, gap, chiral-residue, string-tension, and topology decimals are not required to define the compiler.  Hence the definition/execution layer is closed without claiming that those later computations have already been solved.
\end{proof}

\begin{corollary}[Safe finite-QCD numerical/geometric package]
\label{p10:cor:numerical-finite-qcd-export-package}
The unconditional finite numerical package contains the 24-cell combinatorics and spectra, including the adjacency spectrum
\(8^{(1)}\oplus4^{(4)}\oplus0^{(9)}\oplus(-2)^{(8)}\oplus(-4)^{(2)}\),
the exact face-Hodge spectrum of Theorem~\ref{p10:thm:finite-plaquette-hodge-curvature-hessian}, and \(\deltageom=1/(6\sqrt2)\) as a geometric normalized count.  On the explicit conditional seed/coordinate diagnostic branch, the sharp one-loop crossing is \(0.2843536970\,\mathrm{GeV}\) and the old smooth mixed-clock functional has its zero at \(0.3158981634\,\mathrm{GeV}\).  No \(F_\pi\), \(B_0\), condensate, string tension, mass gap, or topological-susceptibility decimal is included before its correlator/spectral calculation is performed.
\end{corollary}
\begin{proof}
The listed unconditional numbers are obtained directly from the independently replayed 24-cell spectra and finite dimensions.  The comparison values are explicitly tagged by their seed/coordinate branch.  No numerical $F_\pi$, $B_0$, condensate, mass gap, string tension, or susceptibility is included because the preceding theorems leave those as spectral/correlator targets.
\end{proof}

\begin{theorem}[Acyclic Paper-VII-A to Paper-VII-B boundary]
\label{p10:thm:paper-x-xi-boundary}
The dependency direction is
\[
\mathfrak F_{\le6B}\longrightarrow\mathfrak D_{7A}^{\rm QCD,core}
\longrightarrow\text{Paper~VII-B executions}.
\]
Paper~VII-B may evaluate spectra, residues, distributions, thermal traces, transition measures, and multi-hadron blocks from the exported kernel, but comparison data may not flow backward to choose Paper-VII-A coefficients or to disguise an unresolved Paper-VII-A interface as a fitted theorem.
\end{theorem}
\begin{proof}
Paper~VII-A defines the finite matrices, source operators, and their logical statuses before any Paper-VII-B observable is evaluated.  Paper~VII-B operations are functions of that fixed tuple (diagonalization, matrix elements, traces, and cumulants).  Feeding their outputs back into a Paper-VII-A coefficient would reverse this functional dependency and is therefore excluded by construction.
\end{proof}



% Source Paper VII-A, original line 6176; inherited proof bodies unchanged.
\subsection{Earlier theorem synthesis and scope}
Paper~VII-A now separates what the finite Born--Hopf QCD carrier genuinely proves from what must be computed dynamically.  The 24-cell/Hurwitz geometry, exact color representation, singlet projection, finite source/compiler algebra, commutant decomposition, Hodge operators, and fixed-base regulator-chain ancestry are rigorous finite results.  Strong amplitudes, absolute nonperturbative scales, spectral gaps, physical string tension, chiral residues, and topological cumulants are defined as finite spectral/correlator functionals and are exported only with their actual proof status.  The spectral common-heat coefficient does not silently determine the compact-group heat time or the physical transfer time, and exact four-dimensional blocking of a product heat action generally produces multi-face effective interactions.  Center flux does not by itself force area extensivity.  The native finite GW/overlap and index are reused on their proved domain; actual topological action/state sector weights remain readouts of the same packet.  Release v1.64 factorizes the former monolithic \(\mathrm{PAIR}_{xp}^{\rm op}\) item into already generated position geometry/holonomy and the genuinely open positive-area theorem.  The light-front handoff is sharpened as well: a native equal-time Hopf-equivariant good-field seed plus the same finite QCD Hamiltonian/Hopf generator constructs the elapsed null action and forces \(P_{N,L}^+=P_{
m phys,L}^+=(H_L/c_{
m op}+P_{H,L})/\sqrt2\); strong commutation then gives total-null conservation.  The former R1--R2 generator-identification premise is therefore removed, while the native seed, positive-polar inequality, exact physical/Fock species equality, Gauss/zero-mode, color-path and remaining carrier intertwiners stay open.  The remaining frontier is concrete: identify the common-heat matching scale and face heat normalization with any named physical renormalization scale only through the typed metric/quadrature/RG map, transport the occurring G0 quark-source member through the claim-specific G16 RG/threshold map, compute the finite strong commutant/coercivity sequence, prove the chiral thermodynamic/refinement limit, reuse the native finite overlap/index and execute its actual winding weights; retain stronger chiral/topological comparison certificates only where consumed, retain the full \(\LFRep_L\) cross-owner package only for the stronger finite fraction/POVM realization, prove the consumed light-ray/refinement limits, and prove scale-split finite-action universality including generated multi-face interactions only for any claimed representative-independent/canonical action comparison, and discharge the action-preserving bounded-incidence presentation bridge APBI (with APBC as the square-face specialization) before promoting the triality/branch-core area exponent to a positive-tension theorem on the stated pure-glue or certified unscreened sector.  The bounded phase-cluster state and the conditional ordinary local-QFT
reconstruction are supplied by Paper~VII-C; their actual locality and full
covariance tests are retained and are not a G33 action-uniqueness condition.


% Source Paper VII-A, original line 9183; inherited proof bodies unchanged.
\subsection{Native momentum coefficient tests in the finite QCD kernel}
\label{p10:sec:v226-qcd-native-coefficients}
\begin{corollary}[The frozen momentum factor has a complete coordinate basis]
\label{p10:cor:v226-qcd-native-coefficients}
For any actually factorized finite QCD source corner
$A\in\operatorname{End}(\mathcal K)\otimes\mathcal A_D$, its momentum
factor has the exact inverse
$A=\sum_jq_j(A)\otimes D_j$ of
Corollary~\ref{p1b:cor:v226-source-jets}. Every native momentum matrix in
this corner is a degree-six coordinate polynomial; all remaining color,
link, occupation and spatial factors remain in the actual coefficients.
\end{corollary}
\begin{proof}
Apply the tensor inverse to the actual admitted corner, without identifying
its outer carrier with the fifteen-dimensional momentum fibre. It is an
identity on the whole supplied matrix, including every interaction and source
term. Nonreducing truncations retain their actual word defects.
\end{proof}
This propagates the coordinate-derived rotations and protected split to
the existing quark and gauge constructors. It supplies no missing quark
rest dictionary, physical mass gap, chiral order-of-limits prescription,
topological sector weight or Wilson/string-tension limit. The corrected
65-dimensional coordinate-control Lie algebra also prevents inference of
arbitrary Hamiltonian controls from the associative momentum basis alone.

% END FRAMEWORK v2.226: KINEMATIC CLOSURES

% BEGIN FRAMEWORK v2.227: DOWNWARD CLOSURES

% Source Paper VII-A, original line 9211; inherited proof bodies unchanged.
\subsection{Finite gauge-state covariance and complete response comparison}
\label{p10:sec:v227-gauge-covariance}
An actual regulated gauge/rotation symmetry commuting with the finite
Hamiltonian automatically preserves its normalized full-ground or Gibbs
state (Theorem~\ref{p1b:thm:v227-state-symmetry}). On these branches a
finite Ward proof need not assume another invariant vacuum selector.
It still requires the actual source transformation, regulator identities
and any differentiated action/measure word; no chiral measure or
continuum anomaly statement follows from this finite density identity.

For an actual finite kernel/source comparison, the complete response tests
of Theorem~\ref{p1b:thm:v227-finite-response} and its jet corollary close
all energy dependence on the common retained target. The source-visible
minimal environment is then fixed even if a reducing dark extension is
not. Native central-offset retirement applies only after verifying the
actual reducing algebra and even Liouvillian source. Quark/odd source
components crossing its blocks retain their relative energy by
Proposition~\ref{p1b:prop:v227-central-boundary}.

% END FRAMEWORK v2.227: DOWNWARD CLOSURES

% End of predecessor document.
\paragraph{Final theorem status and export boundary.}
Native Hurwitz geometry, the singlet Hamiltonian and finite source operators require actual physical state, scale and limit data before the corresponding experimental claims. The details of the unchanged supporting derivations are preserved in Paper~VII-D.

\paragraph{Technical-reference convention.} References in inherited theorem text to labels formerly described as appendices and now assigned to Paper~VII-D are references to that separate technical companion, not local appendices.



% BEGIN FRAMEWORK v2.338: FULL FINITE PETERSEN OPERATOR AS COMPARISON, NOT QCD KERNEL
\section{Full paired-incidence occupancy operators and exact spectral reduction}
\label{p10:sec:v338-full-petersen-comparison}
Let $K_{\rm inc}=B_uB_u^\dagger$ be the unsigned incidence Gram
from Paper~V-B, acting on the actual 120-dimensional function space
of the given $2I$.  It is not the physical finite color-singlet
Hamiltonian $H_L$, which acts on a different, source-owned state space.

\begin{theorem}[Exact marked-occupancy coupling and soluble full spectrum]
\label{p10:thm:v338-full-occupancy-secular}
For a fixed two-core marking let $D_s$ be multiplication by the
indicator of the 12 elements preserving the pair, and define the
formal geometric family $K(t)=K_{\rm inc}+tD_s$, $t\in\mathbb R$.
Its center-even 60-dimensional characteristic polynomial is
\[
 \begin{split}
 \det(zI-K(t)|_{Z+})={}&z^{13}(z-30)^{10}(z-24)^{20}\\
 &\quad\cdot q_T(z,t)q_S(z,t)q_E(z,t)^4,
 \end{split}
\]
where
\[
 \begin{split}
 q_T(z,t)&=(z-120)(z-30)(z-24)
          -t(z^2-138z+3024),\\
 q_S(z,t)&=z(z-30)-t(z-18),\\
 q_E(z,t)&=z(z-30)(z-24)
          -t(z^2-36z+216).
 \end{split}
\]
On the center-odd 60-dimensional subspace the spectrum is
$\{t^{(6)},0^{(54)}\}$, so
\[
 \det(zI_{120}-K(t))=z^{54}(z-t)^6
            \det(zI-K(t)|_{Z+}).
\]
All 120 modes are included, including the off-sector occupancy
couplings.  At $t=0$ the original incidence spectrum is recovered.
\end{theorem}
\begin{proof}
The pair stabilizer in $A_5$ is $S_3$; the three pair-relation
indicators are constant on its double cosets.  Decompose the even
regular representation into left--right $S_3$ types.  The only
moving blocks for the same-pair indicator reduce to the three
secular polynomials above, with multiplicities $1,1,4$.
The other $13+10+20=43$ even modes remain at $0,30,24$.
The indicator on the odd carrier is a rank-six diagonal projector,
and the incidence kernel vanishes there.  An exact character
calculation and independent 60-by-60 and 120-by-120 numerical
checks certify the factorization; no value of $t$ is selected by a
physical action.
\end{proof}

\begin{proposition}[Diagonal occupancy cannot isolate the new 25-sector]
\label{p10:prop:v338-diagonal-occupancy-leakage}
Let $\Pi_{25}$ be the group-character projector of Paper~V-B.
For any diagonal multiplication $M_f$ on the complete group carrier,
\[
 [\Pi_{25},M_f]=0\quad\Longleftrightarrow\quad
 f\text{ is constant on }2I.
\]
In particular the nonuniform $12,36,72$ occupancy indicators
necessarily mix the 25-sector with the other 17- and 18-dimensional
center-even sectors.  The squared Hilbert--Schmidt block couplings
are
\[
 \begin{array}{c|ccc}
 &17\leftrightarrow25&25\leftrightarrow18&17\leftrightarrow18\\\hline
 D_{12}&13/20&3/5&12/25\\
 D_{36}&5/4&3/2&6/5\\
 D_{72}&2&3/2&6/5
 \end{array}.
\]
None is a physical energy coefficient.
\end{proposition}
\begin{proof}
The character kernel of $\Pi_{25}$ is nonzero whenever the
relative $A_5$ element is a three-cycle (and also for the two
central lifts).  Three-cycles generate $A_5$, so commuting with a
diagonal multiplier forces $f$ constant, including across both
central lifts.  The block norms follow by contracting the exact
central character idempotents against the three stratum indicators;
all displayed values are rational finite sums.  Their integer
certificate is part of this release.
\end{proof}

\begin{remark}[Source-owned Hamiltonian remains separate]
The scalar $t$ in the exploratory operator $K(t)$ is an
unselected mathematical parameter.  The native $H_L$ requires
its actual occupation, color, charge, kinetic norm and source
matrix entries.  The corresponding physical eigengaps and residues
cannot be inferred from $K(t)$ or from its $25$-sector compression.
\end{remark}
% END FRAMEWORK v2.338



% BEGIN S3-CROSS-ACTION-OWNERSHIP v1_178 (18 November 2026)
\subsection{S$^3$ curvature response and the physical-action certificate}
\label{p10:subsec:s3-response-owner-v1178}
\begin{definition}[Physical finite-link action-response certificate]
\label{p10:def:s3-physical-response-certificate-v1178}
For a specified gauge-invariant physical action $S_{\rm phys}$ on the
actual 96-link Hurwitz carrier, or a proved pullback of the generated
position action to this carrier, a \emph{same-owner response certificate}
consists of the complete link configuration domain, spin/color and
Gauss-source maps, the moving spectral/localizer and cutoff data, the
full first and second action jets (including all contact and measure
terms), and the observable/source normalization.
If $P_{\rm blind}$ denotes the rank-eight comparison curvature projector
in Paper~VII-D, the proposed test observable is the restricted bilinear
response
\begin{equation}\label{p10:eq:s3-complete-response-v1178}
\boxed{\mathcal R_{\rm phys}[u,v]
 =D^2 S_{\rm phys}[u,v],\qquad
 u,v\in\operatorname{ran}P_{\rm blind},}
\end{equation}
only \emph{after} a physical meaning for those tangents has been
supplied by the same action owner.  A weaker, directly observable
response need not establish a global position--momentum action isomorphism.
\end{definition}

\begin{remark}[No automatic promotion of comparison curvature zeros]
\label{p10:rem:s3-no-unowned-physical-modes-v1178}
On the massless momentum-shell spin comparison, eight transverse
curvature modes have vanishing quadratic scalar spectral response,
while a declared positive triangular-face comparison has a nonzero
response. This difference is not a contradiction of the conditional
Hurwitz QCD transfer theorem
\ref{p10:thm:hurwitz-shell-qcd-transfer-compiler}: the two tested actions
have not been identified with the same selected physical $H_L$.
Neither a new correction field, a hadron mass, nor a vanishing physical
decay amplitude is inferred.  Spectral poles still require the full
same-owner source resolvent
$\iota^\dagger(z-H_L)^{-1}\iota$, including its complementary return.
\end{remark}
% END S3-CROSS-ACTION-OWNERSHIP v1_178

\printbibliography

% BEGIN FRAMEWORK v2.344: NATIVE C3 EDGE-FAN SOURCE ELIGIBILITY
\section{Character-resolved Hurwitz three-face baryon sources}
\label{p10:sec:v344-threeface-c3-mode}
Retain the gauge-invariant polarized three-face edge determinant
$W_e$ of Proposition~\ref{p10:prop:v2220-edge-fan-baryon}.
By Proposition~\ref{p1:prop:v344-hurwitz-edge-eight-orbits},
there are eight free orbits of oriented Hurwitz edges under left
multiplication by $H=2T$.  Choose representatives $e_a$ and
one nontrivial character $\chi:H\to C_3\subset U(1)$.

\begin{proposition}[Existing determinant admits nonzero $C_3$ modes]
\label{p10:prop:v344-threeface-character-witness}
The character modes
\[
 \boxed{W_{\chi,a}(U)=\frac1{\sqrt{24}}
   \sum_{h\in H}\overline{\chi(h)}\,W_{he_a}(U),
   \qquad a=1,\ldots,8}
\]
are gauge-invariant complex-valued observables built solely from
the pre-existing three-face determinant and fixed shell relabeling.
They are \emph{not identically zero} as algebraic functions of
$SU(3)$ link assignments.  Specifically, put one undirected edge
link equal to $\operatorname{diag}(e^{i\theta},e^{-i\theta},1)$,
its reverse equal to its inverse, and all other links equal to $I$.
For this finite test assignment,
\[
 \boxed{\sum_{a=1}^8|W_{\chi,a}(U)|^2
              =\frac{(1-\cos\theta)^2}{9}.}
\]
In particular it equals $4/9$ at $\theta=\pi$ and is nonzero
for arbitrarily small nonzero $\theta$.
\end{proposition}
\begin{proof}
Each $W_e$ is gauge invariant by the existing polarized determinant
construction, and the Fourier combination preserves that property.
For the one-link assignment, $W_e=1$ except for twelve oriented
neighbouring fan entries, where the color contraction gives
$W_e=(1+2\cos\theta)/3$.  Summing the finite $C_3$ phases over
the eight regular orbits produces the displayed norm identity;
the integer/finitary calculation is recorded in
Proposition~\ref{p1c:prop:v344-hurwitz-petersen-verification}.
\end{proof}

\begin{remark}[Physical admissibility and source/energy distinction]
\label{p10:rem:v344-threeface-physical-boundary}
The one-link witness is an algebraic configuration; membership in
the fully constituted admissible link/action domain is an
independent check.  A nonzero $C_3$ Fourier mode is neither a
physical strange-flavor coupling nor a nonzero vacuum expectation,
spectral residue or hadron mass shift.  An $H$-invariant ensemble
may force its one-point mean to vanish while permitting a
nonzero two-point response.  The mathematical isometry to the
Petersen quintet in
Proposition~\ref{p1:prop:v344-petersen-quintet-intertwiner}
does not insert a new $2I$ term into the physical $H_L$.
\end{remark}
% END FRAMEWORK v2.344


% BEGIN FRAMEWORK v2.345: F40 SOURCE PORT AND LEAKAGE TEST
\clearpage
\section{F40 fivefold-arrow source port for the finite color-singlet action}
\label{p10:sec:v345-f40-source-port}
Let $Z=\ker\partial_{40}\subset\mathbb R^{60}$ and let
$T_R,T_L,A$ be the three intrinsic geometric readouts in
Paper~I-A.  The formal inclusion
$Z\to\mathbb R^{60}\to\ell^2(2I)$ is a comparison, not a map
into the physical $SU(3)$-singlet Fock carrier.

\begin{proposition}[Same-cell fivefold exclusion and conditional source test]
\label{p10:prop:v345-f40-qcd-source-obligation}
In the regular group comparison, if $P_H$ restricts to one
$24$-element Hurwitz cell $H=2T$, then for every fivefold shell
$g\in\mathcal S$,
\[
 \boxed{P_HL_gP_H=0.}
\]
Thus a fivefold permutation cannot be installed as a nonzero
same-cell operator merely by relabeling that one $24$-state
vertex space.  A proposed physical realization must supply a
normalized, color-gauge-compatible isometry
$\iota_L:Z_{\mathbb C}\to\mathcal H_{L,\mathrm{sing}}$, an
actual source $O_L$, and all link/action-domain conditions.
For the already constituted Hamiltonian $H_L$ define
\[
 B_L=(I-\iota_L\iota_L^\dagger)H_L\iota_L,
 \qquad K_L=\iota_L^\dagger H_L\iota_L.
\]
Only $B_L=0$ makes the proposed geometric subspace reducing.
If $B_L\ne0$, physical poles must be extracted from the full
$H_L$ resolvent, not from the spectrum of $K_L$ alone.
\end{proposition}
\begin{proof}
The first equality follows from the disjoint group supports
$H$ and $gH$.  The remaining assertions follow from the
orthogonal decomposition relative to the range of $\iota_L$:
$B_L$ is its off-diagonal Hamiltonian block.  Nonzero leakage
requires the full Schur complement and source-visible residues.
\end{proof}

\begin{corollary}[Coefficient-independent Hodge null test]
\label{p10:cor:v345-f40-symmetric-source-null}
For an $A_5$-equivariant candidate source which factors
\emph{only} through the symmetric quotient readout $T_R+T_L$,
the matrix elements on the two F40 Hodge mirror components
$\mathbf3\oplus\mathbf3'$ vanish identically.  The
antisymmetric readout $T_R-T_L$ is geometrically nonzero on
both components; it is not automatically a term of $H_L$.
\end{corollary}
\begin{proof}
Theorem~\ref{p1:thm:v345-f40-triple-frame} and
Proposition~\ref{p1:prop:v345-f40-balanced-mirror} give the
rank-five symmetric Gram and its missing $\mathbf3\oplus
\mathbf3'$ components.  Multiplication by an overall source
coefficient cannot change the kernel.
\end{proof}
% END FRAMEWORK v2.345

% BEGIN FRAMEWORK v2.349: F40 SOURCE COMPATIBILITY AND SCALAR CONSTRAINT
\clearpage
\section{F40 source consistency and physical scope}
\label{p10:sec:v349-f40-source-filter}
\begin{proposition}[Exact geometric eligibility does not construct the source]
\label{p10:prop:v349-f40-source-eligibility}
If an action-selected color-singlet source is claimed to be a
linear function of \emph{one} underlying F40 cycle through all
three native channels $(T_R,T_L,A)$, its candidate
44-component geometric coefficient packet must satisfy the
23 independent equations of
Theorem~\ref{p1c:thm:v349-f40-23-equations}; equivalently
$P_{44}d=d$.  No diagonal $A_5$-equivariant map from the
complete 21-cycle module to the twisted-25 endomorphism
module alone can be injective; its rank is at most sixteen.
The already existing five-core Petersen quintet is the
minimal missing representation type.  These statements
constrain a \emph{proposed} common-cycle source, not every
possible physical field or source on a different space.
\end{proposition}
\begin{proof}
The native frame compatibility test is necessary and
sufficient exactly for the image of one F40 cycle.  The
multiplicity obstruction is Theorem~\ref{p1:thm:v349-f40-minimal-quintet}.
Both are independent of the unknown physical action coefficient.
\end{proof}
\begin{corollary}[Invariant first-order cycle-only scalar no-go]
\label{p10:cor:v349-f40-linear-scalar-nogo}
An unbroken $A_5$-invariant linear scalar functional on
$\mathcal Z_{40}\cong\mathbf3\oplus\mathbf3'\oplus3\mathbf5$
vanishes identically; the space of invariant symmetric
quadratic forms has dimension eight.  Thus a putative
nonzero \emph{linear, cycle-only} invariant scalar response
cannot be the complete source.  Bilinear matter--cycle
sources, nonlinear sources, and physically broken-symmetry
matrix elements are not excluded.
\end{corollary}
\begin{proof}
No trivial irreducible occurs in $\mathcal Z_{40}$, so
$\operatorname{Hom}_{A_5}(\mathcal Z_{40},\mathbf1)=0$.
The two inequivalent real triplets each contribute one
quadratic invariant and the three quintet copies contribute
six symmetric multiplicity contractions.
\end{proof}
\begin{remark}[Same-action obstruction remains]
\label{p10:rem:v349-f40-source-action-open}
Neither the 23 equations, the harmonic degree-four lift,
nor the golden-ratio Hodge transport constructs the actual
$SU(3)$-singlet isometry, a gauge-compatible link coupling,
the Hamiltonian leakage block, or any correlator pole.
The physical action and source-visible residues remain the
indispensable observable owners.  GP3 is constitutional and is
not being supplemented with a new multiplicity postulate.
\end{remark}
% END FRAMEWORK v2.349

% BEGIN FRAMEWORK v2.350: QCD GAUSS, SOURCE, AND HURWITZ THREE-FACE REPAIR
\clearpage
\section{Repaired fixed-shell QCD compiler and physical gauge/source interfaces}
\label{p10:sec:v350-qcd-exact-repairs}
The exact degree-two scalar Hurwitz spectral intertwiner is
Theorem~\ref{p1:thm:v350-hurwitz-scalar-laplacian}. Its domain is a
$14$-dimensional \emph{scalar} comparison band, distinct from the
$15$-dimensional physical momentum carrier.  The finite spin stencil
of Theorem~\ref{p10:thm:hurwitz-shell-dirac-compression} retains the
fixed geodesic separation $\ell_H=\pi R_{\rm mom}/3$; its Taylor
bound is not a continuum refinement assertion.

\begin{proposition}[Spinorial exactness has a finite residual certificate]
\label{p10:prop:v350-spinor-intertwining-certificate}
Let $\mathcal B_D$ be a specified finite smooth-spinor comparison
subspace, with its actual spin transport and sampling map $S_D$.
Define
\[
 \mathcal E_D:=D_{H,0}S_D-S_D D_{S^3}|_{\mathcal B_D}.
\]
Exact low-band reproduction by this particular Hurwitz stencil holds
if and only if $\mathcal E_D=0$; a nonzero matrix norm
$\|\mathcal E_D\|$ is a finite quantitative obstruction. The
scalar comparison Theorem~\ref{p1:thm:v350-hurwitz-scalar-laplacian}
does not itself discharge this test.
\end{proposition}
\begin{proof}
All maps have the same finite domain and sampled target after a
choice of the genuine spinor basis and connections. Equality of
operators is equivalent to the vanishing of their difference; the
operator norm vanishes exactly when that matrix is zero. No
asymptotic small-spacing limit is involved.
\end{proof}

\begin{theorem}[Exact vertexwise Haar Gauss projector]
\label{p10:thm:v350-local-gauss-projector}
Suppose the actual finite link/matter carrier carries a unitary
representation $\mathcal U$ of the \emph{vertexwise} gauge group
$\mathcal G_H=SU(3)^{V_H}$, with each finite representation truncation
invariant under this action. Then
\begin{equation}
\label{p10:eq:v350-local-gauss-projection}
 P_{\rm Gauss}=\int_{\mathcal G_H}\mathcal U(\{g_v\})
                  \prod_{v\in V_H}d\mu_{\rm Haar}(g_v)
\end{equation}
is the orthogonal projector onto all local Gauss-invariant states.
If the physical finite Hamiltonian obeys
$\mathcal U(g)H_L\mathcal U(g)^*=H_L$ for every vertexwise $g$,
then $[P_{\rm Gauss},H_L]=0$. In the same representation
$\operatorname{Ran}P_{\rm Gauss}$ is contained in the global
$SU(3)$ singlet subspace. Hence the diagonal global singlet
projector alone does not enforce every local Gauss constraint.
\end{theorem}
\begin{proof}
Normalized Haar averaging of any unitary compact-group action
is a self-adjoint idempotent with range the fixed vectors.
The Hamiltonian commutes with the integrand and therefore with
the integral. The diagonal global color group is a subgroup of
$\mathcal G_H$, so every locally fixed vector is globally fixed.
For example a bilocal $\bar q(x)q(y)$ without its Wilson connector
is invariant under diagonal color but generally not under
independent transformations at $x$ and $y$.
\end{proof}

\begin{theorem}[Exact second derivatives of the complete finite source]
\label{p10:thm:v350-dynamical-source-hessian}
Let $H(J)$ be a finite self-adjoint Hamiltonian, twice
differentiable in real sources near zero, and
$W(J)=\log\operatorname{Tr}e^{-\beta H(J)}$ with $\beta>0$.
Write $H_a=\partial_aH$, $H_{ab}=\partial_a\partial_bH$,
$\rho=e^{-\beta H}/\operatorname{Tr}e^{-\beta H}$ and
$\widetilde H_a=H_a-\operatorname{Tr}(\rho H_a)I$. Then
\begin{equation}
\label{p10:eq:v350-source-hessian}
 \partial_aW=-\beta\operatorname{Tr}(\rho H_a),\qquad
 \partial_a\partial_bW=-\beta\operatorname{Tr}(\rho H_{ab})
   +\beta^2\int_0^1\operatorname{Tr}
       (\rho^{1-s}\widetilde H_a\rho^s\widetilde H_b)\,ds.
\end{equation}
Thus nonlinear source contacts and noncommuting insertions
are retained exactly. The finite readout universality assertion
covers observables in the actually supplied Hamiltonian/source
operator dictionary; its extension to scattering or physical
PDFs requires the separate Paper~VII-B/C certificates.
\end{theorem}
\begin{proof}
The Duhamel derivative is
$\partial_a e^{-\beta H}=-\beta\int_0^1
 e^{-\beta(1-s)H}H_a e^{-\beta sH}\,ds$.
Cyclicity of the trace yields the first identity. Differentiate
again, including $\partial_bH_a=H_{ab}$, and subtract the
product of first moments to differentiate the logarithm.
This gives exactly the centered Duhamel correlation in
Eq.~\eqref{p10:eq:v350-source-hessian}.
\end{proof}

\begin{proposition}[Three distinct transports require a polarized determinant]
\label{p10:prop:v350-three-face-polarization}
For three fundamental color transports $A,B,C\in SU(3)$ define
\[
 \mathfrak E(A,B,C)=\frac16\varepsilon_{abc}\varepsilon_{a'b'c'}
 A_{aa'}B_{bb'}C_{cc'}.
\]
This is gauge invariant when the three transports share the stated
common source/target transformation law, and
$|\mathfrak E(A,B,C)|\le1$. In general
$\mathfrak E(A,B,C)\ne\det A$; specifically
\[
 \mathfrak E(U,U,U)=\det U=1,\qquad
 \mathfrak E(A,I,I)=\tfrac13\operatorname{tr}A.
\]
The three-face Hurwitz incidence therefore supplies a nonconstant
baryonic transport invariant rather than a universal unit determinant.
\end{proposition}
\begin{proof}
The contraction is a matrix coefficient of
$A\otimes B\otimes C$ between normalized antisymmetric
three-color singlets, hence has modulus at most one.
Simultaneous endpoint gauge transformations fix the two
antisymmetric tensors because every endpoint color matrix
has determinant one. Equal inputs give the determinant
identity; contraction with two identities gives
$\operatorname{tr}A/3$. Taking
$A=\operatorname{diag}(-1,-1,1)\in SU(3)$ yields $-1/3$,
which is not $\det A=1$.
\end{proof}

\begin{remark}[Physical mass/source and dynamical-gap boundaries]
\label{p10:rem:v350-qcd-physical-firewalls}
The Hamiltonian family remains the whole-family image of the upstream
charged mass/source family; G34 is not reinstated as a new quark-mass
selector. The explicitly realized local gauge representation,
its invariant Hamiltonian, interacting source matrix elements,
and dynamical spectral residues are separate certificates. An excluded
native carrier interval is not a uniform interacting pure-glue gap.
The three-face flux mean and variance are Haar reference statistics,
not baryon mass or spin-splitting predictions; see
Proposition~\ref{p12:prop:v350-flux-reference-witness}.
\end{remark}
% END FRAMEWORK v2.350

% BEGIN FRAMEWORK v2.356: exact finite heat-star moments
\section{Exact three-face heat-star correlations on a declared representative}
\label{p10:sec:v356-heat-star}

\begin{theorem}[Current-free normalized three-face character cumulant]
\label{p10:thm:v356-heat-star-cumulant}
Use normalized Haar measure on $SU(3)$ and the central heat kernel
$K_t$, with character coefficients
$d_Re^{-tC_R/2}$.  Let $U$ be Haar distributed and let
$X_1,X_2,X_3$ be independent heat-distributed elements, each with
density $K_t(X_i)dX_i$.  For the common boundary frame define
$F_i=\chi_{\mathbf3}(X_iU)/3$.  Then
\[
 \mathbb E[F_i]=0,\quad
 \mathbb E[F_iF_j]=0\ (i\ne j),\qquad
 \boxed{\mathbb E[F_1F_2F_3]_c=\frac{e^{-2t}}{27}>0.}
\]
The normalized epsilon-contracted three-fundamental singlet
channel has response $e^{-2t}$.  These are exact boundary-framed
multi-face response observables; an individual $F_i$ is not a
closed gauge-invariant Wilson loop.
\end{theorem}
\begin{proof}
Central heat convolution yields
$\mathbb E_{X_i}[\chi_{\mathbf3}(X_iU)]
=e^{-tC_{\mathbf3}/2}\chi_{\mathbf3}(U)$,
where $C_{\mathbf3}=4/3$.  The Haar integral of one or two
fundamental characters vanishes by center triality, while
$\int\chi_{\mathbf3}(U)^3dU=\dim\operatorname{Inv}
(\mathbf3^{\otimes3})=1$ by the epsilon tensor.  Independence gives
$e^{-3t(4/3)/2}/3^3=e^{-2t}/27$; the vanishing lower moments
make this the connected third cumulant.  The invariant epsilon
contraction removes the normalization $3^{-3}$.
\end{proof}

\begin{proposition}[Strictly positive cubic heat-block diagnostic]
\label{p10:prop:v356-heat-star-cubic}
For every $t>0$ and the same heat convention,
\[
 \boxed{\mathcal C_3(t):=\int_{SU(3)}(K_t(U)-1)^3dU
             \ge54e^{-2t}>0.}
\]
\end{proposition}
\begin{proof}
Expand $K_t-1$ in nontrivial irreducible characters.
Every coefficient $d_Re^{-tC_R/2}$ is positive, and the integral
of three characters is the nonnegative integer
$\dim\operatorname{Inv}(R_1\otimes R_2\otimes R_3)$.
The triples $(\mathbf3,\mathbf3,\mathbf3)$ and
$(\overline{\mathbf3},\overline{\mathbf3},
\overline{\mathbf3})$ each contribute
$3^3e^{-2t}$, yielding the bound.
\end{proof}

\begin{remark}[Heat-time transport is not a theorem of shell incidence]
\label{p10:rem:v356-apbi-frontier}
On the isotropic \emph{reciprocal-position} face presentation
$w_f=A_f^2$ the existing common-heat result fixes
$t_{\rm iso}=5.753295846633\ldots$.  On that explicitly stated
representative the cumulant is $3.72723656463536\times10^{-7}$
and $\mathcal C_3(t_{\rm iso})\simeq5.9022229671513\times10^{-4}$.
The native Hurwitz $24$-cell is a \emph{momentum-shell} complex;
assigning this same face time to every shell triangle requires an
additional action-preserving reciprocal-position/momentum-shell
presentation map preserving the nonlinear face action, its blocking
and source jets (APBI).  Combinatorial incidence alone does not
provide that map, a physically selected gauge action, running
coupling, sea-quark determinant, or a temporal electromagnetic
contact.  The displayed numerics are therefore conditional heat
representative calculations, not measured QCD predictions.
\end{remark}
% END FRAMEWORK v2.356


% BEGIN FRAMEWORK v2.357: exact spin-flux Wilson-character export
\clearpage
\section[Fundamental--adjoint character form of the baryon spin--flux response]{Fundamental--adjoint character form\\of the baryon spin--flux response}
\label{p10:sec:v357-spinflux-wilson}
\begin{proposition}[Exact three-face fundamental--adjoint Wilson contrast]
\label{p10:prop:v357-spinflux-wilson-contrast}
In the fixed-link comparison domain of
Theorem~\ref{p10:thm:v2220-conditional-spinflux}, let
\[
 w_3(U)=\frac13\operatorname{Re}\operatorname{tr}_{\mathbf3}U,
 \qquad w_8(U)=\frac18\operatorname{tr}_{\mathbf8}U.
\]
For each $U\in SU(3)$,
\[
 \left|\frac{\operatorname{tr}_{\mathbf3}U}{3}\right|^2
  -\operatorname{Re}\frac{\operatorname{tr}_{\mathbf3}U}{3}
 =\frac{1+8w_8(U)}9-w_3(U).
\]
Consequently the native edge-fan coefficient appearing in the
comparison spin-flux expansion is exactly
\[
 \boxed{D_e(U)=\sum_{f\supset e}
 \left[\frac{1+8w_8(U_f)}9-w_3(U_f)\right].}
\]
The equality holds pointwise and hence after any common, well-defined
positive gauge expectation, without a new free coefficient.
\end{proposition}
\begin{proof}
The $SU(3)$ tensor identity
$\mathbf3\otimes\overline{\mathbf3}=\mathbf1\oplus\mathbf8$
implies
$|\operatorname{tr}_{\mathbf3}U|^2=1+\operatorname{tr}_{\mathbf8}U$.
Substitution of the normalized characters proves the first identity;
linearity proves its fixed-fan sum and any integrable expectation.
\end{proof}
\begin{remark}[Physical and sign scope]
\label{p10:rem:v357-spinflux-wilson-scope}
The formula reorganizes the coefficient of the already declared
three-face \emph{comparison} transfer. It does not derive a nucleon
mass splitting, identify the comparison stencil with the interacting
QCD Hamiltonian, fix the sign of the gauge expectation, or supply a
source-preserving reciprocal-position/Hurwitz-shell presentation
bridge (APBI). An isolated heat measure is a conditional example,
not a physically selected gauge ensemble.
\end{remark}
% END FRAMEWORK v2.357


% BEGIN FRAMEWORK v2.365: WILSON COLOR-GRAM AND SAME-ACTION FRONTIER
\section{Two-root Wilson color Grams and the baryonic source frontier}
\label{p10:sec:v365-wilson-gram-apbi}
Keep the existing gauge-invariant three-face fan determinant
$W_e$ of Proposition~\ref{p10:prop:v2220-edge-fan-baryon};
its definition is not replaced or rescaled.  Let
$|\epsilon\rangle=6^{-1/2}\sum_{abc}\epsilon_{abc}|abc\rangle$
be the normalized antisymmetric color tensor.  On the selected
edge fan let $\Psi_q=\epsilon$ and
$\Psi_p=(A_1\otimes A_2\otimes A_3)\epsilon$ with the
transporters acting on the corresponding color indices.

\begin{proposition}[Exact normalized color-root overlap]
\label{p10:prop:v365-baryon-color-gram}
For all admissible unitary $SU(3)$ link assignments,
\[
\boxed{W_e=\langle\Psi_q,\Psi_p\rangle,\qquad
|W_e|\le1,\qquad
\det\begin{pmatrix}1&W_e\\\overline{W_e}&1\end{pmatrix}
=1-|W_e|^2.}
\]
The transported quark--antiquark maximally entangled color tensors
similarly give $M_{ij}=\operatorname{tr}(A_iA_j^\dagger)/3$
as a normalized two-root overlap.  These are color-space identities;
$W_e$ contains no fermion insertions, has baryon number zero and is
not itself the physical three-quark interpolator $O_B$.
\end{proposition}
\begin{proof}
The defining epsilon contraction of $W_e$ is the inner product
of the displayed vectors.  Each tensor factor is unitary, hence
both vectors have unit norm and Cauchy--Schwarz proves the bound.
The determinant is the Gram determinant of two normalized vectors.
The meson statement follows from the normalized identity-color
vector $3^{-1/2}\sum_a|a\overline a\rangle$ and the trace pairing.
\end{proof}

\begin{proposition}[Source whitening removes a common raw Gram]
\label{p10:prop:v365-color-whitening}
If the same positive-definite two-root source Gram $G$ is used on
both legs of a color-singlet comparison, its canonically whitened
overlap is $G^{-1/2}GG^{-1/2}=I$.  When $G$ is singular, the
identity holds only on $\operatorname{supp}G$ with its Moore--Penrose
inverse.  Thus the bare $|W_e|^2$ is not a derived weak branching
fraction.  Distinct physical pole states require their actual
same-action cross-Gram and transition current.
\end{proposition}
\begin{proof}
Apply the spectral theorem to $G$, dividing each strictly positive
eigenvalue by its square root on each side.  At zero eigenvalues
the Moore--Penrose definition gives the support projector.
\end{proof}

\begin{definition}[APBI action-preserving hadronic presentation test]
\label{p10:def:v365-apbi-hadron-presentation}
A candidate isometry $\iota:\mathcal H_{\rm fan}\to\mathcal H_L$
into the \emph{same actual} finite QCD color-singlet carrier is
admissible for physical use only after gauge, charge, rotation,
center, fermion-number and source/field placement are matched.
Writing $P=\iota\iota^\dagger$, $Q=I-P$ and
\[
B=QH_L\iota,\qquad
\Sigma(z)=B^\dagger(z-QH_LQ)^{-1}B,
\]
requires the full Feshbach/Schur returned resolvent
\[
\iota^\dagger(z-H_L)^{-1}\iota
=\bigl(z-\iota^\dagger H_L\iota-\Sigma(z)\bigr)^{-1},
\]
whenever both sides are defined.  The $\Sigma$ term may be dropped
only if $B=0$ or a specific approximation error is proven.
Existence and physical actualization of such a source-preserving
$\iota$, and nonzero overlap with the fermionic $O_B$ pole source,
remain unestablished by the finite Hurwitz graph identities.
\end{definition}

\begin{remark}[Reference jet and Cartan candidates remain scoped]
\label{p10:rem:v365-wilson-jets-scope}
The proposal's homogeneous Cartan value $W_e(1/2)=8/9$, the
vanishing $C_3$ Fourier contrast on that one-parameter family,
and the proposed independent-link first-/second-jet ranks $3,8$
are restricted-reference tests.  The current integration does not
promote a formula for that specific Cartan family without its
full source chart and replay.  In particular these candidate ranks
are not physical hadron multiplicities; the source-owned
nonvanishing test of
Proposition~\ref{p10:prop:v344-threeface-character-witness}
is preserved independently.  No new $SU(3)_c$ gauge connection,
Hamiltonian eigenvalue, mass, width or $8/9$ rate multiplier is added.
\end{remark}
% END FRAMEWORK v2.365


% BEGIN FRAMEWORK v2.366: HURWITZ CHAMBER/PHYSICAL SOURCE INTERFACE
\section{Character-resolved chamber sources and action-owned hadronic tests}
\label{p10:sec:v366-hurwitz-physical-source-interface}
The preceding full QCD Hamiltonian $H_L$, its genuine fermionic
interpolators, Wilson determinant $W_e$, and the D22 heat-action
Hessian are retained.  The following finite character identities
\emph{do not} add an interaction or create a new baryon species.

\begin{proposition}[Hurwitz triality diagonalizes the existing D22 label Hessian]
\label{p10:prop:v366-hurwitz-character-hessian}
For the three normal $Q_8$ octets of one marked $2T$, let
$\omega=e^{2\pi i/3}$ and $P_{\chi},P_{\bar\chi}$ project onto
the nonconstant characters of the cyclic permutation of the octets.
The three-face relative-incidence label matrix has the identity
\[
 \boxed{3I_3-J_3=3(P_{\chi}+P_{\bar\chi})}.
\]
Thus the \emph{already defined} D22 heat-action Hessian's
relative-label factor $k(t)(3I_3-J_3)$ has zero uniform eigenvalue
and two character eigenvalues $3k(t)$.  Independently, for the
already defined polarized three-face determinant
$W_e=D(e^{i\varepsilon X_1},e^{i\varepsilon X_2},
 e^{i\varepsilon X_3})$ with traceless Hermitian $X_i$,
\[
 \boxed{W_e=1-\frac{\varepsilon^2}{12}
  \sum_{i,j}(3\delta_{ij}-1)\operatorname{tr}(X_iX_j)
  +O(\varepsilon^3)}.
\]
Equality of the \emph{label quadratic form} does not identify the
Wilson Taylor coefficient with the complete heat-action strength.
\end{proposition}
\begin{proof}
The discrete Fourier projectors are
$P_1=J_3/3$ and $P_{\chi}+P_{\bar\chi}=I_3-P_1$.
The determinant's linear term vanishes by tracelessness, and its
second derivative is the displayed relative-face quadratic form,
as evaluated by the exact symbolic Taylor check in the supplied
\texttt{v1\_11} finite certificate.  D22's original $k(t)$ is
unchanged by this Fourier basis change.
\end{proof}

\begin{proposition}[Eight Wilson edge modes admit a four-plus-four conjugation grading]
\label{p10:prop:v366-hurwitz-wilson-conjugation-grading}
On the previously established eight inverse-paired nearest-edge
orbits and their $C_3$-character Wilson functions $W_{\chi,r}$,
let $\mathsf C_0$ be the \emph{gauge-link} conjugation/inverse-edge
operation.  The existing three-face determinant obeys
\[
 \mathsf C_0 W_{\chi,r}=\chi(r)W_{\chi,r^{-1}},\qquad
 W^{(\pm)}_{\chi,r}=
 \frac{W_{\chi,r}\pm\chi(r)W_{\chi,r^{-1}}}{\sqrt2},
\]
with four $+$ and four $-$ eigenfunctions.
The explicit eight-function one-link evaluation matrix has
nonzero determinant $2^{20}3^2$ in the normalization of the
verifier, proving the eight functions are \emph{linearly
independent}.  For a strictly positive full-support pure-gauge
link measure their centered covariance is positive-definite;
if that measure is $\mathsf C_0$-invariant the covariance has
independent $4+4$ blocks.  This mathematical gauge-function
gradation does not assert the full interacting Hamiltonian or
physical vacuum possesses charge-conjugation symmetry.
\end{proposition}
\begin{proof}
Inverting the oriented edge and complex-conjugating all color
links sends each twisted orbit to its inverse with character phase
$\chi(r)$; each inverse pair diagonalizes into the displayed two
vectors.  Eight explicit link evaluations have determinant
$2^{20}3^2\ne0$, establishing linear independence.  A centered
linear combination with zero variance under a strictly positive
full-support measure would be constant on the whole link domain;
it vanishes at the all-identity configuration and linear
independence forces all coefficients to vanish.  Covariance
between odd and even graded functions vanishes under an invariant
measure.  No polynomial algebraic independence is claimed.
\end{proof}

\begin{proposition}[Hurwitz spin-two character tensor and factorized branch norms]
\label{p10:prop:v366-hurwitz-spin2-character-source}
The nontrivial local $C_3$ character has a normalized tensor line
in the spin-two symmetric-traceless three-space,
\[
 Q_{\chi}=3^{-1/2}\operatorname{diag}(1,\omega,\omega^2),\qquad
 \operatorname{tr}Q_{\chi}=
 \operatorname{tr}Q_{\chi}^{2}=0,\qquad
 \operatorname{tr}Q_{\chi}^{3}=1/\sqrt3.
\]
It has $\langle\mathbf J\rangle=0$ and isotropic quadratic spin
moments.  On the \emph{abstract factorized} carrier
$Q_{\chi}\otimes\psi$, where $\psi$ is any normalized spinor,
\[
 \boxed{\|P_{4'}(Q_{\chi}\otimes\psi)\|^2=\tfrac25,
 \qquad\|P_6(Q_{\chi}\otimes\psi)\|^2=\tfrac35}.
\]
These are exact normalized source-projection weights for the
induced $\mathbf5\otimes\mathbf2=\mathbf4'\oplus\mathbf6$
representation, not physical hadron decay fractions or an assertion
that the geometric rotation action equals the physical $SU(2)$ one.
\end{proposition}
\begin{proof}
The invariant $Q_{\chi}$ line has vanishing dipole by
its tetrahedral character covariance.  The projectors
$P_{4'}=(2I-X)/5$ and $P_6=(X+3I)/5$ of
Proposition~\ref{p1:prop:v366-twisted-quintet-spinorial-branch}
then have expectations $2/5$ and $3/5$ because
$\langle X\rangle=2\langle\mathbf J\rangle\cdot
\langle\mathbf S\rangle=0$.  The exact tensor/cubic and spin
matrix identities are checked separately in the $v1\_10$
certificate.
\end{proof}

\begin{definition}[Required same-action Hurwitz--QCD mixed-source test]
\label{p10:def:v366-hurwitz-stress-baryon-interface}
An optional chamber, rank-15 Petersen or induced-quintet source is
\emph{physically eligible} only after an explicit isometry
$\iota_L$ into the \emph{actual} color-singlet $\mathcal H_L$
intertwines gauge transformations, physical rotations, center,
baryon number, charge and flavor.  Under a genuinely
$\mathsf C$-invariant actual action/state define the connected
mixed correlator for each Wilson eigenmode and each native gluonic
stress spin-two component $O_{G,m}^{2^{++}}$ by
\[
 \Gamma^{(\pm)}_{r,m}(t)=
 \langle\Omega|W^{(\pm)}_{\chi,r}(t)
  (O_{G,m}^{2^{++}})^{\dagger}(0)|\Omega\rangle_c.
\]
If the actual $O_G$ is $\mathsf C$-even and $\mathsf C$ acts as the
verified gauge-link conjugation, then $\Gamma^{(-)}_{r,m}(t)=0$;
no symmetry guarantees $\Gamma^{(+)}_{r,m}(t)\ne0$.
The gauge-only $W_e$ and $W_{\chi,r}$ have baryon number zero;
a baryon-number-one candidate must also use the established
genuine three-quark operator $O_B$ on a common gauge-compatible
carrier (for example a well-defined projected product
$P_{\rm sing}W_{\chi,r}O_BP_{\rm sing}$).
Finally compute the actual leakage
$(I-\iota_L\iota_L^\dagger)H_L\iota_L$ and retain the Schur
self-energy of Definition~\ref{p10:def:v365-apbi-hadron-presentation}.
Neither a nonzero $\Gamma^{(+)}$, an $O_B$ pole nor a
new decay transition is asserted here.
\end{definition}
% END FRAMEWORK v2.366


% BEGIN PAPER VII-A v1.173: TWO/THREE CORE SOURCE PORTS, NATIVE QCD UNCHANGED
\section{Local $2O$ sources and mobile $2I$ multi-core comparison ports}
\label{p10:sec:v1173-octahedral-icosahedral-source-ports}
The native QCD gauge group, its $24$-vertex Hurwitz stencil, the
color-singlet Hamiltonian $H_L$, the Wilson/source operators and
the charge/kinetic ports above remain unmodified.  The following
are typed comparison-source results: a five-core $2I$ completion
is not an ordinary automorphism of the same isolated Hurwitz graph.

\begin{proposition}[Formal nine-dimensional meson flavor-source port]
\label{p10:prop:v1173-formal-two-core-meson-port}
Let a supplied five-core $2I$ packet provide the $20$ directed core
pair labels and let $\mathcal N_9$ be the $A_5$ submodule
$\mathbf1\oplus\mathbf3\oplus\mathbf5$ formed by the unoriented
singlet and quintet and the momentum-compatible oriented triplet.
The embedded $A_5$ action on $V=\operatorname{Im}\mathbb H$ supplies
an equivariant real isometry
$\Upsilon:\mathcal N_9\to\operatorname{Herm}_3(\mathbb C)$.
For the already defined meson source $O_{\Gamma,T}^{(L)}$, the
linear word $O_{\Gamma,\Upsilon(\xi)}^{(L)}$ is an admissible
\emph{formal candidate} source label for every $\xi\in\mathcal N_9$.
In the finite free color-singlet one-pair CAR test, orthonormal
$\xi_1,\ldots,\xi_9$ have identity Gram
$\langle\xi_i,\xi_j\rangle=\delta_{ij}$ under the normalized
matrix-trace pairing.  This is not a claim about a physical
interacting-vacuum pole residue or an actual $A_5$ action on flavor.
\end{proposition}
\begin{proof}
The pair module has $A_5$ character
$\mathbf1\oplus2\mathbf4\oplus\mathbf5\oplus\mathbf3\oplus\mathbf3'$.
The chosen nine has the same irreducible character as
$\operatorname{Herm}_3(\mathbb C)$ under the real
three-dimensional icosahedral representation, namely
$\mathbb RI\oplus i\mathfrak{so}(3)\oplus
\operatorname{Sym}_0^2(\mathbb R^3)$.
Orthonormal transport gives the stated candidate isometry.
The color-singlet one-pair normalization
$3^{-1/2}\sum_c q_{fc}^\dagger\bar q_{gc}^\dagger|0\rangle$
gives the Hilbert--Schmidt Gram $\operatorname{tr}(T^\dagger S)$
by the finite CAR relations.  None of these algebraic operations
supplies physical source occupancy.
\end{proof}

\begin{proposition}[Three-face color-current relative Gram]
\label{p10:prop:v1173-three-core-color-current-gram}
In the already derived normalized antisymmetric three-quark color
singlet $|\epsilon\rangle$, let $T_i^A$ be the color generator
on constituent slot $i$, normalized by
$\sum_A(T_i^A)^2=4I/3$, and put
$C_{12}^A=T_1^A-T_2^A$, $C_{23}^A=T_2^A-T_3^A$,
$C_{31}^A=T_3^A-T_1^A$.  Then
\[
 \boxed{\sum_A\langle\epsilon|C_p^AC_q^A|\epsilon\rangle
       =6\left(I_3-\frac13J_3\right)_{pq}.}
\]
This positive relative-label Gram and the existing D22 label
Hessian share the rank-two $C_3$ projector.  A colored composite
junction current nevertheless requires a generated same-action
source between the actual gauge-face, quark and mobile-core owners;
this equality of labels alone is not such an interaction.
\end{proposition}
\begin{proof}
In the color singlet $\sum_A\langle T_i^AT_j^A\rangle=-2/3$
for $i\ne j$, since $(T_1+T_2+T_3)|\epsilon\rangle=0$.
The diagonal current Gram is $4$ and the off-diagonal is $-2$;
this is exactly the displayed matrix.
\end{proof}
% END PAPER VII-A v1.173


% BEGIN FRAMEWORK v3.023: MARKED 2I/HURWITZ TO NATIVE QCD SOURCE MAP
\section{Marked Hurwitz fan as a reference QCD quark-source map}
\label{p10:sec:v3023-marked-fan-source-map}
This is a map on the \emph{existing} $24$-vertex Hurwitz link/CAR
reference owner, not a second color gauge action or a claim that all
mobile $2I$ charts share one physical position-space realisation.
Retain the three face vertices $r_1,r_2,r_3$ of the oriented Hurwitz
edge $e=(p,q)$ and the existing two-edge transports
$A_i=U_{qr_i}U_{r_ip}$ from
Proposition~\ref{p10:prop:v2220-edge-fan-baryon}.
For a fixed gauge-link configuration define $U_{r_iv}$ ($v=p,q$)
as the already generated root-to-site parallel transports; here
$U_{r_iq}^{\dagger}U_{r_ip}=A_i$.

\begin{proposition}[Rooted three-quark CAR isometry and polarized Wilson overlap]
\label{p10:prop:v3023-marked-baryon-source-overlap}
Let $W=\mathbb C^2_{\rm spin}\otimes\mathbb C^3_{\rm flavor}$,
$\mathcal B_{56}=\operatorname{Sym}^3W$, and let
$b^\dagger_{r_i,d,\alpha}$ be three distinct-site quark creation
operators in the native reference CAR ledger. For $v=p,q$ put
\[
\mathcal I_{v,e}^B(\chi;U)=\frac1{\sqrt6}
\sum_{a,b,c;\,d_i,\alpha_i}
\epsilon_{abc}\chi_{\alpha_1\alpha_2\alpha_3}
\prod_{i=1}^3(U_{r_iv})_{d_ia_i}
 b^\dagger_{r_i,d_i,\alpha_i}|0\rangle,
\qquad(a_1,a_2,a_3)=(a,b,c),
\]
with factors ordered by the distinct site labels $r_1,r_2,r_3$.
For the normalized symmetric-tensor inner product,
\[
\boxed{(\mathcal I_{v,e}^B)^\dagger\mathcal I_{v,e}^B=I_{56},
\qquad(\mathcal I_{q,e}^B)^\dagger\mathcal I_{p,e}^B
 =W_e(U)I_{56},}
\]
where $W_e$ is \emph{the existing} polarized determinant in
\eqref{p10:eq:v2220-baryon-wilson}. The displayed scalar overlap
assumes a common aligned spin frame and a fixed ordered three-site
orbital reference. The operator-valued, simultaneously transformed
link/quark section is locally $SU(3)$ gauge invariant on the
admitted Gauss-compatible source domain; frozen-link source vectors
alone are not asserted to be physical Gauss states.
\end{proposition}
\begin{proof}
For each root, unitarity of the three color transporters and
the distinct-site CAR yield the normalized color contraction
$\sum_{abc}|\epsilon_{abc}|^2/6=1$ and the symmetric spin--flavor
inner product. There are $\binom{6+3-1}{3}=56$ such states. The
cross-root contraction factors over the three fixed sites and
produces $U_{r_iq}^\dagger U_{r_ip}=A_i$ on each color index.
The two epsilon tensors give precisely $W_e$, with the same
orientation as \eqref{p10:eq:v2220-baryon-wilson}. Under a local
gauge change the site modes and link transporters transform together;
the color epsilon cancels the common root $SU(3)$ matrix. This
uses only the native color connection and CAR ledger.
\end{proof}

\begin{proposition}[Spin-transport correction and two-leg meson overlap]
\label{p10:prop:v3023-spin-meson-source-overlap}
Let $S_{r_iv}$ be the native unitary spin parallel transport, acting
on the spin factor of $W$, and let $\Pi_{56}$ be the orthogonal
symmetric-cube projector. With the same color dressing, the full
cross-root baryon source Gram is
\[
\boxed{(\mathcal I_{q,e}^B)^\dagger\mathcal I_{p,e}^B
 =W_e(U)\,\Pi_{56}\Bigl[\bigotimes_{i=1}^3
 (S_{r_iq}^\dagger S_{r_ip})\Bigr]\Pi_{56}
 \Big|_{\operatorname{Sym}^3W}.}
\]
The simple $W_eI_{56}$ result is restricted to trivial compressed
relative spin transport. Further, with the already defined real
isometry $\Upsilon:\mathcal N_9\to\operatorname{Herm}_3(\mathbb C)$
from Proposition~\ref{p10:prop:v1173-formal-two-core-meson-port},
a normalized $q\bar q$ reference source at two distinct face sites
$r_i,r_j$ transported to $v=p,q$ has, in a common spin/Dirac
channel, the exact real-linear flavor-module overlap
\[
\boxed{(\mathcal I_{q,e}^{M,ij})^\dagger
 \mathcal I_{p,e}^{M,ij}
 =\frac13\operatorname{tr}(A_iA_j^\dagger)\,I_9.}
\]
No physical $A_5$ flavor-symmetry occurrence, interacting pole,
quark mass, decay width or S-wave spatial symmetrizer follows from
these reference CAR maps.
\end{proposition}
\begin{proof}
The spin parallel transports act independently on the three
ordered sites, and their overlap is compressed to $\mathcal B_{56}$.
The independent color contraction is the same $W_e$ as above.
For the meson, the normalized transported color singlet
$3^{-1/2}\sum_a(U_{r_iv}e_a)\otimes
\overline{(U_{r_jv}e_a)}$ has root-to-root overlap
$\operatorname{tr}(A_iA_j^\dagger)/3$; the Hilbert--Schmidt
orthonormality of the nine $\Upsilon$ matrices supplies $I_9$.
\end{proof}

\begin{proposition}[Exact Wilson presentation defect and its nonzero witness]
\label{p10:prop:v3023-two-presentation-holonomy-defect}
Suppose a second, 600-cell-dual path triple $A_i^{600}$ and the
Hurwitz triple $A_i^H$ have both been defined in one common gauge
owner, with the same endpoint transformation law. Set
$R_i=A_i^{600}(A_i^H)^{-1}\in SU(3)$ and
$d_i=\|R_i-I\|_{\rm op}$. Then each $d_i$ is gauge invariant and
\[
\boxed{|W(A_1^{600},A_2^{600},A_3^{600})
 -W(A_1^H,A_2^H,A_3^H)|\le\sum_{i=1}^3d_i.}
\]
The combinatorial duality alone does not make the defect zero:
for $A_i^H=I$, $A_2^{600}=A_3^{600}=I$ and
$A_1^{600}=\operatorname{diag}(e^{i\pi/3},e^{-i\pi/3},1)$,
$W_H=1$ but $W_{600}=2/3$.
\end{proposition}
\begin{proof}
Write $W(A_1,A_2,A_3)=\langle\epsilon|
A_1\otimes A_2\otimes A_3|\epsilon\rangle$ and telescope
the difference one tensor factor at a time. The other factors are
unitary and $\|\epsilon\|=1$, so each term is bounded by
$\|A_i^{600}-A_i^H\|_{\rm op}=d_i$. The witness follows from
$W(A,I,I)=\operatorname{tr}(A)/3$. The common-endpoint gauge
transformation conjugates $R_i$ and preserves its norm.
\end{proof}

\begin{remark}[Source--action and physical-limit firewall]
\label{p10:rem:v3023-global-apbi-remains-open}
The marked-fan isometries do not identify the distinct left-coset
600-cell nearest edges with intra-Hurwitz links, nor prove a
simultaneous five-core/25-transporter same-action APBI. A complete
physical map still needs compatible gauge and spin transitions,
quark-source grading, position/clock assignment, actual nonlinear
action and measure, Gauss quotient, and nonzero interacting pole
residues with the full Schur/Feshbach leakage term. The new source
operators are not extra elementary species; hence their existence
alone cannot change the one-loop strong beta coefficient.
\end{remark}
% END FRAMEWORK v3.023


% BEGIN FRAMEWORK v3.027: THREE-PATH CARTAN AND SOURCE-OWNED CHARACTER LEAKAGE
\section{Geometric three-path comparison and the native character source test}
\label{p10:sec:v3027-600-hurwitz-cartans}
The $600$-cell incidence and five Hurwitz cosets in Paper~I-B are
mathematical finite geometry. The native spatial QCD links here live
on the eight-neighbour Hurwitz $24$-cell, not on the $600$-cell's
nearest edges. The following \emph{comparison extension} applies
the already defined Cartan map
$g(q)=\operatorname{diag}(\rho(q),1)$ to the distinct $600$-cell edges;
it is not an independently generated gauge field on those edges.

\begin{proposition}[Three-path polar color and Clifford-tangent reconstruction]
\label{p10:prop:v3027-600-three-path-cartan}
Fix one Hurwitz nearest pair $q,q'$ in a left coset, and let
$\mathcal I(q,q')=\{r:q\sim_{600}r\sim_{600}q'\}$.
The set has exactly three elements. For the ordered comparison links
$U^{600}_{qr}(s)$ and $U^{600}_{rq'}(s)$, define
\[
 \mathscr T_{qq'}(s)=\frac13\sum_{r\in\mathcal I(q,q')}
 U^{600}_{qr}(s)U^{600}_{rq'}(s).
\]
At the three source profile values $s\in\{0,\tfrac12,1\}$,
\[
\boxed{\operatorname{polar}_{\rm unitary}(\mathscr T_{qq'}(s))
     =U^{24}_{qq'}(s).}
\]
At $s=\tfrac12$ the singular values of $\mathscr T$ are
$(\kappa,\kappa,1)$, where
$\kappa=\sqrt3/2+\sqrt{15}/30>0$.
The native unit spatial Clifford tangent also obeys
\[
\boxed{\mathbf n_{qq'}=
 \frac{\sin(\pi/5)}{\varphi\sqrt3}
 \sum_{r\in\mathcal I(q,q')}
   (\mathbf n_{qr}+\mathbf n_{rq'}).}
\]
The polar operation cannot be replaced by one common scalar
normalization since the nontrivial singular values differ from $1$.
\end{proposition}
\begin{proof}
The exact $600$-cell two-step return identity in Paper~I-B
supplies the three intermediates. In its $\mathbb Z[\varphi]$
quaternionic chart their coordinate sum is
$\sum_{r\in\mathcal I}r=\varphi(q+q')$.
Taking tangent components gives the displayed unit-vector identity;
ordered quaternionic products with $g(q)$ at the three named profile
values yield the positive two-color factor $\kappa$ and unit third
color. All links transform with the same endpoint gauge ownership,
so the unitary polar factor has the native Hurwitz link transformation
law; the finite group action transports this representative
calculation to the $480$ native links in the five cosets.
The construction remains a comparison, not a source-licensed
nonlinear polar assembly in the action.
\end{proof}

\begin{theorem}[Cartan-dependent lifting of geometric dark directions]
\label{p10:thm:v3027-cartan-dark-rank}
Let $\mathcal B_s$ be the $24\times96$ inter-coset incidence
extended to site$\otimes\mathbb C^3$ by the preceding \emph{unitary}
comparison links, and take source-side nullity as
$\dim\ker\mathcal B_s^*$. For the three source profile values,
\[
\boxed{\dim\ker\mathcal B_s^*=
  \begin{cases}6&s=0,\\2&s=\tfrac12,\\6&s=1.\end{cases}}
\]
At $s=\tfrac12$ the two-color $48\times48$ source Gram matrix has
exact squared singular spectrum
\[
\boxed{\left\{
 (3-3/\sqrt5)^{(8)},
 (9-1/\sqrt5)^{(24)},
 (18+2/\sqrt5)^{(12)},
 (30+6/\sqrt5)^{(4)}\right\}.}
\]
The untouched third color retains a two-dimensional kernel.
At every named profile $\|\mathcal B_s\|_{\rm HS}^2=864$;
thus the change is one of rank/coherent cancellation rather than
increased total squared coupling. These are fixed-Cartan comparison
ranks, not physical color-singlet decay projectors.
\end{theorem}
\begin{proof}
At the two endpoints the links are identity or pure-gauge, so their
Gram is unitarily equivalent to
$(12I+3N)\otimes I_3$, of rank $3(24-2)=66$.
At the middle profile the two-color Gram equals
$12I_{48}+\beta\mathscr M$ with
$\beta=(15+\sqrt5)/10$ and
$\mathscr M=N\otimes I_2+D_\rho(N\otimes I_2)D_\rho^*$.
Its exact Gaussian-integer minimal-polynomial certificate,
traces and eigenvalue multiplicities are given in Paper~VII-D
and imply the four strictly positive values displayed here.
The other color has the scalar Gram $12I+3N$, hence kernel rank $2$.
There are $288$ directed coset-to-complement edges and every
comparison link is unitary on three colors; the squared
Hilbert--Schmidt norm is therefore $288\cdot3=864$.
\end{proof}

\begin{theorem}[Character conservation in the native spatial kernel and exact temporal leakage]
\label{p10:thm:v3027-native-character-clutch}
Now use the \emph{actual source-owned} three-clock-layer,
$864$-dimensional Hermitian Hurwitz regulator $H_*$ of this paper,
not the extended $600$-cell comparison. On site space put
$E_{\rm char}=E_{-4}$, the rank-two projector onto the nontrivial
characters of $2T/Q_8$, and define
\[
 \Pi=I_{\rm clock,3}\otimes E_{\rm char}
          \otimes I_{\rm spin,4}\otimes I_{\rm color,3},
 \quad \operatorname{rank}\Pi=72.
\]
Each native spatial Cartan block, the native diagonal spin block,
and the two interior identity temporal links preserve the range
of $\Pi$. The closing temporal gauge clutch does not:
\[
\boxed{\|(I-\Pi)H_*\Pi\|_{\rm HS}^2=16,
\qquad\|[H_*,\Pi]\|_{\rm HS}^2=32.}
\]
The projector $\Pi$ is a fixed-chart mathematical character test,
not a licensed gauge-invariant hadron preparation or decay projector.
\end{theorem}
\begin{proof}
Right Hurwitz multiplication preserves the native edge tangent and
the right-frame Cartan link because
$\rho(qh)\rho(rh)^*=\rho(q)\rho(r)^*$.
Consequently the native spatial kernels commute with the right
$2T$ regular representation and preserve both one-dimensional
quotient-character sectors. Diagonal terms and interior temporal
identity links do likewise. The closing sitewise color transporter
is $\mathcal G=\operatorname{diag}_{q\in2T}(\rho(q)\oplus1)$.
Its compression on $E_{\rm char}\otimes\mathbb C^2$ is zero:
character products factor through $2T/Q_8$ while
$\sum_{h\in Q_8}\rho(h)=0$ (already the central pair
$h,-h$ cancels). The color singlet component $1$ remains internal.
Thus the two-character, two-color source subspace of dimension $4$
leaks completely on the closing link. The existing clock spin factor
$F=\frac12\Gamma_5(\Gamma_4-I_4)$ obeys
$\|F\|_{\rm HS}^2=2$; one closing orientation contributes $4\cdot2=8$
and its Hermitian opposite contributes another $8$.
There is no other leakage, so its squared norm is $16$.
For any Hermitian $H_*$ and orthogonal $\Pi$,
$\|[H_*,\Pi]\|_{\rm HS}^2=2\|(I-\Pi)H_*\Pi\|_{\rm HS}^2$,
which gives $32$.
\end{proof}

\begin{remark}[Finite source response is not yet a physical decay law]
\label{p10:rem:v3027-character-leakage-firewall}
The ten scalar $600$-cell zero modes, the rank change in an extended
Cartan edge comparison, and the nonzero native temporal-clutch leakage
are distinct operators with distinct domains. They are not identified
as the same physical transition matrix. A gauge-invariant physical
Fock projector, a same-action intermediate Clifford propagator,
and complete daughter-channel residues remain independently required
before extracting any decay selection rule or branching fraction.
\end{remark}
% END FRAMEWORK v3.027 VII-A MODULE


% BEGIN GAUSS-REDUCED SHAPE RESPONSE INTERFACE v1.17
\section{Same-action QCD shear Hessian and electromagnetic current source}
\label{bh:berger:v117:p7a:shear-source}
\begin{definition}[Metric-dependent physical finite source member]
\label{bh:berger:v117:p7a:source-data}
Given an \emph{actually completed} finite QCD transfer member with its
local Gauss projection, specified physical finite source state,
normalized clock, sources, link/occupation data and complete matter
metric variation, let $H_L(B)$ denote the resulting physical
Hamiltonian under the same volume-preserving position-space shear
$B\in V_5$. Its derivatives $H_a$, $H_{ab}$ include variation of the
local Gauss projector, finite embedding, kinetic normalization,
localizers and contact terms when these depend on the metric.
This is a typed request to the existing compiler, not an appended
phenomenological shape Hamiltonian.
\end{definition}
\begin{proposition}[Finite ground-state source Hessian and current f-sum]
\label{bh:berger:v117:p7a:source-hessian}
On a gapped, nondegenerate, locally Gauss-invariant ground of the
above physically specified member, with $H_L(0)|r\rangle=E_r|r\rangle$,
its complete stationary ground-energy shape Hessian is
\[
 K_{ab}=\langle0|H_{ab}|0\rangle
 -2\operatorname{Re}\sum_{r>0}
   \frac{\langle0|H_a|r\rangle
         \langle r|H_b|0\rangle}{E_r-E_0}.
\]
For a self-adjoint physical electromagnetic-current component $J_u$
acting on the same Gauss-invariant finite Hilbert space,
\[
 \sum_{r>0}(E_r-E_0)|\langle r|J_u|0\rangle|^2
 =\tfrac12\langle0|[J_u,[H_L(0),J_u]]|0\rangle\ge0.
\]
Neither formula assigns a physical transition energy or nonzero
current residue without the actual interacting Hamiltonian.
For a degenerate ground use the correct spectral projection and
reduced resolvent; for thermal states use the finite Duhamel kernel.
\end{proposition}
\begin{proof}
Differentiate the simple ground eigenvalue twice and resolve
$|\partial_a0\rangle$ over its orthogonal excited states; include
the direct derivative $H_{ab}$. The second identity follows by
inserting a complete finite energy eigenbasis into the double
commutator. Local electromagnetic color-singlet Noether bilinears
commute with the local color Gauss action on the properly realized
finite gauge carrier; this does not by itself make a residue nonzero.
\end{proof}
\begin{remark}[No automatic octahedral instability or radiative decay]
\label{bh:berger:v117:p7a:instability}
Only if the actual physical source state retains one octahedral
frame may the computed Hessian be written $k_2P_2+k_3P_3$.
The sign of $k_2-k_3$ and whether the complete \emph{constrained}
gravitational operator develops a growing mode are not determined by
the representation labels. A finite nonzero current matrix element
is not a decay width: require energy support, gauge/photonic
constraints, physical poles, complete channels and observation.
\end{remark}
% END GAUSS-REDUCED SHAPE RESPONSE INTERFACE v1.17




% BEGIN FRAMEWORK v3.068 GEOMETRY BRIDGE SOURCE INTEGRATION
\section{Source-scoped regulator transitions, Wilson root, and complete orbital weak derivative}
\noindent\emph{Scope.} The statements below retain their declared finite geometric, source-chart, and action-relative hypotheses. They do not assert physical particle decay or unconditional source occupation.
\begin{theorem}[Exact native character-clutch bright projector]
\label{p10:thm:v368-character-bright-projector}
In the certified unit-background $864$-dimensional Hermitian Hurwitz
regulator $H_*$ of Theorem~\ref{p10:thm:v199-native-unit-component},
let $\Pi=I_{3}\otimes E_{-4}\otimes I_{4}\otimes I_{3}$ and
$Q=I-\Pi$. Define $B=QH_*\Pi$ and $P_\pm=(I_4\pm\Gamma_4)/2$.
With $P_2=\operatorname{diag}(1,1,0)$ in the fixed Cartan color chart,
\[
\boxed{B^\dagger B=
 |0\rangle\langle0|\otimes E_{-4}\otimes P_-\otimes P_2
 +|2\rangle\langle2|\otimes E_{-4}\otimes P_+\otimes P_2.}
\]
In particular, $B^\dagger B$ is a rank-$16$ orthogonal projector on
$\operatorname{ran}\Pi$; its spectrum is
$\{1^{(16)},0^{(56)}\}$ and $\|B\|_{\rm HS}^{2}=16$.
\end{theorem}
\begin{proof}
The closing temporal transport is
$\mathcal G(q)=\operatorname{diag}(\rho(q),1)$ on the source's
$2T$ Hurwitz sites. The nontrivial quotient characters of
$2T/Q_8\cong C_3$ are constant on each $Q_8$ coset and
$\sum_{h\in Q_8}\rho(h)=0$ by the $h,-h$ pairing. Therefore the
rank-two character space has zero compression of the two-color
closing transporter. Since the transporter is unitary, it maps
the resulting four-dimensional character--two-color input
isometrically into its orthogonal complement. The temporal spin word
$F=\Gamma_5(\Gamma_4-I_4)/2$ obeys
$F^\dagger F=P_-$ and $FF^\dagger=P_+$.
The source's spatial blocks and two interior clock links preserve
$\operatorname{ran}\Pi$ by the existing
Theorem~\ref{p10:thm:v3027-native-character-clutch}; hence the
closing link and its Hermitian opposite are the only leakage blocks.
Their input clock supports are orthogonal, yielding exactly the
displayed Gram. Both summands have rank $2\cdot2\cdot2=8$.
\end{proof}

\begin{theorem}[Three exact native spectral-reachability steps]
\label{p10:thm:v368-character-onset-filtration}
In the same fixed unit regulator, let
$D=\Pi H_*\Pi\vert_{\operatorname{ran}\Pi}$, and put
\[
\mathcal O_k=\begin{bmatrix}B\\BD\\\vdots\\BD^{k-1}\end{bmatrix},
\qquad
\mathcal N_k=\bigcap_{j=1}^{k}
\ker(QH_*^j\vert_{\operatorname{ran}\Pi}).
\]
Then
\[
\boxed{\operatorname{rank}\mathcal O_1=16,\quad
       \operatorname{rank}\mathcal O_2=32,\quad
       \operatorname{rank}\mathcal O_3=48,}
\]
\[
\boxed{\dim\mathcal N_0=72,\quad\dim\mathcal N_1=56,
\quad\dim\mathcal N_2=40,\quad\dim\mathcal N_3=24.}
\]
The final $24$-dimensional common kernel is the reducing Cartan
third-color line in the selected $SU(2)\subset SU(3)$ background.
It is not a physical $SU(3)$ color-singlet particle sector.
\end{theorem}
\begin{proof}
The preceding theorem implies $\operatorname{rank}\mathcal O_k\le16k$.
The Cartan third-color line is fixed by all source-owned spatial,
interior temporal and closing links and reduces the Hermitian
$H_*$ on the two-character sector, giving an upper bound $48$
on the total observable rank.
The native site/clock/spin/color stencil compresses on the remaining
$48$ character--two-color states to a matrix $D_{(2)}$ over
$\mathbb Q(i,\sqrt3)$, with the closing-link compression zero by
the $Q_8$ character sum.
The corresponding exact matrices
\[
P_{\rm bright}^{(2)},\qquad
\begin{bmatrix}P_{\rm bright}^{(2)}\\
P_{\rm bright}^{(2)}D_{(2)}\end{bmatrix},\qquad
\begin{bmatrix}P_{\rm bright}^{(2)}\\
P_{\rm bright}^{(2)}D_{(2)}\\
P_{\rm bright}^{(2)}D_{(2)}^2\end{bmatrix}
\]
have nonzero minors of orders $16,32,48$ respectively, under the
good field reduction
$i\mapsto22$, $\sqrt3\mapsto10$, and
$(-1+i\sqrt3)/2\mapsto61$ in $\mathbb F_{97}$.
The minor determinants there are $35,16,61$, respectively, and
all numerator/denominator operations are licensed in the reduction.
Their exact source-derived entries and their row/column witness
indices are supplied by the companion
\texttt{Paper00-GeometryObligationBridges-v1\_7-exact-ranks.py}
certificate. Nonzero modular minors imply rank at least $16,32,48$
in characteristic zero; the preceding upper bounds prove equality.
Finally, for any Hermitian block operator
$\left(\begin{smallmatrix}D&B^\dagger\\B&C\end{smallmatrix}\right)$,
the stacked maps $QH_*^j\Pi$ and $BD^{j-1}$ have the same common
kernel by triangular elimination of the complementary return terms.
Thus the stated filtration dimensions follow.
\end{proof}

\begin{corollary}[Rank and kernel of the full native Schur self-energy]
\label{p10:cor:v368-native-character-schur-support}
With $C=QH_*Q\vert_{\operatorname{ran}Q}$, the exact self-energy
$\Sigma_{\Pi}(z)=B^\dagger(z-C)^{-1}B$ obeys
\[
\boxed{\operatorname{rank}\Sigma_{\Pi}(z)=16,
\quad\ker\Sigma_{\Pi}(z)=\ker B\quad(\operatorname{Im}z\ne0),}
\]
and $\Sigma_\Pi(z)=z^{-1}P_{\rm bright}+O(z^{-2})$ as
$|z|\to\infty$. The fixed kernel belongs to the self-energy,
not necessarily to the full projected resolvent.
\end{corollary}
\begin{proof}
For $z=x+iy$ with $y\ne0$ and any $v$ with $Bv\ne0$,
\[
\operatorname{Im}\langle v,\Sigma_{\Pi}(z)v\rangle
=-y\langle Bv,[(x-C)^2+y^2]^{-1}Bv\rangle\ne0.
\]
The bright projector theorem makes $B$ injective on its
$16$-dimensional bright source space and zero on its $56$-dimensional
kernel. The large-$z$ identity is the ordinary resolvent expansion.
\end{proof}

\begin{remark}[Action ownership and nonpromotion to a decay width]
\label{p10:rem:v368-native-onset-scope}
These theorems concern the action-derived \emph{Hermitian Wilson
regulator} at a fixed native Hurwitz background. In particular
$e^{-isH_*}$ is an auxiliary spectral evolution, not the proper
time evolution of a parent-particle Fock preparation. The invariant
Cartan third-color line is not a physical $SU(3)$ gauge singlet.
Nothing here identifies the comparison $600$-chamber face adjacency
with a QCD operator, replaces the action-owned full Fock propagator,
or supplies the parent/daughter state map and physical width conditions
in Definition~\ref{p5c:def:v3027-geometric-decay-certificate}.
\end{remark}
\begin{theorem}[Fourth-winding Wilson root and native full-rank transfer]
\label{p10:thm:v368-wilson-root-clutch}
On the certified native three-clock Cartan background of
Theorem~\ref{p10:thm:v199-native-unit-component}, let
$\Omega(q)=U_{01}(q)U_{12}(q)U_{20}(q)=\mathrm{diag}(\rho(q),1)$,
$q\in2T$. The locally gauge-invariant normalized fourth temporal
Wilson winding satisfies
\[
 w_4(q)=\tfrac13\operatorname{tr}\Omega(q)^4
 =\mathbf1_{Q_8}(q),\qquad
 P_4=I_{\rm clock,3}\otimes M_{w_4}\otimes I_{\rm spin,4}
 \otimes I_{\rm color,3},\quad\operatorname{rank}P_4=288.
\]
The three $Q_8$-coset averages of the winding traces agree at
$n=1,2,3$ but differ first at $n=4$, with averages $(1,0,0)$.
For the genuine $864$-dimensional Hermitian native source regulator
$H_*$, set $B_4=(I-P_4)H_*P_4$. Then
\[
 \operatorname{rank}B_4=288,\qquad
 \|B_4\|_{\rm HS}^2=648,\qquad
 \|[H_*,P_4]\|_{\rm HS}^2=1296.
\]
For every nonreal $z$, the exact Schur self-energy
$\Sigma_4(z)=B_4^\dagger[z-(I-P_4)H_*(I-P_4)]^{-1}B_4$
on the rooted space has rank $288$.
\end{theorem}
\begin{proof}
On $Q_8$, $q_0\in\{0,\pm1\}$; on the 16 complementary
Hurwitz sites $q_0=\pm\frac12$. Since the SU(2) trace gives
$\frac13\operatorname{tr}\Omega(q)^n=\frac13[1+2T_n(q_0)]$,
$T_4(x)=8x^4-8x^2+1$ gives the pointwise root selector.
Each $Q_8$ coset is a rotated cross-polytope and hence a
spherical three-design, implying equality of all degree-at-most-three
coset averages. There are no $Q_8$--$Q_8$ native nearest edges;
each of the eight root sites has eight complement neighbors.
Each native Clifford--color edge block
$\frac38\Gamma_5(\Gamma\cdot n-I_4)\otimes U_{qr}$ has
Hilbert--Schmidt square $(3/8)^2(8)(3)=27/8$.
Clock and site-diagonal terms preserve $P_4$, yielding
$3\cdot64\cdot27/8=648$ and the Hermitian commutator identity.
Full column rank in each clock is certified over
$\mathbb Q(i,\sqrt3)$ by nonzero $96$-minors with residues
$22,86,22$ modulo $97$ under $i\mapsto22$,
$\sqrt3\mapsto10$, with all denominators invertible.
Thus $\ker B_4=0$. For nonreal $z=x+iy$, $y\ne0$,
$\operatorname{Im}\langle v,\Sigma_4(z)v\rangle
=-y\langle B_4v,[(x-C)^2+y^2]^{-1}B_4v\rangle\ne0$
for all $v\ne0$, where $C=(I-P_4)H_*(I-P_4)$ on its range.
\end{proof}
\begin{proposition}[Root Wilson character-trine and Gauss-compatible source]
\label{p10:prop:v368-wilson-root-fock-source}
Let $E$ be the normalized two-character $C_3=2T/Q_8$
site embedding. Then
\[
 E^\dagger M_{w_4}E=\tfrac13\begin{pmatrix}1&1\\1&1\end{pmatrix},\qquad
 \operatorname{spec}=\{2/3,0\}.
\]
Its three coset translates form an exact trine POVM on $\mathbb C^2$.
On any Gauss-compatible finite Fock chart, the composite of this
closed-Wilson trace and the already defined local $O_B$ (respectively
local meson bilinear) is gauge invariant and preserves the appropriate
baryon number of the source. At the reference background it restricts
the 24-site local baryon source to the eight rooted sites; the
$56$-dimensional $\operatorname{Sym}^3(\mathbb C^2\otimes\mathbb C^3)$
spin--flavor factor is unchanged. This source construction establishes
neither physical parent/daughter pole projectors nor partial widths.
\end{proposition}
\begin{remark}[Reference/background and action/decay firewalls]
$M_{w_4}$ is a numerical orthogonal projection only on the declared
Cartan unit reference configuration. For a general dynamical link
configuration, the fourth Wilson trace is a gauge-invariant operator
but not in general a projection. The current full-column-rank theorem
is about $H_*$, not the Gauss-reduced interacting QCD Fock Hamiltonian.
No $600$-chamber action edge or proper decay clock is supplied by it.
\end{remark}

\begin{proposition}[Center-neutral Wilson root effect on a fluctuating gauge chart]
\label{p10:prop:v368-center-neutral-Wilson-effect}
For any closed temporal $SU(3)$ holonomy $\Omega$, the effect
$\mathcal R_2=|\operatorname{tr}\Omega^2|^2/9$
is gauge invariant, temporal-center neutral, and satisfies
$0\le\mathcal R_2\le1$. On the source's fixed Hurwitz Cartan
reference clutch it equals $1$ at $q=\pm1$, $1/9$ at the six
noncentral $Q_8$ elements, and $0$ outside $Q_8$.
On general fluctuating backgrounds it is an effect, not a projector;
its finite pure-gauge heat-chart positivity does not prove an
interacting physical baryon-state occupation.
\end{proposition}
\begin{proof}
Color conjugation leaves the closed trace invariant, and the
center phase cancels against its complex conjugate. The bound
$|\operatorname{tr}\Omega^2|\le3$ follows from unitarity.
On the reference clutch $\Omega=\operatorname{diag}(\rho(q),1)$,
$\operatorname{tr}\Omega^2=1+2(2q_0^2-1)=4q_0^2-1$.
The stated support values follow from the Hurwitz coordinates.
\end{proof}
\begin{proposition}[Complete first-harmonic weak--Weyl derivative of the quadratic Cartan Petersen mode]
\label{p10:prop:v368-complete-first-shell-weak-derivative}
Fix a marked Hopf axis $u\in S^2$ and the existing Petersen flux $K$,
$KK^{\mathsf T}=5I_3$. Put $v_e=-2P_uKe/3$ and
$h_v(q)=q_0(v\cdot\operatorname{Im}q)$. Let
$a_v^T=P_T[h_vq^{-1}dq/2]$ be the background-gauge-transverse
Born connection variation, with the degree-two ghost inverse
$\chi_v=-\sum_j (X_jh_v)e_j/12$.
On the physical fixed-frame orbital restriction containing
$\varphi_a=2q_a$, $a=0,1,2,3$, define
\[
 N_{ac}(v)_{\alpha b}=
 \int_{S^3}\varphi_a\varphi_c\,
 a_v^T(X_b)_{\alpha}\,d\mu,
 \qquad
 F(v)_{ac}=\sum_{\alpha=1}^{2}\sum_{b=1}^3
 N_{ac}(v)_{\alpha b}\,\frac{\sigma_\alpha}{2}
 \otimes(-\sigma_b).
\]
Then each entry of $N_{ac}$ is an exact uniform 24-point Hurwitz
average. For every unit real $v\perp u$, $F(v)$ is Hermitian, has
rank $16$, and has spectrum
\[
 \left\{\pm\frac5{144},\ \pm\frac1{144},\
 \pm\frac{3+2\sqrt3}{144},\
 \pm\frac{2\sqrt3-3}{144}\right\},
\]
with each signed eigenvalue twice. For real transverse $v,w$,
$\operatorname{Tr}_{16}(F(v)F(w))=17v\cdot w/1296$.
Consequently the pulled Petersen connection-coordinate Gram in the
radius-$R$ orthonormal frame is
\[
 G_E=\frac{17}{2916R^2}K^{\mathsf T}P_uK,
\quad
 \operatorname{spec}G_E=\{(85/(2916R^2))^{(2)},0^{(13)}\}.
\]
On the complete left-chiral three-family, three-color quark first-harmonic
source block the spectrum is repeated ninefold for any unitary $C_{ud}$.
A rank-eight cross-block obtained by keeping only $\varphi_0\leftrightarrow
\varphi_i$ is not the complete first-shell operator.
\end{proposition}
\begin{proposition}[Naked weak adjoint is killed by Gauss averaging;
conditional static Cartan source equation]
\label{p10:prop:v368-naked-Gauss-and-frozen-source}
For the complete operator $F(v)$ above, the already present quaternionic
weak subgroup $Q_8\subset SU(2)$ satisfies
\[
 \frac18\sum_{g\in Q_8}
 (I_4\otimes\rho_2(g)\otimes I_2)F(v)
 (I_4\otimes\rho_2(g)^\dagger\otimes I_2)=0.
\]
This is a no-go for a *naked* weak-adjoint insertion between local weak
Gauss singlets, not for a jointly gauge-dressed field--current operator.
For the existing static nonconstant Cartan energy in the two Petersen
coordinates $\xi=P_uKe$, one has
\[
 \mathcal E(\xi)=\frac14+\frac5{324}\|\xi\|^2
 +\frac1{4320}\|\xi\|^4.
\]
For an externally specified, correctly normalized *frozen* real source
$-j\cdot\xi$, the unique static stationary solution is
\[
 \bigl(5/162+\|\xi\|^2/1080\bigr)\xi=j.
\]
Neither proposition selects a physical excitation, interacting pole or
physical partial width. The gauge-dressed complete source response must
be evaluated on the same-action physical Fock/Gauss carrier.
\end{proposition}
% END FRAMEWORK v3.068 GEOMETRY BRIDGE SOURCE INTEGRATION


% BEGIN PAPER00 v1.5 -- exact rooted and dressed action currents
% New source-owned theorems for Paper VII-A (proposed v1.185).
% Derived from exact rooted CAR source p10:prop:v3023-marked-baryon-source-overlap,
% the Gauss projector p10:thm:v350-local-gauss-projector,
% and the finite Schur source p10:eq:finite-schur-source-block.
\section{Rooted axial insertions and all-energy Schur transition vertices}
\label{p10:sec:v375-rooted-dressed-current}

\begin{theorem}[Exact rooted three-quark current and Wilson polarization]
\label{p10:thm:v375-rooted-CAR-axial-insertion}
Let $W=\mathbb C^2_{\rm spin}\otimes\mathbb C^3_{\rm flavor}$ and let
$\mathcal I_{v,e}^{B}:\operatorname{Sym}^3W\to\mathcal H_{\rm CAR}$
be the normalized three-distinct-site, link-dressed source of
Proposition~\ref{p10:prop:v3023-marked-baryon-source-overlap}.
At site $r_j$ let $\mathcal J_j(A)$ be the colour-identity
normal-ordered one-body insertion of an endomorphism $A\in\operatorname{End}(W)$,
and put $\mathcal J(A)=\sum_{j=1}^3\mathcal J_j(A)$.
Write $T_{jv}$ for the unitary spin--flavour transport in a declared
source frame, and set
$M_{j}=T_{jq}^{\dagger}T_{jp}$ and
$D_j(A)=T_{jq}^{\dagger}AT_{jp}$.
Then on the fixed gauge-link/source reference chart,
\[
\boxed{
 (\mathcal I_{q,e}^{B})^\dagger\mathcal J(A)\mathcal I_{p,e}^{B}
 =W_e(U)\,\Pi_{\rm sym}
   \Big[\sum_{j=1}^{3}
   M_1\otimes\cdots\otimes D_j(A)\otimes\cdots\otimes M_3
   \Big]\Pi_{\rm sym}\Big|_{\operatorname{Sym}^3W}.}
\]
Here the $j$th factor is $D_j(A)$ and all others are $M_k$.
For aligned common spin frames $T_{jp}=T_{jq}=I_W$, this reduces to
\[
\boxed{(\mathcal I_{q,e}^{B})^\dagger\mathcal J(A)\mathcal I_{p,e}^{B}
 =W_e(U)\,d\Gamma_3(A),\qquad
 d\Gamma_3(A)=\sum_{j=1}^3 A_{(j)}\big|_{\operatorname{Sym}^3W}.}
\]
In particular for $A^{ia}=(\sigma^i/2)\otimes(\lambda^a/2)$
this is exactly $W_e(U)G^{ia}$ for the \emph{existing}
$56$-state primitive axial packet; with $p=q$ it is $G^{ia}$.
The same fixed-link identity is pointwise in a dynamical gauge-link
configuration and therefore holds as a fibrewise operator identity
on a well-defined Gauss-compatible link/CAR chart.
It makes no claim that the entire physical axial source derivative is
only this local one-body insertion; mode/spin reduction, moving sources,
link-source derivatives, and interacting pole dressings retain their
same-action conditions.
\end{theorem}
\begin{proof}
Use the three distinct-site CAR to identify the source space with
$(\mathbb C^3\otimes W)^{\otimes3}$ on its one-fermion-per-site block.
The colour component is the normalized antisymmetric tensor
$\epsilon/\sqrt6$ transported by $U_{r_jv}$, and the $W$ component
is the normalized symmetric tensor transported by $T_{jv}$.
Because $\mathcal J_j(A)$ is the colour identity on every site,
its colour contraction is the \emph{same} polarized determinant
$W_e(U)$ for each $j$, while its $W$ contraction places $D_j(A)$
in the $j$th slot and $M_k$ in the others. Compress by
$\Pi_{\rm sym}$ and sum. The aligned specialization is the
ordinary symmetric-cube differential of $A$. The $24$ Pauli--Gell-Mann
instances are exactly the previously defined $G^{ia}$.
The argument is pointwise in the link variables and uses the existing
root-to-site gauge transformation rule, so it survives the admissible
operator-valued link multiplication without a frozen-link ensemble
assumption. No expectation value or physical pole is used.
\end{proof}

% Proposed Paper VII-A insertion v1.188. This is a finite heat-reference theorem,
% not an occurrence claim about the full interacting QCD Hamiltonian.
\section{The minimal full-Hurwitz baryon heat membrane and its Gauss source}
\label{p10:sec:v378-baryon-heat-membrane}

\begin{theorem}[Unique three-face triality membrane of a Hurwitz baryon edge]
\label{p10:thm:v378-unique-baryon-triality-membrane}
For each oriented edge $e=(p,q)$ of the native $24$-cell, denote by
$r_1,r_2,r_3$ the third vertices of its three incident triangular
faces, and orient the three peripheral quark paths from $q$ through
$r_i$ to $p$.  In the oriented cellular chain complex with $\mathbb Z_3$
coefficients, the negative of the resulting six-edge baryon boundary
has no filling by one or two triangular faces, and has exactly one
filling by three \emph{distinct} triangular faces with coefficients
$\pm1$: the three faces $(p,q,r_i)$ with their cancellation
orientations. This holds for all $96$ edges.
\end{theorem}
\begin{proof}
The $24$ native Hurwitz units have adjacency $\langle v,w\rangle=1/2$.
By edge transitivity take $p=(1,0,0,0)$ and
$q=(1,1,1,1)/2$.  The three common adjacent vertices can be chosen
as $(1,1,-1,-1)/2$, $(1,-1,1,-1)/2$ and
$(1,-1,-1,1)/2$; their pairwise inner products vanish, so no two
are adjacent. The six peripheral edges are the distinct spokes
$(q,r_i)$ and $(r_i,p)$.  Any triangular face contains at most two
of these spokes. It contains two only if it is one of the three
$(p,q,r_i)$ faces: a face joining spokes with different $r_i$
would require two of the mutually nonadjacent third vertices to be
adjacent. Hence fewer than three faces cannot cancel all six spokes.
Exactly three faces must each cancel two spokes, forcing all three
incident faces and their unique cancelling orientations. The
remaining common edge has multiplicity three and vanishes in
$\mathbb Z_3$. Edge transitivity proves the result everywhere.
The complete rational-coordinate $96$-edge check is recorded in
\texttt{Paper00-HurwitzBaryonMembraneEvidence-v1\_1.json}.
\end{proof}

\begin{theorem}[Fixed leading baryon overlap on the full 96-face heat reference]
\label{p10:thm:v378-full-heat-baryon-leading}
On the native Hurwitz link configuration space take independent
normalized Haar measures $dU_e$ on its $96$ unoriented edges, the
existing $SU(3)$ heat kernel convention
$K_t(U)=\sum_R d_R e^{-tC_R/2}\chi_R(U)$,
and the \emph{homogeneous product-heat reference representative}
\[
 d\mu_t(U)=\frac{1}{Z_t}\prod_{f\in F_H}K_t(U_f)
                       \prod_{e\in E_H}dU_e,\qquad t>0.
\]
For the existing gauge-invariant baryon polarized determinant
$W_e$ of Proposition~\ref{p10:prop:v2220-edge-fan-baryon},
\[
\boxed{\langle W_e\rangle_{\mu_t}
=e^{-2t}+O(e^{-8t/3})\quad(t\longrightarrow\infty).}
\]
The leading coefficient is exactly one and is independent of the
oriented edge. In particular this reference expectation is real
and strictly positive for sufficiently large $t$.  The asymptotic
is \emph{not} an exact equality at arbitrary $t$, not the interacting
fermion/gluon vacuum expectation and not a physical decay width.
\end{theorem}
\begin{proof}
For large $t$, the heat-character series converges uniformly and
its lowest nontrivial terms are
$3e^{-2t/3}(\chi_{\mathbf3}+\chi_{\bar{\mathbf3}})$.
Every one- or two-face insertion leaves at least one of the six
peripheral edges of $W_e$ unsaturated and has zero Haar integral.
By Theorem~\ref{p10:thm:v378-unique-baryon-triality-membrane}, a
three-face insertion can satisfy the necessary center-triality
condition only on the three adjacent faces, each with the
cancelling orientation. Since each peripheral edge then carries one
source fundamental and one inserted representation, Haar
orthogonality forces the inserted representation to be exactly
$\bar{\mathbf3}$, rather than a higher triality-equivalent one.
Condition on the common edge link $U_{pq}$. The three two-edge
peripheral transports $A_i$ have independent Haar laws, and
\[
 \int_{SU(3)} A_{aa'}\,
   \overline{\operatorname{tr}(U_{pq}A)}\,dA
       =\frac{(U_{pq}^{\dagger})_{aa'}}{3}.
\]
The two epsilon contractions give
$\det(U_{pq}^{\dagger})/27=1/27$. Multiplication by the three
fundamental heat coefficients gives
$(3e^{-2t/3})^3/27=e^{-2t}$, without a free edge or
channel weight. Every other surviving insertion has at least four
nontrivial faces; the smallest nontrivial $SU(3)$ heat exponent per
face is $2t/3$, giving $O(e^{-8t/3})$. The normalized partition
function tends to one and differs from one by at most
$O(e^{-2t/3})$ from uniform heat-kernel convergence, so
normalization preserves the same remainder order. Charge
conjugation of all links makes the expectation real, and symmetry
of the homogeneous measure makes it edge-independent.
\end{proof}

\begin{corollary}[Normalized locally Gauss-invariant heat-reference baryon source]
\label{p10:cor:v378-gauss-rooted-heat-axial}
On the declared link Hilbert space
$L^2(SU(3)^{E_H},\prod_e dU_e)$ tensored with the existing
three-site color CAR carrier, let $\mathcal I^B_{v,e}(\chi;U)$
be the rooted source of
Proposition~\ref{p10:prop:v3023-marked-baryon-source-overlap}.
The square-root heat density times that source,
\[
 \Phi_{v,t}(\chi;U)=Z_t^{-1/2}
       \left(\prod_{f\in F_H}K_t(U_f)\right)^{1/2}
       \mathcal I^B_{v,e}(\chi;U),\qquad v=p,q,
\]
is a normalized linear isometry of the symmetric $56$-state
spin--flavor carrier into a \emph{locally Gauss-invariant reference}
subspace under the simultaneous link/quark gauge action. For
aligned spin frames, its root-to-root overlap and color-identity
one-body axial insertion are exactly
\[
\boxed{
 \Phi_{q,t}^{\dagger}\Phi_{p,t}=b_t I_{56},\qquad
 \Phi_{q,t}^{\dagger}\mathcal J^{ia}\Phi_{p,t}
       =b_t G^{ia},\quad b_t=\langle W_e\rangle_{\mu_t}.}
\]
Thus both have the action-owned leading coefficient
$b_t=e^{-2t}+O(e^{-8t/3})$ on this \emph{reference branch}.
The construction produces a nonzero Gauss-compatible finite source,
not a selected interacting QCD ground state, physical baryon pole
or proof of full source occurrence under G0.
\end{corollary}
\begin{proof}
The product heat density is invariant under independent local
$SU(3)$ transformations because each face holonomy transforms by
conjugation and $K_t$ is central. The existing link-dressed
three-quark epsilon source transforms invariantly with it.
Pointwise unitarity of its color transports gives
$(\mathcal I^B_{v,e})^\dagger\mathcal I^B_{v,e}=I_{56}$;
normalization of $\mu_t$ therefore gives the isometry. The
existing cross-root CAR contraction is
$(\mathcal I^B_{q,e})^\dagger\mathcal I^B_{p,e}=W_eI_{56}$.
The color-identity axial insertion similarly contracts to
$W_eG^{ia}$ by
Theorem~\ref{p10:thm:v375-rooted-CAR-axial-insertion}.
Integrating both identities against $\mu_t$ proves the result.
\end{proof}

\begin{proposition}[Scalar Wilson source overlap cancels from a canonically normalized bare current]
\label{p10:prop:v378-root-gram-current-cancellation}
The preceding cancellation holds more generally for any normalized,
spin--flavor-independent gauge-link reference density, not only the
heat representative.  Write $\Phi_p,\Phi_q$ for its two normalized
aligned-spin source maps and $b=\Phi_p^\dagger\Phi_q|_{\mathbf1_{56}}$.
The $112$-source Gram and the projected color-identity one-body
current factor \emph{in the same nonorthonormal source frame} as
\[
 S=\begin{pmatrix}1&b\\\overline b&1\end{pmatrix}\otimes I_{56},
 \qquad
 J_{\rm bare}^{ia}=\begin{pmatrix}1&b\\\overline b&1\end{pmatrix}
                      \otimes G^{ia}.
\]
For $|b|<1$, canonical Gram orthonormalization gives exactly
\[
 \boxed{S^{-1/2}J_{\rm bare}^{ia}S^{-1/2}
             =I_2\otimes G^{ia}.}
\]
For $|b|=1$, the Moore--Penrose version is the projector onto the
supported root space tensored with $G^{ia}$. In particular the
symmetric/antisymmetric normalized root sectors have no additional
primitive root-changing axial vertex. Any nontrivial additional
current on this source support must appear as a deviation of the
actual same-action current matrix from this source-Gram factorization.
\end{proposition}
\begin{proof}
At each fixed link configuration the color-identity one-body current
acts on the symmetric spin--flavor coordinate as $G^{ia}$ and leaves
the rooted color dressing unchanged. Its contraction between any two
root maps therefore equals their color Gram times $G^{ia}$. Integration
against any scalar gauge-link measure preserves this factorization.
Because $S$ is positive semidefinite, conjugating $S\otimes G^{ia}$
by its inverse square root yields the identity on the full root
space when $S$ is invertible and the support projector otherwise.
This is an identity for the bare source operator, not for a fully
dressed physical axial current.
\end{proof}


\begin{corollary}[Exact root-changing current residual after Gram normalization]
\label{p10:cor:v378-root-current-residual-certificate}
In the preceding nonsingular two-root source construction let
$\mathcal J_{\mathrm{src}}^{ia}$ denote \emph{any} Hermitian
source-projected current from a specified action, and set
$D^{ia}=\mathcal J_{\mathrm{src}}^{ia}-S\otimes G^{ia}$,
where now $S=\left(\begin{smallmatrix}1&b\\\bar b&1\end{smallmatrix}\right)$
acts on the root pair alone. For $0<|b|<1$, let $E_\pm$ denote its
rank-one spectral projectors. Then, writing
$\widehat{\mathcal J}^{ia}= (S^{-1/2}\otimes I_{56})
\mathcal J_{\mathrm{src}}^{ia}(S^{-1/2}\otimes I_{56})$,
\[
 \boxed{(E_-\otimes I_{56})\widehat{\mathcal J}^{ia}
 (E_+\otimes I_{56})=\frac{(E_-\otimes I_{56})D^{ia}
 (E_+\otimes I_{56})}{\sqrt{1-|b|^2}}.}
\]
In particular a nonzero normalized primitive root-changing current
requires a nonzero \emph{Gram-subtracted} action-current residual.
For this isolated pair the equality gives the exact Hilbert--Schmidt
norm amplification factor $(1-|b|^2)^{-1/2}$; it is not a physical
channel-weight prediction and does not apply at $|b|=1$ without
restricting to the one-dimensional root support.
\end{corollary}
\begin{proof}
The two eigenvalues of $S$ are $1\pm|b|$ and its positive spectral
square root commutes with $E_\pm$. The identity contribution
$I_2\otimes G^{ia}$ has vanishing $-+$ block. Left and right
multiplication by $S^{-1/2}$ on that block contribute precisely
$(1-|b|)^{-1/2}(1+|b|)^{-1/2}$. The residual identity follows.
\end{proof}

% END proposed v1.188 heat reference closure.

\begin{theorem}[Two-energy dressed-current Schur identity and pole vertices]
\label{p10:thm:v375-two-energy-dressed-transition}
Let $H=H^*$ be a fixed, fully specified finite Hamiltonian on the
same Gauss-compatible action domain as its self-adjoint source derivative
$J=\partial_\phi H$. For a \emph{fixed} orthogonal source/complement
splitting write
\[
 H=\begin{pmatrix}K&B^\dagger\\B&C\end{pmatrix},
 \qquad
 J=\begin{pmatrix}J_{PP}&J_{PQ}\\J_{QP}&J_{QQ}\end{pmatrix},
 \quad R_Q(z)=(z-C)^{-1},
 \quad F(z)=[z-K-B^\dagger R_Q(z)B]^{-1}.
\]
For $z_i,z_f$ in the common resolvent domain, define
\[
\boxed{\begin{aligned}
 \mathcal V_J(z_f,z_i)={}&J_{PP}
 +J_{PQ}R_Q(z_i)B+B^\dagger R_Q(z_f)J_{QP}\\
 &+B^\dagger R_Q(z_f)J_{QQ}R_Q(z_i)B.
\end{aligned}}
\]
Then the exact off-shell three-point resolvent identity is
\[
\boxed{P(z_f-H)^{-1}J(z_i-H)^{-1}P
 =F(z_f)\mathcal V_J(z_f,z_i)F(z_i).}
\]
Let $E_i,E_f$ be real isolated simple eigenvalues of $H$ outside
$\operatorname{spec}(C)$, whose source components are nonzero.
Choose unit source vectors $u_a$ spanning
$\ker[E_a-K-B^\dagger R_Q(E_a)B]$ and define
$Z_a=I+B^\dagger R_Q(E_a)^2B$.
The uniquely normalized lifted eigenvectors are
$\Psi_a=\binom{I}{R_Q(E_a)B}u_a/(u_a^\dagger Z_a u_a)^{1/2}$,
and their same-action transition is exactly
\[
\boxed{\langle\Psi_f,J\Psi_i\rangle=
\frac{u_f^\dagger\mathcal V_J(E_f,E_i)u_i}
{(u_f^\dagger Z_f u_f)^{1/2}(u_i^\dagger Z_i u_i)^{1/2}}.}
\]
Thus no additional \emph{finite} transition-vertex or pole-residue
normalization is permissible once $H$ and $J$ are fixed.
No real-pole identity alone supplies an absorptive width, daughter
scattering states, a current type, or a physical action member.
\end{theorem}
\begin{proof}
The two block columns of the full resolvent are
$(z-H)^{-1}P=\binom{I}{R_Q(z)B}F(z)$, and the corresponding
block row is $P(z-H)^{-1}=F(z)(I, B^\dagger R_Q(z))$.
Multiply with the complete $J$ between distinct spectral arguments;
the four blocks yield $\mathcal V_J$ exactly. At a real pole, the
complement eigenvector equation gives $v=R_Q(E)Bu$ and its squared
norm is $u^\dagger Z(E)u$. The first block equation gives the
source pole condition. Normalizing both lifted states and inserting
$J$ gives the claimed expression. No perturbation-series truncation
or chosen Schur path coefficients occur.
\end{proof}

\begin{corollary}[Isolated degenerate multiplets have a canonical dressed-current frame]
\label{p10:cor:v375-degenerate-multiplet-vertex}
The preceding theorem extends without a simplicity hypothesis to
an isolated real eigenvalue $E\notin\operatorname{spec}(C)$ of any
finite multiplicity $r_E$. Put
$N_E=\ker[E-K-B^\dagger R_Q(E)B]$ and
$Z_E=I+B^\dagger R_Q(E)^2B\succ0$.
Choose any column map $U_E:\mathbb C^{r_E}\to N_E$
with $U_E^\dagger Z_EU_E=I_{r_E}$.
Then $\Psi_E=\binom{I}{R_Q(E)B}U_E$ is an isometry onto the
\emph{full} eigenvalue-$E$ subspace, and
\[
\boxed{
 \operatorname*{Res}_{z=E}F(z)=U_EU_E^\dagger,
 \qquad \mathcal M_{fi}=U_{E_f}^\dagger
 \mathcal V_J(E_f,E_i)U_{E_i}.}
\]
The only remaining choice is a unitary basis change within each
actual degenerate eigenspace. Consequently the dressed-current theorem
applies to the degenerate $SU(3)\times SU(2)$ octet/decuplet
multiplets needed by Paper~VII-B, without selecting separate pole
vectors or inserting independent reduced coefficients.
\end{corollary}
\begin{proof}
The complement eigenvalue equation gives the one-to-one lifting
$u\mapsto (u,R_Q(E)Bu)$ from $N_E$ onto the full eigenspace.
The metric induced by this lift is $Z_E$, which is positive definite,
so $U_E$ can be chosen $Z_E$-orthonormal. The full eigenprojector
is $\Psi_E\Psi_E^\dagger$; compressing to $P$ gives its
source resolvent residue $U_EU_E^\dagger$. Inserting the full
current $J$ between the two lifted isometries yields
$U_{E_f}^\dagger\mathcal V_J(E_f,E_i)U_{E_i}$.
Unitary changes of basis within each eigenspace change only the
matrix presentation of the same physical current tensor.
\end{proof}

\begin{corollary}[Moment-certified bound on physical-current dressing]
\label{p10:cor:v375-schur-current-error-bound}
In Theorem~\ref{p10:thm:v375-two-energy-dressed-transition}, let
$M_2=\iota^\dagger H^2\iota$, $K=\iota^\dagger H\iota$,
$\delta_a=\operatorname{dist}(z_a,\operatorname{spec}C)>0$, and
\[
 \beta_a=\frac{\|B\|_{\rm op}}{\delta_a}
 =\frac{\|M_2-K^2\|_{\rm op}^{1/2}}{\delta_a},
 \qquad a\in\{i,f\}.
\]
Then the exact projected vertex obeys the quantitative bound
\[
\boxed{\|\mathcal V_J(z_f,z_i)-J_{PP}\|_{\rm op}
 \le\|J\|_{\rm op}(\beta_f+\beta_i+\beta_f\beta_i).}
\]
For real isolated poles outside $\operatorname{spec}C$,
$ I\preceq Z(E_a)\preceq(1+\beta_a^2)I$.
If the pole source vectors $u_a$ are unit-normalized, the full
normalized pole current differs from the corresponding bare
source contraction by at most
\[
\boxed{\begin{aligned}
\big|\langle\Psi_f,J\Psi_i\rangle-u_f^\dagger J_{PP}u_i\big|
\le\|J\|_{\rm op}\Bigl(\beta_f+\beta_i+\beta_f\beta_i
 +1-\frac1{\sqrt{(1+\beta_f^2)(1+\beta_i^2)}}\Bigr).
\end{aligned}}
\]
These are \emph{computable error certificates}, not assertions that
$\beta_a$ is small for a selected physical baryon pole.
\end{corollary}
\begin{proof}
Self-adjointness of $C$ gives $\|R_Q(z_a)\|=1/\delta_a$.
Apply submultiplicativity separately to the three non-bare terms
of $\mathcal V_J$, noting $\|J_{PQ}\|,\|J_{QP}\|,\|J_{QQ}\|\le\|J\|$.
The existing compression--leakage identity yields
$M_2-K^2=B^\dagger B$, giving the moment form of $\beta$.
On real resolvent gaps, $R_Q(E)^2\succeq0$, so
$1\le u^\dagger Z(E)u\le1+\beta(E)^2$ for every unit $u$.
Insert and subtract the bare source contraction in the normalized
pole formula, use the preceding vertex estimate, and bound the
normalization factor by its indicated interval.
\end{proof}

\begin{proposition}[Moving source-chart current and quantitative leakage bound]
\label{p10:prop:v375-moving-source-current}
Let $H(\phi)$ be a differentiable self-adjoint finite action family
on one fixed Hilbert-space trivialization, and
$\iota(\phi)^\dagger\iota(\phi)=I$ a differentiable, possibly
field-dependent source isometry. Put
$P=\iota\iota^\dagger$, $Q=I-P$,
$K=\iota^\dagger H\iota$, $B=QH\iota$,
$J=\partial_\phi H$, $A=\iota^\dagger\partial_\phi\iota$
(anti-Hermitian), and $\eta=Q\partial_\phi\iota$.
At the declared field value,
\[
\boxed{\partial_\phi K=\iota^\dagger J\iota
 +[K,A]+\eta^\dagger B+B^\dagger\eta,}
\]
with the sharp general estimate
\[
\boxed{\|\partial_\phi K-[K,A]-\iota^\dagger J\iota\|_{\rm HS}
 \le2\|B\|_{\rm op}\|\eta\|_{\rm HS}.}
\]
Here $\|B\|_{\rm op}^2=\|\iota^\dagger H^2\iota-K^2\|_{\rm op}$
by Theorem~\ref{p10:thm:v327-compression-leakage-pythagoras}.
A fixed source chart has $A=\eta=0$; if $B=0$ or $\eta=0$,
there is no extrinsic current correction beyond frame transport.
For a moving source-visible resolvent the complete derivative is
\[
 \partial_\phi[\iota^\dagger(z-H)^{-1}\iota]
 =\iota'^\dagger(z-H)^{-1}\iota
 +\iota^\dagger(z-H)^{-1}J(z-H)^{-1}\iota
 +\iota^\dagger(z-H)^{-1}\iota'.
\]
This is a source/field-placement certificate, not a license to
ignore derivatives of a physical Wilson, spin, measure, or localizer.
\end{proposition}
\begin{proof}
Differentiate $\iota^\dagger\iota=I$ to see $A^\dagger=-A$.
Decompose $\iota'=\iota A+\eta$ with $\iota^\dagger\eta=0$.
The product rule for $K$ gives the commutator and the two leakage
terms. Use the triangle inequality and
$\|\eta^\dagger B\|_{\rm HS}\le
\|\eta\|_{\rm HS}\|B\|_{\rm op}$.
The norm identity is the existing compression--leakage theorem.
Differentiating the full resolvent uses
$\partial_\phi(z-H)^{-1}=(z-H)^{-1}J(z-H)^{-1}$.
\end{proof}

\begin{theorem}[Exact projected Noether--Ward identity, charge residues and broken-current defect]
\label{p10:thm:v375-projected-Ward-current}
In the two-energy Schur theorem, let $Q_0=Q_0^*$ be a
microscopic charge/generator on the full finite action Hilbert
space such that the source projection preserves its charge sectors,
$[P,Q_0]=0$. Write $q=P Q_0 P|_{\operatorname{Ran}P}$ and
$D=[H,Q_0]$; $D$ is allowed to be anti-Hermitian.
The four-block vertex formula is algebraic and applies to $D$ as well
as to field derivatives. For any distinct $z_f,z_i$ outside the
relevant spectra, the exact projected Ward identity is
\[
\boxed{
 (z_f-z_i)\mathcal V_{Q_0}(z_f,z_i)
  =F(z_f)^{-1}q-qF(z_i)^{-1}
   +\mathcal V_D(z_f,z_i).}
\]
For isolated pole multiplets $E_f,E_i$ outside $\operatorname{spec}(C)$,
with the normalized degenerate source frames $U_{E_a}$ of
Corollary~\ref{p10:cor:v375-degenerate-multiplet-vertex}, this becomes
\[
 \boxed{(E_f-E_i)\,U_{E_f}^\dagger
 \mathcal V_{Q_0}(E_f,E_i)U_{E_i}
 =U_{E_f}^\dagger\mathcal V_D(E_f,E_i)U_{E_i}.}
\]
Hence every transition of a nearly conserved charge between
distinct-energy normalized poles satisfies the non-adjustable bound
\[
\boxed{|\langle\Psi_f,Q_0\Psi_i\rangle|
 \le\frac{\|[H,Q_0]\|_{\rm op}}{|E_f-E_i|}
 \qquad(E_f\ne E_i).}
\]
When $D=0$, the entire source-visible charge vertex is
\[
\boxed{
 \mathcal V_{Q_0}(z_f,z_i)
 =\frac{F(z_f)^{-1}-F(z_i)^{-1}}{z_f-z_i}\,q,
 \qquad \mathcal V_{Q_0}(z,z)=Z(z)q,}
\]
where $Z(z)=I+B^\dagger(z-C)^{-2}B$ and the coincident relation
is read in a resolvent gap. In particular the $Z$ factors cancel
exactly in the normalized pole charge: the full conserved charge
on each pole multiplet is precisely the representation of $q$;
no additional charge residue renormalization is possible. Distinct
charge sectors have zero charge-changing matrix elements.
This is a vector/symmetry Ward normalization and a quantitative
broken-current test, not an axial-Goldstone residue or a decay-width
formula. The source-invariance hypothesis $[P,Q_0]=0$ must be
verified for the actual action source.
\end{theorem}
\begin{proof}
Use $R(z)=(z-H)^{-1}$ and $R(z)H=zR(z)-I$ to compute
\[
 R(z_f)[H,Q_0]R(z_i)
 =(z_f-z_i)R(z_f)Q_0R(z_i)
   -Q_0R(z_i)+R(z_f)Q_0.
\]
Compress to $P$, where $Q_0$ is block diagonal, and use the exact
two-energy Schur factorization for $Q_0$ and for $D$. Multiplication
by $F(z_f)^{-1}$ and $F(z_i)^{-1}$ yields the first boxed identity.
Inserting the normalized pole source frames kills the inverse
Schur factors and proves the on-shell relation; Cauchy--Schwarz
bounds its right-hand side by $\|D\|_{\rm op}$.
If $D=0$, the charge commutes with $H$, $C$, and $F(z)$, so
$qF(z)^{-1}=F(z)^{-1}q$. The two-resolvent identity and its
coincident limit give $\mathcal V_{Q_0}(z,z)=Z(z)q$.
On an isolated eigenspace, its $Z$-orthonormal source frame is an
isometric copy of the full pole, and thus its charge is $q$
restricted to that eigenspace. This statement agrees with the
already established vector-flavour commutator conservation
(Theorem~\ref{p10:thm:vector-flavor-charge-conservation}); its new
content is the explicit all-energy Schur-vertex and residue identity.
\end{proof}

\begin{theorem}[Finite mixed-moment reconstruction of the entire source-current three-point function]
\label{p10:thm:v375-mixed-moment-current-reconstruction}
Let $H=H^*$ act on a $d$-dimensional finite Gauss-compatible
state space, $V:\mathbb C^s\to\mathcal H$ the selected action source,
and $J=J^*$ its selected field derivative. Define the exact
\emph{source moments} and \emph{mixed current moments}
\[
 M_k=V^\dagger H^kV\quad(0\le k\le2d-1),\qquad
 N_{mn}=V^\dagger H^mJH^nV\quad(0\le m,n<d).
\]
On block indices $(m,a),(n,b)$ for $0\le m,n<d$ form the
$d s\times d s$ moment matrices
\[
 S=(M_{m+n})_{m,n},\qquad S_1=(M_{m+n+1})_{m,n},
 \qquad T=(N_{mn})_{m,n}.
\]
Let $r=\operatorname{rank}S$, and choose a positive-rank eigendecomposition
$S=U_r\Lambda_rU_r^\dagger$, $\Lambda_r\succ0$.
Let $E_0:\mathbb C^s\to\mathbb C^{ds}$ include the $m=0$
source block. Define the fully moment-reconstructed matrices
\[
\boxed{\begin{aligned}
 h&=\Lambda_r^{-1/2}U_r^\dagger S_1U_r\Lambda_r^{-1/2},\\
 j&=\Lambda_r^{-1/2}U_r^\dagger TU_r\Lambda_r^{-1/2},\\
 v&=\Lambda_r^{1/2}U_r^\dagger E_0.
\end{aligned}}
\]
Then $r$ is exactly the action-source cyclic dimension and, for
all $z_f,z_i\in\operatorname{res}(H)$,
\[
\boxed{V^\dagger(z_f-H)^{-1}J(z_i-H)^{-1}V
 =v^\dagger(z_f-h)^{-1}j(z_i-h)^{-1}v.}
\]
Thus the complete finite \emph{one-current insertion} three-point
function, including every pole, residue and off-shell normalization
visible from $V$, is determined by finitely many exact same-action
source/current moments, uniquely up to an $r$-dimensional unitary
change of frame. No independent 56-to-240 embedding is required
for this observable. The theorem does not assert that $M_k,N_{mn}$
have already been computed from the action, and a two-current
four-point function may require additional higher mixed moments.
\end{theorem}
\begin{proof}
Let $Y=[V,HV,\ldots,H^{d-1}V]$. Then
$S=Y^\dagger Y$, $S_1=Y^\dagger HY$ and $T=Y^\dagger JY$.
The map $W=Y U_r\Lambda_r^{-1/2}$ is an isometry onto
$\mathcal K=\operatorname{ran}Y$, the source-cyclic space.
By Cayley--Hamilton, $H\mathcal K\subseteq\mathcal K$;
self-adjointness makes $\mathcal K$ reducing.
Therefore $h=W^\dagger HW$, $j=W^\dagger JW$ and
$v=W^\dagger V$ have the stated moment expressions. Since
$(z-H)^{-1}V=W(z-h)^{-1}v$ and
$V^\dagger(z-H)^{-1}=v^\dagger(z-h)^{-1}W^\dagger$,
compressing a single $J$ insertion gives the claimed identity.
Any other isometry onto $\mathcal K$ differs by a unitary, which
leaves the three-point response unchanged.
\end{proof}

% END PAPER00 v1.5

% BEGIN PAPER00 v1.6 -- transfer inversion and relative spectral port
% Proposed v1.6 VII-A addition. Exactly scoped finite positive-transfer theorem.
% Proposed v1.6 VII-A addition. Exactly scoped finite positive-transfer theorem.
\section{Transfer-source inversion and mass-to-current moment interface}
\label{p10:sec:v376-transfer-moment-inversion}
\begin{theorem}[Unique finite transfer-source inversion at every positive transfer time]
\label{p10:thm:v376-positive-transfer-current-inversion}
Let $H(\phi)=H(\phi)^\dagger$ be a differentiable finite
Gauss-compatible action member on a fixed carrier, and let
$c(\phi)\in\mathbb R$ be a differentiable scalar offset.  For
$t>0$ put $T_t(\phi)=\exp[-t(H(\phi)-c(\phi)I)]$.
At a fixed source value, diagonalize $H-cI=\sum_a\epsilon_a
|a\rangle\langle a|$ and put $J=H'$, $\widetilde J=J-c'I$.
Then the complete noncommuting source derivative has entries
\[
 \boxed{ (T_t')_{ab}=-\ell_{ab}(t)\widetilde J_{ab},\quad
 \ell_{ab}(t)=\int_0^t e^{-(t-s)\epsilon_a-s\epsilon_b}\,ds
 =t e^{-t(\epsilon_a+\epsilon_b)/2}
 \frac{\sinh[t(\epsilon_a-\epsilon_b)/2]}
 {t(\epsilon_a-\epsilon_b)/2}>0. }
\]
The quotient has its continuous value $1$ at zero, including
all degenerate eigenvalues. Consequently $T_t'$ determines
$\widetilde J$ \emph{uniquely} for every $t>0$, and the
reconstruction $-(T_t')_{ab}/\ell_{ab}(t)$ is independent of
$t$ for a common action source family. If
$\epsilon_{\min}\le\epsilon_a\le\epsilon_{\max}$, inversion
from a perturbation of the derivative satisfies the sharp uniform
Hilbert--Schmidt bound
\[
 \boxed{\|\Delta\widetilde J\|_{\rm HS}
 \le \frac{e^{t\epsilon_{\max}}}{t}\,\|\Delta T_t'\|_{\rm HS}.}
\]
For any two \emph{orthogonal} pole states, the scalar offset
$c'I$ does not contribute to their transition matrix element.
Therefore the normalized transfer $e^{-t(H-E_0)I}$, when its
source derivative exists, already fixes every off-diagonal
finite current without an independently chosen vacuum-energy
source derivative. A numerical current requires the actual
source-dependent $T_t(\phi)$; knowing $T_t(0)$ alone is
insufficient.
\end{theorem}
\begin{proof}
Duhamel's differentiation formula gives
$T_t'=-\int_0^t e^{-(t-s)(H-cI)}\widetilde J
 e^{-s(H-cI)}ds$.
Evaluate it in an eigenbasis of $H-cI$. The integral is strictly
positive for every pair of real eigenvalues; it has the displayed
hyperbolic-sine form, including a continuous diagonal limit.
Thus entrywise inversion is unique, even if $H$ and $J$ do not
commute. Since the integrand lies between
$e^{-t\epsilon_{\max}}$ and $e^{-t\epsilon_{\min}}$,
$\ell_{ab}\ge t e^{-t\epsilon_{\max}}$. Apply this entrywise
to the Hilbert--Schmidt norm; the bound is attained in an
extremal diagonal eigenspace. A scalar identity has vanishing
matrix element between orthogonal states. All transfer times
recover the same original $\widetilde J$ by the identity.
\end{proof}

\begin{proposition}[First mixed moment and the native leakage blocks]
\label{p10:prop:v376-first-mixed-moment-leakage}
For a fixed source projector $P$ with complement $Q$, write
$H=\left(\begin{smallmatrix}K&B^\dagger\\ B&C\end{smallmatrix}\right)$
and
$J=\left(\begin{smallmatrix}j&X\\ Y&D\end{smallmatrix}\right)$.
The action-owned mixed Krylov moments
$N_{mn}=P H^m J H^nP|_{P}$ satisfy
\[
\boxed{\begin{aligned}
N_{00}&=j,\\
N_{10}&=Kj+B^\dagger Y,\\
N_{01}&=jK+XB,\\
N_{11}&=KjK+KXB+B^\dagger YK+B^\dagger DB.
\end{aligned}}
\]
Moreover $M_2-K^2=B^\dagger B$ as already proved in
Theorem~\ref{p10:thm:v327-compression-leakage-pythagoras}, and
\[
 \boxed{\|N_{10}-Kj\|_{\rm op},\;
 \|N_{01}-jK\|_{\rm op}
 \le \|J\|_{\rm op}\,\|M_2-K^2\|_{\rm op}^{1/2}.}
\]
Thus the rooted CAR current fixes $N_{00}$ but higher mixed
moments require the actual Hamiltonian leakage and the same-action
current on its complement. These terms are computed matrix words,
not freely chosen decay strengths.
\end{proposition}
\begin{proof}
Block multiply $PH^m JH^nP$ for $(m,n)=(0,0),(1,0),
(0,1),(1,1)$. The leakage identity supplies
$B^\dagger B=M_2-K^2$, and
$\|B^\dagger Y\|,\|XB\|\le\|B\|\|J\|$.
\end{proof}

\begin{corollary}[Exact centered mixed moment isolates the doubly leaked current]
\label{p10:cor:v376-centered-double-leakage-current}
Under the preceding source/complement hypotheses define the
\emph{centered mixed-current moment}
\[
 \boxed{\mathcal X_J
 :=N_{11}-K N_{01}-N_{10}K+K N_{00}K
 =B^\dagger J_{QQ}B.}
\]
The bare source current and both single-leakage current blocks
cancel \emph{exactly}; no commutation of $H$ and $J$ or spin--flavor
symmetry is assumed. In particular,
\[
\boxed{\|\mathcal X_J\|_{\rm op}
 \le \|J_{QQ}\|_{\rm op}\,\|M_2-K^2\|_{\rm op}.}
\]
The centered expression is invariant under any common scalar
energy-origin change $H\mapsto H+cI$ at the same source value.
For an exactly spin--flavor-covariant $H$ with one octet and one
decuplet source copy, the $8\leftarrow10$ component is
proportional to $P_8G^{ia}P_{10}$ and its coefficient can be
extracted by a single $24$-current Hilbert--Schmidt contraction
with denominator $40$. The identity makes the \emph{auxiliary}
part of the transition computable from the action, but does not
supply its value before the action/current member is fixed.
\end{corollary}
\begin{proof}
Insert the four block-multiplication identities of
Proposition~\ref{p10:prop:v376-first-mixed-moment-leakage}
and cancel the terms involving $KjK$, $KXB$ and $B^\dagger YK$;
only $B^\dagger D B$ remains. The compression--leakage identity
and submultiplicativity give the norm bound. A scalar energy
origin changes $K$ by the same scalar but leaves the off-diagonal
leakage block $B$ unchanged. Exact spin--flavor covariance and
Schur's lemma imply the last reduced-tensor statement, with the
existing full-packet normalization
$\sum_{ia}\|P_8G^{ia}P_{10}\|_{\rm HS}^2=40$.
\end{proof}

% END PAPER00 v1.6
% BEGIN PAPER00 v1.7 -- all-order centered positive-moment Schur closure
\section{Positive moment certificates for the source-visible auxiliary current}
\label{p10:sec:v377-positive-moment-Schur}
\begin{theorem}[All-order centered auxiliary moments and source-cyclic Schur reconstruction]
\label{p10:thm:v377-auxiliary-krylov-centering}
Let $H=H^*$ act on a finite $d$-dimensional Gauss-compatible member,
$V:\mathbb C^s\to\mathcal H$ an isometric, \emph{fixed} source embedding,
$P=VV^\dagger$, $Q=I-P$, and $J=J^*$ the same-member field
Hamiltonian derivative.  For $m,n\ge0$ put
$M_n=V^\dagger H^nV$ and $N_{mn}=V^\dagger H^mJH^nV$,
with $M_0=I_s$, and for $m,n\ge1$ define
\[
\boxed{\begin{aligned}
S_{mn}&=M_{m+n}-M_mM_n,\\
\Xi_{mn}&=N_{mn}-M_mN_{0n}-N_{m0}M_n+M_mN_{00}M_n,\\
C_{mn}&=M_{m+n+1}-M_mM_{n+1}-M_{m+1}M_n+M_mM_1M_n.
\end{aligned}}
\]
Let $E_n=QH^nV$, $C=QHQ|_{Q\mathcal H}$, and $J_Q=QJQ|_{Q\mathcal H}$. Then
\[
 \boxed{S_{mn}=E_m^\dagger E_n,\qquad
 \Xi_{mn}=E_m^\dagger J_QE_n,\qquad
 C_{mn}=E_m^\dagger CE_n.}
\]
The span $\mathcal K_Q=\operatorname{span}\{E_n\mathbb C^s:1\le n<d\}$
is the smallest $C$-invariant auxiliary subspace containing $B=QHV$;
it is reducing for $C$. Therefore the \emph{complete source-visible Schur
Hamiltonian and current} are determined, up to unitary on $\mathcal K_Q$,
by $S_{mn},C_{mn},\Xi_{mn}$, $M_1,N_{00}$ and the cross-moment blocks
\[
 V^\dagger J E_n=N_{0n}-N_{00}M_n,\qquad
 E_n^\dagger J V=N_{n0}-M_nN_{00}.
\]
In particular there is no independent choice of a source-visible
auxiliary representation once these \emph{actual same-action moments}
are evaluated. This theorem does not compute the moments from the
Born--Hopf QCD action or identify a physical source-dependent embedding.
\end{theorem}
\begin{proof}
Write $E_n=QH^nV$. Inserting $Q=I-P$ into
$E_m^\dagger E_n$ and $E_m^\dagger JE_n$ gives the
first two identities, including all four centered terms.  Inserting
$Q$ on either side of $H$ gives the third.  The cross blocks follow
similarly.  Since $P H^nV=VM_n$,
\[
 CE_n=QHQH^nV=E_{n+1}-E_1M_n.
\]
Cayley--Hamilton expresses $E_d$ in the span of the preceding $E_n$,
so $\mathcal K_Q$ is invariant for the self-adjoint $C$, hence reducing.
Let $Y_Q=[E_1,\ldots,E_{d-1}]$ and take a positive-rank
factorization $S=U_r\Lambda_rU_r^\dagger=Y_Q^\dagger Y_Q$.
Then $W_Q=Y_QU_r\Lambda_r^{-1/2}$ is an isometry onto
$\mathcal K_Q$.  Its auxiliary Hamiltonian and current are respectively
$\Lambda_r^{-1/2}U_r^\dagger(C_{mn})U_r\Lambda_r^{-1/2}$ and
$\Lambda_r^{-1/2}U_r^\dagger(\Xi_{mn})U_r\Lambda_r^{-1/2}$.
The leakage and cross blocks are recovered from the first $E_1$ block
and the displayed cross moments.  Self-adjoint reduction makes
$P\mathcal H\oplus\mathcal K_Q$ reducing for $H$; compressing $J$ onto
that same space leaves every source-to-source two-energy current
correlator unchanged.  Different factorization frames differ by a unitary.
\end{proof}

\begin{corollary}[Sharp semidefinite feasibility envelope for auxiliary currents]
\label{p10:cor:v377-auxiliary-current-gram-envelope}
Under Theorem~\ref{p10:thm:v377-auxiliary-krylov-centering},
for any $c\ge\|J_Q\|_{\rm op}$, the complete centered moment matrix
$\Xi=(\Xi_{mn})_{m,n\ge1}$ obeys
\[
 \boxed{-cS\preceq\Xi\preceq cS.}
\]
Conversely, for a fixed positive $S$ and a Hermitian $\Xi$ satisfying
these two inequalities, there exists a Hermitian auxiliary current of
norm at most $c$ on the minimal $S$-rank carrier realizing exactly
$\Xi$. Thus the interval is sharp when \emph{only} the auxiliary-current
norm is bounded. It is not an independent physical current-selection law;
a specified full action may restrict the interval further through its
cross blocks, symmetry, and source jets.
\end{corollary}
\begin{proof}
For $Y_Q$ as above, $S=Y_Q^\dagger Y_Q$ and
$\Xi=Y_Q^\dagger J_QY_Q$. Positivity of $cI\pm J_Q$
gives both inequalities by congruence. Conversely they force
$\ker S\subseteq\ker\Xi$. On the positive eigenspace
$S=U_r\Lambda_rU_r^\dagger$ define
$j=\Lambda_r^{-1/2}U_r^\dagger\Xi U_r\Lambda_r^{-1/2}$.
Then $-cI\preceq j\preceq cI$, and the isometry
$W_Q=Y_QU_r\Lambda_r^{-1/2}$ yields the realizing current
$W_QjW_Q^\dagger$, extended by zero outside the source-cyclic
auxiliary carrier. Taking $j=\pm cI$ saturates either bound.
\end{proof}

\begin{theorem}[Four-moment positive flatness and exact single-pole current]
\label{p10:thm:v377-four-moment-flatness}
Keep the fixed finite source and write
$H=\left(\begin{smallmatrix}K&B^\dagger\\ B&C\end{smallmatrix}\right)$.
The following positive auxiliary moment matrices are determined
\emph{only} by $M_1,M_2,M_3,M_4$:
\[
\begin{aligned}
T_0&=B^\dagger B=M_2-K^2,\\
T_1&=B^\dagger CB=M_3-KM_2-M_2K+K^3,\\
T_2&=B^\dagger C^2B
 =S_{22}-S_{21}K-KS_{12}+KS_{11}K.
\end{aligned}
\]
Assume $\operatorname{tr}T_0>0$ and define the \emph{moment-selected},
not independently chosen, auxiliary energy
$\Delta_*:=\operatorname{tr}T_1/\operatorname{tr}T_0$.
Then
\[
\boxed{\mathcal V_*:=T_2-2\Delta_*T_1+\Delta_*^2T_0
 =B^\dagger(C-\Delta_*I)^2B\succeq0.}
\]
The following statements are equivalent: (i) $\mathcal V_*=0$;
(ii) $(C-\Delta_*I)B=0$; (iii) the complete source-visible
Schur self-energy has the exact all-energy form
$\Sigma(z)=T_0/(z-\Delta_*)$.
If they hold, the \emph{complete} current vertex of
Theorem~\ref{p10:thm:v375-two-energy-dressed-transition} is
\[
\boxed{\begin{aligned}
\mathcal V_J(z_f,z_i)={}&N_{00}
 +\frac{N_{01}-N_{00}K}{z_i-\Delta_*}
 +\frac{N_{10}-KN_{00}}{z_f-\Delta_*}\\
 &+\frac{\Xi_{11}}{(z_f-\Delta_*)(z_i-\Delta_*)}.
\end{aligned}}
\]
No unlisted auxiliary pole or channel coefficient may be appended to
these \emph{source-visible finite} responses. Without exact flatness,
for $z\notin\operatorname{spec}C\cup\{\Delta_*\}$,
\[
\boxed{\left\|\Sigma(z)-\frac{T_0}{z-\Delta_*}\right\|_{\rm op}
\le\frac{\sqrt{\|T_0\|_{\rm op}}\,
\sqrt{\operatorname{tr}\mathcal V_*}}
{|z-\Delta_*|\,\operatorname{dist}(z,\operatorname{spec}C)}.}
\]
The exact certificate and the error bound are finite mathematical
statements. The physical QCD branch must first supply its own moments;
no claim is made that its auxiliary response is single-pole.
\end{theorem}
\begin{proof}
The first two formulas follow by multiplying $H^2,H^3$.
For $E_1=B$ and $E_2=BK+CB$, the all-order Gram identity gives
$B^\dagger C^2B=(E_2-E_1K)^\dagger(E_2-E_1K)$,
which expands to the stated $T_2$. Positivity gives the equivalence
between (i) and (ii). Condition (ii) means the leakage is supported
entirely on the $\Delta_*$ eigenspace of the self-adjoint $C$;
hence (iii). Conversely equality of the resolvents near infinity
fixes $T_1=\Delta_*T_0$ and $T_2=\Delta_*^2T_0$, giving (i).
Insert $(z-C)^{-1}B=B/(z-\Delta_*)$ and its adjoint-side
algebraic identity into all four blocks of the complete Schur
vertex, then use $N_{01}-N_{00}K=J_{PQ}B$,
$N_{10}-KN_{00}=B^\dagger J_{QP}$, and
$\Xi_{11}=B^\dagger J_{QQ}B$.
Finally use the resolvent identity
\[
(z-C)^{-1}B-\frac{B}{z-\Delta_*}
 =(z-C)^{-1}(C-\Delta_*)B\,(z-\Delta_*)^{-1}
\]
and the norm bounds
$\|B\|_{\rm op}=\sqrt{\|T_0\|_{\rm op}}$,
$\|(C-\Delta_*)B\|_{\rm HS}=\sqrt{\operatorname{tr}\mathcal V_*}$.
\end{proof}

\begin{theorem}[Positive polynomial flatness identifies every visible auxiliary pole]
\label{p10:thm:v377-multipole-flatness}
For the same finite self-adjoint $C$, define
$T_n=B^\dagger C^nB$ for $n\ge0$. All $T_n$ are recursively
computable from the exact mass moments via the large-$z$ identity
\[
 \Sigma(z)=zI-K-\left(\sum_{n\ge0}M_nz^{-n-1}\right)^{-1}
 =\sum_{n\ge0}T_nz^{-n-1}.
\]
For any nonzero real polynomial $p(x)=\sum_{j=0}^r a_jx^j$,
\[
 \boxed{\mathcal F_p:=\sum_{m,n=0}^r a_ma_nT_{m+n}
 =B^\dagger p(C)^2B\succeq0.}
\]
The finite exact condition $\mathcal F_p=0$ is equivalent to
$p(C)B=0$. It proves that all source-visible auxiliary poles lie
among the real zeros of $p$, and for distinct visible zeros
$\lambda_1,\ldots,\lambda_q$ it gives the unique positive return
\[
 \boxed{\Sigma(z)=\sum_{j=1}^q\frac{W_j}{z-\lambda_j},
 \qquad W_j\succeq0,\quad
 T_n=\sum_j\lambda_j^nW_j.}
\]
The weights $W_j$ follow from the first $q$ auxiliary moments by
Vandermonde inversion. If the actual auxiliary current moments
$L_{mn}=B^\dagger C^mJ_{QQ}C^nB$ are also supplied for
$0\le m,n<q$, the entire doubly leaked two-energy current is
uniquely reconstructed:
\[
 \boxed{B^\dagger(z_f-C)^{-1}J_{QQ}(z_i-C)^{-1}B
 =\sum_{j,k=1}^q\frac{R_{jk}}{(z_f-\lambda_j)(z_i-\lambda_k)},}
\]
where the $R_{jk}$ are obtained by two-sided Vandermonde inversion.
These are \emph{finite exact} source responses; no physical continuum
branching width is inferred from discrete real poles.
\end{theorem}
\begin{proof}
The Schur identity and large-$z$ resolvent expansion determine every
$T_n$ recursively from the source moments. Self-adjointness yields
$\mathcal F_p=(p(C)B)^\dagger(p(C)B)$ as a positive operator.
Zero therefore means $p(C)B=0$. The spectral theorem then
restricts the cyclic subspace generated by $B$ to eigenspaces
at the real roots of $p$, with positive weights
$W_j=B^\dagger P_{\lambda_j}(C)B$.
The Vandermonde matrix on distinct $\lambda_j$ is invertible.
For a current insertion insert the same spectral resolution on
both sides of $J_{QQ}$; its coefficient matrix is recovered from
$L_{mn}=\sum_{j,k}\lambda_j^m\lambda_k^nR_{jk}$ by
left and right Vandermonde inversion.
\end{proof}

\begin{corollary}[Exact local source return for the native Hurwitz adjacency]
\label{p10:cor:v377-hurwitz-local-spectral-return}
For the already established $24$-vertex Hurwitz adjacency $A_H$ of
Theorem~\ref{p10:thm:hurwitz-shell-adjacency-spectrum}, a normalized
vertex source $e_q$ has the exact spectral weights
\[
\boxed{\langle e_q,P_8e_q\rangle=\tfrac1{24},\quad
 \langle e_q,P_4e_q\rangle=\tfrac16,\quad
 \langle e_q,P_0e_q\rangle=\tfrac38,\quad
 \langle e_q,P_{-2}e_q\rangle=\tfrac13,\quad
 \langle e_q,P_{-4}e_q\rangle=\tfrac1{12}.}
\]
Consequently the complete finite local graph return is
\[
 \boxed{g_q(z)=\frac{1}{24(z-8)}+\frac{1}{6(z-4)}
 +\frac{3}{8z}+\frac{1}{3(z+2)}+\frac{1}{12(z+4)}.}
\]
For distinct vertices the off-diagonal residues depend only on the
native quaternionic inner product, and are exactly:
\[
\begin{array}{c|rrrrr}
\langle q,q'\rangle&P_8&P_4&P_0&P_{-2}&P_{-4}\\\hline
1&1/24&1/6&3/8&1/3&1/12\\
1/2&1/24&1/12&0&-1/12&-1/24\\
0&1/24&0&-1/8&0&1/12\\
-1/2&1/24&-1/12&0&1/12&-1/24\\
-1&1/24&-1/6&3/8&-1/3&1/12
\end{array}
\]
These are exact \emph{geometric adjacency} source-to-source kernels.
They are not interacting QCD pole energies, quark masses, or decay widths.
\end{corollary}
\begin{proof}
The $24$ vertices comprise the eight axial unit quaternions and
the sixteen half-unit sign quaternions; adjacency is the exact
inner-product relation $\langle q,q'\rangle=1/2$.
For eigenvalues $8,4,0,-2,-4$, evaluate the exact spectral
polynomials $P_\lambda=\prod_{\mu\ne\lambda}
(A_H-\mu I)/(\lambda-\mu)$ on the $24\times24$ adjacency.
Their ranks are $1,4,9,8,2$ as in the existing spectrum theorem,
and their entries are constant on each indicated inner-product
relation, giving the displayed rational matrix. The diagonal
entries also follow by vertex transitivity as rank$/24$.
The spectral resolution of $(z-A_H)^{-1}$ gives the return.
\end{proof}
% END PAPER00 v1.7

\end{document}
